% In this file you should put the actual content of the blueprint.
% It will be used both by the web and the print version.
% It should *not* include the \begin{document}
%
% If you want to split the blueprint content into several files then
% the current file can be a simple sequence of \input. Otherwise It
% can start with a \section or \chapter for instance.

\part{The power of adaptivity for \texorpdfstring{$\eps$}{epsilon}-best arm identification in linear bandits}
\label{part:main}

\textbf{Overview}

This blueprint guides the Lean formalization of the paper
\emph{On the Power of Adaptivity for $\eps$-Best Arm Identification in Linear Bandits}
by Arnab Maiti, Yunbei Xu and Kevin Jamieson (COLT 2026). The paper studies the minimax
sample complexity of $\eps$-best arm identification in linear bandits with Gaussian noise:
a compact action set $\Xs \subset \R^d$ spanning $\R^d$, an unknown reward vector $\theta$,
observations $y_t = \ip{x_t}{\theta} + \eta_t$, and the goal of returning $\rec \in \Xs$ with
$\max_{x \in \Xs} \ip{x - \rec}{\theta} \le \eps$ with probability at least $1 - \delta$.
Its results are:
\begin{enumerate}
  \item[(i)] a non-adaptive fixed-design algorithm with sample complexity
        $O\big((d\log(1/\delta) + \gw(\Xs)^2)/\eps^2\big)$, where $\gw(\Xs)$ is a Gaussian
        width term of the action set (Theorem~\ref{thm:upper});
  \item[(ii)] matching lower bounds: $\Omega(d\log(1/\delta)/\eps^2)$ for every algorithm,
        adaptive or not (Theorem~\ref{thm:lower_adaptive}), and
        $\Omega(\gw(\Xs)^2/\eps^2)$ for non-adaptive fixed designs
        (Theorem~\ref{thm:lower_nonadaptive});
  \item[(iii)] properties of the Gaussian width term and its computation for structured
        action sets (Chapter~\ref{chap:width});
  \item[(iv)] structured action sets on which adaptivity yields at most logarithmic
        improvements (Chapter~\ref{chap:log_gains}), through a Median-Elimination based
        adaptive algorithm and adaptive lower bounds;
  \item[(v)] an $\ell_2$-norm estimation procedure with sample complexity
        $O(d\log(1/\delta)/\eps^2)$ (Theorem~\ref{thm:norm_estimation});
  \item[(vi)] an action set on which adaptivity yields a polynomial-factor improvement over
        every non-adaptive algorithm (Theorem~\ref{thm:separation}).
\end{enumerate}

\textbf{Goals of the formalization.} Two goals drive the organization of this document.
First, formalize the main results of the paper. Second, build a reusable library, to be
contributed to \emph{Lean Machine Learning} (LML) and Mathlib: the objects (linear bandits,
PAC identification algorithms, fixed designs, optimal experimental design, Gaussian width,
Gaussian concentration, sub-exponential variables, divergence decompositions, Median
Elimination) are therefore introduced in the generality in which they are useful, even when
the paper only needs a special case, and extra work is done when needed to reach that
generality.

\textbf{Deviations from the paper.} The blueprint departs from the paper whenever a
different route is simpler to formalize while giving the same (or a stronger) statement;
every such point is recorded in a remark next to the statement it concerns. The main ones:
\begin{itemize}
  \item The rounding of a design distribution to a fixed design of size $O(d)$ is proved by
        a one-sided barrier argument in the style of Batson--Spielman--Srivastava
        (Section~\ref{sec:rounding}) instead of the cited rounding procedures; the constant
        $360$ of the paper becomes another absolute constant.
  \item The Sudakov--Fernique inequality is only used in its ``covariance domination''
        special case, which has a one-line proof by Jensen's inequality
        (Lemma~\ref{lem:gauss_comparison}); the general inequality is only needed for the
        secondary lower bound $\gw(\Xs) \ge \Omega(\sqrt{d \log d})$
        (Chapter~\ref{chap:pre_sudakov}).
  \item John's theorem in the proof of the adaptive lower bound is replaced by the
        Kiefer--Wolfowitz theorem on $G$-optimal designs, which is needed anyway.
  \item The Bayesian argument of the non-adaptive lower bound is done without Gaussian
        conditioning, through an orthogonal decomposition and independence of uncorrelated
        jointly Gaussian vectors.
  \item Borell--TIS is obtained from the Maurey--Pisier Gaussian concentration argument,
        with the constant $\cG = \pi^2/4$ in the exponent instead of $1$.
  \item The large-norm regime of the $\ell_2$-norm estimation uses a deterministic
        basis design instead of random Rademacher rows, which removes the need for
        random-matrix concentration.
  \item The lower bounds for structured sets are proved directly for randomized algorithms
        (no Yao principle) and with a single argument for all block sizes.
  \item Non-adaptive algorithms are deterministic fixed designs with an arbitrary (possibly
        randomized) recommendation rule, the class of the paper's Theorem~3; the
        $\Omega(kd^2/\eps^2)$ lower bound of the polynomial separation
        (Theorem~\ref{thm:separation_lower_nonadaptive}) is stated for this class, whereas the
        paper's Theorem~7 says ``possibly randomized non-adaptive''. Budgets are
        deterministic (Remark~\ref{rem:stopping_times}).
  \item The rank-one Hanson--Wright inequality used for the Rademacher directions is replaced
        by an MGF bound for the square of a sub-Gaussian variable
        (Lemma~\ref{lem:rademacher_square_subexp}), and the Gaussian Hanson--Wright inequality
        by diagonalization (Lemma~\ref{lem:gauss_quadratic_form_subexp}); the paper's
        tail-to-MGF lemma is not needed.
\end{itemize}

\textbf{Reading guide.}
Part~\ref{part:main} follows the paper: Chapter~\ref{chap:model} sets up the model,
Chapter~\ref{chap:design} the experimental design tools, Chapters~\ref{chap:upper}
to~\ref{chap:lower_adaptive} the fixed-design upper bound and the lower bounds,
Chapter~\ref{chap:width} the Gaussian width, Chapter~\ref{chap:log_gains} the structured
sets, Chapter~\ref{chap:norm} the $\ell_2$-norm estimation and Chapter~\ref{chap:separation}
the polynomial separation.
Part~\ref{part:prereq} collects the probabilistic and information-theoretic prerequisites
that are not in Mathlib or LML, decomposed into small lemmas: Gaussian vectors and linear
Gaussian processes, Gaussian concentration, sub-exponential variables, divergences of
algorithm-environment sequences, the Sudakov--Fernique inequality, Median Elimination and
Gaussian posteriors of adaptively collected data.

\textbf{Conventions.} An \emph{action set} is a nonempty compact subset of $\R^d$; the
paper's standing assumption $\Span(\Xs) = \R^d$ (``spanning'') is stated explicitly in the
results that need it (designs, widths, upper bounds), because the structured sets of
Chapter~\ref{chap:log_gains} are not spanning. Time is $0$-indexed, as in LML: an algorithm with budget $T$ plays the
actions $x_0, \dots, x_{T-1}$ and sums over $t < T$. All random objects live on a probability
space $(\Omega, \mathcal{F}, \Prob)$ and $\E$ is the expectation under $\Prob$; when the
reward vector matters we write $\Prob_\theta$, $\E_\theta$. Norms are Euclidean unless stated
otherwise, and for a positive definite matrix $W$ we write $\norm{x}_W^2 := x^\top W x$.

\textbf{Status.} Nothing is formalized yet. Each statement will carry a \verb+\lean+ tag naming
its Lean declaration once formalized.

\textbf{Blueprint author.} Rémy Degenne. Anyone is welcome to contribute!

\chapter{Model: linear bandits, identification algorithms and the Gaussian width}
\label{chap:model}

This chapter fixes the objects of the paper (Section~1 and the notations of Section~1.1),
in the form in which they will be formalized on top of the LML framework of sequential
algorithms and environments (\texttt{Learning.Algorithm}, \texttt{Learning.Environment},
\texttt{Learning.IsAlgEnvSeq}, \texttt{Learning.stationaryEnv}).

\section{Action sets and the linear Gaussian environment}
\label{sec:model_env}

\begin{definition}[Action set]\label{def:arm_set}
Let $d \ge 1$. An \emph{action set} (or arm set) is a nonempty compact subset $\Xs \subset \R^d$.
It is \emph{spanning} if $\Span(\Xs) = \R^d$. The paper assumes spanning throughout; in this
blueprint the spanning assumption is stated explicitly in every statement that needs it (the
existence of positive definite designs, the Gaussian width, the upper bounds), and the
definitions of environments, algorithms, regret and PAC guarantees below make sense for every
action set, spanning or not (the structured sets of Chapter~\ref{chap:log_gains} are not
spanning). We write $\ball_d := \{x \in \R^d : \norm{x} \le 1\}$ for the unit
Euclidean ball and $\sphere^{d-1}$ for the unit sphere. For $\theta \in \R^d$ the
\emph{value} of an arm $x$ is $\ip{x}{\theta}$ and an \emph{optimal arm} is
$x_\star(\theta) \in \argmax_{x \in \Xs} \ip{x}{\theta}$ (it exists by compactness).
\end{definition}

\textbf{Lean remark.} The ambient space is \texttt{EuclideanSpace ℝ (Fin d)}. The action
type of the LML algorithms is the subtype $\Xs$ with its Borel $\sigma$-algebra, so that
algorithms play in $\Xs$ by typing. Compactness is only used through the existence of
maximizers of continuous functions and the finiteness of $\sup_{x \in \Xs}\norm{x}$; the
spanning assumption is used through the existence of positive definite design matrices
(Lemma~\ref{lem:exists_posdef_design}) and is a separate hypothesis
(\texttt{Submodule.span ℝ 𝒳 = ⊤}) of the statements that need it. Finite action sets are the important special case,
and several statements are first proved for finite $\Xs$ and then extended by density
(Lemma~\ref{lem:sup_countable_dense}).

\begin{definition}[Linear Gaussian environment]\label{def:env}
\uses{def:arm_set}
For $\theta \in \R^d$, the \emph{linear Gaussian environment} $\nu_\theta$ is the stationary
environment (LML \texttt{stationaryEnv}) with reward kernel
$x \mapsto \Normal(\ip{x}{\theta}, 1)$, $x \in \Xs$: at each round $t$ the learner plays
$x_t \in \Xs$ and observes $y_t = \ip{x_t}{\theta} + \eta_t$ where the $\eta_t$ are i.i.d.\
standard Gaussian, independent of the past and of the learner's randomness.
\end{definition}

\begin{lemma}[Noise representation]\label{lem:noise_representation}
\uses{def:env}
Let $(x_t, y_t)_{t < T}$ be an algorithm-environment sequence for an algorithm with actions
in $\Xs$ and the environment $\nu_\theta$, on $(\Omega, \Prob)$. Then
$\eta_t := y_t - \ip{x_t}{\theta}$, $t < T$, are i.i.d.\ $\Normal(0,1)$ and, for every $t$,
$\eta_t$ is independent of the pair $\big((x_s, y_s)_{s < t},\, x_t\big)$ (jointly, i.e.\ of the
$\sigma$-algebra generated by the history up to $t-1$ together with $x_t$).
\end{lemma}

\begin{proof}
\uses{def:env}
The LML statement \texttt{IsAlgEnvSeq.hasCondDistrib\_feedback} gives that the conditional law
of $y_t$ given the history up to $t-1$ and $x_t$ is $\Normal(\ip{x_t}{\theta},1)$; shifting by
the measurable function $\ip{x_t}{\theta}$ of the conditioning variables gives that the
conditional law of $\eta_t$ is $\Normal(0,1)$, a constant kernel, which is the claimed
independence. Joint independence of the $\eta_t$ follows by induction on $t$.
\end{proof}

\section{Identification algorithms and PAC guarantees}
\label{sec:model_pac}

\begin{definition}[Identification algorithm with budget $T$]\label{def:bai_alg}
\uses{def:env}
An \emph{identification algorithm with budget $T \ge 0$} on $\Xs$ is a pair
$\mathsf{Alg} = (\mathrm{alg}, \rho)$ of an LML algorithm $\mathrm{alg}$ with actions in $\Xs$
and observations in $\R$ (a sampling rule), and a \emph{recommendation rule}: a Markov kernel
$\rho$ from histories of length $T$ (the sequence $(x_t, y_t)_{t<T}$) to $\Xs$. Run against
$\nu_\theta$, it produces the history $\hist_T = (x_t,y_t)_{t<T}$ and then the recommendation
$\rec \sim \rho(\hist_T)$. We write $\Prob_\theta$ for the law of $(\hist_T, \rec)$ (unique by
\texttt{Learning.isAlgEnvSeq\_unique}).
\end{definition}

\textbf{Lean remark.} The randomness of the recommendation is allowed on purpose (the lower
bounds are stated for randomized procedures). Budget $T = 0$ is allowed: the history of length
$0$ lives in a one-point space (\texttt{Fin 0 → 𝒳 × ℝ}) and the recommendation kernel is then
just a probability measure on $\Xs$; this degenerate case is needed by the singular-design lemma
(Lemma~\ref{lem:singular_design_fails}) and its use in Chapter~\ref{chap:separation}. Histories
of length $T \ge 1$ are LML's \texttt{Iic (T-1) → 𝒳 × ℝ}. A deterministic recommendation is the kernel
\texttt{Kernel.deterministic}. Stopping times are not modelled: the paper's PAC algorithms use
a deterministic budget $T = (H_1 \log(1/\delta) + H_2)/\eps^2$ and all its lower bounds are
stated for such budgets; a randomized budget can be handled by padding
(Remark~\ref{rem:stopping_times}).

\begin{definition}[Simple regret]\label{def:regret}
\uses{def:arm_set}
The \emph{simple regret} of $\rec \in \Xs$ for the reward vector $\theta$ is
\[
  r(\rec, \theta) := \max_{x \in \Xs} \ip{x}{\theta} - \ip{\rec}{\theta} \;\ge 0 .
\]
We also set $Z(\theta) := \max_{x, y \in \Xs} \ip{x - y}{\theta}$, so that
$r(\rec,\theta) \le Z(\theta)$ for every $\rec \in \Xs$.
\end{definition}

\begin{definition}[$(\eps,\delta)$-PAC algorithm]\label{def:pac}
\uses{def:bai_alg, def:regret}
Let $\eps > 0$ and $\delta \in (0,1)$. An identification algorithm with budget $T$ is
\emph{$(\eps,\delta)$-PAC on $\Xs$} if for every $\theta \in \R^d$,
\[
  \Prob_\theta\big( r(\rec, \theta) \le \eps \big) \ge 1 - \delta .
\]
\end{definition}

\begin{lemma}[Expected regret of a PAC algorithm under a prior]\label{lem:pac_expected_regret}
\uses{def:pac}
Let $\mathsf{Alg}$ be $(\eps,\delta)$-PAC with budget $T$ and let $\pi$ be a probability
measure on $\R^d$ such that $\int Z(\theta)\,\pi(d\theta) < \infty$. Let $\kappa$ be a Markov
kernel from $\R^d$ to the space of pairs $(\hist_T, \rec)$ such that $\kappa(\theta) = \Prob_\theta$
for $\pi$-almost every $\theta$ (``run $\mathsf{Alg}$ against $\nu_\theta$ conditionally on
$\theta$''), and let $(\theta, \hist_T, \rec) \sim \pi \otimes \kappa$. Then
\[
  \E\big[ r(\rec, \theta) \big] \le \eps + \delta \int Z(\theta)\, \pi(d\theta) .
\]
\end{lemma}

\begin{proof}
\uses{def:regret}
Pathwise, $r \le \eps + Z(\theta)\indic\{r > \eps\}$. Taking expectations under
$\pi \otimes \kappa$ (Fubini for the composition-product, Mathlib
\texttt{MeasureTheory.Measure.integral\_compProd}),
$\E[Z(\theta)\indic\{r>\eps\}] = \int Z(\theta)\,\kappa(\theta)(r > \eps)\,\pi(d\theta)
= \int Z(\theta)\,\Prob_\theta(r > \eps)\,\pi(d\theta) \le \delta \int Z\,d\pi$.

\textbf{Lean remark.} The kernel $\kappa$ is a hypothesis, not a construction: measurability of
$\theta \mapsto \Prob_\theta$ for LML trajectory measures is not available off the shelf. Both
uses of the lemma supply $\kappa$ explicitly: for a fixed design
$\kappa(\theta) = \Normal(X\theta, I_T) \otimes \rho$ (Lemma~\ref{lem:fixed_design_law}), and
Chapter~\ref{chap:pre_bayes} builds $\kappa$ from the likelihood ratio of
Lemma~\ref{lem:env_likelihood_ratio}.
\end{proof}

\begin{remark}[Stopping times]\label{rem:stopping_times}
An algorithm with a stopping time $\tau$ and $\E[\tau] \le T_0$ can be turned into one with
deterministic budget $T = T_0/\delta$ by padding with dummy rounds: by Markov's inequality
$\Prob(\tau > T) \le \delta$, so an $(\eps,\delta)$-PAC algorithm with random budget becomes
$(\eps, 2\delta)$-PAC with budget $T_0/\delta$. All lower bounds below therefore transfer to
expected budgets up to the factor $1/\delta$, which is what the paper does in
Appendix~I. This reduction is not part of the formalization plan.
\end{remark}

\section{Non-adaptive fixed designs}
\label{sec:model_fixed}

\begin{definition}[Fixed-design algorithm]\label{def:fixed_design}
\uses{def:bai_alg}
A \emph{fixed design} of length $T$ is a deterministic sequence $x_0, \dots, x_{T-1} \in \Xs$.
The associated \emph{fixed-design (non-adaptive) algorithm} is the identification algorithm
whose sampling rule is deterministic and ignores the observations, playing $x_t$ at round
$t$ (LML \texttt{detAlgorithm} with a history-independent next-action function), together with
an arbitrary recommendation kernel $\rho$ from $(y_t)_{t<T} \in \R^T$ to $\Xs$ (budget
$T \ge 0$; for $T = 0$ the design is empty and $\rho$ is a probability measure on $\Xs$). We write
$X \in \R^{T \times d}$ for the matrix with rows $x_t^\top$ and
$\Sigma_T := X^\top X = \sum_{t<T} x_t x_t^\top$ for its \emph{design matrix}. Since the actions
are deterministic, the history $\hist_T$ is the image of $y = (y_t)_{t<T}$ under the measurable
bijection $y \mapsto ((x_t, y_t))_{t<T}$; we identify $(\hist_T, \rec)$ with $(y, \rec)$ and
write $\Prob_\theta$ for the law of either.
\end{definition}

\begin{lemma}[Law of the observations under a fixed design]\label{lem:fixed_design_law}
\uses{def:fixed_design, lem:noise_representation}
Under a fixed design run against $\nu_\theta$, the observation vector
$y = (y_t)_{t<T}$ satisfies $y = X\theta + \eta$ with $\eta \sim \Normal(0, I_T)$, and
the recommendation is $\rec \sim \rho(y)$. In particular the law of $(y, \rec)$ is
$\Normal(X\theta, I_T) \otimes \rho$.
\end{lemma}

\begin{proof}
\uses{lem:noise_representation, lem:gauss_pi}
Lemma~\ref{lem:noise_representation} with deterministic actions: the $\eta_t$ are i.i.d.\
standard Gaussians, so $\eta \sim \Normal(0,I_T)$ (Lemma~\ref{lem:gauss_pi}).
\end{proof}

\section{Design matrices and the Gaussian width}
\label{sec:model_width}

\begin{definition}[Design distributions and design matrices]\label{def:design_matrix}
\uses{def:arm_set}
A \emph{design distribution} on $\Xs$ is a finitely supported probability vector
$\lambda = (\lambda_x)_{x \in S}$, $S \subset \Xs$ finite, $\lambda_x \ge 0$,
$\sum_{x \in S} \lambda_x = 1$; we write $\lambda \in \simplex_{\Xs}$ and
$\supp(\lambda) := \{x \in S : \lambda_x > 0\}$. Its \emph{design matrix} is
\[
  A(\lambda) := \sum_{x \in S} \lambda_x\, x x^\top \in \PSD .
\]
The \emph{set of design matrices} is $\designset(\Xs) := \{A(\lambda) : \lambda \in
\simplex_{\Xs}\} = \conv\{x x^\top : x \in \Xs\}$. Finally, for a fixed design
$x_0,\dots,x_{T-1}$ the \emph{empirical design distribution} is
$\lambda_T := \frac1T \sum_{t<T} \delta_{x_t}$, so that $\Sigma_T = T\, A(\lambda_T)$.
\end{definition}

\textbf{Lean remark.} A design distribution is a finitely supported function
$\Xs \to [0,1]$ summing to $1$ (\texttt{Finsupp}, or a \texttt{Finset} with a weight function;
zero weights are allowed and $\supp$ is the set of positive weights). Restricting to finitely
supported designs loses nothing: by
Carathéodory's theorem (Lemma~\ref{lem:caratheodory_design}) every barycenter of
$\{xx^\top : x \in \Xs\}$ under a probability measure on $\Xs$ is also $A(\lambda)$ for a
$\lambda$ supported on at most $d(d+1)/2 + 1$ points. Finitely supported designs are what
the fixed-design construction consumes.

\begin{lemma}[Basic properties of the set of design matrices]\label{lem:designset_basic}
\uses{def:design_matrix}
$\designset(\Xs)$ is a convex compact subset of the symmetric matrices, contained in $\PSD$,
and every $A \in \designset(\Xs)$ satisfies $\tr(A) \le \sup_{x\in\Xs}\norm{x}^2$.
\end{lemma}

\begin{proof}
\uses{def:design_matrix}
The map $x \mapsto xx^\top$ is continuous, so $\{xx^\top : x\in\Xs\}$ is compact; the convex
hull of a compact subset of a finite-dimensional space is compact
(Mathlib \texttt{Set.Finite.isCompact\_convexHull} / \texttt{isCompact\_convexHull}).
Positive semidefiniteness and the trace bound are preserved by convex combinations.
\end{proof}

\begin{lemma}[Existence of a positive definite design]\label{lem:exists_posdef_design}
\uses{def:design_matrix}
If $\Xs$ is spanning, there exists $\lambda \in \simplex_{\Xs}$ with $A(\lambda)$ positive
definite. More precisely, if $b_1, \dots, b_d \in \Xs$ is a basis of $\R^d$ (which exists since
$\Span(\Xs) = \R^d$),
then the uniform design on $\{b_1,\dots,b_d\}$ has $A(\lambda) \succ 0$.
\end{lemma}

\begin{proof}
\uses{def:design_matrix}
For $v \ne 0$, $v^\top A(\lambda) v = \frac1d \sum_i \ip{b_i}{v}^2 > 0$ since some
$\ip{b_i}{v} \neq 0$ (the $b_i$ span $\R^d$).
\end{proof}

\begin{definition}[Gaussian width of a compact set for a matrix]\label{def:width_matrix}
\uses{def:arm_set}
For a nonempty compact set $K \subseteq \R^d$ (typically $K = \Xs$, or a region $K = R \subseteq \Xs$
in Chapter~\ref{chap:log_gains}), a positive definite matrix $A \in \PD$ and
$g \sim \Normal(0, I_d)$, the \emph{Gaussian width of $K$ for $A$} is
\[
  \gwidth{K}{A} := \E\Big[ \sup_{x \in K} \ip{x}{A^{-1/2} g} \Big]
  = \E\Big[ \sup_{x \in K} \ip{A^{-1/2}x}{g} \Big].
\]
Here $A^{-1/2}$ is the positive definite square root of $A^{-1}$ (Mathlib \texttt{CFC.sqrt}).
The lemmas of this section that do not involve the design set $\designset(\Xs)$
(Lemmas~\ref{lem:width_matrix_wd}, \ref{lem:width_homog}, \ref{lem:width_mono} and
\ref{lem:width_symm}) are stated and proved for an arbitrary nonempty compact index set $K$;
they are written with $K = \Xs$ for readability.
\end{definition}

\begin{lemma}[The width for a matrix is well defined]\label{lem:width_matrix_wd}
\uses{def:width_matrix, lem:sup_countable_dense}
(For an arbitrary nonempty compact index set $\Xs$.) The random variable $\sup_{x\in\Xs}\ip{A^{-1/2}x}{g}$ is measurable, integrable, and
$0 \le \gwidth{\Xs}{A} \le \sup_{x\in\Xs}\norm{A^{-1/2}x}\cdot \sqrt d$.
Moreover $\gwidth{\Xs}{A}$ only depends on the law of the Gaussian vector
$G := A^{-1/2}g \sim \Normal(0, A^{-1})$: $\gwidth{\Xs}{A} = \E[\sup_{x\in\Xs}\ip{x}{G}]$ for
any $G \sim \Normal(0,A^{-1})$.
\end{lemma}

\begin{proof}
\uses{lem:sup_countable_dense, lem:gauss_norm_moments, lem:gauss_linear_image}
Measurability: the supremum can be taken over a countable dense subset of $\Xs$
(Lemma~\ref{lem:sup_countable_dense}). Integrability and the upper bound:
$\abs{\sup_x \ip{A^{-1/2}x}{g}} \le \sup_x\norm{A^{-1/2}x}\,\norm{g}$ and
$\E\norm{g} \le \sqrt{\E\norm{g}^2} = \sqrt d$ (Lemma~\ref{lem:gauss_norm_moments}).
Nonnegativity: for any $x_0\in\Xs$, $\sup_x \ip{A^{-1/2}x}{g} \ge \ip{A^{-1/2}x_0}{g}$, whose
expectation is $0$. The last claim is that $A^{-1/2}g \sim \Normal(0,A^{-1})$
(Lemma~\ref{lem:gauss_linear_image}) and that the expectation of a measurable function only
depends on the law.
\end{proof}

\begin{definition}[Gaussian width of $\Xs$]\label{def:width}
\uses{def:width_matrix, def:design_matrix, lem:exists_posdef_design}
Let $\Xs$ be spanning. The \emph{Gaussian width term} of the action set is
\[
  \gw(\Xs) := \inf\big\{ \gwidth{\Xs}{A} : A \in \designset(\Xs) \cap \PD \big\}
  = \inf_{\lambda \in \simplex_{\Xs},\ A(\lambda)\succ 0}
    \E\Big[\max_{x\in\Xs}\ip{x}{A(\lambda)^{-1/2}g}\Big].
\]
The set over which the infimum is taken is nonempty by Lemma~\ref{lem:exists_posdef_design},
so $0 \le \gw(\Xs) < \infty$.
\end{definition}

\begin{lemma}[Homogeneity]\label{lem:width_homog}
\uses{def:width_matrix, lem:sqrt_props}
(For an arbitrary nonempty compact index set $\Xs$.) For $A \in \PD$ and $c > 0$,
$\gwidth{\Xs}{cA} = c^{-1/2}\,\gwidth{\Xs}{A}$.
\end{lemma}

\begin{proof}
\uses{def:width_matrix}
$(cA)^{-1/2} = c^{-1/2}A^{-1/2}$ and the supremum is positively homogeneous.
\end{proof}

\begin{lemma}[Monotonicity in the Loewner order]\label{lem:width_mono}
\uses{def:width_matrix, lem:gauss_comparison, lem:loewner_inv}
(For an arbitrary nonempty compact index set $\Xs$.) If $A, B \in \PD$ and $A \preceq B$, then $\gwidth{\Xs}{B} \le \gwidth{\Xs}{A}$.
\end{lemma}

\begin{proof}
\uses{lem:loewner_inv, lem:gauss_comparison, lem:width_matrix_wd}
By Lemma~\ref{lem:loewner_inv}, $B^{-1} \preceq A^{-1}$. The processes
$x \mapsto \ip{x}{G_B}$ and $x\mapsto\ip{x}{G_A}$ with $G_B\sim\Normal(0,B^{-1})$,
$G_A \sim \Normal(0,A^{-1})$ are centered linear Gaussian processes whose covariance matrices
are ordered, so Lemma~\ref{lem:gauss_comparison} (covariance domination) gives
$\E\sup_x\ip{x}{G_B} \le \E\sup_x\ip{x}{G_A}$; conclude with Lemma~\ref{lem:width_matrix_wd}.
\end{proof}

\begin{lemma}[Invariance under linear isomorphisms]\label{lem:width_iso}
\uses{def:width, lem:width_matrix_wd, lem:gauss_linear_image, lem:loewner_congruence}
Let $\Xs$ be spanning and let $M \in \R^{d\times d}$ be invertible and $\Xs' := M\Xs$. Then for every $A \in \PD$,
$\gwidth{\Xs'}{MAM^\top} = \gwidth{\Xs}{A}$, the map $A \mapsto MAM^\top$ is a bijection
from $\designset(\Xs)\cap\PD$ onto $\designset(\Xs')\cap\PD$, and $\gw(\Xs') = \gw(\Xs)$.
\end{lemma}

\begin{proof}
\uses{lem:gauss_law_of_cov, lem:width_matrix_wd}
The process $x \mapsto \ip{Mx}{G'}$ with $G' \sim \Normal(0,(MAM^\top)^{-1})$ is centered
Gaussian with covariance $x^\top M^\top (MAM^\top)^{-1} M y = x^\top A^{-1} y$, the same as
$x \mapsto \ip{x}{G}$, $G \sim \Normal(0,A^{-1})$. Two centered Gaussian vectors with the same
covariance have the same law (Lemma~\ref{lem:gauss_law_of_cov}), so the two suprema have the
same expectation by Lemma~\ref{lem:width_matrix_wd}. The design matrix of the pushed-forward
design is $A(M\lambda) = \sum \lambda_x (Mx)(Mx)^\top = MA(\lambda)M^\top$, and congruence by
an invertible matrix preserves positive definiteness.
\end{proof}

\begin{lemma}[Width of an empirical design]\label{lem:width_empirical}
\uses{def:width, def:design_matrix, lem:width_homog}
Let $\Xs$ be spanning and let $x_0,\dots,x_{T-1} \in \Xs$ with $\Sigma_T = \sum_{t<T} x_tx_t^\top$ positive definite.
Then $\gwidth{\Xs}{\Sigma_T} \ge \gw(\Xs)/\sqrt T$.
\end{lemma}

\begin{proof}
\uses{lem:width_homog, def:width}
$\Sigma_T = T A(\lambda_T)$ with $A(\lambda_T) \in \designset(\Xs)\cap\PD$, so
$\gwidth{\Xs}{\Sigma_T} = T^{-1/2}\gwidth{\Xs}{A(\lambda_T)} \ge T^{-1/2}\gw(\Xs)$.
\end{proof}

\begin{lemma}[Symmetry of the width]\label{lem:width_symm}
\uses{def:width_matrix}
(For an arbitrary nonempty compact index set $\Xs$.) For $A \in \PD$ and $G \sim \Normal(0,A^{-1})$,
$\E[\max_{x\in\Xs}\ip{x}{-G}] = \E[\max_{x\in\Xs}\ip{x}{G}] = \gwidth{\Xs}{A}$ and
$\E[Z(G)] = \E[\max_{x,y\in\Xs}\ip{x-y}{G}] = 2\,\gwidth{\Xs}{A}$.
\end{lemma}

\begin{proof}
\uses{lem:gauss_law_of_cov}
$-G \sim \Normal(0,A^{-1})$ as well (Lemma~\ref{lem:gauss_law_of_cov}), and
$\max_{x,y}\ip{x-y}{G} = \max_x\ip{x}{G} + \max_y\ip{y}{-G}$.
\end{proof}

\section{Sequential composition of algorithms}
\label{sec:model_composition}

The adaptive algorithms of the paper (Chapters~\ref{chap:log_gains},~\ref{chap:norm}
and~\ref{chap:separation}) are built by running a first algorithm for $T_1$ rounds and then
a second algorithm whose parameters depend on the first history. The following framework
lemma isolates what is needed.

\begin{lemma}[Sequential composition]\label{lem:composition}
\uses{def:bai_alg}
Let $\mathrm{alg}_1$ be an algorithm with actions in $\Xs$, $T_1 \ge 1$, and for each history
$h$ of length $T_1$ let $\mathrm{alg}_2(h)$ be an algorithm, measurably in $h$. There is an
algorithm $\mathrm{alg}$ (the \emph{composition}) whose policy coincides with that of
$\mathrm{alg}_1$ at rounds $t < T_1$ and with that of $\mathrm{alg}_2(\hist_{T_1})$ at rounds
$t \ge T_1$ (applied to the history after time $T_1$). For every environment $\nu$, under the
law of the composition against $\nu$: the history $\hist_{T_1}$ has the law of
$\mathrm{alg}_1$ against $\nu$, and conditionally on $\hist_{T_1} = h$ the shifted process
$(x_{T_1+s}, y_{T_1+s})_{s\ge 0}$ is an algorithm-environment sequence for
$(\mathrm{alg}_2(h), \nu)$. Consequently, if $(E_h)_h$ is a family of sets of trajectories, jointly measurable in
$(h, \cdot)$, with
$\Prob^{\mathrm{alg}_2(h),\nu}(E_h) \ge 1-\delta$ for every $h$ in a set of $\mathrm{alg}_1$-probability
at least $1 - \delta'$, then the composition satisfies
$\Prob\big((x_{T_1+s},y_{T_1+s})_{s \ge 0} \in E_{\hist_{T_1}}\big) \ge 1 - \delta - \delta'$.
\end{lemma}

\begin{proof}
\uses{def:bai_alg}
Define the policy of $\mathrm{alg}$ at time $t\ge T_1$ from the history $h'$ of length $t$ by
splitting $h' = (h, h'')$ with $h$ of length $T_1$ and applying the policy of
$\mathrm{alg}_2(h)$ at time $t - T_1$ to $h''$ (with the initial action distribution of
$\mathrm{alg}_2(h)$ at $t = T_1$); measurability in $h'$ holds by assumption. The trajectory
measure of the composition (\texttt{Learning.trajMeasure}, an Ionescu--Tulcea construction)
is, by construction of the kernels, the composition-product of the law of $\hist_{T_1}$ under
$\mathrm{alg}_1$ with the kernel $h \mapsto$ trajectory measure of $(\mathrm{alg}_2(h),\nu)$
shifted by $T_1$. The probability bound is then Fubini for the composition-product.
\end{proof}

\textbf{Lean remark.} This lemma is a general statement about the LML framework, to be
contributed to LML. The measurability of $h \mapsto \mathrm{alg}_2(h)$ means that the policy
kernels and initial distributions, as functions of $(h, \cdot)$, are kernels; in all uses
below $\mathrm{alg}_2(h)$ is a fixed algorithm reparametrized by a measurable function of $h$
(a finite set of candidate arms, or a scale $r_0$).

\chapter{Optimal experimental design and rounding}
\label{chap:design}

This chapter collects the deterministic (linear-algebraic and convex-analytic) tools on which
the fixed-design algorithm of Chapter~\ref{chap:upper} and the lower bounds of
Chapters~\ref{chap:lower_nonadaptive} and~\ref{chap:lower_adaptive} rest. It proves:
\begin{itemize}
  \item the facts about the Loewner order, eigenvalues and matrix square roots that are used
        throughout (Section~\ref{sec:design_loewner}); Mathlib already provides the order
        (\texttt{Matrix.instPartialOrder}, scoped in \texttt{MatrixOrder}), positive
        (semi)definiteness (\texttt{Matrix.PosDef}, \texttt{Matrix.PosSemidef}), the spectral
        theorem (\texttt{Matrix.IsHermitian.spectral\_theorem}) and the square root through the
        continuous functional calculus (\texttt{CFC.sqrt});
  \item the existence of a $G$-optimal design, i.e.\ the ``$D$-optimal implies $G$-optimal''
        half of the Kiefer--Wolfowitz theorem (Section~\ref{sec:design_kw}). This replaces the
        use of John's theorem in the paper (Appendix~B.2) and gives the lemma stated without
        proof in the paper's Appendix~B.1;
  \item Carath\'eodory's theorem for design matrices, and the continuity and coercivity of
        $A \mapsto \gwidth{\Xs}{A}$, which give a minimizer of the Gaussian width
        (Section~\ref{sec:design_caratheodory});
  \item a rounding theorem (Section~\ref{sec:rounding}): every design distribution $\lambda$ can
        be rounded, for every budget $T \ge 4d$, to a multiset of $T$ points of its support whose
        design matrix dominates $(T/4)A(\lambda)$ in the Loewner order. This replaces the rounding
        procedures cited by the paper \cite{katz2020empirical,allen2021near}; the proof is the
        one-sided barrier argument of Batson, Spielman and Srivastava \cite{batson2012twice}
        with unit weights.
\end{itemize}
Nothing in this chapter is probabilistic except the two lemmas on the width, which only use
the Gaussian facts of Chapter~\ref{chap:pre_gaussian}. Since action sets
(Definition~\ref{def:arm_set}) are not assumed to span $\R^d$, every statement of this chapter
that relies on the existence of a positive definite design
(Lemma~\ref{lem:exists_posdef_design}) or on the width $\gw(\Xs)$ (Definition~\ref{def:width})
carries an explicit hypothesis ``$\Xs$ is spanning''; the linear-algebra facts,
Carath\'eodory's theorem, the continuity of the width and the rounding theorem hold for every
action set.

\textbf{Notation.} $\Sym$ is the space of symmetric real $d \times d$ matrices, $\PSD \subset
\Sym$ the positive semidefinite ones and $\PD \subset \PSD$ the positive definite ones. For
$A, B \in \Sym$ we write $A \preceq B$ if $B - A \in \PSD$ and $A \prec B$ if $B - A \in \PD$
(the \emph{Loewner order}). For $A \in \Sym$ we write $\lambda_1(A) \le \dots \le \lambda_d(A)$
for its eigenvalues with multiplicity, $\lambda_{\min}(A) := \lambda_1(A)$,
$\lambda_{\max}(A) := \lambda_d(A)$, and $\norm{A}_{\mathrm{op}}$ for the operator norm,
$\norm{A}_F$ for the Frobenius norm. For $W \in \PD$, $\norm{x}_W^2 := x^\top W x$.

\section{The Loewner order and matrix square roots}
\label{sec:design_loewner}

\begin{lemma}[Loewner order and quadratic forms]\label{lem:loewner_quadratic}
Let $A, B \in \Sym$. Then $A \preceq B$ if and only if $x^\top A x \le x^\top B x$ for every
$x \in \R^d$, and $A \prec B$ if and only if $x^\top A x < x^\top B x$ for every $x \ne 0$.
In particular $\preceq$ is a partial order on $\Sym$ compatible with addition and with
multiplication by nonnegative scalars, and $A \preceq B$ together with $B \preceq A$ implies
$A = B$.
\end{lemma}

\begin{proof}
By definition $B - A \in \PSD$ means that $x^\top(B-A)x \ge 0$ for all $x$, which is the first
claim; the second is the definition of $\PD$. Compatibility with addition and nonnegative
scalars is immediate on quadratic forms. Antisymmetry: if $x^\top(B-A)x = 0$ for all $x$ then
the symmetric matrix $B - A$ vanishes (polarization: $2\,x^\top(B-A)y = (x+y)^\top(B-A)(x+y) -
x^\top(B-A)x - y^\top(B-A)y = 0$ for all $x, y$).
\end{proof}

\textbf{Lean remark.} Mathlib's \texttt{Matrix.PosSemidef} is defined through
\texttt{Matrix.posSemidef\_iff\_dotProduct\_mulVec} (Hermitian plus nonnegative quadratic form),
and the order instance is \texttt{Matrix.le\_iff : A ≤ B ↔ (B - A).PosSemidef} in
\texttt{Mathlib/Analysis/Matrix/Order.lean} (\texttt{open scoped MatrixOrder}). Over $\R$ the
Hermitian condition is symmetry. Antisymmetry is part of the \texttt{PartialOrder} instance.

\begin{lemma}[Congruence preserves the Loewner order]\label{lem:loewner_congruence}
\uses{lem:loewner_quadratic}
Let $A, B \in \Sym$ with $A \preceq B$ and let $M \in \R^{k \times d}$. Then
$M A M^\top \preceq M B M^\top$ in $\mathcal{S}^k$; in particular (take $A = 0$) congruence by
an arbitrary matrix $M$ maps positive semidefinite matrices to positive semidefinite matrices.
If moreover $k = d$, $M$ is invertible and
$A \prec B$, then $M A M^\top \prec M B M^\top$; in particular congruence by an invertible
matrix maps $\PD$ into $\PD$.
\end{lemma}

\begin{proof}
\uses{lem:loewner_quadratic}
For $z \in \R^k$, $z^\top M (B-A) M^\top z = (M^\top z)^\top (B - A) (M^\top z) \ge 0$ by
Lemma~\ref{lem:loewner_quadratic}; with $A = 0$ this says that $MBM^\top$ is positive
semidefinite whenever $B$ is, for every $M$. If $M$ is invertible and $z \ne 0$ then $M^\top z \ne 0$, so
the quantity is $> 0$ when $A \prec B$.
\end{proof}

\textbf{Lean remark.} \texttt{Matrix.PosSemidef.mul\_mul\_conjTranspose\_same} and
\texttt{Matrix.PosDef.mul\_mul\_conjTranspose\_same} (the latter needs injectivity of
$M^\top$, i.e.\ \texttt{Function.Injective M.vecMul}) applied to $B - A$.

\begin{lemma}[Eigenvalues and quadratic forms]\label{lem:eigen_quadratic_bounds}
\uses{lem:loewner_quadratic}
Let $A \in \Sym$. There is an orthonormal basis $u_1, \dots, u_d$ of $\R^d$ with
$A u_j = \lambda_j(A) u_j$, i.e.\ $A = U \operatorname{diag}(\lambda_1(A),\dots,\lambda_d(A))
U^\top = \sum_j \lambda_j(A) u_j u_j^\top$ with $U = [u_1 \cdots u_d]$ orthogonal. Consequently:
\begin{enumerate}
  \item[(i)] $\lambda_{\min}(A)\norm{x}^2 \le x^\top A x \le \lambda_{\max}(A)\norm{x}^2$ for all
        $x$, with equality on the left at $x = u_1$ and on the right at $x = u_d$; hence
        $\lambda_{\min}(A) = \min_{\norm{x} = 1} x^\top A x$, $A \in \PSD$ iff
        $\lambda_{\min}(A) \ge 0$, $A \in \PD$ iff $\lambda_{\min}(A) > 0$, and for $c \in \R$,
        $cI \preceq A$ iff $\lambda_{\min}(A) \ge c$.
  \item[(ii)] If $A \preceq B$ then $\lambda_{\min}(A) \le \lambda_{\min}(B)$; in particular
        $\lambda_{\min}(A + P) \ge \lambda_{\min}(A)$ for $P \in \PSD$.
  \item[(iii)] $\norm{A}_{\mathrm{op}} = \max_j \abs{\lambda_j(A)}$, which equals
        $\lambda_{\max}(A)$ when $A \in \PSD$.
  \item[(iv)] $\tr A = \sum_j \lambda_j(A)$, $\det A = \prod_j \lambda_j(A)$,
        $\tr(A^2) = \sum_j \lambda_j(A)^2 \le d\,\max_j\lambda_j(A)^2$, and for every $c \in \R$
        the matrix $A - cI$ has the same eigenvectors with eigenvalues $\lambda_j(A) - c$.
  \item[(v)] If $A \in \PD$ then $A^{-1} = U \operatorname{diag}(1/\lambda_j(A)) U^\top \in \PD$,
        $A^{-2} = U\operatorname{diag}(1/\lambda_j(A)^2)U^\top$, $\tr(A^{-1}) = \sum_j 1/\lambda_j(A)$,
        $\tr(A^{-2}) = \sum_j 1/\lambda_j(A)^2$ and $\norm{A^{-1}}_{\mathrm{op}} =
        1/\lambda_{\min}(A)$.
\end{enumerate}
\end{lemma}

\begin{proof}
\uses{lem:loewner_quadratic}
The existence of the orthonormal eigenbasis is the spectral theorem for real symmetric
matrices. Writing $x = \sum_j \ip{x}{u_j} u_j$ gives $x^\top A x = \sum_j \lambda_j \ip{x}{u_j}^2$
and $\norm{x}^2 = \sum_j \ip{x}{u_j}^2$, from which (i) follows by bounding each $\lambda_j$ by
$\lambda_{\min}$ and $\lambda_{\max}$. (ii) is (i) combined with
Lemma~\ref{lem:loewner_quadratic}: $\lambda_{\min}(B) = \min_{\norm{x}=1} x^\top B x \ge
\min_{\norm{x}=1} x^\top A x$. (iii): $\norm{Ax}^2 = \sum_j \lambda_j^2 \ip{x}{u_j}^2 \le
\max_j \lambda_j^2 \norm{x}^2$ with equality at the corresponding $u_j$. (iv): trace and
determinant are invariant under conjugation by $U$; $A^2 = U\operatorname{diag}(\lambda_j^2)U^\top$;
$A - cI = U\operatorname{diag}(\lambda_j - c)U^\top$. (v): $U\operatorname{diag}(1/\lambda_j)U^\top$
is a two-sided inverse of $A$ since $U^\top U = I$, and the remaining identities are (iii) and
(iv) applied to $A^{-1}$.
\end{proof}

\textbf{Lean remark.} \texttt{Matrix.IsHermitian.spectral\_theorem},
\texttt{Matrix.IsHermitian.eigenvalues}, \texttt{Matrix.IsHermitian.eigenvectorUnitary},
\texttt{Matrix.IsHermitian.det\_eq\_prod\_eigenvalues},
\texttt{Matrix.IsHermitian.trace\_eq\_sum\_eigenvalues} (in
\texttt{Mathlib/Analysis/Matrix/Spectrum.lean}); \texttt{Matrix.PosDef.eigenvalues\_pos},
\texttt{Matrix.PosSemidef.eigenvalues\_nonneg}. Mathlib's eigenvalues are indexed by the
index type of the matrix (not sorted), so $\lambda_{\min}$ is
\texttt{Finset.univ.inf'} of the eigenvalue function; statement (i) is the only place where
sorting matters and it is best stated as ``$\min_j \lambda_j \norm{x}^2 \le x^\top A x$''.
The identity $x^\top A x = \sum_j \lambda_j \ip{x}{u_j}^2$ is the coordinate form of the
spectral theorem in the eigenbasis (\texttt{OrthonormalBasis} of \texttt{EuclideanSpace}).

\begin{lemma}[Square roots of positive definite matrices]\label{lem:sqrt_props}
\uses{lem:eigen_quadratic_bounds}
Let $A \in \PD$ with spectral decomposition $A = U\operatorname{diag}(\lambda_j)U^\top$. The
matrix $A^{1/2} := U\operatorname{diag}(\sqrt{\lambda_j})U^\top$ is the unique $S \in \PSD$ with
$S^2 = A$ (for $A \in \PSD$ the same formula defines the unique positive semidefinite square
root). Moreover:
\begin{enumerate}
  \item[(i)] $A^{1/2} \in \PD$, it commutes with $A$, its eigenvalues are $\sqrt{\lambda_j(A)}$
        and $\lambda_{\min}(A^{1/2}) = \sqrt{\lambda_{\min}(A)}$;
  \item[(ii)] $A^{-1/2} := (A^{1/2})^{-1} = (A^{-1})^{1/2} = U\operatorname{diag}(\lambda_j^{-1/2})U^\top
        \in \PD$, $(A^{-1/2})^2 = A^{-1}$, $\norm{A^{-1/2}}_{\mathrm{op}} = \lambda_{\min}(A)^{-1/2}$;
  \item[(iii)] $A^{1/2} A^{-1/2} = A^{-1/2} A^{1/2} = I$ and $A^{-1/2} A A^{-1/2} = I$,
        $A^{1/2} A^{-1} A^{1/2} = I$;
  \item[(iv)] $x^\top A x = \norm{A^{1/2}x}^2$ and $x^\top A^{-1} x = \norm{A^{-1/2}x}^2$ for all
        $x$; in particular $\norm{x}_{A^{-1}} = \norm{A^{-1/2}x}$;
  \item[(v)] $\ip{u}{v} = \ip{A^{1/2}u}{A^{-1/2}v}$ for all $u, v \in \R^d$;
  \item[(vi)] for $c > 0$, $(cA)^{1/2} = \sqrt c\, A^{1/2}$ and $(cA)^{-1/2} = c^{-1/2}A^{-1/2}$;
        and $I^{1/2} = I$.
\end{enumerate}
\end{lemma}

\begin{proof}
\uses{lem:eigen_quadratic_bounds}
$S_0 := U\operatorname{diag}(\sqrt{\lambda_j})U^\top$ is symmetric with $S_0^2 =
U\operatorname{diag}(\lambda_j)U^\top = A$ and is positive definite by
Lemma~\ref{lem:eigen_quadratic_bounds}(i) since $\sqrt{\lambda_j} > 0$. Uniqueness: if $S \in
\PSD$ and $S^2 = A$, then $S$ commutes with $A = S^2$, hence preserves the eigenspaces of $A$;
on the eigenspace of $A$ for $\lambda$, $S$ is a positive semidefinite operator whose square is
$\lambda\,\mathrm{id}$, hence (diagonalizing $S$ there) equals $\sqrt\lambda\,\mathrm{id}$; so
$S = S_0$. (i) is read off the decomposition. (ii): $U\operatorname{diag}(\lambda_j^{-1/2})U^\top$
is the inverse of $A^{1/2}$ and the positive square root of $A^{-1} =
U\operatorname{diag}(\lambda_j^{-1})U^\top$ (Lemma~\ref{lem:eigen_quadratic_bounds}(v)), and its
largest eigenvalue is $\lambda_{\min}(A)^{-1/2}$. (iii) follows from (ii). (iv): $x^\top A x =
x^\top A^{1/2} A^{1/2} x = \norm{A^{1/2}x}^2$ since $A^{1/2}$ is symmetric, and likewise for
$A^{-1}$. (v): $\ip{A^{1/2}u}{A^{-1/2}v} = u^\top A^{1/2}A^{-1/2} v = \ip{u}{v}$ by (iii).
(vi): $cA = U\operatorname{diag}(c\lambda_j)U^\top$ and $\sqrt{c\lambda_j} = \sqrt c\sqrt{\lambda_j}$.
\end{proof}

\textbf{Lean remark.} Mathlib defines the square root through the continuous functional
calculus: \texttt{CFC.sqrt} (\texttt{Mathlib/Analysis/SpecialFunctions/ContinuousFunctionalCalculus/Rpow/Basic.lean}),
with \texttt{CFC.sqrt\_mul\_sqrt\_self}, \texttt{CFC.sqrt\_nonneg}, \texttt{CFC.sqrt\_unique},
\texttt{CFC.sqrt\_eq\_iff}, and for matrices \texttt{Matrix.PosSemidef.inv\_sqrt :
(CFC.sqrt A)⁻¹ = CFC.sqrt A⁻¹} and \texttt{Matrix.PosSemidef.det\_sqrt} in
\texttt{Mathlib/Analysis/Matrix/Order.lean}. The explicit spectral formula is available through
\texttt{Matrix.IsHermitian.cfc} and \texttt{Matrix.IsHermitian.cfc\_eq}
(\texttt{Mathlib/Analysis/Matrix/HermitianFunctionalCalculus.lean}), which state that \texttt{cfc f A} equals $U\operatorname{diag}(f(\lambda_j))U^\top$; this is
the bridge between the abstract \texttt{CFC.sqrt} and the concrete computations of this chapter.
Since Lemma~\ref{lem:sqrt_props} is the only place where the definition of $A^{1/2}$ is unfolded,
the rest of the blueprint can use whichever definition is convenient. Note that
\texttt{CFC.sqrt A = 0} for $A$ not positive semidefinite; all uses below are for $A \in \PD$.

\begin{lemma}[Inversion reverses the Loewner order]\label{lem:loewner_inv}
\uses{lem:loewner_congruence, lem:eigen_quadratic_bounds, lem:sqrt_props}
Let $A \in \PD$ and $B \in \Sym$ with $A \preceq B$. Then $B \in \PD$ and $B^{-1} \preceq A^{-1}$.
\end{lemma}

\begin{proof}
\uses{lem:loewner_congruence, lem:loewner_quadratic, lem:eigen_quadratic_bounds, lem:sqrt_props}
$B \in \PD$: for $x \ne 0$, $x^\top B x \ge x^\top A x > 0$ (Lemma~\ref{lem:loewner_quadratic}).
Let $C := A^{-1/2} B A^{-1/2} \in \Sym$. Congruence of $A \preceq B$ by $M = A^{-1/2}$
(Lemma~\ref{lem:loewner_congruence}) gives $I = A^{-1/2} A A^{-1/2} \preceq C$
(Lemma~\ref{lem:sqrt_props}(iii)). By Lemma~\ref{lem:eigen_quadratic_bounds}(i) applied to
$C - I$, every eigenvalue $\mu_j$ of $C$ satisfies $\mu_j \ge 1$; in particular $C \in \PD$.
Write $C = V\operatorname{diag}(\mu_j)V^\top$ with $V$ orthogonal. Then
$I - C^{-1} = V\operatorname{diag}(1 - 1/\mu_j)V^\top$ has nonnegative eigenvalues
$1 - 1/\mu_j \ge 0$ (Lemma~\ref{lem:eigen_quadratic_bounds}(v)), so $C^{-1} \preceq I$. Now
$C^{-1} = A^{1/2} B^{-1} A^{1/2}$ (indeed $A^{-1/2}BA^{-1/2}\cdot A^{1/2}B^{-1}A^{1/2} = I$ by
Lemma~\ref{lem:sqrt_props}(iii)), and congruence of $C^{-1} \preceq I$ by $A^{-1/2}$ gives
$B^{-1} = A^{-1/2} C^{-1} A^{-1/2} \preceq A^{-1/2} I A^{-1/2} = A^{-1}$.
\end{proof}

\textbf{Lean remark.} With the C$^*$-algebra structure on $\texttt{Matrix (Fin d) (Fin d) ℝ}$
(scoped instances of \texttt{Mathlib/Analysis/CStarAlgebra/Matrix.lean}, operator norm), this is
\texttt{CStarAlgebra.inv\_le\_inv} in
\texttt{Mathlib/Analysis/CStarAlgebra/ContinuousFunctionalCalculus/Order.lean}, stated for units
$a, b$ with $0 \le a \le b$. The proof above is the elementary route in case the instance
plumbing (the order on matrices versus the C$^*$-algebra order,
\texttt{Matrix.instStarOrderedRing}) turns out to be painful.

\begin{lemma}[Scalar domination and inverses]\label{lem:loewner_scalar}
\uses{lem:loewner_inv, lem:loewner_quadratic}
Let $c > 0$, $A \in \PD$ and $B \in \Sym$ with $cA \preceq B$. Then $B \in \PD$,
$B^{-1} \preceq c^{-1}A^{-1}$, and $\norm{x}_{B^{-1}}^2 \le c^{-1}\norm{x}_{A^{-1}}^2$ for every
$x \in \R^d$.
\end{lemma}

\begin{proof}
\uses{lem:loewner_inv, lem:loewner_quadratic}
$cA \in \PD$, so Lemma~\ref{lem:loewner_inv} gives $B \in \PD$ and
$B^{-1} \preceq (cA)^{-1} = c^{-1}A^{-1}$; the quadratic-form inequality is
Lemma~\ref{lem:loewner_quadratic}.
\end{proof}

\begin{lemma}[Trace of the inverse: AM--HM]\label{lem:trace_inv_amhm}
\uses{lem:eigen_quadratic_bounds}
For $A \in \PD$, $\tr(A^{-1}) \ge d^2/\tr(A)$.
\end{lemma}

\begin{proof}
\uses{lem:eigen_quadratic_bounds}
By Lemma~\ref{lem:eigen_quadratic_bounds}(iv)--(v), $\tr A = \sum_j \lambda_j$ and $\tr(A^{-1}) =
\sum_j 1/\lambda_j$ with all $\lambda_j > 0$. By the Cauchy--Schwarz inequality,
$d^2 = \big(\sum_j \sqrt{\lambda_j}\cdot\lambda_j^{-1/2}\big)^2 \le \sum_j\lambda_j\cdot\sum_j
\lambda_j^{-1}$.
\end{proof}

\section{Kiefer--Wolfowitz: existence of a \texorpdfstring{$G$}{G}-optimal design}
\label{sec:design_kw}

Recall (Definition~\ref{def:design_matrix}) that design distributions $\lambda \in \simplex_{\Xs}$
are finitely supported and that $\designset(\Xs) = \{A(\lambda)\} = \conv\{xx^\top : x \in \Xs\}$.
The Kiefer--Wolfowitz theorem \cite{kiefer1960equivalence} (see also
\cite[Chapter~9]{pukelsheim2006optimal} and \cite[Chapter~21]{lattimore2020bandit}) states that
a design maximizes $\det A(\lambda)$ if and only if it minimizes $\max_{x \in \Xs}
\norm{x}_{A(\lambda)^{-1}}^2$, and that the minimal value is $d$. We only need the implication
``$D$-optimal $\Rightarrow$ $G$-optimal'' together with the value $d$.

\begin{definition}[$G$-value of a matrix]\label{def:g_value}
\uses{def:arm_set, lem:eigen_quadratic_bounds}
For $A \in \PD$ the \emph{$G$-value} of $A$ on $\Xs$ is
\[
  g_{\Xs}(A) := \max_{x \in \Xs} x^\top A^{-1} x = \max_{x\in\Xs}\norm{x}_{A^{-1}}^2 .
\]
The maximum exists since $x \mapsto x^\top A^{-1} x$ is continuous and $\Xs$ is compact and
nonempty.
\end{definition}

\begin{lemma}[The $G$-value of a design matrix is at least $d$]\label{lem:g_ge_d}
\uses{def:g_value, def:design_matrix}
Let $\lambda \in \simplex_{\Xs}$ with $A(\lambda) \in \PD$. Then
\[
  \sum_{x \in \supp\lambda} \lambda_x\, x^\top A(\lambda)^{-1} x = d ,
\]
hence $g_{\Xs}(A(\lambda)) \ge \max_{x\in\supp\lambda} x^\top A(\lambda)^{-1}x \ge d$. If
$g_{\Xs}(A(\lambda)) = d$, then $x^\top A(\lambda)^{-1} x = d$ for every $x \in \supp\lambda$.
\end{lemma}

\begin{proof}
\uses{def:g_value, def:design_matrix}
Write $A := A(\lambda)$. For each $x$, $x^\top A^{-1} x = \tr(A^{-1} x x^\top)$ (cyclicity of the
trace). By linearity of the trace,
$\sum_x \lambda_x x^\top A^{-1} x = \tr\big(A^{-1}\sum_x \lambda_x xx^\top\big) = \tr(A^{-1}A) =
\tr(I_d) = d$. Since the $\lambda_x$ are nonnegative and sum to one, a weighted average equal to
$d$ forces the maximum over the support to be $\ge d$, and $g_{\Xs}(A) \ge$ that maximum because
$\supp\lambda \subset \Xs$. If $g_{\Xs}(A) = d$, all terms $x^\top A^{-1}x$, $x \in \supp\lambda$,
are $\le d$ and their $\lambda$-average is $d$, so each equals $d$ (a nonnegative combination
$\sum_x \lambda_x (d - x^\top A^{-1} x) = 0$ of nonnegative terms with $\lambda_x > 0$ vanishes
termwise).
\end{proof}

\begin{definition}[$D$-optimal design matrix]\label{def:d_optimal}
\uses{def:design_matrix}
A \emph{$D$-optimal design matrix} is a matrix $A_D \in \designset(\Xs)$ maximizing $\det$ over
$\designset(\Xs)$: $\det A_D = \max_{A \in \designset(\Xs)} \det A$.
\end{definition}

\begin{lemma}[Existence and positivity of a $D$-optimal design]\label{lem:exists_d_optimal}
\uses{def:arm_set, def:d_optimal, lem:designset_basic, lem:exists_posdef_design, lem:eigen_quadratic_bounds}
A $D$-optimal design matrix $A_D$ exists. Assume moreover that $\Xs$ is spanning. Then
$\det A_D > 0$; consequently $A_D \in \PD$.
\end{lemma}

\begin{proof}
\uses{lem:designset_basic, lem:exists_posdef_design, lem:eigen_quadratic_bounds}
$\designset(\Xs)$ is compact and nonempty (Lemma~\ref{lem:designset_basic}) and $\det$ is
continuous (a polynomial in the entries), so a maximizer exists. If $\Xs$ is spanning,
Lemma~\ref{lem:exists_posdef_design} gives $A' \in \designset(\Xs) \cap \PD$, and
$\det A' = \prod_j \lambda_j(A') > 0$ (Lemma~\ref{lem:eigen_quadratic_bounds}(i),(iv)); hence
$\det A_D \ge \det A' > 0$. Finally $A_D \in \PSD$ (Lemma~\ref{lem:designset_basic}), so its
eigenvalues are $\ge 0$ with product $\det A_D > 0$, so they are all $> 0$ and $A_D \in \PD$.
\end{proof}

\textbf{Lean remark.} \texttt{IsCompact.exists\_isMaxOn} with \texttt{Continuous.matrix\_det};
the last step is \texttt{Matrix.PosSemidef.posDef\_iff\_det\_ne\_zero}
(\texttt{Mathlib/Analysis/Matrix/Order.lean}).

\begin{lemma}[Directional derivative of $\log\det$]\label{lem:logdet_directional}
\uses{lem:eigen_quadratic_bounds, lem:sqrt_props}
Let $A \in \PD$ and $M \in \Sym$. There is $t_0 > 0$ such that $A + tM \in \PD$ for all
$\abs{t} < t_0$, and
\[
  \lim_{t \to 0,\ t \ne 0} \frac{\log\det(A + tM) - \log\det A}{t} = \tr(A^{-1}M) .
\]
\end{lemma}

\begin{proof}
\uses{lem:eigen_quadratic_bounds, lem:sqrt_props, lem:loewner_congruence}
Let $C := A^{-1/2} M A^{-1/2} \in \Sym$, with spectral decomposition
$C = V\operatorname{diag}(\mu_1,\dots,\mu_d)V^\top$ (Lemma~\ref{lem:eigen_quadratic_bounds}). Set
$t_0 := 1/(1 + \max_j\abs{\mu_j})$. For $\abs{t} < t_0$ all $1 + t\mu_j$ are $> 0$, so
$I + tC = V\operatorname{diag}(1+t\mu_j)V^\top \in \PD$, and
$A + tM = A^{1/2}(I + tC)A^{1/2}$ (Lemma~\ref{lem:sqrt_props}(iii)) is in $\PD$ by congruence
(Lemma~\ref{lem:loewner_congruence}). Taking determinants,
$\det(A + tM) = \det(A^{1/2})^2\det(I + tC) = \det A\cdot\prod_j(1 + t\mu_j)$ (using
$\det(A^{1/2})^2 = \det((A^{1/2})^2) = \det A$ and $\det(I+tC) = \det(\operatorname{diag}(1+t\mu_j))$
by invariance under conjugation by $V$). Hence, for $0 < \abs{t} < t_0$,
\[
  \frac{\log\det(A+tM) - \log\det A}{t} = \sum_{j} \frac{\log(1 + t\mu_j)}{t}
  \xrightarrow[t \to 0]{} \sum_j \mu_j ,
\]
since $\log(1 + t\mu)/t \to \mu$ as $t \to 0$ (derivative of $\log$ at $1$; the term is $0$ when
$\mu = 0$). Finally $\sum_j\mu_j = \tr C = \tr(A^{-1/2}MA^{-1/2}) = \tr(A^{-1}M)$ by cyclicity of
the trace and $A^{-1/2}A^{-1/2} = A^{-1}$ (Lemma~\ref{lem:sqrt_props}(ii)).
\end{proof}

\textbf{Lean remark.} The limit of $\log(1+t\mu)/t$ is \texttt{HasDerivAt.log} composed with
the affine map, or \texttt{Real.hasDerivAt\_log}; the sum of limits is \texttt{Filter.Tendsto}
over \texttt{Finset.sum}. The statement is deliberately about a one-dimensional limit along
$t \to 0$ (\texttt{nhdsWithin 0 \{0\}ᶜ}) rather than the Fr\'echet derivative of $\det$, which
keeps the proof inside elementary calculus. Only the one-sided limit $t \to 0^+$ is used below.

\begin{theorem}[Kiefer--Wolfowitz: existence of a $G$-optimal design]\label{thm:kiefer_wolfowitz}
\uses{def:arm_set, def:g_value, def:design_matrix, lem:g_ge_d}
Assume that $\Xs$ is spanning. There is a design distribution $\lambda_G \in \simplex_{\Xs}$
(finitely supported, as every design distribution) such that $A_G := A(\lambda_G)
\in \PD$ and $g_{\Xs}(A_G) = d$, i.e.\ $\max_{x\in\Xs}\norm{x}_{A_G^{-1}}^2 = d$; in particular
$A_G$ minimizes $g_{\Xs}$ over $\designset(\Xs)\cap\PD$ (Lemma~\ref{lem:g_ge_d}). Moreover every
$x \in \supp\lambda_G$ satisfies $x^\top A_G^{-1} x = d$. (A numerical bound on
$\abs{\supp\lambda_G}$ is available but deliberately not part of the statement; see
Remark~\ref{rem:design_kw_support}.)
\end{theorem}

\begin{proof}
\uses{lem:exists_d_optimal, def:design_matrix, lem:caratheodory_design, lem:designset_basic, lem:logdet_directional, lem:g_ge_d}
Let $A_D$ be a $D$-optimal design matrix; since $\Xs$ is spanning, $A_D \in \PD$
(Lemma~\ref{lem:exists_d_optimal}). By definition of $\designset(\Xs)$
(Definition~\ref{def:design_matrix}) there is a finitely supported $\lambda_G \in \simplex_{\Xs}$
with $A(\lambda_G) = A_D$ (Lemma~\ref{lem:caratheodory_design} shows that it can even be chosen
with $\abs{\supp\lambda_G} \le d(d+1)/2+1$, which we do not need); write $A_G := A_D$.

Fix $x \in \Xs$ and let $B := xx^\top = A(\delta_x) \in \designset(\Xs)$ and $M := B - A_G \in
\Sym$. By convexity of $\designset(\Xs)$ (Lemma~\ref{lem:designset_basic}),
$A_G + tM = (1-t)A_G + tB \in \designset(\Xs)$ for all $t \in [0,1]$, so by $D$-optimality
$\det(A_G + tM) \le \det A_G$ for $t \in [0,1]$. Let $t_0$ be given by
Lemma~\ref{lem:logdet_directional}; for $0 < t < \min(t_0, 1)$ both determinants are positive and
$\log$ is increasing, so $\big(\log\det(A_G + tM) - \log\det A_G\big)/t \le 0$. Letting $t \to
0^+$, Lemma~\ref{lem:logdet_directional} yields $\tr(A_G^{-1}M) \le 0$, i.e.
\[
  \tr(A_G^{-1}xx^\top) - \tr(A_G^{-1}A_G) = x^\top A_G^{-1}x - d \le 0 .
\]
As $x \in \Xs$ was arbitrary, $g_{\Xs}(A_G) \le d$. Lemma~\ref{lem:g_ge_d} gives
$g_{\Xs}(A_G) \ge d$, hence $g_{\Xs}(A_G) = d$, and the same lemma gives
$x^\top A_G^{-1}x = d$ on the support.
\end{proof}

\begin{remark}[Support size of the $G$-optimal design]\label{rem:design_kw_support}
The proof of Theorem~\ref{thm:kiefer_wolfowitz} may produce $\lambda_G$ through
Lemma~\ref{lem:caratheodory_design}, hence with $\abs{\supp\lambda_G} \le d(d+1)/2 + 1$; the
paper (Appendix~B.1) states the existence of a $G$-optimal design with
$\abs{\supp\lambda_*} \le d(d+1)/2$ (the sharper bound uses that the optimal matrices lie on
the boundary of the hull). This numerical bound is deliberately kept out of the statements of
Theorem~\ref{thm:kiefer_wolfowitz} and Lemma~\ref{lem:width_attained}: only the finiteness of
the support is ever used, and the formalization may prove Lemma~\ref{lem:caratheodory_design}
with the weaker bound $d^2 + 1$ (see the Lean remark there) without changing either statement. In
Chapter~\ref{chap:lower_adaptive}, $\lambda_G$ replaces the John ellipsoid of the paper:
Corollary~\ref{cor:g_optimal_normalized} provides exactly the normalized points and the isotropy
identity that the paper extracts from John's theorem.
\end{remark}

\begin{corollary}[Whitened arms have norm at most $\sqrt d$]\label{cor:g_optimal_norm_bound}
\uses{thm:kiefer_wolfowitz, lem:sqrt_props}
Let $\Xs$ be spanning and let $A_G$ be as in Theorem~\ref{thm:kiefer_wolfowitz}. Every $x \in \Xs$ satisfies
$\norm{A_G^{-1/2}x}^2 = x^\top A_G^{-1}x \le d$.
\end{corollary}

\begin{proof}
\uses{thm:kiefer_wolfowitz, lem:sqrt_props}
$\norm{A_G^{-1/2}x}^2 = x^\top A_G^{-1}x$ by Lemma~\ref{lem:sqrt_props}(iv), and
$x^\top A_G^{-1} x \le g_{\Xs}(A_G) = d$.
\end{proof}

\begin{corollary}[Normalized $G$-optimal design]\label{cor:g_optimal_normalized}
\uses{thm:kiefer_wolfowitz, cor:g_optimal_norm_bound, lem:sqrt_props}
Let $\Xs$ be spanning, let $\lambda_G$, $A_G$ be as in Theorem~\ref{thm:kiefer_wolfowitz} and set $M := d^{-1/2}A_G^{-1/2}$,
$v_x := Mx = A_G^{-1/2}x/\sqrt d$ for $x \in \supp\lambda_G$, and
$\Xs' := M\Xs = \{Mx : x \in \Xs\}$. Then:
\begin{enumerate}
  \item[(i)] $\norm{v_x} = 1$ for every $x \in \supp\lambda_G$;
  \item[(ii)] $\sum_{x \in \supp\lambda_G}\lambda_{G,x}\, v_xv_x^\top = I_d/d$;
  \item[(iii)] $\Xs' \subset \ball_d$, and $v_x \in \Xs'$ for every $x \in \supp\lambda_G$;
        $\Xs'$ is compact and spans $\R^d$.
\end{enumerate}
\end{corollary}

\begin{proof}
\uses{thm:kiefer_wolfowitz, cor:g_optimal_norm_bound, lem:sqrt_props}
(i): $\norm{v_x}^2 = x^\top A_G^{-1}x/d = 1$ by Lemma~\ref{lem:sqrt_props}(iv) and the support
property of Theorem~\ref{thm:kiefer_wolfowitz}. (ii):
$\sum_x\lambda_{G,x}v_xv_x^\top = \frac1d A_G^{-1/2}\big(\sum_x\lambda_{G,x}xx^\top\big)A_G^{-1/2}
= \frac1d A_G^{-1/2}A_GA_G^{-1/2} = I_d/d$ (Lemma~\ref{lem:sqrt_props}(iii)). (iii): for
$x \in \Xs$, $\norm{Mx}^2 = x^\top A_G^{-1}x/d \le 1$ by Corollary~\ref{cor:g_optimal_norm_bound};
$v_x = Mx$ with $x \in \supp\lambda_G \subset \Xs$; $M$ is invertible, so $M\Xs$ is compact and
spanning.
\end{proof}

\section{Carath\'eodory, continuity and attainment of the width}
\label{sec:design_caratheodory}

\begin{lemma}[Carath\'eodory for design matrices]\label{lem:caratheodory_design}
\uses{def:design_matrix}
Every $A \in \designset(\Xs)$ is of the form $A = A(\lambda)$ for a design distribution
$\lambda \in \simplex_{\Xs}$ with $\abs{\supp\lambda} \le d(d+1)/2 + 1$.
\end{lemma}

\begin{proof}
\uses{def:design_matrix}
$\designset(\Xs) = \conv S$ with $S := \{xx^\top : x \in \Xs\} \subset \Sym$, and $\Sym$ is a real
vector space of dimension $n := d(d+1)/2$. By Carath\'eodory's theorem, every point of $\conv S$
is a convex combination of an affinely independent family of points of $S$, which has at most
$n + 1$ elements: $A = \sum_{i=1}^k \alpha_i s_i$ with $k \le n+1$, $\alpha_i > 0$,
$\sum_i\alpha_i = 1$, $s_i \in S$ pairwise distinct. Choose $x_i \in \Xs$ with $s_i = x_ix_i^\top$;
the $x_i$ are pairwise distinct since the $s_i$ are. Define $\lambda_{x_i} := \alpha_i$; then
$\lambda \in \simplex_{\Xs}$, $A(\lambda) = A$ and $\abs{\supp\lambda} = k \le n + 1$.
\end{proof}

\textbf{Lean remark.} \texttt{convexHull\_eq\_union} (\texttt{Mathlib/Analysis/Convex/Caratheodory.lean})
writes the hull as a union over affinely independent finite subsets, and
\texttt{AffineIndependent.card\_le\_finrank\_succ} bounds their cardinality by
$\texttt{finrank} + 1$. The ambient space must be the submodule of symmetric matrices to get
$n = d(d+1)/2$ (its \texttt{finrank} is not in Mathlib and would have to be computed from the
basis of matrices $E_{ii}$ and $E_{ij} + E_{ji}$, $i < j$). Applying Carath\'eodory in the full
matrix space $\texttt{Matrix (Fin d) (Fin d) ℝ}$ instead gives the weaker bound $d^2 + 1$, which
is sufficient for every use in this blueprint (the support size of $\lambda_G$ and $\lambda_1$
only needs to be finite); the formalization may take this shortcut. Note also that if design
distributions are encoded as finitely supported functions $\Xs \to [0,1]$ (\texttt{Finsupp}),
the definition of $\designset(\Xs)$ as $\{A(\lambda)\}$ and the equality with the convex hull is
itself a small lemma (\texttt{Finset.centerMass\_mem\_convexHull} and
\texttt{convexHull\_eq\_union\_convexHull\_finite\_subsets}).

\begin{lemma}[Local Lipschitz continuity of the square root]\label{lem:design_sqrt_continuous}
\uses{lem:sqrt_props, lem:eigen_quadratic_bounds}
For $B, B' \in \PD$,
\[
  \norm{B^{1/2} - B'^{1/2}}_{\mathrm{op}} \le \norm{B^{1/2} - B'^{1/2}}_F \le
  \frac{d\,\norm{B - B'}_{\mathrm{op}}}{\sqrt{\lambda_{\min}(B)} + \sqrt{\lambda_{\min}(B')}} .
\]
Consequently $B \mapsto B^{1/2}$ is continuous on $\PD$, and $A \mapsto A^{-1/2}$ is continuous
on $\PD$.
\end{lemma}

\begin{proof}
\uses{lem:sqrt_props, lem:eigen_quadratic_bounds}
Let $X := B^{1/2} - B'^{1/2}$. Then $B^{1/2}X + XB'^{1/2} = B - B'$ (expand, using
$(B^{1/2})^2 = B$ and $(B'^{1/2})^2 = B'$). Let $(u_i, a_i)_i$ be an orthonormal eigenbasis of
$B^{1/2}$ and $(v_j, b_j)_j$ one of $B'^{1/2}$; by Lemma~\ref{lem:sqrt_props}(i),
$a_i \ge \alpha := \sqrt{\lambda_{\min}(B)}$ and $b_j \ge \beta := \sqrt{\lambda_{\min}(B')}$.
Since $B^{1/2}$ is symmetric, $u_i^\top B^{1/2} = a_iu_i^\top$, hence
$u_i^\top(B - B')v_j = u_i^\top B^{1/2}Xv_j + u_i^\top XB'^{1/2}v_j = (a_i + b_j)\,u_i^\top Xv_j$, so
$\abs{u_i^\top X v_j} \le \norm{B - B'}_{\mathrm{op}}/(\alpha + \beta)$. The Frobenius norm is
invariant under orthogonal changes of basis, so $\norm{X}_F^2 = \sum_{i,j}(u_i^\top Xv_j)^2 \le
d^2\norm{B-B'}_{\mathrm{op}}^2/(\alpha+\beta)^2$. The operator norm is bounded by the Frobenius
norm. Continuity at $B_0 \in \PD$: the bound with $B' = B_0$ gives
$\norm{B^{1/2} - B_0^{1/2}}_{\mathrm{op}} \le d\,\norm{B - B_0}_{\mathrm{op}}/\sqrt{\lambda_{\min}(B_0)}$
for every $B \in \PD$, with a constant independent of $B$. Finally $A \mapsto A^{-1}$ is
continuous on the (open) set of invertible matrices (entries of $A^{-1} = \operatorname{adj}(A)/\det A$
are rational functions with nonvanishing denominator), maps $\PD$ to $\PD$, and
$A^{-1/2} = (A^{-1})^{1/2}$ (Lemma~\ref{lem:sqrt_props}(ii)).
\end{proof}

\textbf{Lean remark.} This lemma is not in the outline; it is stated locally because the
continuity of \texttt{CFC.sqrt} in its argument is not readily available in Mathlib for
matrices, whereas the Sylvester-equation argument above is elementary. Continuity of the inverse
is \texttt{continuousAt\_matrix\_inv} (\texttt{Mathlib/Topology/Instances/Matrix.lean}) or
\texttt{NormedRing.inverse\_continuousAt}. All norms on the finite-dimensional space of matrices
are equivalent, so the choice of norm in the statement is immaterial for continuity.

\begin{lemma}[Continuity of the width in the matrix]\label{lem:width_matrix_continuous}
\uses{def:width_matrix, lem:width_matrix_wd, lem:design_sqrt_continuous, lem:gauss_norm_moments}
Let $R := \sup_{x\in\Xs}\norm{x}$. For $A, A' \in \PD$,
\[
  \abs{\gwidth{\Xs}{A} - \gwidth{\Xs}{A'}} \le R\sqrt d\,\norm{A^{-1/2} - A'^{-1/2}}_{\mathrm{op}} .
\]
Consequently $A \mapsto \gwidth{\Xs}{A}$ is continuous on $\PD$.
\end{lemma}

\begin{proof}
\uses{lem:width_matrix_wd, lem:design_sqrt_continuous, lem:gauss_norm_moments}
For every $g \in \R^d$, two suprema of functions on $\Xs$ differ by at most the supremum of the
difference:
$\abs{\sup_x\ip{A^{-1/2}x}{g} - \sup_x\ip{A'^{-1/2}x}{g}} \le \sup_x\abs{\ip{(A^{-1/2} - A'^{-1/2})x}{g}}
\le R\,\norm{A^{-1/2} - A'^{-1/2}}_{\mathrm{op}}\norm{g}$ (Cauchy--Schwarz). Both suprema are
integrable (Lemma~\ref{lem:width_matrix_wd}); taking expectations and using
$\E\norm{g} \le \sqrt d$ (Lemma~\ref{lem:gauss_norm_moments}) gives the inequality. Continuity
follows from the continuity of $A \mapsto A^{-1/2}$ on $\PD$
(Lemma~\ref{lem:design_sqrt_continuous}).
\end{proof}

\begin{lemma}[Coercivity of the width on design matrices]\label{lem:width_coercive}
\uses{def:arm_set, def:width_matrix, def:design_matrix, lem:sqrt_props, lem:gauss_1d_moments, lem:gauss_linear_image}
Let $\Xs$ be spanning and let $(A_n)_n$ be a sequence in $\designset(\Xs) \cap \PD$ converging to a singular matrix
$A_\infty$. Then $\gwidth{\Xs}{A_n} \to +\infty$. More precisely, for every $A \in \PD$, every
unit vector $u$ and all $x, x' \in \Xs$,
\[
  \gwidth{\Xs}{A} \ \ge\ \frac{\norm{A^{-1/2}(x - x')}}{\sqrt{2\pi}}
  \ \ge\ \frac{\abs{\ip{u}{x - x'}}}{\sqrt{2\pi}\,\sqrt{u^\top A u}} ,
\]
and if $A_n \in \designset(\Xs)$ with $u^\top A_n u \to 0$ then there are $x, x' \in \Xs$ with
$\ip{u}{x - x'} \ne 0$.
\end{lemma}

\begin{proof}
\uses{lem:sqrt_props, lem:gauss_1d_moments, lem:gauss_linear_image, lem:width_matrix_wd, lem:designset_basic}
\emph{Lower bound.} Let $a := A^{-1/2}x$, $b := A^{-1/2}x'$. For all $g$,
$\sup_{z\in\Xs}\ip{A^{-1/2}z}{g} \ge \max(\ip{a}{g},\ip{b}{g}) = \tfrac12\ip{a+b}{g} +
\tfrac12\abs{\ip{a-b}{g}}$. Taking expectations, $\E\ip{a+b}{g} = 0$ and $\ip{a-b}{g} \sim
\Normal(0,\norm{a-b}^2)$ (Lemma~\ref{lem:gauss_linear_image}) has $\E\abs{\ip{a-b}{g}} =
\norm{a-b}\sqrt{2/\pi}$ (Lemma~\ref{lem:gauss_1d_moments}), so $\gwidth{\Xs}{A} \ge
\norm{a-b}/\sqrt{2\pi} = \norm{A^{-1/2}(x-x')}/\sqrt{2\pi}$. Next, by
Lemma~\ref{lem:sqrt_props}(v) and Cauchy--Schwarz,
$\abs{\ip{u}{x-x'}} = \abs{\ip{A^{1/2}u}{A^{-1/2}(x-x')}} \le \norm{A^{1/2}u}\,\norm{A^{-1/2}(x-x')}$
with $\norm{A^{1/2}u}^2 = u^\top A u$ (Lemma~\ref{lem:sqrt_props}(iv)).

\emph{A separating pair exists.} Suppose $A_n = A(\lambda^{(n)}) \in \designset(\Xs)$ with
$u^\top A_n u \to 0$, and suppose for contradiction that $\ip{u}{x - x'} = 0$ for all
$x, x' \in \Xs$, i.e.\ $\ip{u}{x} = c$ for all $x \in \Xs$ and some constant $c$. Then
$u^\top A_nu = \sum_x\lambda^{(n)}_x\ip{u}{x}^2 = c^2$ for all $n$, so $c = 0$, so $u \perp \Xs$
and hence $u \perp \Span(\Xs) = \R^d$ (this is where the spanning hypothesis is used),
contradicting $\norm{u} = 1$.

\emph{Divergence.} $A_\infty \in \PSD$ since $\PSD$ is closed and $A_n \in \PSD$, and $A_\infty$
is singular, so there is a unit vector $u$ with $A_\infty u = 0$; then $u^\top A_nu \to
u^\top A_\infty u = 0$, and $u^\top A_n u > 0$ for each $n$. Pick $x, x'$ as in the previous
paragraph. The lower bound gives $\gwidth{\Xs}{A_n} \ge \abs{\ip{u}{x-x'}}/(\sqrt{2\pi}\sqrt{u^\top A_nu})
\to +\infty$.
\end{proof}

\begin{remark}
The restriction to $A_n \in \designset(\Xs)$ (rather than arbitrary $A_n \in \PD$, as the
outline of the chapter first suggested) is necessary: for $\Xs = \{e_1 + e_2, e_1 - e_2\}
\subset \R^2$ and $A_n = \operatorname{diag}(1/n, 1)$, one has $\gwidth{\Xs}{A_n} =
\E[\sqrt n g_1 + \abs{g_2}] = \sqrt{2/\pi}$ for all $n$, although $A_n \to \operatorname{diag}(0,1)$
is singular. The point is that a design matrix built from $\Xs$ cannot degenerate in a direction
$u$ along which $\Xs$ is not constant.
\end{remark}

\begin{lemma}[The width is attained on a positive definite design]\label{lem:width_attained}
\uses{def:arm_set, def:width, def:design_matrix, lem:designset_basic, lem:width_matrix_continuous, lem:width_coercive}
Let $\Xs$ be spanning. There is $A_1 \in \designset(\Xs)\cap\PD$ with $\gwidth{\Xs}{A_1} = \gw(\Xs)$,
i.e.\ a (finitely supported) design distribution $\lambda_1 \in \simplex_{\Xs}$ with
$A_1 = A(\lambda_1) \in \PD$ and $\gwidth{\Xs}{A(\lambda_1)} = \gw(\Xs)$.
\end{lemma}

\begin{proof}
\uses{def:width, def:design_matrix, lem:designset_basic, lem:width_matrix_continuous, lem:width_coercive, lem:exists_posdef_design}
Since $\Xs$ is spanning, the infimum defining $\gw(\Xs)$ is over a nonempty set (Lemma~\ref{lem:exists_posdef_design})
and is finite, so there is a sequence $A_n \in \designset(\Xs)\cap\PD$ with
$\gwidth{\Xs}{A_n} \to \gw(\Xs)$. By compactness of $\designset(\Xs)$
(Lemma~\ref{lem:designset_basic}) a subsequence $A_{n_k}$ converges to some
$A_\infty \in \designset(\Xs)$. If $A_\infty$ were singular, Lemma~\ref{lem:width_coercive}
would give $\gwidth{\Xs}{A_{n_k}} \to +\infty$, contradicting convergence to the finite value
$\gw(\Xs)$. Hence $A_\infty \in \PD$ ($A_\infty \in \PSD$ and nonsingular), and by continuity
(Lemma~\ref{lem:width_matrix_continuous}) $\gwidth{\Xs}{A_\infty} = \lim_k\gwidth{\Xs}{A_{n_k}} =
\gw(\Xs)$. Take $A_1 := A_\infty$; the representation $A_1 = A(\lambda_1)$ with a finitely
supported $\lambda_1 \in \simplex_{\Xs}$ is the definition of $\designset(\Xs)$
(Definition~\ref{def:design_matrix}).
\end{proof}

\textbf{Lean remark.} The argument is a standard ``minimize a continuous coercive function on a
compact set'': it can be organized as follows. The set $K_c := \{A \in \designset(\Xs) :
\gwidth{\Xs}{A} \le c\}$ for $c > \gw(\Xs)$ is nonempty, contained in $\PD$ (by coercivity and
closedness of $\designset$), and closed in $\designset(\Xs)$ (sequential closedness: a limit of
elements of $K_c$ is nonsingular by coercivity, and then the width is continuous there), hence
compact, and a minimizer of the continuous function $\gwidth{\Xs}{\cdot}$ on $K_c$
(\texttt{IsCompact.exists\_isMinOn}) is a global minimizer. Subsequence extraction is
\texttt{IsCompact.tendsto\_subseq}. As for Theorem~\ref{thm:kiefer_wolfowitz},
Lemma~\ref{lem:caratheodory_design} shows that $\lambda_1$ can be chosen with
$\abs{\supp\lambda_1} \le d(d+1)/2+1$; this numerical bound is kept out of the statement
(Remark~\ref{rem:design_kw_support}).

\section{Rounding a design distribution to a fixed design}
\label{sec:rounding}

Throughout this section, $\lambda \in \simplex_{\Xs}$ is a design distribution with
$A := A(\lambda) \in \PD$, support $\supp\lambda = \{v_1,\dots,v_m\}$ and weights
$p_i := \lambda_{v_i} > 0$, $\sum_{i=1}^m p_i = 1$. The \emph{whitened} vectors are
$u_i := A^{-1/2}v_i$, so that (Lemma~\ref{lem:sqrt_props}(iii))
\[
  \sum_{i=1}^m p_i\,u_iu_i^\top = A^{-1/2}\Big(\sum_i p_iv_iv_i^\top\Big)A^{-1/2} = A^{-1/2}AA^{-1/2} = I_d .
\]
The goal is Theorem~\ref{thm:rounding}: for every $T \ge 4d$ one can pick $T$ points of the
support (with repetitions) whose empirical design matrix is at least $(T/4)A$. This is the
lower-barrier half of the argument of \cite{batson2012twice}, with all added rank-one
matrices given weight one (no upper barrier is needed since only a lower bound on
$\Sigma_T$ is used).

\begin{definition}[Lower barrier potential]\label{def:barrier_potential}
\uses{lem:eigen_quadratic_bounds}
For $B \in \Sym$ and $l \in \R$ with $l < \lambda_{\min}(B)$, so that $B - lI \in \PD$, the
\emph{lower barrier potential} is
\[
  \Phi_l(B) := \tr\big((B - lI)^{-1}\big) = \sum_{j=1}^d \frac{1}{\lambda_j(B) - l} .
\]
\end{definition}

The second expression is Lemma~\ref{lem:eigen_quadratic_bounds}(iv)--(v). Note that
$\Phi_l(B) \ge 1/(\lambda_{\min}(B) - l)$ (keep only the smallest eigenvalue), that $\Phi_l(B)$ is
increasing in $l$, and that $\Phi_{-d}(0) = \tr((dI)^{-1}) = 1$.

\begin{lemma}[Sherman--Morrison, trace form]\label{lem:sherman_morrison_trace}
\uses{lem:loewner_quadratic, lem:eigen_quadratic_bounds}
Let $M \in \PD$ and $u \in \R^d$. Then $M + uu^\top \in \PD$, $1 + u^\top M^{-1}u > 0$,
\[
  (M + uu^\top)^{-1} = M^{-1} - \frac{M^{-1}uu^\top M^{-1}}{1 + u^\top M^{-1}u},
  \qquad
  \tr\big((M + uu^\top)^{-1}\big) = \tr(M^{-1}) - \frac{u^\top M^{-2}u}{1 + u^\top M^{-1}u} .
\]
\end{lemma}

\begin{proof}
\uses{lem:loewner_quadratic, lem:eigen_quadratic_bounds}
$M + uu^\top \in \PD$ since $x^\top(M+uu^\top)x = x^\top Mx + \ip{u}{x}^2 > 0$ for $x \ne 0$; and
$u^\top M^{-1}u \ge 0$ since $M^{-1} \in \PD$ (Lemma~\ref{lem:eigen_quadratic_bounds}(v)). Set
$s := u^\top M^{-1}u$ and $N := M^{-1} - M^{-1}uu^\top M^{-1}/(1+s)$. Then
\[
  (M + uu^\top)N = I - \frac{uu^\top M^{-1}}{1+s} + uu^\top M^{-1} - \frac{u\,(u^\top M^{-1}u)\,u^\top M^{-1}}{1+s}
  = I + uu^\top M^{-1}\Big(1 - \frac{1}{1+s} - \frac{s}{1+s}\Big) = I ,
\]
so $N = (M+uu^\top)^{-1}$ (a right inverse of a square matrix is the inverse). Taking traces,
$\tr(M^{-1}uu^\top M^{-1}) = \tr(u^\top M^{-1}M^{-1}u) = u^\top M^{-2}u$ by cyclicity.
\end{proof}

\textbf{Lean remark.} This is the special case $U = u$ (a $d \times 1$ matrix), $V = u^\top$,
$C = (1)$ of the Woodbury identity \texttt{Matrix.add\_mul\_mul\_inv\_eq\_sub} in
\texttt{Mathlib/LinearAlgebra/Matrix/NonsingularInverse.lean}; the direct verification above is
probably shorter than instantiating it (\texttt{Matrix.inv\_eq\_right\_inv} turns the one-sided
computation $(M + uu^\top)N = I$ into $(M+uu^\top)^{-1} = N$). The rank-one
matrix $uu^\top$ is \texttt{Matrix.vecMulVec u u}.

\begin{lemma}[One barrier step]\label{lem:barrier_step}
\uses{def:barrier_potential, lem:sherman_morrison_trace, lem:eigen_quadratic_bounds}
Let $l \in \R$, $l' := l + 1/2$, and $B \in \Sym$ with $\lambda_{\min}(B) > l'$. Set
$M := B - l'I \in \PD$, $\Delta := \Phi_{l'}(B) - \Phi_l(B)$, and for $u \in \R^d$
\[
  L_B(u) := \frac{u^\top M^{-2}u}{\Delta} - u^\top M^{-1}u .
\]
Then $\Delta > 0$, and if $L_B(u) \ge 1$ then
$\lambda_{\min}(B + uu^\top) > l'$ and $\Phi_{l'}(B + uu^\top) \le \Phi_l(B)$.
\end{lemma}

\begin{proof}
\uses{def:barrier_potential, lem:sherman_morrison_trace, lem:eigen_quadratic_bounds}
$\Delta > 0$: with $\lambda_j := \lambda_j(B)$, $\Phi_{l'}(B) - \Phi_l(B) = \sum_j\big(\tfrac{1}{\lambda_j - l'} -
\tfrac{1}{\lambda_j - l}\big)$ and each term is positive since $0 < \lambda_j - l' < \lambda_j - l$.
Next, $uu^\top \in \PSD$, so $\lambda_{\min}(B + uu^\top) \ge \lambda_{\min}(B) > l'$
(Lemma~\ref{lem:eigen_quadratic_bounds}(ii)), and $\Phi_{l'}(B + uu^\top)$ is defined. Since
$(B + uu^\top) - l'I = M + uu^\top$, Lemma~\ref{lem:sherman_morrison_trace} gives
\[
  \Phi_{l'}(B + uu^\top) = \tr(M^{-1}) - \frac{u^\top M^{-2}u}{1 + u^\top M^{-1}u}
  = \Phi_{l'}(B) - \frac{u^\top M^{-2}u}{1 + u^\top M^{-1}u} .
\]
The hypothesis $L_B(u) \ge 1$ reads $u^\top M^{-2}u/\Delta \ge 1 + u^\top M^{-1}u$; multiplying by
$\Delta/(1 + u^\top M^{-1}u) > 0$ gives $u^\top M^{-2}u/(1 + u^\top M^{-1}u) \ge \Delta$, hence
$\Phi_{l'}(B + uu^\top) \le \Phi_{l'}(B) - \Delta = \Phi_l(B)$.
\end{proof}

\begin{lemma}[Averaging the barrier criterion]\label{lem:barrier_average}
\uses{def:barrier_potential, lem:barrier_step, lem:eigen_quadratic_bounds}
Let $u_1,\dots,u_m \in \R^d$ and $p_1,\dots,p_m \ge 0$ with $\sum_ip_i = 1$ and
$\sum_i p_iu_iu_i^\top = I_d$. Let $l \in \R$, $l' := l + 1/2$, and $B \in \Sym$ with
$\lambda_{\min}(B) > l$ and $\Phi_l(B) \le 1$. Then $\lambda_{\min}(B) \ge l + 1 > l'$, and with
$L_B$ as in Lemma~\ref{lem:barrier_step},
\[
  \sum_{i=1}^m p_i\,L_B(u_i) \ \ge\ 2 - \Phi_l(B) \ \ge\ 1 ;
\]
in particular there is an index $i$ with $L_B(u_i) \ge 1$.
\end{lemma}

\begin{proof}
\uses{def:barrier_potential, lem:barrier_step, lem:eigen_quadratic_bounds}
Since $1/(\lambda_{\min}(B) - l) \le \Phi_l(B) \le 1$ and $\lambda_{\min}(B) - l > 0$, we get
$\lambda_{\min}(B) - l \ge 1$, so $\lambda_{\min}(B) \ge l + 1 > l + 1/2 = l'$ and
Lemma~\ref{lem:barrier_step} applies: $M = B - l'I \in \PD$ and $\Delta > 0$.

Let $\lambda_j := \lambda_j(B)$, $a_j := \lambda_j - l \ge 1$ and $b_j := a_j - 1/2 = \lambda_j - l'
\ge 1/2$, so that $M$ has eigenvalues $b_j$ (Lemma~\ref{lem:eigen_quadratic_bounds}(iv)). Define
\[
  P := \tr(M^{-2}) = \sum_j\frac{1}{b_j^2},\qquad
  Q := \sum_j\frac{1}{a_jb_j} > 0,\qquad
  \Phi_l(B) = \sum_j\frac1{a_j},\qquad \Phi_{l'}(B) = \tr(M^{-1}) = \sum_j\frac1{b_j} .
\]
By linearity of the trace and $\sum_ip_iu_iu_i^\top = I$: $\sum_ip_i\,u_i^\top M^{-2}u_i =
\tr(M^{-2}\sum_ip_iu_iu_i^\top) = P$ and $\sum_ip_i\,u_i^\top M^{-1}u_i = \tr(M^{-1}) =
\Phi_{l'}(B)$. Hence $\sum_ip_iL_B(u_i) = P/\Delta - \Phi_{l'}(B)$. Now
\[
  \Delta = \sum_j\Big(\frac1{b_j} - \frac1{a_j}\Big) = \sum_j\frac{a_j - b_j}{a_jb_j} = \frac{Q}{2},
  \qquad
  \Phi_{l'}(B) = \Phi_l(B) + \Delta = \Phi_l(B) + \frac Q2 ,
\]
so $\sum_ip_iL_B(u_i) = 2P/Q - \Phi_l(B) - Q/2$, and the claim $\sum_ip_iL_B(u_i) \ge 2 - \Phi_l(B)$
is equivalent to $2P/Q - Q/2 \ge 2$, i.e.\ (multiplying by $Q/2 > 0$) to $P - Q \ge Q^2/4$.
We have $P - Q = \sum_j\big(\tfrac1{b_j^2} - \tfrac1{a_jb_j}\big) = \sum_j\tfrac{a_j-b_j}{a_jb_j^2} =
\tfrac12\sum_j\tfrac1{a_jb_j^2} \ge 0$, and by Cauchy--Schwarz,
\[
  Q = \sum_j\frac{1}{\sqrt{a_j}\,b_j}\cdot\frac{1}{\sqrt{a_j}}
  \le \Big(\sum_j\frac1{a_jb_j^2}\Big)^{1/2}\Big(\sum_j\frac1{a_j}\Big)^{1/2}
  = \sqrt{2(P-Q)}\,\sqrt{\Phi_l(B)} \le \sqrt{2(P-Q)} ,
\]
using $\Phi_l(B) \le 1$. Thus $Q^2 \le 2(P - Q)$ and $Q^2/4 \le (P-Q)/2 \le P - Q$. Finally
$2 - \Phi_l(B) \ge 1$ because $\Phi_l(B) \le 1$, and a $p$-weighted average of the numbers
$L_B(u_i)$ being $\ge 1$ forces $\max_iL_B(u_i) \ge 1$.
\end{proof}

\begin{theorem}[Rounding to a fixed design of any size $T \ge 4d$]\label{thm:rounding}
\uses{def:design_matrix, lem:sqrt_props, def:barrier_potential}
Let $\lambda \in \simplex_{\Xs}$ with $A := A(\lambda) \in \PD$ and let $T \ge 4d$ be an integer.
There exist $x_0, \dots, x_{T-1} \in \supp\lambda$ (with repetitions allowed) such that
\[
  \Sigma_T := \sum_{t<T}x_tx_t^\top \ \succeq\ \frac T4\,A(\lambda) .
\]
\end{theorem}

\begin{proof}
\uses{lem:barrier_average, lem:barrier_step, lem:loewner_congruence, lem:eigen_quadratic_bounds, lem:sqrt_props, def:barrier_potential}
Use the notation of the beginning of the section: $v_1,\dots,v_m$, $p_i$, $u_i = A^{-1/2}v_i$
with $\sum_ip_iu_iu_i^\top = I$. We construct by induction on $t \le T$ indices
$i_0,\dots,i_{t-1} \in \{1,\dots,m\}$ such that $B_t := \sum_{s<t}u_{i_s}u_{i_s}^\top$ and
$l_t := -d + t/2$ satisfy the invariant
\[
  \lambda_{\min}(B_t) > l_t \qquad\text{and}\qquad \Phi_{l_t}(B_t) \le 1 .
\]
\emph{Base case} $t = 0$: $B_0 = 0$ has $\lambda_{\min}(B_0) = 0 > -d = l_0$ and
$\Phi_{-d}(0) = \tr((dI)^{-1}) = d\cdot\frac1d = 1$.
\emph{Step:} given the invariant at $t$, Lemma~\ref{lem:barrier_average} (with $l = l_t$,
$l' = l_t + 1/2 = l_{t+1}$, $B = B_t$) provides an index $i_t$ with $L_{B_t}(u_{i_t}) \ge 1$, and
Lemma~\ref{lem:barrier_step} then gives, for $B_{t+1} = B_t + u_{i_t}u_{i_t}^\top$,
$\lambda_{\min}(B_{t+1}) > l_{t+1}$ and $\Phi_{l_{t+1}}(B_{t+1}) \le \Phi_{l_t}(B_t) \le 1$.

After $T$ steps, $\lambda_{\min}(B_T) > l_T = -d + T/2 \ge T/4$, where the last inequality is
$T/4 \ge d$, i.e.\ $T \ge 4d$. By Lemma~\ref{lem:eigen_quadratic_bounds}(i), $B_T \succeq (T/4)I$.
Set $x_t := v_{i_t} \in \supp\lambda$. Since $v_i = A^{1/2}u_i$ (Lemma~\ref{lem:sqrt_props}(iii)),
\[
  \sum_{t<T}x_tx_t^\top = \sum_{t<T}A^{1/2}u_{i_t}u_{i_t}^\top A^{1/2} = A^{1/2}B_TA^{1/2}
  \ \succeq\ A^{1/2}\Big(\frac T4 I\Big)A^{1/2} = \frac T4 A
\]
by congruence with $M = A^{1/2}$ (Lemma~\ref{lem:loewner_congruence}) and $A^{1/2}A^{1/2} = A$.
\end{proof}

\textbf{Lean remark.} The proof is an existence statement proved by induction on $t$:
``for every $t \le T$ there is a list of $t$ indices whose $B_t$ satisfies the invariant''.
The greedy choice at each step is \texttt{Classical.choose} applied to
Lemma~\ref{lem:barrier_average}; no explicit sequence needs to be defined. The design matrix
$\Sigma_T$ is $\sum_{t<T}$ \texttt{vecMulVec (x t) (x t)} for \texttt{x : Fin T → Xs}. Note that
the theorem produces a design of length exactly $T$ with points in the support of $\lambda$,
which is what Definition~\ref{def:fixed_design} consumes.

\begin{corollary}[Fixed design with controlled inverse and width]\label{cor:fixed_design_from_distribution}
\uses{thm:rounding, def:width_matrix, lem:loewner_scalar, lem:width_mono, lem:width_homog}
In the setting of Theorem~\ref{thm:rounding}, let $x_0,\dots,x_{T-1}$ be the design it provides
and $\Sigma_T := \sum_{t<T}x_tx_t^\top$. Then $\Sigma_T \in \PD$,
\[
  \norm{x}_{\Sigma_T^{-1}}^2 \le \frac 4T\,\norm{x}_{A(\lambda)^{-1}}^2 \quad\text{for all } x \in \R^d,
  \qquad\text{and}\qquad
  \gwidth{\Xs}{\Sigma_T} \le \frac{2}{\sqrt T}\,\gwidth{\Xs}{A(\lambda)} .
\]
\end{corollary}

\begin{proof}
\uses{thm:rounding, lem:loewner_scalar, lem:width_mono, lem:width_homog}
Theorem~\ref{thm:rounding} gives $(T/4)A \preceq \Sigma_T$ with $A = A(\lambda) \in \PD$.
Lemma~\ref{lem:loewner_scalar} with $c = T/4$ yields $\Sigma_T \in \PD$ and
$\norm{x}_{\Sigma_T^{-1}}^2 \le (4/T)\norm{x}_{A^{-1}}^2$. For the width,
Lemma~\ref{lem:width_mono} applied to $(T/4)A \preceq \Sigma_T$ (both in $\PD$) gives
$\gwidth{\Xs}{\Sigma_T} \le \gwidth{\Xs}{(T/4)A} = (T/4)^{-1/2}\gwidth{\Xs}{A} =
(2/\sqrt T)\gwidth{\Xs}{A}$ by Lemma~\ref{lem:width_homog}.
\end{proof}

\begin{lemma}[Naive rounding by ceilings]\label{lem:naive_rounding}
\uses{def:design_matrix, lem:loewner_quadratic}
Let $\lambda \in \simplex_{\Xs}$ with support $\{v_1,\dots,v_m\}$ and weights $p_i > 0$, and let
$T \ge 2m$ be an integer. Set $\kappa_i := \lceil (T-m)p_i \rceil \in \N$. Then
$\sum_i\kappa_i \le T$, and any sequence $x_0,\dots,x_{T-1}$ consisting of $\kappa_i$ copies of
$v_i$ for each $i$, padded with $T - \sum_i\kappa_i$ further copies of $v_1$, satisfies
\[
  \sum_{t<T}x_tx_t^\top \ \succeq\ \sum_i\kappa_iv_iv_i^\top \ \succeq\ (T-m)\,A(\lambda) \ \succeq\ \frac T2\,A(\lambda) .
\]
\end{lemma}

\begin{proof}
\uses{def:design_matrix, lem:loewner_quadratic}
Since $\lceil y\rceil < y + 1$, $\sum_i\kappa_i < (T-m)\sum_ip_i + m = T$, so $\sum_i\kappa_i \le T$
(integers) and the padding is well defined. Each $v_iv_i^\top \in \PSD$, so adding the padding
copies only increases the sum in the Loewner order (Lemma~\ref{lem:loewner_quadratic}). Next,
$\sum_i\kappa_iv_iv_i^\top - (T-m)A(\lambda) = \sum_i(\kappa_i - (T-m)p_i)v_iv_i^\top$ is a
combination of positive semidefinite matrices with nonnegative coefficients
$\kappa_i - (T-m)p_i \ge 0$, hence positive semidefinite. Finally $T - m \ge T/2$ because
$T \ge 2m$, and $A(\lambda) \in \PSD$, so $(T-m)A(\lambda) \succeq (T/2)A(\lambda)$.
\end{proof}

\begin{remark}
Lemma~\ref{lem:naive_rounding} is a one-line substitute for Theorem~\ref{thm:rounding} with a
better constant ($T/2$ instead of $T/4$), but it requires $T \ge 2m$ where $m = \abs{\supp\lambda}$
is only known to be finite (at most $d(d+1)/2+1$ when $\lambda$ is produced by
Lemma~\ref{lem:caratheodory_design}); it therefore only gives
fixed designs with budget $T = \Omega(d^2)$, which is not enough for the $d\log(1/\delta)/\eps^2$
term of Theorem~\ref{thm:upper} when $\log(1/\delta) \ll d$. Theorem~\ref{thm:rounding} removes
this restriction. Either lemma may be used in a first version of the formalization of
Chapter~\ref{chap:upper}, at the price of a weaker theorem in the naive case.
\end{remark}

\chapter{The non-adaptive fixed-design algorithm}
\label{chap:upper}

This chapter proves Theorem~1 of the paper (Section~2.1 and Appendix~B.1 there): a
non-adaptive fixed-design algorithm with budget
$T = \lceil 600(\gw(\Xs)^2 + d\log(2/\delta))/\eps^2\rceil$ is $(\eps,\delta)$-PAC on every
spanning action set $\Xs$ (Theorem~\ref{thm:upper}); the spanning hypothesis is what makes the
width $\gw(\Xs)$ and the $G$-optimal design exist. The design is obtained by rounding
(Theorem~\ref{thm:rounding}) the mixture of a width-minimizing design
(Lemma~\ref{lem:width_attained}) and a $G$-optimal design (Theorem~\ref{thm:kiefer_wolfowitz});
the estimator is least squares; the recommendation is an (approximate) maximizer of the
estimated value. The probabilistic input is the law of the least-squares error
(Chapter~\ref{chap:pre_gaussian}) and the Borell--TIS inequality for linear Gaussian processes in
the form of Corollary~\ref{cor:sup_bound_symmetric} (Chapter~\ref{chap:pre_concentration}).
Compared with the paper: the constant $360$ becomes $600$ because of the concentration
constant $\cG = \pi^2/4$ of Theorem~\ref{thm:borell_tis} and because the recommendation is only
an $\eps/4$-approximate maximizer (a measurable exact argmax selector on a compact set is not
available without extra work; Lemma~\ref{lem:approx_argmax_selector}); the rounding step uses
Theorem~\ref{thm:rounding} instead of the cited rounding procedure. The chapter ends with a
direct treatment of the unit ball (Lemma~\ref{lem:ball_identification}), which needs none of
the width machinery and is used in Chapter~\ref{chap:separation}.

\section{The mixed design and its rounding}
\label{sec:upper_design}

\begin{definition}[Mixed design]\label{def:mixed_design}
\uses{def:arm_set, def:design_matrix, lem:width_attained, thm:kiefer_wolfowitz}
Let $\Xs$ be spanning. Let $\lambda_1 \in \simplex_{\Xs}$ with $A_1 := A(\lambda_1) \in \PD$ and
$\gwidth{\Xs}{A_1} = \gw(\Xs)$ (Lemma~\ref{lem:width_attained}), and let $\lambda_G \in \simplex_{\Xs}$
with $A_G := A(\lambda_G) \in \PD$ and $g_{\Xs}(A_G) = d$ (Theorem~\ref{thm:kiefer_wolfowitz}). The
\emph{mixed design distribution} is
\[
  \lambda_0 := \tfrac12\lambda_1 + \tfrac12\lambda_G \in \simplex_{\Xs},
  \qquad
  A_0 := A(\lambda_0) = \tfrac12A_1 + \tfrac12A_G \in \PD ,
\]
where $\lambda_0$ is supported on $\supp\lambda_1\cup\supp\lambda_G$ (finite) with weights
$\lambda_{0,x} = \frac12\lambda_{1,x} + \frac12\lambda_{G,x}$ (a weight being $0$ outside the
corresponding support).
\end{definition}

\begin{lemma}[Bounds for the mixed design]\label{lem:mixed_design_bounds}
\uses{def:mixed_design, def:width_matrix, def:g_value}
In the setting of Definition~\ref{def:mixed_design} ($\Xs$ spanning), for every $x \in \Xs$,
$\norm{x}_{A_0^{-1}}^2 \le 2d$, and $\gwidth{\Xs}{A_0} \le \sqrt2\,\gw(\Xs)$.
\end{lemma}

\begin{proof}
\uses{lem:loewner_scalar, cor:g_optimal_norm_bound, lem:width_mono, lem:width_homog, lem:loewner_quadratic}
Since $\frac12A_1 \in \PSD$, $\frac12A_G \preceq A_0$ (Lemma~\ref{lem:loewner_quadratic}), and
Lemma~\ref{lem:loewner_scalar} with $c = 1/2$ gives $\norm{x}_{A_0^{-1}}^2 \le 2\norm{x}_{A_G^{-1}}^2
\le 2d$ for $x \in \Xs$, by Corollary~\ref{cor:g_optimal_norm_bound}. Similarly
$\frac12A_1 \preceq A_0$ with $\frac12A_1, A_0 \in \PD$, so by monotonicity of the width
(Lemma~\ref{lem:width_mono}) and homogeneity (Lemma~\ref{lem:width_homog} with $c = 1/2$,
$c^{-1/2} = \sqrt2$),
$\gwidth{\Xs}{A_0} \le \gwidth{\Xs}{\tfrac12A_1} = \sqrt2\,\gwidth{\Xs}{A_1} = \sqrt2\,\gw(\Xs)$.
\end{proof}

\section{Least squares and the difference process}
\label{sec:upper_ls}

\begin{definition}[Least-squares estimator]\label{def:least_squares}
\uses{def:fixed_design}
For a fixed design $x_0,\dots,x_{T-1}$ with $\Sigma_T = X^\top X \in \PD$ and observations
$y \in \R^T$, the \emph{least-squares estimator} is
\[
  \hat\theta := \Sigma_T^{-1}X^\top y = \Sigma_T^{-1}\sum_{t<T}y_tx_t \in \R^d ,
\]
a linear (hence measurable) function of $y$.
\end{definition}

\begin{lemma}[Law of the least-squares error]\label{lem:least_squares_law}
\uses{def:least_squares, lem:fixed_design_law}
Under the fixed design run against $\nu_\theta$,
$\hat\theta - \theta = \Sigma_T^{-1}X^\top\eta \sim \Normal(0,\Sigma_T^{-1})$, where
$\eta \sim \Normal(0,I_T)$ is the noise vector. In particular
$\hat\theta - \theta \eqd \Sigma_T^{-1/2}g$ with $g \sim \Normal(0,I_d)$.
\end{lemma}

\begin{proof}
\uses{lem:fixed_design_law, lem:gauss_linear_image, lem:gauss_law_of_cov, lem:sqrt_props}
By Lemma~\ref{lem:fixed_design_law}, $y = X\theta + \eta$ with $\eta \sim \Normal(0,I_T)$. Then
$\hat\theta = \Sigma_T^{-1}X^\top X\theta + \Sigma_T^{-1}X^\top\eta = \theta + \Sigma_T^{-1}X^\top\eta$.
Lemma~\ref{lem:gauss_linear_image} with the matrix $M := \Sigma_T^{-1}X^\top \in \R^{d\times T}$
gives $M\eta \sim \Normal(0, MM^\top)$ with $MM^\top = \Sigma_T^{-1}X^\top X\Sigma_T^{-1} =
\Sigma_T^{-1}$ (using that $\Sigma_T^{-1}$ is symmetric). The last claim is
Lemma~\ref{lem:gauss_law_of_cov} together with $(\Sigma_T^{-1/2})^2 = \Sigma_T^{-1}$
(Lemma~\ref{lem:sqrt_props}(ii)).
\end{proof}

\begin{lemma}[Measurable approximate argmax selector]\label{lem:approx_argmax_selector}
\uses{def:arm_set, lem:sup_countable_dense}
Let $\Xs \subset \R^d$ be compact and nonempty and let $\gamma > 0$. There is a Borel measurable
map $s_\gamma : \R^d \to \Xs$ such that for every $v \in \R^d$,
\[
  \ip{s_\gamma(v)}{v} \ \ge\ \max_{x\in\Xs}\ip{x}{v} - \gamma .
\]
If $\Xs$ is finite, there is a measurable $s_0 : \R^d \to \Xs$ with $\ip{s_0(v)}{v} =
\max_{x\in\Xs}\ip{x}{v}$ for all $v$ (an exact argmax selector).
\end{lemma}

\begin{proof}
\uses{lem:sup_countable_dense}
Let $M(v) := \max_{x\in\Xs}\ip{x}{v}$; it is continuous in $v$ (Lipschitz with constant
$R := \sup_{x\in\Xs}\norm{x}$, Lemma~\ref{lem:sup_countable_dense}). Let $(q_n)_{n\in\N}$ be a
sequence that is dense in $\Xs$ ($\Xs$ is a separable metric space). For each $n$ the set
$U_n := \{v \in \R^d : \ip{q_n}{v} > M(v) - \gamma\}$ is open, as a strict inequality between
continuous functions. The $U_n$ cover $\R^d$: given $v$, let $x^\star \in \Xs$ attain $M(v)$ and
pick $n$ with $\norm{q_n - x^\star} < \gamma/(1 + \norm{v})$; then
$\ip{q_n}{v} \ge \ip{x^\star}{v} - \norm{q_n - x^\star}\norm{v} > M(v) - \gamma$. Define
$n(v) := \min\{n : v \in U_n\}$ and $s_\gamma(v) := q_{n(v)}$. Then $s_\gamma(v) \in \Xs$ and
$\ip{s_\gamma(v)}{v} > M(v) - \gamma$ by construction. Measurability: $s_\gamma$ takes countably
many values and $\{v : s_\gamma(v) = q_n\} \supseteq \{v : n(v) = n\} = U_n \setminus \bigcup_{k<n}U_k$;
more precisely $s_\gamma^{-1}(\{q\}) = \bigcup_{n : q_n = q}\{n(\cdot) = n\}$ is a countable union of
Borel sets, so every preimage of a subset of the countable range is Borel.

For finite $\Xs = \{q_1,\dots,q_m\}$ the same construction with $\gamma = 0$ works: the sets
$F_n := \{v : \ip{q_n}{v} = M(v)\}$ are closed (hence Borel), cover $\R^d$ since the maximum is
attained, and $s_0(v) := q_{\min\{n : v \in F_n\}}$ is measurable for the same reason.
\end{proof}

\textbf{Lean remark.} $n(v)$ is \texttt{Nat.find} applied to the covering property; the
measurability of a map with countable range whose fibers are measurable is
\texttt{measurable\_to\_countable'} (\texttt{Mathlib/MeasureTheory/MeasurableSpace/Constructions.lean}),
and \texttt{TopologicalSpace.exists\_countable\_dense} gives the dense sequence (turn the
countable dense set into a sequence with \texttt{Set.Countable.exists\_eq\_range} after
handling emptiness). The codomain is the subtype $\Xs$ with the Borel $\sigma$-algebra of the
subspace topology, as in Definition~\ref{def:arm_set}. The recommendation kernel of
Definition~\ref{def:fixed_design_algorithm} is then \texttt{Kernel.deterministic} of the
measurable map $y \mapsto s_\gamma(\hat\theta(y))$. The paper uses an exact argmax; we use
$\gamma = \eps/4$ and absorb the loss in the constant (Theorem~\ref{thm:upper}).

\begin{lemma}[Regret is controlled by the difference process]\label{lem:regret_le_difference_process}
\uses{def:regret}
Let $\theta, \hat\theta \in \R^d$, $\gamma \ge 0$, and $\rec \in \Xs$ with
$\ip{\rec}{\hat\theta} \ge \max_{x\in\Xs}\ip{x}{\hat\theta} - \gamma$. Then
\[
  r(\rec,\theta) \le D + \gamma,
  \qquad
  D := \max_{x,x'\in\Xs}\ip{x - x'}{\hat\theta - \theta} .
\]
\end{lemma}

\begin{proof}
\uses{def:regret}
Let $x^\star \in \argmax_{x\in\Xs}\ip{x}{\theta}$. Then
\[
  r(\rec,\theta) = \ip{x^\star}{\theta} - \ip{\rec}{\theta}
  = \ip{x^\star - \rec}{\theta - \hat\theta} + \big(\ip{x^\star}{\hat\theta} - \ip{\rec}{\hat\theta}\big) .
\]
The second bracket is at most $\max_x\ip{x}{\hat\theta} - \ip{\rec}{\hat\theta} \le \gamma$, and the
first term equals $\ip{\rec - x^\star}{\hat\theta - \theta} \le D$ (take $x = \rec$, $x' = x^\star$).
\end{proof}

\begin{lemma}[Concentration of the difference process]\label{lem:difference_process_concentration}
\uses{def:width_matrix, lem:width_symm, cor:sup_bound_symmetric}
Let $\Sigma \in \PD$, $\Delta \sim \Normal(0,\Sigma^{-1})$, $D := \max_{x,x'\in\Xs}\ip{x-x'}{\Delta}$
and $\sigma^2 := \max_{x\in\Xs}\norm{x}_{\Sigma^{-1}}^2$. Then $\E D = 2\,\gwidth{\Xs}{\Sigma}$, and
for every $\delta \in (0,1)$,
\[
  \Prob\Big( D \le 2\,\gwidth{\Xs}{\Sigma} + 2\sigma\sqrt{2\cG\log(1/\delta)} \Big) \ge 1 - \delta .
\]
\end{lemma}

\begin{proof}
\uses{lem:width_symm, lem:width_matrix_wd, cor:sup_bound_symmetric, cor:sup_bound_whp, lem:sqrt_props}
The expectation: $D = \max_x\ip{x}{\Delta} + \max_{x'}\ip{x'}{-\Delta}$ and both maxima have
expectation $\gwidth{\Xs}{\Sigma}$ by Lemma~\ref{lem:width_symm} (the width only depends on the
law $\Normal(0,\Sigma^{-1})$ of $\Delta$, Lemma~\ref{lem:width_matrix_wd}). The deviation bound is
Corollary~\ref{cor:sup_bound_symmetric} applied to the centered Gaussian vector $G = \Delta$ with
covariance $S = \Sigma^{-1}$ and index set $\Xs$: there $\sigma^2 = \sup_{x\in\Xs}x^\top Sx =
\max_x\norm{x}_{\Sigma^{-1}}^2$, and the statement is exactly
$\Prob(D \le \E D + 2\sigma\sqrt{2\cG\log(1/\delta)}) \ge 1-\delta$. (Equivalently, apply
Corollary~\ref{cor:sup_bound_whp} to the compact index set $\Xs - \Xs = \{x - x' : x,x'\in\Xs\}$,
whose variance proxy is $\max_{x,x'}\norm{x-x'}_{\Sigma^{-1}}^2 \le (2\sigma)^2$ by the triangle
inequality for the norm $\norm{\cdot}_{\Sigma^{-1}} = \norm{\Sigma^{-1/2}\cdot}$,
Lemma~\ref{lem:sqrt_props}(iv).)
\end{proof}

\section{The algorithm and its guarantee}
\label{sec:upper_theorem}

\begin{definition}[The fixed-design identification algorithm]\label{def:fixed_design_algorithm}
\uses{def:arm_set, def:fixed_design, def:mixed_design, thm:rounding, def:least_squares, lem:approx_argmax_selector, def:width}
Let $\Xs$ be spanning, $\eps \in (0,1]$ and $\delta \in (0,1)$, and set
\[
  T := \Big\lceil \frac{600\,\big(\gw(\Xs)^2 + d\log(2/\delta)\big)}{\eps^2} \Big\rceil .
\]
Then $T \ge 4d$, as required by Theorem~\ref{thm:rounding}: indeed $\log(2/\delta) > \log 2 > 0.69$
and $\eps \le 1$, so $T \ge 600\cdot0.69\,d/\eps^2 \ge 414\,d \ge 4d$.
The \emph{fixed-design identification algorithm} $\mathsf{Alg}_{\mathrm{fd}}(\eps,\delta)$
is the fixed-design algorithm (Definition~\ref{def:fixed_design}) with
\begin{enumerate}
  \item[(i)] design $x_0,\dots,x_{T-1} \in \supp\lambda_0$ given by Theorem~\ref{thm:rounding}
        applied to the mixed design distribution $\lambda_0$ of Definition~\ref{def:mixed_design}
        with budget $T$; its design matrix $\Sigma_T$ is positive definite
        (Corollary~\ref{cor:fixed_design_from_distribution});
  \item[(ii)] estimator $\hat\theta = \Sigma_T^{-1}X^\top y$ (Definition~\ref{def:least_squares});
  \item[(iii)] recommendation $\rec := s_{\eps/4}(\hat\theta)$, with $s_{\eps/4}$ the measurable
        $\eps/4$-approximate argmax selector of Lemma~\ref{lem:approx_argmax_selector}; the
        recommendation kernel is the Dirac kernel at the measurable map $y \mapsto
        s_{\eps/4}(\Sigma_T^{-1}X^\top y)$.
\end{enumerate}
\end{definition}

\begin{remark}[The range of $\eps$ and the constant]\label{rem:upper_eps_range}
The definition is restricted to $\eps \in (0,1]$ (the regime of the paper), which matches
Theorem~\ref{thm:upper} and makes the requirement $T \ge 4d$ of Theorem~\ref{thm:rounding}
automatic. For $\eps > 1$ one can define the algorithm in the same way with
$T := \max\big(\lceil 600(\gw(\Xs)^2 + d\log(2/\delta))/\eps^2\rceil, 4d\big)$; the proof of
Theorem~\ref{thm:upper} goes through unchanged, since it only uses $T \ge 4d$ and
$T \ge 600(\gw(\Xs)^2 + d\log(2/\delta))/\eps^2$. If $\Xs$ is finite, the exact selector $s_0$
can be used in (iii); then $D \le \eps$ suffices in the proof below, and since that proof gives
$D^2 \le 320\,(\gw(\Xs)^2 + d\log(2/\delta))/T$ on the good event, the constant $600$ can be
reduced to $320$. The paper's constant is $360$.
\end{remark}

\begin{theorem}[Fixed-design upper bound (Theorem~1 of the paper)]\label{thm:upper}
\uses{def:arm_set, def:fixed_design_algorithm, def:pac, def:width}
Let $\Xs$ be spanning, $\eps \in (0,1]$ and $\delta \in (0,1)$. The algorithm
$\mathsf{Alg}_{\mathrm{fd}}(\eps,\delta)$
of Definition~\ref{def:fixed_design_algorithm} is $(\eps,\delta)$-PAC on $\Xs$, with budget
\[
  T = \Big\lceil \frac{600\,(\gw(\Xs)^2 + d\log(2/\delta))}{\eps^2} \Big\rceil
  \le \frac{600\,(\gw(\Xs)^2 + d\log(2/\delta))}{\eps^2} + 1 .
\]
\end{theorem}

\begin{corollary}[Theorem~1 of the paper, existential form]\label{cor:upper_exists}
\lean{COLT83.exists_isFixedDesign_isPAC}\leanok
\uses{thm:upper, def:fixed_design, def:pac, def:width}
Let $\Xs$ be spanning, $\eps \in (0,1]$ and $\delta \in (0,1)$. There exist a budget
$T \le 600\,(\gw(\Xs)^2 + d\log(2/\delta))/\eps^2 + 1$ and a fixed-design identification
algorithm with budget $T$ (Definition~\ref{def:fixed_design}) that is $(\eps,\delta)$-PAC on
$\Xs$.
\end{corollary}
\begin{proof}
\uses{thm:upper}
Take $\mathsf{Alg}_{\mathrm{fd}}(\eps,\delta)$ of Theorem~\ref{thm:upper}, which is a fixed-design
algorithm with the required budget.
\end{proof}

\begin{proof}
\uses{cor:fixed_design_from_distribution, lem:mixed_design_bounds, lem:least_squares_law, lem:difference_process_concentration, lem:regret_le_difference_process, lem:approx_argmax_selector, def:pac}
Fix $\theta \in \R^d$ and write $L := \log(1/\delta) > 0$, $\gw := \gw(\Xs)$.

\emph{Design quantities.} By Corollary~\ref{cor:fixed_design_from_distribution} (applied to
$\lambda_0$, $A_0 = A(\lambda_0)$) and Lemma~\ref{lem:mixed_design_bounds}, for every $x \in \Xs$,
\[
  \norm{x}_{\Sigma_T^{-1}}^2 \le \frac4T\norm{x}_{A_0^{-1}}^2 \le \frac{8d}{T},
  \qquad
  \gwidth{\Xs}{\Sigma_T} \le \frac{2}{\sqrt T}\gwidth{\Xs}{A_0} \le \frac{2\sqrt2\,\gw}{\sqrt T} .
\]
Hence $\sigma_T^2 := \max_{x\in\Xs}\norm{x}_{\Sigma_T^{-1}}^2 \le 8d/T$.

\emph{Concentration.} By Lemma~\ref{lem:least_squares_law}, $\Delta := \hat\theta - \theta \sim
\Normal(0,\Sigma_T^{-1})$ under $\Prob_\theta$. Let $D := \max_{x,x'\in\Xs}\ip{x-x'}{\Delta}$, a
measurable function of $y$. Lemma~\ref{lem:difference_process_concentration} gives an event
$E$ with $\Prob_\theta(E) \ge 1-\delta$ on which
\[
  D \le 2\,\gwidth{\Xs}{\Sigma_T} + 2\sigma_T\sqrt{2\cG L}
  \le \frac{4\sqrt2\,\gw}{\sqrt T} + 2\sqrt{\frac{8d}{T}}\sqrt{2\cG L}
  = \frac{4\sqrt2\,\gw + 8\sqrt{\cG dL}}{\sqrt T} .
\]

\emph{Arithmetic.} Using $(a+b)^2 \le 2a^2 + 2b^2$, $(4\sqrt2)^2 = 32$, $8^2 = 64$,
\[
  \big(4\sqrt2\,\gw + 8\sqrt{\cG dL}\big)^2 \le 64\,\gw^2 + 128\,\cG\,dL
  \le 64\,\gw^2 + 320\,dL \le 320\,\big(\gw^2 + d\log(2/\delta)\big),
\]
where $128\cG = 32\pi^2 < 32\cdot 9.87 < 320$ and $L = \log(1/\delta) \le \log(2/\delta)$.
Since $T \ge 600(\gw^2 + d\log(2/\delta))/\eps^2$, on $E$,
\[
  D^2 \le \frac{320\,(\gw^2 + d\log(2/\delta))}{T} \le \frac{320}{600}\,\eps^2 = \frac{8}{15}\eps^2
  \le \frac{9}{16}\eps^2 ,
\]
because $8/15 = 0.5\overline{3} \le 0.5625 = 9/16$; hence $D \le 3\eps/4$ on $E$.

\emph{Regret.} The recommendation satisfies $\ip{\rec}{\hat\theta} \ge \max_x\ip{x}{\hat\theta} - \eps/4$
(Lemma~\ref{lem:approx_argmax_selector}), so by Lemma~\ref{lem:regret_le_difference_process}
with $\gamma = \eps/4$, on $E$: $r(\rec,\theta) \le D + \eps/4 \le 3\eps/4 + \eps/4 = \eps$.
Thus $\Prob_\theta(r(\rec,\theta) \le \eps) \ge \Prob_\theta(E) \ge 1 - \delta$, for every $\theta$,
which is Definition~\ref{def:pac}. The bound $\lceil u\rceil \le u + 1$ gives the budget bound.
\end{proof}

\textbf{Lean remark.} The theorem is an existence statement whose witness is built from
non-constructive choices: $\lambda_1$ (Lemma~\ref{lem:width_attained}), $\lambda_G$
(Theorem~\ref{thm:kiefer_wolfowitz}), the rounded design (Theorem~\ref{thm:rounding}) and the
selector (Lemma~\ref{lem:approx_argmax_selector}); each is obtained by \texttt{Classical.choose}
and Definition~\ref{def:fixed_design_algorithm} packages them into a
\texttt{Learning.detAlgorithm} with a history-independent policy plus a deterministic
recommendation kernel. The good event $E$ is defined in terms of $y$ only, so the PAC
probability is computed under the law $\Normal(X\theta,I_T)$ of $y$
(Lemma~\ref{lem:fixed_design_law}) and the law of $\Delta = \Sigma_T^{-1}X^\top(y - X\theta)$; no
conditioning is needed. The constant $\cG = \pi^2/4$ enters only through the numerical
inequality $32\pi^2 < 320$, i.e.\ $\pi^2 < 10$; the crude bound \texttt{Real.pi\_lt\_four}
is not enough ($16 \not< 10$), but \texttt{Real.pi\_lt\_d2} ($\pi < 3.15$) gives
$\pi^2 < 9.93 < 10$.

\begin{remark}[Asymptotic form]\label{rem:upper_asymptotic}
In the asymptotic form of the paper's Theorem~1: for every spanning action set $\Xs$, every
$\eps \in (0,1]$ and every $\delta \in (0,1/2]$ there is a non-adaptive $(\eps,\delta)$-PAC
algorithm with budget $O\big((d\log(1/\delta) + \gw(\Xs)^2)/\eps^2\big)$, since
$\log(2/\delta) \le 2\log(1/\delta)$ for $\delta \le 1/2$. This asymptotic form is not meant to
be formalized; Theorem~\ref{thm:upper} is the explicit statement.
\end{remark}

\begin{remark}[Comparison with the paper]
The paper's algorithm and analysis (Section~2.1) are the same, with three differences: (a) the
paper's rounding lemma (Appendix~B.1, cited from \cite{katz2020empirical,allen2021near})
requires $T \ge 180d$ and yields $\tau(A_T) \le 2\tau(A(\lambda_0))$, whereas
Theorem~\ref{thm:rounding} yields the Loewner bound $\Sigma_T \succeq (T/4)A_0$ for $T \ge 4d$,
which implies the same kind of control with the factor $4$ instead of $2$; (b) the paper uses
the Borell--TIS inequality with constant $1$ in the exponent, we have $\cG = \pi^2/4$; (c) the
paper's recommendation is an exact maximizer. Items (a)--(c) explain the change of constant from
$360$ to $600$. The dependence on $\eps$, $\delta$, $d$ and $\gw(\Xs)$ is unchanged.
\end{remark}

\section{The unit ball: a direct algorithm}
\label{sec:upper_ball}

\begin{lemma}[Identification on the unit ball by a basis design]\label{lem:ball_identification}
\uses{def:pac, def:fixed_design, def:least_squares}
Let $\Xs = \ball_d$, $\eps > 0$, $\delta \in (0,1)$, and
\[
  n := \Big\lceil \frac{8d + 48\log(1/\delta)}{\eps^2}\Big\rceil, \qquad T := d\,n .
\]
Consider the fixed design that plays each standard basis vector $e_1,\dots,e_d$ exactly $n$
times ($x_t := e_{1 + \lfloor t/n\rfloor}$ for $t < T$), the estimator $\hat\theta \in \R^d$ whose
$i$-th coordinate is the empirical mean of the $n$ observations made at $e_i$ (this is the
least-squares estimator, since $\Sigma_T = nI_d$), and the recommendation
$\rec := \hat\theta/\norm{\hat\theta}$ if $\hat\theta \ne 0$ and $\rec := e_1$ otherwise. This
fixed-design algorithm with budget $T$ is $(\eps,\delta)$-PAC on $\ball_d$, and
$T \le (8d^2 + 48d\log(1/\delta))/\eps^2 + d$.
\end{lemma}

\begin{proof}
\uses{lem:least_squares_law, lem:fixed_design_law, lem:noise_representation, lem:gauss_pi, lem:gauss_law_of_cov, lem:chi_square_tail, def:regret}
\emph{Law of the error.} $\Sigma_T = \sum_{t<T}x_tx_t^\top = n\sum_ie_ie_i^\top = nI_d \in \PD$ and
$\Sigma_T^{-1}X^\top y$ has $i$-th coordinate $\frac1n\sum_{t : x_t = e_i}y_t$, so $\hat\theta$ is
the least-squares estimator and Lemma~\ref{lem:least_squares_law} gives
$\Delta := \hat\theta - \theta \sim \Normal(0, I_d/n)$; by Lemma~\ref{lem:gauss_law_of_cov},
$\Delta \eqd g/\sqrt n$ with $g \sim \Normal(0,I_d)$, so $\norm{\Delta}^2 \eqd \norm{g}^2/n$.
(Directly: by Lemmas~\ref{lem:fixed_design_law} and~\ref{lem:noise_representation},
$\hat\theta_i = \theta_i + \frac1n\sum_{t : x_t = e_i}\eta_t$ with the $\eta_t$ i.i.d.\ standard
Gaussian, and the $d$ block averages are independent $\Normal(0,1/n)$ variables,
Lemma~\ref{lem:gauss_pi}.)

\emph{Good event.} By Lemma~\ref{lem:chi_square_tail},
$\Prob(\norm{g}^2 \ge 2d + 12\log(1/\delta)) \le \delta$. Let $E := \{\norm{\Delta}^2 <
(2d + 12\log(1/\delta))/n\}$, so $\Prob_\theta(E) \ge 1 - \delta$. On $E$, by the choice of $n$,
\[
  \norm{\Delta}^2 < \frac{2d + 12\log(1/\delta)}{n}
  \le \frac{(2d + 12\log(1/\delta))\,\eps^2}{8d + 48\log(1/\delta)} = \frac{\eps^2}{4},
\]
so $\norm{\Delta} \le \eps/2$.

\emph{Regret on the good event.} On $\ball_d$, $\max_{x\in\ball_d}\ip{x}{\theta} = \norm{\theta}$
(Cauchy--Schwarz, attained at $\theta/\norm\theta$ or at any point if $\theta = 0$), so
$r(\rec,\theta) = \norm{\theta} - \ip{\rec}{\theta}$. If $\hat\theta \ne 0$:
\[
  \ip{\rec}{\theta} = \ip{\rec}{\hat\theta} - \ip{\rec}{\Delta} = \norm{\hat\theta} - \ip{\rec}{\Delta}
  \ge \norm{\hat\theta} - \norm{\Delta} \ge \norm{\theta} - 2\norm{\Delta},
\]
using $\norm{\rec} = 1$ and $\norm{\hat\theta} \ge \norm{\theta} - \norm{\Delta}$; hence
$r(\rec,\theta) \le 2\norm{\Delta} \le \eps$. If $\hat\theta = 0$: $\norm{\theta} = \norm{\Delta}
\le \eps/2$ and $r(\rec,\theta) = \norm\theta - \ip{e_1}{\theta} \le 2\norm{\theta} \le \eps$.
In both cases $r(\rec,\theta) \le \eps$ on $E$, so
$\Prob_\theta(r(\rec,\theta) \le \eps) \ge 1-\delta$ for every $\theta$.

\emph{Budget.} $T = dn \le d\big((8d + 48\log(1/\delta))/\eps^2 + 1\big)$.
\end{proof}

\begin{remark}
Lemma~\ref{lem:ball_identification} is a special case of Theorem~\ref{thm:upper} ($\ball_d$ is
spanning, so the theorem applies; the lemma itself uses no spanning hypothesis): since
$\gw(\ball_d) \le d$ (Proposition~\ref{prop:width_le_d}), Theorem~\ref{thm:upper} gives a budget of
order $(d^2 + d\log(1/\delta))/\eps^2$ on the ball, the same order as here. The direct argument is
recorded because it needs neither the width nor the rounding machinery, only the
$\chi^2$ tail bound, and because Chapter~\ref{chap:separation} composes it with the norm
estimation algorithm on each block of the separation instance. The map $y \mapsto \rec$ is
measurable as it is continuous on $\{\hat\theta \ne 0\}$ and constant on the closed set
$\{\hat\theta = 0\}$ (\texttt{Measurable.piecewise}).
\end{remark}

\chapter{Lower bound for non-adaptive fixed designs}
\label{chap:lower_nonadaptive}

This chapter proves Theorem~3 of the paper (Section~2.2 and Appendices~B.3 and~B.4): on a
spanning action set $\Xs$, every non-adaptive fixed-design algorithm that is $(\eps,\delta)$-PAC
with $\delta \le 1/10$ uses a budget $T \ge \gw(\Xs)^2/(70\,\eps^2)$
(Theorem~\ref{thm:lower_nonadaptive}). The argument is
Bayesian. Under the prior $\theta \sim \Normal(0, \tau^2 \Sigma_T^{-1})$, the expected simple
regret of \emph{any} recommendation rule is at least a constant times
$\gwidth{\Xs}{\Sigma_T}$ (Lemma~\ref{lem:bayes_regret_lower}), while the expected regret of an
$(\eps,\delta)$-PAC algorithm is at most $\eps + \delta\,\E Z(\theta)$
(Lemma~\ref{lem:pac_expected_regret}). Comparing the two bounds and using
$\gwidth{\Xs}{\Sigma_T} \ge \gw(\Xs)/\sqrt T$ (Lemma~\ref{lem:width_empirical}) gives the
theorem. Singular design matrices are handled separately
(Lemma~\ref{lem:singular_design_fails}).

From the prerequisite chapters we use the linear theory of Gaussian vectors of
Chapter~\ref{chap:pre_gaussian}: linear images (Lemma~\ref{lem:gauss_linear_image}), joint
Gaussianity of linear functionals of a Gaussian vector (Lemma~\ref{lem:gauss_joint_linear}),
independence of uncorrelated jointly Gaussian vectors
(Lemma~\ref{lem:gauss_uncorrelated_indep}, Mathlib
\texttt{HasGaussianLaw.indepFun\_of\_covariance\_eval}) and the vanishing of
$\E\ip{f(V)}{W}$ for independent $V$, $W$ with $W$ centered
(Lemma~\ref{lem:indep_integral_zero}).

What is new with respect to the paper. The paper uses the Gaussian posterior mean
$\E[\theta \mid y] = c\,\Sigma_T^{-1}X^\top y$; conditional expectations of Gaussian vectors
are not convenient to formalize, so this step is replaced by the orthogonal decomposition
$\theta = W + c\,\Sigma_T^{-1}X^\top y$ with $W$ independent of $y$
(Lemma~\ref{lem:posterior_decomposition}), which yields the same identity for
$\E\ip{\rec}{\theta}$ (Lemma~\ref{lem:bayes_regret_identity}). The recommendation is an
arbitrary Markov kernel, so all statements hold for randomized recommendation rules. The
singular case (Appendix~B.4 of the paper, stated there for finite $\Xs$) is proved for every
compact spanning $\Xs$ and every recommendation kernel.

\section{The Bayesian model and the orthogonal decomposition}
\label{sec:lower_nonadaptive_bayes}

Throughout this chapter $x_0, \dots, x_{T-1} \in \Xs$ is a fixed design of length $T$,
$X \in \R^{T\times d}$ is the matrix with rows $x_t^\top$, $\Sigma_T = X^\top X$ is its design
matrix (Definition~\ref{def:fixed_design}), and $\rho$ is a recommendation kernel from
$\R^T$ to $\Xs$. In Sections~\ref{sec:lower_nonadaptive_bayes}
and~\ref{sec:lower_nonadaptive_regret} we assume $T \ge 1$ and $\Sigma_T \succ 0$.

\begin{definition}[Gaussian prior on the reward vector]\label{def:bayes_prior}
\uses{def:fixed_design, lem:fixed_design_law}
Let $\tau > 0$. In the \emph{Bayesian model with parameter $\tau$}, the reward vector
$\theta$ is a random vector with law $\pi := \Normal(0, \tau^2\Sigma_T^{-1})$, the noise
$\eta \sim \Normal(0, I_T)$ is independent of $\theta$, the observation vector is
$y := X\theta + \eta \in \R^T$ and the recommendation is $\rec \sim \rho(y)$. That is, the
joint law of $(\theta, \eta, \rec)$ is
\[
  \big(\Normal(0,\tau^2\Sigma_T^{-1}) \otimes \Normal(0, I_T)\big)
  \otimes \big( (\theta,\eta) \mapsto \rho(X\theta + \eta) \big).
\]
Equivalently, the joint law of $(\theta, y, \rec)$ is $\pi \otimes \kappa$ for the Markov
kernel $\kappa(\theta) := \Normal(X\theta, I_T) \otimes \rho$, and by
Lemma~\ref{lem:fixed_design_law} $\kappa(\theta)$ is the law $\Prob_\theta$ of $(y, \rec)$ under
the fixed-design algorithm run against $\nu_\theta$ (recall from
Definition~\ref{def:fixed_design} that $(y, \rec)$ is identified with $(\hist_T, \rec)$). This
is the setting of Lemma~\ref{lem:pac_expected_regret}. We write $\E$ for
the expectation under this joint law.
\end{definition}

\textbf{Lean remark.} A randomized recommendation rule $\rec = f(y, U)$ using an internal
seed $U \sim \mu$ independent of $(\theta, \eta)$, as in the paper, is the same thing as the
Markov kernel $\rho(y) := f(y,\cdot)_\#\mu$; only the kernel form is used. The Bayesian model
should be set up so that $(\theta,\eta)$ is visibly a linear image of a standard Gaussian
vector of $\R^{d+T}$: take $(g, g') \sim \Normal(0, I_{d+T})$ (\texttt{stdGaussian}) and
$\theta := \tau\,\Sigma_T^{-1/2} g$, $\eta := g'$; then every linear functional of
$(\theta,\eta)$ has \texttt{HasGaussianLaw}, which is what
Lemma~\ref{lem:posterior_decomposition} needs. The law of $\theta$ is Mathlib's
\texttt{multivariateGaussian} with mean $0$ and covariance $\tau^2\Sigma_T^{-1}$ by Lemma~\ref{lem:gauss_law_of_cov}.

\begin{lemma}[Orthogonal decomposition of the reward vector]\label{lem:posterior_decomposition}
\uses{def:bayes_prior}
In the Bayesian model with parameter $\tau$, let
\[
  c := \frac{\tau^2}{1+\tau^2} \in (0,1), \qquad
  W := \theta - c\,\Sigma_T^{-1}X^\top y \in \R^d .
\]
Then $(W, y)$ is a centered jointly Gaussian vector of $\R^d \times \R^T$ with
$\mathrm{Cov}(W, y) = \E[W y^\top] = 0 \in \R^{d \times T}$. Consequently $W$ and $y$ are
independent, and $W$ is integrable with $\E W = 0$.
\end{lemma}

\begin{proof}
\uses{lem:gauss_pi, lem:gauss_linear_image, lem:gauss_law_of_cov, lem:gauss_joint_linear, lem:gauss_uncorrelated_indep, lem:gauss_norm_moments}
Since $\Sigma_T^{-1}X^\top X = I_d$,
\[
  W = \theta - c\,\Sigma_T^{-1}X^\top(X\theta + \eta) = (1-c)\,\theta - c\,\Sigma_T^{-1}X^\top\eta .
\]
\emph{Joint Gaussianity.} Let $V := (\theta, \eta) \in \R^{d+T}$. Let
$(g, g') \sim \Normal(0, I_{d+T})$; its two blocks $g \sim \Normal(0,I_d)$ and
$g' \sim \Normal(0, I_T)$ are independent (Lemma~\ref{lem:gauss_pi}), and
$\tau\Sigma_T^{-1/2}g \sim \Normal(0, \tau^2\Sigma_T^{-1})$
(Lemmas~\ref{lem:gauss_linear_image} and~\ref{lem:gauss_law_of_cov}). The law of a pair of
independent random vectors is the product of the marginals, so
$V \eqd D\,(g, g')$ with $D := \mathrm{diag}(\tau\Sigma_T^{-1/2}, I_T)$, and hence
$V \sim \Normal(0, S)$ with $S := DD^\top = \mathrm{diag}(\tau^2\Sigma_T^{-1}, I_T)$
(Lemma~\ref{lem:gauss_linear_image}). Now $W = L_1 V$ and $y = L_2 V$ for the linear maps
\[
  L_1 := \big[\, (1-c) I_d \;\; -c\,\Sigma_T^{-1}X^\top \,\big] \in \R^{d \times (d+T)}, \qquad
  L_2 := \big[\, X \;\; I_T \,\big] \in \R^{T \times (d+T)},
\]
so by Lemma~\ref{lem:gauss_joint_linear} the pair $(W, y)$ is centered jointly Gaussian with
cross-covariance $L_1 S L_2^\top$.

\emph{Cross-covariance.} Using the block structure of $S$,
\[
  L_1 S L_2^\top
  = (1-c)\,\tau^2\Sigma_T^{-1} X^\top \;-\; c\,\Sigma_T^{-1}X^\top I_T
  = \big((1-c)\tau^2 - c\big)\,\Sigma_T^{-1}X^\top = 0,
\]
because $(1-c)\tau^2 = \tau^2/(1+\tau^2) = c$. (Equivalently, in the form of the outline:
$\mathrm{Cov}(\theta, y) = \tau^2\Sigma_T^{-1}X^\top$,
$\mathrm{Cov}(y,y) = \tau^2 X\Sigma_T^{-1}X^\top + I_T$, and
$\Sigma_T^{-1}X^\top(\tau^2 X\Sigma_T^{-1}X^\top + I_T) = (1+\tau^2)\Sigma_T^{-1}X^\top$, so that
$\mathrm{Cov}(W,y) = \tau^2\Sigma_T^{-1}X^\top - c(1+\tau^2)\Sigma_T^{-1}X^\top = 0$.)

\emph{Independence and integrability.} Uncorrelated jointly Gaussian vectors are independent
(Lemma~\ref{lem:gauss_uncorrelated_indep}). A Gaussian vector has finite second moments
(Lemma~\ref{lem:gauss_norm_moments}), hence is integrable, and it is centered.
\end{proof}

\begin{lemma}[Bayes identity for the value of the recommendation]\label{lem:bayes_regret_identity}
\uses{def:bayes_prior, lem:posterior_decomposition}
In the Bayesian model with parameter $\tau$ and with $c = \tau^2/(1+\tau^2)$, the random
variables $\ip{\rec}{\theta}$ and $\ip{\rec}{\Sigma_T^{-1}X^\top y}$ are integrable and
\[
  \E\ip{\rec}{\theta} = c\,\E\ip{\rec}{\Sigma_T^{-1}X^\top y} .
\]
\end{lemma}

\begin{proof}
\uses{lem:posterior_decomposition, lem:indep_integral_zero, lem:gauss_norm_moments}
Let $R := \sup_{x\in\Xs}\norm{x} < \infty$ ($\Xs$ is compact). Integrability:
$\abs{\ip{\rec}{\theta}} \le R\norm{\theta}$ and
$\abs{\ip{\rec}{\Sigma_T^{-1}X^\top y}} \le R\norm{\Sigma_T^{-1}X^\top y}$, and Gaussian
vectors are integrable (Lemma~\ref{lem:gauss_norm_moments}). By
Lemma~\ref{lem:posterior_decomposition}, $\theta = W + c\,\Sigma_T^{-1}X^\top y$, so
\[
  \ip{\rec}{\theta} = \ip{\rec}{W} + c\,\ip{\rec}{\Sigma_T^{-1}X^\top y},
\]
and it remains to prove $\E\ip{\rec}{W} = 0$. Let
$\bar\rho(y) := \int_{\Xs} x\,\rho(y, dx) \in \R^d$; it is a measurable function of $y$ with
$\norm{\bar\rho(y)} \le R$. Since, given $(\theta, \eta)$, the recommendation has law
$\rho(y)$ and $W$ is a function of $(\theta,\eta)$, Fubini for the composition-product gives
$\E\ip{\rec}{W} = \E\ip{\bar\rho(y)}{W}$. Now $W$ is integrable and centered, $W$ is
independent of $y$ (Lemma~\ref{lem:posterior_decomposition}) and $\bar\rho$ is bounded
measurable, so Lemma~\ref{lem:indep_integral_zero} (with $V = y$ and $f = \bar\rho$) gives
$\E\ip{\bar\rho(y)}{W} = 0$.
\end{proof}

\textbf{Lean remark.} The seed of a randomized rule never appears: integrating the kernel
$\rho$ first reduces the claim to the deterministic bounded function $\bar\rho$ of $y$, and
the only independence needed is $W \perp y$, which comes from Mathlib's
\texttt{HasGaussianLaw.indepFun\_of\_covariance\_eval} through
Lemma~\ref{lem:gauss_uncorrelated_indep}. No conditional expectation is used anywhere in
this chapter.

\section{Lower bound on the Bayesian expected regret}
\label{sec:lower_nonadaptive_regret}

\begin{lemma}[Bayesian lower bound on the expected simple regret]\label{lem:bayes_regret_lower}
\uses{def:bayes_prior, def:regret, def:width_matrix}
In the Bayesian model with parameter $\tau > 0$, for every recommendation kernel $\rho$,
$r(\rec,\theta)$ is integrable and
\[
  \E\big[ r(\rec, \theta) \big] \;\ge\; \frac{\tau(1-\tau)}{1+\tau^2}\,\gwidth{\Xs}{\Sigma_T} .
\]
\end{lemma}

\begin{proof}
\uses{lem:bayes_regret_identity, lem:width_matrix_wd, lem:gauss_linear_image, lem:gauss_law_of_cov}
Write $w_T := \gwidth{\Xs}{\Sigma_T}$ and $c := \tau^2/(1+\tau^2)$.
\begin{enumerate}
  \item[(i)] \emph{Decomposition.} $r(\rec,\theta) = \max_{x\in\Xs}\ip{x}{\theta} - \ip{\rec}{\theta}$.
    Both terms are integrable: the first by Lemma~\ref{lem:width_matrix_wd} (applied below)
    and the second by Lemma~\ref{lem:bayes_regret_identity}. Hence
    $\E[r(\rec,\theta)] = \E\max_x\ip{x}{\theta} - \E\ip{\rec}{\theta}$.
  \item[(ii)] \emph{The prior term.} $G := \tau^{-1}\theta \sim \Normal(0, \Sigma_T^{-1})$
    (Lemma~\ref{lem:gauss_linear_image}), and $\max_x\ip{x}{\theta} = \tau\max_x\ip{x}{G}$
    since $\tau > 0$. By Lemma~\ref{lem:width_matrix_wd} (the width only depends on the law
    of $G$), $\E\max_x\ip{x}{\theta} = \tau\,w_T$.
  \item[(iii)] \emph{The recommendation term.} By Lemma~\ref{lem:bayes_regret_identity} and
    $\Sigma_T^{-1}X^\top y = \Sigma_T^{-1}X^\top X\theta + \Sigma_T^{-1}X^\top\eta
    = \theta + \Sigma_T^{-1}X^\top\eta$,
    \[
      \E\ip{\rec}{\theta} = c\,\E\ip{\rec}{\theta + \Sigma_T^{-1}X^\top\eta} .
    \]
    Pointwise, $\ip{\rec}{v} \le \max_{x\in\Xs}\ip{x}{v}$ for every $v \in \R^d$ because
    $\rec \in \Xs$, and $\max_x\ip{x}{a+b} \le \max_x\ip{x}{a} + \max_x\ip{x}{b}$. Therefore
    \[
      \E\ip{\rec}{\theta} \le c\Big( \E\max_x\ip{x}{\theta} + \E\max_x\ip{x}{\Sigma_T^{-1}X^\top\eta} \Big).
    \]
  \item[(iv)] \emph{The noise term.} $\Sigma_T^{-1}X^\top\eta$ is a linear image of
    $\eta \sim \Normal(0, I_T)$, so $\Sigma_T^{-1}X^\top\eta \sim
    \Normal(0, \Sigma_T^{-1}X^\top X \Sigma_T^{-1}) = \Normal(0, \Sigma_T^{-1})$
    (Lemmas~\ref{lem:gauss_linear_image} and~\ref{lem:gauss_law_of_cov}), and
    Lemma~\ref{lem:width_matrix_wd} gives $\E\max_x\ip{x}{\Sigma_T^{-1}X^\top\eta} = w_T$.
  \item[(v)] \emph{Conclusion.} Combining,
    \[
      \E[r(\rec,\theta)] \ge \tau w_T - c(\tau w_T + w_T)
      = \Big(\tau - \frac{\tau^2(1+\tau)}{1+\tau^2}\Big) w_T
      = \frac{\tau + \tau^3 - \tau^2 - \tau^3}{1+\tau^2}\, w_T
      = \frac{\tau(1-\tau)}{1+\tau^2}\, w_T .
    \]
\end{enumerate}
\end{proof}

\begin{lemma}[Optimized constant]\label{lem:bayes_regret_lower_opt}
\uses{lem:bayes_regret_lower}
With $\tau := \sqrt2 - 1$, for every recommendation kernel $\rho$,
\[
  \E\big[ r(\rec,\theta) \big] \ge \frac{\sqrt2 - 1}{2}\,\gwidth{\Xs}{\Sigma_T}
  \ge 0.207\,\gwidth{\Xs}{\Sigma_T} .
\]
\end{lemma}

\begin{proof}
\uses{lem:bayes_regret_lower}
For $\tau = \sqrt2 - 1$: $\tau^2 = 3 - 2\sqrt2$, so $1 + \tau^2 = 4 - 2\sqrt2 = 2(2-\sqrt2)$,
and $1 - \tau = 2 - \sqrt2$. Hence
$\frac{\tau(1-\tau)}{1+\tau^2} = \frac{(\sqrt2-1)(2-\sqrt2)}{2(2-\sqrt2)} = \frac{\sqrt2-1}{2}$.
Finally $1.414^2 = 1.999396 \le 2$ gives $\sqrt2 \ge 1.414$ and
$(\sqrt2-1)/2 \ge 0.207$. Apply Lemma~\ref{lem:bayes_regret_lower} (the width is
nonnegative by Lemma~\ref{lem:width_matrix_wd}).
\end{proof}

\begin{remark}[Choice of $\tau$]\label{rem:lower_nonadaptive_tau}
The function $\tau \mapsto \tau(1-\tau)/(1+\tau^2)$ has derivative
$(1 - 2\tau - \tau^2)/(1+\tau^2)^2$, which vanishes on $(0,1)$ exactly at $\tau = \sqrt2 - 1$;
this is the maximizer, as stated in the paper. Only the value at $\sqrt2 - 1$ is used.
\end{remark}

\section{Singular designs}
\label{sec:lower_nonadaptive_singular}

\begin{lemma}[Singular designs fail]\label{lem:singular_design_fails}
\uses{def:arm_set, def:fixed_design, lem:fixed_design_law, def:regret}
Let $\Xs$ be spanning and let $x_0,\dots,x_{T-1} \in \Xs$ be a fixed design of length $T \ge 0$
(Definition~\ref{def:fixed_design}) whose design matrix $\Sigma_T = \sum_{t<T}x_tx_t^\top$ is
singular; if $T = 0$ (so that $\Sigma_0 = 0$), assume moreover that $\Xs$ has at least two points
(for $T \ge 1$ this follows from the spanning hypothesis, see the proof). Let
$\eps > 0$ and let $\rho$ be any recommendation kernel from $\R^T$ to $\Xs$. Then there are
$\theta_+, \theta_- \in \R^d$ such that
\begin{enumerate}
  \item[(i)] the laws of $(y, \rec)$ under $\theta_+$ and under $\theta_-$ coincide:
    $\Prob_{\theta_+} = \Prob_{\theta_-} = \Normal(0, I_T) \otimes \rho$;
  \item[(ii)] $r(\rec, \theta_+) + r(\rec,\theta_-) = 4\eps$ for every $\rec \in \Xs$.
\end{enumerate}
Consequently there is $\theta \in \{\theta_+,\theta_-\}$ with
$\Prob_\theta\big(r(\rec,\theta) > \eps\big) \ge 1/2$; in particular the fixed-design
algorithm is not $(\eps,\delta)$-PAC for any $\delta < 1/2$.
\end{lemma}

\begin{proof}
\uses{lem:fixed_design_law, def:arm_set}
\emph{Step 1: a direction invisible to the design but not to $\Xs$.} We claim there is
$v \in \R^d$ such that $\ip{x_t}{v} = 0$ for all $t < T$ and $x \mapsto \ip{x}{v}$ is not
constant on $\Xs$.
If $T \ge 1$: since $\Sigma_T$ is singular there is $v \ne 0$ with $\Sigma_T v = 0$, and then
$0 = v^\top\Sigma_T v = \sum_{t<T}\ip{x_t}{v}^2$ forces $\ip{x_t}{v} = 0$ for every $t$. If
$\ip{\cdot}{v}$ were constant on $\Xs$, its value would be $\ip{x_0}{v} = 0$, so
$v \perp \Xs$, hence $v \perp \Span(\Xs) = \R^d$ (spanning) and $v = 0$, a contradiction. (This also
shows that $\Xs$ has at least two points when $T \ge 1$.)
If $T = 0$: pick $x \ne x'$ in $\Xs$ and $v := x - x'$; then
$\ip{x}{v} - \ip{x'}{v} = \norm{v}^2 > 0$, and the condition on the $x_t$ is vacuous.

\emph{Step 2: the two instances.} Let $a := \max_{x\in\Xs}\ip{x}{v}$ and
$b := \min_{x\in\Xs}\ip{x}{v}$ (compactness); by Step~1, $a > b$. Put
$\alpha := 4\eps/(a-b) > 0$ and $\theta_\pm := \pm\alpha v$. Since $\ip{x_t}{v} = 0$ for all
$t$, $X\theta_+ = X\theta_- = 0$, so by Lemma~\ref{lem:fixed_design_law} the law of
$(y,\rec)$ is $\Normal(0, I_T)\otimes\rho$ under both instances. This is (i).

\emph{Step 3: the regrets.} For every $\rec\in\Xs$,
\[
  r(\rec,\theta_+) = \max_x \alpha\ip{x}{v} - \alpha\ip{\rec}{v} = \alpha\big(a - \ip{\rec}{v}\big),
  \qquad
  r(\rec,\theta_-) = \max_x \big(-\alpha\ip{x}{v}\big) + \alpha\ip{\rec}{v} = \alpha\big(\ip{\rec}{v} - b\big),
\]
so $r(\rec,\theta_+) + r(\rec,\theta_-) = \alpha(a-b) = 4\eps$. This is (ii).

\emph{Step 4: conclusion.} By (ii), for every $\rec \in \Xs$ at least one of
$r(\rec,\theta_+)$, $r(\rec,\theta_-)$ is $\ge 2\eps > \eps$. With $\mu := \Normal(0,I_T)\otimes\rho$
and $E_\pm := \{(y, x) \in \R^T\times\Xs : r(x,\theta_\pm) > \eps\}$ (measurable, since
$x \mapsto r(x,\theta)$ is continuous), $E_+ \cup E_- = \R^T\times\Xs$, hence
$\mu(E_+) + \mu(E_-) \ge 1$ and one of them is at least $1/2$. By (i),
$\Prob_{\theta_\pm}(r(\rec,\theta_\pm) > \eps) = \mu(E_\pm)$.
\end{proof}

\begin{remark}[The case $T = 0$]\label{rem:lower_nonadaptive_singleton}
Budget $T = 0$ is a legitimate fixed design (Definitions~\ref{def:bai_alg}
and~\ref{def:fixed_design}): the design is empty, $\Sigma_0 = 0$ is singular and $\rho$ is a
probability measure on $\Xs$. For $T = 0$ the hypothesis that $\Xs$ has at least two points
cannot be dropped: if $\Xs = \{x_0\}$ with $x_0 \ne 0$ (a spanning action set when $d = 1$),
then every recommendation equals $x_0$ and the regret is identically $0$, so the length-$0$
design is $(\eps,\delta)$-PAC. For $T \ge 1$ no such assumption is needed. This is harmless for
Theorem~\ref{thm:lower_nonadaptive}, which assumes $T \ge 1$; its conclusion also holds for
$T = 0$, since either $\Xs$ has two points and no length-$0$ design is $(\eps,\delta)$-PAC with
$\delta \le 1/10$, or $\Xs = \{x_0\}$ and $\gw(\{x_0\}) = 0$ (as $\E\ip{x_0}{A^{-1/2}g} = 0$).
\end{remark}

\section{The lower bound}
\label{sec:lower_nonadaptive_theorem}

\begin{theorem}[Lower bound for non-adaptive fixed designs]\label{thm:lower_nonadaptive}
\lean{COLT83.le_budget_of_isFixedDesign_of_isPAC}\leanok
\uses{def:arm_set, def:fixed_design, def:pac, def:width}
Let $\Xs$ be spanning, $\eps > 0$ and $\delta \in (0, 1/10]$. If a fixed-design algorithm with budget $T \ge 1$
(design $x_0,\dots,x_{T-1} \in \Xs$ and any recommendation kernel $\rho$) is
$(\eps,\delta)$-PAC on $\Xs$, then
\[
  T \;\ge\; \frac{\gw(\Xs)^2}{70\,\eps^2} .
\]
\end{theorem}

\begin{proof}
\uses{lem:singular_design_fails, def:bayes_prior, def:fixed_design, lem:fixed_design_law, lem:pac_expected_regret, lem:width_symm, lem:width_matrix_wd, lem:gauss_linear_image, lem:bayes_regret_lower_opt, lem:width_empirical}
\emph{Step 1: $\Sigma_T$ is positive definite.} Otherwise Lemma~\ref{lem:singular_design_fails}
(with $T \ge 1$ and $\Xs$ spanning) gives $\theta$ with $\Prob_\theta(r(\rec,\theta) > \eps) \ge 1/2 > \delta$,
contradicting the PAC property. So $\Sigma_T \succ 0$, and we write
$w_T := \gwidth{\Xs}{\Sigma_T} \ge 0$.

\emph{Step 2: upper bound on the Bayesian regret.} Let $\tau := \sqrt2 - 1$ and consider the
Bayesian model of Definition~\ref{def:bayes_prior} with prior $\pi = \Normal(0,\tau^2\Sigma_T^{-1})$.
Its joint law is $\pi \otimes \kappa$ with the Markov kernel
$\kappa(\theta) = \Normal(X\theta, I_T) \otimes \rho$ (a Gaussian kernel whose mean depends
linearly on $\theta$, composed with $\rho$), and $\kappa(\theta) = \Prob_\theta$ for every
$\theta$ by Lemma~\ref{lem:fixed_design_law}, with $(\hist_T, \rec)$ identified with
$(y, \rec)$ as in Definition~\ref{def:fixed_design}. Hence Lemma~\ref{lem:pac_expected_regret}
applies with this $\kappa$, provided $\int Z\,d\pi < \infty$. With $G := \tau^{-1}\theta \sim \Normal(0,\Sigma_T^{-1})$
(Lemma~\ref{lem:gauss_linear_image}) and $Z(\theta) = \tau Z(G)$, Lemma~\ref{lem:width_symm}
gives $\int Z\,d\pi = \E Z(\theta) = \tau\,\E Z(G) = 2\tau w_T < \infty$. Hence
\[
  \E\big[r(\rec,\theta)\big] \le \eps + \delta \cdot 2\tau\, w_T .
\]

\emph{Step 3: lower bound on the Bayesian regret.} By Lemma~\ref{lem:bayes_regret_lower_opt},
$\E[r(\rec,\theta)] \ge 0.207\, w_T$.

\emph{Step 4: arithmetic.} Combining, $(0.207 - 2\tau\delta)\, w_T \le \eps$. Since
$1.4143^2 = 2.00024 \ge 2$ we have $\tau = \sqrt2 - 1 \le 0.4143$, so with $\delta \le 1/10$,
$2\tau\delta \le 0.08286 \le 0.083$, and $0.207 - 0.083 = 0.124 \ge 0.12$. Therefore
$0.12\, w_T \le \eps$.

\emph{Step 5: from $w_T$ to $\gw(\Xs)$.} By Lemma~\ref{lem:width_empirical},
$w_T \ge \gw(\Xs)/\sqrt T$, so $0.12\,\gw(\Xs) \le \eps\sqrt T$; squaring (both sides are
nonnegative), $0.0144\,\gw(\Xs)^2 \le \eps^2 T$. Since $0.0144 \cdot 70 = 1.008 \ge 1$, this
gives $T \ge \gw(\Xs)^2/(70\,\eps^2)$.
\end{proof}

\begin{corollary}[The paper's form of the lower bound]\label{cor:lower_nonadaptive_H}
\uses{thm:lower_nonadaptive}
Let $\eps > 0$ and $0 \le H_1 \le H_2$. Suppose that for every $\delta \in (0,1)$ there is a
fixed-design algorithm that is $(\eps,\delta)$-PAC on $\Xs$ with budget
$T_\delta \le (H_1\log(1/\delta) + H_2)/\eps^2$. Then
\[
  H_2 \;\ge\; \frac{\gw(\Xs)^2}{70\,(\log 10 + 1)} .
\]
\end{corollary}

\begin{proof}
\uses{thm:lower_nonadaptive}
Only $\delta = 1/10$ is used. Theorem~\ref{thm:lower_nonadaptive} gives
$T_{1/10} \ge \gw(\Xs)^2/(70\eps^2)$, while
$T_{1/10} \le (H_1 \log 10 + H_2)/\eps^2 \le H_2(\log 10 + 1)/\eps^2$ because $H_1 \le H_2$.
Hence $H_2(\log 10 + 1) \ge \gw(\Xs)^2/70$.
\end{proof}

\begin{remark}[Comparison with the paper]\label{rem:lower_nonadaptive_paper}
The paper obtains $H_2 \ge 0.0144\,\gw(\Xs)^2/(\log 10 + 1)$; we lose slightly by rounding
$0.0144$ down to $1/70$. The statement of Theorem~\ref{thm:lower_nonadaptive} is for a single
algorithm and a single $\delta$, which is the form used in Chapter~\ref{chap:separation}.
The lower bound $\Omega(d\log(1/\delta)/\eps^2)$ that the paper states alongside
Theorem~3 is the adaptive lower bound of Chapter~\ref{chap:lower_adaptive}, which applies in
particular to fixed designs.
\end{remark}

\chapter{Lower bound for adaptive algorithms}
\label{chap:lower_adaptive}

This chapter proves Theorem~2 of the paper (Appendix~B.2, Theorem~\emph{main} and
Lemma~\emph{baseline} there): for $d \ge 2$ and $\delta < 1/16$, every $(\eps,\delta)$-PAC
identification algorithm on a spanning action set $\Xs$, adaptive or not, with a randomized
recommendation or not,
uses a budget $n \ge c_{\mathrm{ad}}\, d\log(1/\delta)/\eps^2$ with the explicit constant
$c_{\mathrm{ad}} = 1/(11556 + 5760\sqrt2) \ge 1/20000$ (Theorem~\ref{thm:lower_adaptive}).
The proof has three steps. First, a linear change of coordinates
(Lemma~\ref{lem:transform_preserves_pac}) reduces to a \emph{normalized instance}
$\Xs' \subseteq \ball_d$ containing unit vectors $v_1,\dots,v_m$ with
$\sum_j \lambda_j v_jv_j^\top = I_d/d$ (Definition~\ref{def:normalized_instance}). Second, a
two-point argument gives the baseline $n \ge c_0\log(1/(4\delta))/\eps^2$
(Lemma~\ref{lem:baseline_lower}). Third, a PAC algorithm is turned into a test between
$\theta = 0$ and the mixture of the $\theta^{(j)} = 3\eps v_j$
(Lemma~\ref{lem:test_from_alg}), and the mixture's KL to $\Prob_0$ is at most
$9N\eps^2/(2d)$ (Lemma~\ref{lem:mixture_kl}), which forces $N \ge (2d/(9\eps^2))\log(1/(8\delta))$
(Lemma~\ref{lem:test_lower}).

From the prerequisite chapters we use the divergence decomposition for algorithm-environment
sequences (Theorem~\ref{thm:divergence_decomposition},
Corollary~\ref{cor:divergence_decomposition_gaussian}), the fact that a common recommendation
kernel does not change the KL (Lemma~\ref{lem:kl_data_processing_recommendation}), the
Bretagnolle--Huber inequality (Lemma~\ref{lem:bretagnolle_huber}) and the convexity of the KL
in its second argument for finite mixtures (Lemma~\ref{lem:kl_mixture_convex}), all from
Chapter~\ref{chap:pre_info} (Mathlib \texttt{ProbabilityTheory.klDiv}; see
\cite[Chapters~14 and~15]{lattimore2020bandit}), together with Gaussian mean concentration
(Lemma~\ref{lem:gauss_mean_concentration}, in the form of Lemma~\ref{lem:noise_azuma} for the
formalization) and the sequential composition of algorithms (Lemma~\ref{lem:composition}; the
formalization only needs the special case of Lemma~\ref{lem:prefix_transfer}).

What is new with respect to the paper. The paper normalizes the instance with the L\"owner
ellipsoid of $\conv(\Xs\cup(-\Xs))$ and John's theorem \cite{ball1997elementary}; we use
instead the Kiefer--Wolfowitz theorem on $G$-optimal designs
(Theorem~\ref{thm:kiefer_wolfowitz}, \cite{kiefer1960equivalence,pukelsheim2006optimal}),
which is needed anyway for the upper bound and gives exactly the structure the argument
uses (Remark~\ref{rem:lower_adaptive_john}). All constants are explicit, the recommendation
is a Markov kernel, and the test built from the algorithm is an honest sequential
composition in the sense of Lemma~\ref{lem:composition}. Throughout, $d \ge 2$; the spanning
hypothesis on $\Xs$ enters only through Theorem~\ref{thm:kiefer_wolfowitz} and is stated where
it is needed.

\section{Linear transformations and the normalized instance}
\label{sec:lower_adaptive_transform}

\begin{lemma}[Linear transformation of an instance]\label{lem:transform_preserves_pac}
\uses{def:arm_set, def:env, def:bai_alg, def:regret, def:pac}
Let $M \in \R^{d\times d}$ be invertible and $\Xs' := M\Xs = \{Mx : x \in \Xs\}$. Then $\Xs'$
is an action set (nonempty and compact), and it is spanning if $\Xs$ is. For every identification algorithm
$\mathsf{Alg} = (\mathrm{alg}, \rho)$ on $\Xs$ with budget $n$ there is an identification
algorithm $\mathsf{Alg}' = (\mathrm{alg}',\rho')$ on $\Xs'$ with budget $n$, the
\emph{transformed algorithm}, with the following property. For $\theta \in \R^d$ put
$\theta' := M^{-\top}\theta$ and let
$\Phi\big((x_t,y_t)_{t<n}, \rec\big) := \big((Mx_t, y_t)_{t<n}, M\rec\big)$. Then the law
$\Prob'_{\theta'}$ of $(\hist'_n, \rec')$ for $\mathsf{Alg}'$ against the environment
$\nu'_{\theta'}$ on $\Xs'$ is the image of $\Prob_\theta$ by $\Phi$:
\[
  \Prob'_{\theta'} = \Phi_\#\Prob_\theta .
\]
Moreover $\ip{Mx}{\theta'} = \ip{x}{\theta}$ for all $x \in \R^d$, hence the simple regrets
satisfy $r'(M\rec, \theta') = r(\rec,\theta)$ (with $r'$ the simple regret on $\Xs'$), and
$\mathsf{Alg}'$ is $(\eps,\delta)$-PAC on $\Xs'$ if and only if $\mathsf{Alg}$ is
$(\eps,\delta)$-PAC on $\Xs$.
\end{lemma}

\begin{proof}
\uses{def:env, def:bai_alg, def:pac, lem:noise_representation}
\emph{$\Xs'$ is an action set.} $\Xs'$ is the nonempty image of a compact set by a continuous map,
and if $\Span(\Xs) = \R^d$ then $\Span(M\Xs) = M\,\Span(\Xs) = M\R^d = \R^d$ since $M$ is invertible.

\emph{Values.} $\ip{Mx}{M^{-\top}\theta} = x^\top M^\top M^{-\top}\theta = \ip{x}{\theta}$.
Hence $\max_{x'\in\Xs'}\ip{x'}{\theta'} = \max_{x\in\Xs}\ip{Mx}{\theta'} = \max_{x\in\Xs}\ip{x}{\theta}$
and $\ip{M\rec}{\theta'} = \ip{\rec}{\theta}$, so $r'(M\rec,\theta') = r(\rec,\theta)$.

\emph{The transformed algorithm.} Let $\phi := (x \mapsto Mx)$, a homeomorphism from $\Xs$
onto $\Xs'$, and let $\Phi_t$ be the induced bijection of histories of length $t$,
$\Phi_t((x_s,y_s)_{s<t}) := ((Mx_s,y_s)_{s<t})$. Define $\mathrm{alg}'$ by pushing forward the
initial action distribution of $\mathrm{alg}$ by $\phi$, and, at time $t \ge 1$, given a
history $h'$ of length $t$ on $\Xs'$, by playing the image by $\phi$ of the action drawn by
the policy of $\mathrm{alg}$ given $\Phi_t^{-1}(h')$. Define
$\rho'(h') := \phi_\#\,\rho(\Phi_n^{-1}(h'))$. These are Markov kernels since $\phi$ and
$\Phi_t^{-1}$ are measurable.

\emph{Identification of the law.} Let $(\hist_n,\rec)$ have law $\Prob_\theta$, i.e.\ be an
algorithm-environment sequence for $(\mathrm{alg},\nu_\theta)$ followed by
$\rec \sim \rho(\hist_n)$. We check that $\Phi(\hist_n,\rec)$ is an algorithm-environment
sequence for $(\mathrm{alg}',\nu'_{\theta'})$ followed by $\rho'$. Conditionally on the past
$\Phi_t(\hist_t)$, the action $Mx_t$ has the image by $\phi$ of the conditional law of $x_t$
given $\hist_t$, which is the policy of $\mathrm{alg}$ at $\Phi_t^{-1}\Phi_t(\hist_t) = \hist_t$,
i.e.\ the policy of $\mathrm{alg}'$ at $\Phi_t(\hist_t)$. Conditionally on the past and on
$Mx_t$, the observation $y_t$ has law $\Normal(\ip{x_t}{\theta}, 1) = \Normal(\ip{Mx_t}{\theta'},1)$,
which is the reward kernel of $\nu'_{\theta'}$ evaluated at $Mx_t$
(Definition~\ref{def:env}). Finally, conditionally on $\Phi_n(\hist_n)$, $M\rec$ has law
$\phi_\#\rho(\hist_n) = \rho'(\Phi_n(\hist_n))$. By uniqueness of the law of an
algorithm-environment sequence (\texttt{Learning.isAlgEnvSeq\_unique},
Definition~\ref{def:bai_alg}), $\Phi_\#\Prob_\theta = \Prob'_{\theta'}$.

\emph{PAC equivalence.} For every $\theta$,
$\Prob'_{\theta'}(r'(\rec',\theta') \le \eps) = \Prob_\theta(r'(M\rec,\theta') \le \eps)
= \Prob_\theta(r(\rec,\theta)\le\eps)$. Since $\theta \mapsto M^{-\top}\theta$ is a bijection
of $\R^d$, the PAC condition for all $\theta'$ on $\Xs'$ is equivalent to the PAC condition
for all $\theta$ on $\Xs$.
\end{proof}

\textbf{Lean remark.} The formalization does \emph{not} transform the algorithm: it works in
the original coordinates and only uses the map $M$ inside the analysis. Indeed every quantity of
Sections~\ref{sec:lower_adaptive_baseline} and~\ref{sec:lower_adaptive_test} is of the form
$\ip{Mx}{u}$ for an action $x \in \Xs$ and a fixed vector $u$, and $\ip{Mx}{u} = \ip{x}{Mu}$
since $M$ is symmetric (\texttt{COLT83.inner\_toEuclideanCLM\_normalizeMat}); the reward vectors
of the normalized instance are therefore replaced by their images $Mu$ on the original instance
(e.g.\ $\theta_\pm = \pm\frac{4\eps}{\varrho}Mu$ in
Definition~\ref{def:lower_adaptive_two_point_instance}). This lemma is thus not needed for
Theorem~\ref{thm:lower_adaptive} and is not formalized; the transport of a run along a
measurable map of the \emph{sample space} (as opposed to the action type) is
\texttt{Learning.IdentAlg.IsRun.comp\_hasLaw}.

\begin{definition}[Normalized instance]\label{def:normalized_instance}
\lean{COLT83.normalizeMat, COLT83.IsGOptimalDesign.norm_normalizeMat_le, COLT83.IsGOptimalDesign.norm_normalizeMat_eq_one, COLT83.IsGOptimalDesign.designMatrix_mapDomain_normalizeMat, COLT83.transpose_normalizeMat, COLT83.inner_toEuclideanCLM_normalizeMat}\leanok
\uses{def:arm_set, thm:kiefer_wolfowitz, cor:g_optimal_norm_bound, cor:g_optimal_normalized, lem:transform_preserves_pac}
Let $\Xs$ be spanning and let $\lambda_G \in \simplex_{\Xs}$ be a $G$-optimal design with
$A_G := A(\lambda_G) \succ 0$ (Theorem~\ref{thm:kiefer_wolfowitz}), with support $\{x_1, \dots, x_m\}$ ($m \ge 1$ distinct
points of $\Xs$) and weights $\lambda_j := \lambda_{G,x_j} > 0$, $\sum_{j=1}^m \lambda_j = 1$.
Set
\[
  M := d^{-1/2} A_G^{-1/2}, \qquad \Xs' := M\Xs, \qquad v_j := M x_j \quad (1 \le j \le m).
\]
By Corollaries~\ref{cor:g_optimal_norm_bound} and~\ref{cor:g_optimal_normalized}:
$\Xs' \subseteq \ball_d$ is a spanning action set (Lemma~\ref{lem:transform_preserves_pac}),
$v_j \in \Xs'$, $\norm{v_j} = 1$ for every $j$, and
\[
  \sum_{j=1}^m \lambda_j\, v_j v_j^\top = \frac{1}{d} I_d .
\]
In particular $m \ge d \ge 2$ (the left-hand side has rank at most $m$). We let $J$ be a random
index in $\{1,\dots,m\}$ with $\Prob(J = j) = \lambda_j$, so that
$\E[v_J v_J^\top] = I_d/d$. The statements of Sections~\ref{sec:lower_adaptive_baseline}
and~\ref{sec:lower_adaptive_test} are about algorithms on the normalized instance $\Xs'$; there
we write $\Prob_\theta$, $\E_\theta$ for the law of $(\hist, \rec)$ of the algorithm under
consideration against the environment $\nu'_\theta$ of $\Xs'$ (Definition~\ref{def:bai_alg}),
and $r$ for the simple regret on $\Xs'$.
\end{definition}

\begin{remark}[Kiefer--Wolfowitz instead of John]\label{rem:lower_adaptive_john}
The paper obtains points $u_j \in \Xs\cap\sphere^{d-1}$ (after normalization by the L\"owner
ellipsoid of $\conv(\Xs\cup(-\Xs))$) and weights $c_j > 0$ with $\sum_j c_j u_ju_j^\top = I_d$,
$\sum_j c_j = d$, by John's theorem applied to the polar body. Definition~\ref{def:normalized_instance}
provides the same data with $c_j = d\lambda_j$: after the linear map $M = d^{-1/2}A_G^{-1/2}$
the support points of the $G$-optimal design are unit vectors whose weighted second-moment
matrix is $I_d/d$, and $\Xs' \subseteq \ball_d$ because $\norm{A_G^{-1/2}x}^2 \le d$ for
all $x \in \Xs$. Only these three facts are used below.
\end{remark}

\section{The baseline lower bound}
\label{sec:lower_adaptive_baseline}

\begin{lemma}[A well-separated pair of unit arms]\label{lem:separated_pair}
\lean{COLT83.IsGOptimalDesign.sum_mul_inner_normalizeMat_sq, COLT83.IsGOptimalDesign.exists_inner_normalizeMat_le, COLT83.IsGOptimalDesign.norm_sub_sq_normalizeMat, COLT83.IsGOptimalDesign.exists_separated_pair}\leanok
\uses{def:normalized_instance}
In the setting of Definition~\ref{def:normalized_instance},
\[
  \E\big[\ip{v_1}{v_J}^2\big] = \sum_{j=1}^m \lambda_j \ip{v_1}{v_j}^2 = \frac1d .
\]
Consequently there is $k \in \{1,\dots,m\}$ with $\ip{v_1}{v_k} \le 1/\sqrt d \le 1/\sqrt2$.
Moreover, every $k$ with $\ip{v_1}{v_k} \le 1/\sqrt2$ satisfies $k \ne 1$ and
\[
  \norm{v_1 - v_k}^2 = 2 - 2\ip{v_1}{v_k} \ge 2 - \sqrt2 > 0 .
\]
\end{lemma}

\begin{proof}
\leanok
\uses{def:normalized_instance}
$\sum_j\lambda_j\ip{v_1}{v_j}^2 = v_1^\top\big(\sum_j\lambda_jv_jv_j^\top\big)v_1
= v_1^\top (I_d/d) v_1 = \norm{v_1}^2/d = 1/d$. A weighted average with weights
$\lambda_j > 0$, $\sum_j\lambda_j = 1$, is at least its smallest term, so some $k$ has
$\ip{v_1}{v_k}^2 \le 1/d$, hence $\ip{v_1}{v_k} \le \abs{\ip{v_1}{v_k}} \le 1/\sqrt d$, and
$1/\sqrt d \le 1/\sqrt2$ because $d \ge 2$. Now let $k$ be any index with
$\ip{v_1}{v_k} \le 1/\sqrt2$. Since $\ip{v_1}{v_1} = 1 > 1/\sqrt2$, $k \ne 1$.
Finally $\norm{v_1 - v_k}^2 = \norm{v_1}^2 + \norm{v_k}^2 - 2\ip{v_1}{v_k} = 2 - 2\ip{v_1}{v_k}
\ge 2 - \sqrt2$.
\end{proof}

\begin{definition}[The two-point instance]\label{def:lower_adaptive_two_point_instance}
\uses{def:normalized_instance, lem:separated_pair}
In the setting of Definition~\ref{def:normalized_instance}, let $\eps > 0$ and let
$k \in \{1,\dots,m\}$ be an index with $\ip{v_1}{v_k} \le 1/\sqrt2$, which exists by
Lemma~\ref{lem:separated_pair}; by the same lemma $k \ne 1$, and the distance
$\varrho := \norm{v_1 - v_k}$ satisfies $\varrho^2 \ge 2 - \sqrt2 > 0$. Define
\[
  u := \frac{v_1 - v_k}{\varrho} \in \sphere^{d-1}, \qquad
  s := \frac{\ip{v_1}{u} + \ip{v_k}{u}}{2}, \qquad
  \theta_+ := \frac{4\eps}{\varrho}\, u, \qquad \theta_- := -\frac{4\eps}{\varrho}\, u .
\]
(The letter $\varrho$ is used for this distance to avoid a clash with the recommendation
kernel $\rho$.) Since $\ip{v_1 - v_k}{u} = \norm{v_1-v_k}^2/\varrho = \varrho$, we have
$\ip{v_1}{u} = s + \varrho/2$ and $\ip{v_k}{u} = s - \varrho/2$.
\end{definition}

\begin{lemma}[Two-point instances and the induced test]\label{lem:two_point_test}
\uses{def:lower_adaptive_two_point_instance, def:regret, def:pac}
Let $\eps > 0$ and let $\varrho$, $u$, $s$, $\theta_\pm$ be as in
Definition~\ref{def:lower_adaptive_two_point_instance}. For every $\rec \in \Xs'$:
\begin{enumerate}
  \item[(i)] if $r(\rec,\theta_+) \le \eps$ then $\ip{\rec}{u} \ge s + \varrho/4$;
  \item[(ii)] if $r(\rec,\theta_-) \le \eps$ then $\ip{\rec}{u} \le s - \varrho/4$.
\end{enumerate}
Consequently, if an identification algorithm on $\Xs'$ with budget $n$ is
$(\eps,\delta)$-PAC, then the $\{0,1\}$-valued measurable function
$\hat H := \indic\{\ip{\rec}{u} \ge s\}$ of the recommendation satisfies
\[
  \Prob_{\theta_+}(\hat H = 0) \le \delta, \qquad \Prob_{\theta_-}(\hat H = 1) \le \delta .
\]
\end{lemma}

\begin{proof}
\uses{def:lower_adaptive_two_point_instance, def:regret, def:pac}
Recall from Definition~\ref{def:lower_adaptive_two_point_instance} that
$\ip{v_1}{u} = s + \varrho/2$ and $\ip{v_k}{u} = s - \varrho/2$.

(i) $r(\rec,\theta_+) \le \eps$ means $\ip{\rec}{\theta_+} \ge \max_{x\in\Xs'}\ip{x}{\theta_+} - \eps
\ge \ip{v_1}{\theta_+} - \eps$ (as $v_1 \in \Xs'$), i.e.\
$\frac{4\eps}{\varrho}\ip{\rec}{u} \ge \frac{4\eps}{\varrho}\big(s + \frac\varrho2\big) - \eps$.
Multiplying by $\varrho/(4\eps) > 0$: $\ip{\rec}{u} \ge s + \varrho/2 - \varrho/4 = s + \varrho/4$.

(ii) Similarly $\ip{\rec}{\theta_-} \ge \ip{v_k}{\theta_-} - \eps$ reads
$-\frac{4\eps}{\varrho}\ip{\rec}{u} \ge -\frac{4\eps}{\varrho}\big(s - \frac\varrho2\big) - \eps$, i.e.\
$\ip{\rec}{u} \le s - \varrho/2 + \varrho/4 = s - \varrho/4$.

For the last claim: by (i), $\{\hat H = 0\} = \{\ip{\rec}{u} < s\} \subseteq \{r(\rec,\theta_+) > \eps\}$,
whose $\Prob_{\theta_+}$-probability is at most $\delta$ by the PAC property; by (ii),
$\{\hat H = 1\} = \{\ip{\rec}{u} \ge s\} \subseteq \{r(\rec,\theta_-) > \eps\}$, of
$\Prob_{\theta_-}$-probability at most $\delta$.
\end{proof}

\textbf{Lean remark.} $\hat H$ is the composition of the recommendation kernel $\rho$ with the
measurable map $x \mapsto \indic\{\ip{x}{u}\ge s\}$ from $\Xs'$ to \texttt{Bool}; the pair
(sampling rule, this composed kernel) is a \emph{test} in the sense of
Definition~\ref{def:mixture_test} below. The events $\{\hat H = 0\}$, $\{\hat H = 1\}$ are
events of $(\hist_n,\rec)$, as required by Lemma~\ref{lem:kl_data_processing_recommendation}.
In the formalization the two implications (i), (ii) and the error bounds are established inside
the proof of Lemma~\ref{lem:baseline_lower} (\texttt{COLT83.baseline\_le\_budget\_of\_isPAC}),
in the original coordinates (the set $B := \{x : \ip{Mx}{u} \ge s\}$ of recommendations), and the
error probabilities are those of the canonical runs of Definition~\ref{def:bai_alg}
(\texttt{IdentAlg.IsFixedBudget.isRun\_fixedBudgetRunMeasure}).

\begin{lemma}[KL divergence between the two-point instances]\label{lem:kl_two_instances}
\lean{Learning.LinearBandit.IsRun.klDiv_map_out_le_of_sq_le}\leanok
\uses{def:normalized_instance, def:lower_adaptive_two_point_instance, def:bai_alg}
Let $\eps > 0$ and let $\varrho$, $\theta_\pm$ be as in
Definition~\ref{def:lower_adaptive_two_point_instance}. For every identification algorithm on
the normalized instance $\Xs' \subseteq \ball_d$ (Definition~\ref{def:normalized_instance})
with budget $n$ (any sampling rule, any recommendation kernel),
\[
  \KL\big(\Prob_{\theta_+} \,\big\|\, \Prob_{\theta_-}\big) \le \frac{32\, n\, \eps^2}{\varrho^2} .
\]
\end{lemma}

\begin{proof}
\leanok
\uses{lem:kl_data_processing_recommendation, cor:divergence_decomposition_gaussian, def:normalized_instance}
By Lemma~\ref{lem:kl_data_processing_recommendation}, the KL between the laws of
$(\hist_n,\rec)$ equals the KL between the laws of $\hist_n$ (the recommendation kernel is
the same under both instances). By Corollary~\ref{cor:divergence_decomposition_gaussian} with
$\theta_+ - \theta_- = \frac{8\eps}{\varrho}u$,
\[
  \KL\big(\Prob_{\theta_+}\big\|\Prob_{\theta_-}\big)
  = \frac12\sum_{t<n}\E_{\theta_+}\big[\ip{x_t}{\theta_+ - \theta_-}^2\big]
  = \frac{32\eps^2}{\varrho^2}\sum_{t<n}\E_{\theta_+}\big[\ip{x_t}{u}^2\big]
  \le \frac{32\, n\,\eps^2}{\varrho^2},
\]
using $\abs{\ip{x_t}{u}} \le \norm{x_t}\norm{u} \le 1$ since $x_t \in \Xs' \subseteq \ball_d$
(Definition~\ref{def:normalized_instance}) and $\norm{u} = 1$.
\end{proof}

\begin{lemma}[Baseline lower bound]\label{lem:baseline_lower}
\lean{COLT83.baseline_le_budget_of_isPAC}\leanok
\uses{def:normalized_instance, def:pac}
Let $\eps > 0$ and $\delta \in (0, 1/4)$. Every $(\eps,\delta)$-PAC identification algorithm
on $\Xs'$ (equivalently, by the Lean remark after Lemma~\ref{lem:transform_preserves_pac}, on the
spanning action set $\Xs$ itself, $d \ge 2$) with budget $n$ satisfies
\[
  n \;\ge\; \frac{2 - \sqrt2}{32\,\eps^2}\,\log\frac{1}{4\delta} .
\]
In particular this holds for $\delta < 1/16$, the range of Theorem~\ref{thm:lower_adaptive}.
\end{lemma}

\begin{proof}
\leanok
\uses{def:lower_adaptive_two_point_instance, lem:two_point_test, lem:kl_two_instances, lem:bretagnolle_huber, lem:separated_pair}
Fix $k$, $\varrho$, $u$, $s$ and $\theta_\pm$ as in
Definition~\ref{def:lower_adaptive_two_point_instance} (the index $k$ exists by
Lemma~\ref{lem:separated_pair}). Let $\hat H$ be the test of Lemma~\ref{lem:two_point_test} and $A := \{\hat H = 0\}$, an
event of $(\hist_n,\rec)$. The Bretagnolle--Huber inequality (Lemma~\ref{lem:bretagnolle_huber})
applied to $\mu = \Prob_{\theta_+}$, $\nu = \Prob_{\theta_-}$ gives
\[
  \Prob_{\theta_+}(\hat H = 0) + \Prob_{\theta_-}(\hat H = 1)
  \ge \frac12\exp\big(-\KL(\Prob_{\theta_+}\|\Prob_{\theta_-})\big).
\]
The left-hand side is at most $2\delta$ by Lemma~\ref{lem:two_point_test}, so
$\exp(-\KL) \le 4\delta$, i.e.\ $\KL(\Prob_{\theta_+}\|\Prob_{\theta_-}) \ge \log\frac{1}{4\delta}$,
which is positive since $\delta < 1/4$. With Lemma~\ref{lem:kl_two_instances},
$\frac{32n\eps^2}{\varrho^2} \ge \log\frac{1}{4\delta}$, i.e.\
$n \ge \frac{\varrho^2}{32\eps^2}\log\frac1{4\delta}$, and $\varrho^2 \ge 2 - \sqrt2$
(Lemma~\ref{lem:separated_pair}) together with $\log\frac{1}{4\delta} > 0$ gives the claim.
\end{proof}

\section{The mixture testing problem}
\label{sec:lower_adaptive_test}

\begin{definition}[Tests and the mixture testing problem]\label{def:mixture_test}
\lean{COLT83.mixtureParam, COLT83.designWeight, COLT83.histLaw, COLT83.mixtureHistLaw, COLT83.IsDesign.isProbabilityMeasure_mixtureHistLaw, COLT83.IsGOptimalDesign.inner_mixtureParam_self}\leanok
\uses{def:normalized_instance, def:bai_alg}
Let $\eps > 0$ and, for $1 \le j \le m$, $\theta^{(j)} := 3\eps\, v_j$. A \emph{test with budget
$N \ge 1$} on $\Xs'$ is a pair $\mathsf{Test} = (\mathrm{alg}, \kappa)$ of an LML algorithm
$\mathrm{alg}$ with actions in $\Xs'$ and observations in $\R$, and a Markov kernel $\kappa$
from histories of length $N$ to $\{0,1\}$; run against $\nu'_\theta$ it produces
$\hist_N$ and the \emph{decision} $\hat H \sim \kappa(\hist_N)$, and we write
$\Prob^{\mathrm{test}}_\theta$ for the law of $(\hist_N,\hat H)$ (a probability measure on the
product of the history space with $\{0,1\}$). A test is thus an identification algorithm in the
sense of Definition~\ref{def:bai_alg} whose recommendation type is $\{0,1\}$ instead of $\Xs'$,
and $\Prob^{\mathrm{test}}_\theta$ is its law in the sense of that definition; when the test is
obtained from an identification algorithm $(\mathrm{alg},\rho)$ with budget $N$ by composing
$\rho$ with a measurable map $f : \Xs' \to \{0,1\}$ (as in Lemma~\ref{lem:two_point_test}),
$\Prob^{\mathrm{test}}_\theta$ is the image of the law $\Prob_\theta$ of $(\hist_N,\rec)$ under
$(h, x) \mapsto (h, f(x))$. The two hypotheses are
\[
  H_0 :\ \theta = 0, \qquad H_1 :\ \theta = \theta^{(J)} \text{ with } J \sim \lambda_G ,
\]
with the corresponding laws
\[
  \Prob_0 := \Prob^{\mathrm{test}}_0 \quad\text{and}\quad
  \Prob_1 := \sum_{j=1}^m \lambda_j\,\Prob^{\mathrm{test}}_{\theta^{(j)}} ,
\]
the latter being the finite mixture (with the weights $\lambda_j$ of
Definition~\ref{def:normalized_instance}) of the laws of $(\hist_N,\hat H)$; it is a
probability measure. We write $\E_0$ for the expectation under $\Prob_0$. The \emph{errors} of
the test are $\Prob_0(\hat H = 1)$ and
$\Prob_1(\hat H = 0) = \sum_j\lambda_j\Prob^{\mathrm{test}}_{\theta^{(j)}}(\hat H = 0)$.
\end{definition}

\textbf{Lean remark.} A test is an identification algorithm whose recommendation type is
\texttt{Bool} instead of $\Xs'$; both are instances of ``LML algorithm plus a recommendation
kernel into a measurable type'', and the KL statements of Chapter~\ref{chap:pre_info}
(Lemma~\ref{lem:kl_data_processing_recommendation}) are stated for an arbitrary
recommendation type. The mixture $\Prob_1$ is the finite weighted sum of measures
$\sum_j \lambda_j \bullet \Prob^{\mathrm{test}}_{\theta^{(j)}}$ (\texttt{Finset.sum} of \texttt{Measure}s with
\texttt{ENNReal} weights); it is not the law of an algorithm-environment sequence for a
single $\theta$, which is why the KL to it is controlled by convexity
(Lemma~\ref{lem:kl_mixture_convex}) rather than by the divergence decomposition directly.
In the formalization the alternatives are written in the original coordinates
(\texttt{COLT83.mixtureParam}: $\theta^{(j)} = 3\eps M(Mx_j)$, so that
$\ip{x}{\theta^{(j)}} = 3\eps\ip{Mx}{v_j}$), the laws $\Prob^{\mathrm{test}}_\theta$ are the laws
\texttt{COLT83.histLaw} of the history of length $N$ under the canonical trajectory measure of
the test's sampling rule (the decision is a measurable function of this history, so nothing is
lost, cf.\ Lemma~\ref{lem:kl_data_processing_recommendation}), and the mixture is
\texttt{COLT83.mixtureHistLaw}.

\begin{lemma}[KL divergence to the mixture]\label{lem:mixture_kl}
\lean{COLT83.IsGOptimalDesign.sum_mul_inner_mixtureParam_sq, COLT83.IsGOptimalDesign.sum_mul_inner_mixtureParam_sq_le, COLT83.IsGOptimalDesign.klDiv_histLaw_zero_mixtureHistLaw_le}\leanok
\uses{def:mixture_test}
For every test with budget $N$ on $\Xs'$,
\[
  \KL(\Prob_0\|\Prob_1)
  \le \sum_{j=1}^m \lambda_j\, \KL\big(\Prob_0 \big\| \Prob^{\mathrm{test}}_{\theta^{(j)}}\big)
  = \frac{9\eps^2}{2}\sum_{t<N}\E_0\Big[ x_t^\top \Big(\sum_{j=1}^m\lambda_j v_jv_j^\top\Big) x_t \Big]
  = \frac{9\eps^2}{2d}\sum_{t<N}\E_0\big[\norm{x_t}^2\big]
  \le \frac{9 N \eps^2}{2d} .
\]
\end{lemma}

\begin{proof}
\leanok
\uses{lem:kl_mixture_convex, lem:kl_data_processing_recommendation, cor:divergence_decomposition_gaussian, def:normalized_instance}
The first inequality is the convexity of $\nu \mapsto \KL(\mu\|\nu)$ for the finite mixture
$\Prob_1 = \sum_j\lambda_j\Prob^{\mathrm{test}}_{\theta^{(j)}}$ (Lemma~\ref{lem:kl_mixture_convex}). For each
$j$, by Lemma~\ref{lem:kl_data_processing_recommendation} the KL between the laws of
$(\hist_N,\hat H)$ under $0$ and $\theta^{(j)}$ equals the KL between the laws of $\hist_N$,
and by Corollary~\ref{cor:divergence_decomposition_gaussian}
\[
  \KL\big(\Prob_0\big\|\Prob^{\mathrm{test}}_{\theta^{(j)}}\big)
  = \frac12\sum_{t<N}\E_0\big[\ip{x_t}{0 - \theta^{(j)}}^2\big]
  = \frac{9\eps^2}{2}\sum_{t<N}\E_0\big[\ip{x_t}{v_j}^2\big] .
\]
Note that the expectation is under $\Prob_0$ for every $j$. Multiplying by $\lambda_j$,
summing over $j$ and exchanging the finite sum with the expectation,
$\sum_j\lambda_j\ip{x_t}{v_j}^2 = x_t^\top\big(\sum_j\lambda_jv_jv_j^\top\big)x_t = \norm{x_t}^2/d$
by Definition~\ref{def:normalized_instance}, which gives the two equalities. Finally
$\norm{x_t} \le 1$ since $x_t \in \Xs' \subseteq \ball_d$.
\end{proof}

\begin{lemma}[Lower bound for the testing problem]\label{lem:test_lower}
\lean{COLT83.IsGOptimalDesign.exp_neg_le_of_mixture}\leanok
\uses{def:mixture_test}
Let $\delta \in (0,1/8)$. Every test with budget $N$ on $\Xs'$ such that
$\Prob_0(\hat H = 1) \le 2\delta$ and $\Prob_1(\hat H = 0) \le 2\delta$ satisfies
\[
  N \;\ge\; \frac{2d}{9\,\eps^2}\,\log\frac{1}{8\delta} .
\]
\end{lemma}

\begin{proof}
\leanok
\uses{lem:bretagnolle_huber, lem:mixture_kl}
Apply the Bretagnolle--Huber inequality (Lemma~\ref{lem:bretagnolle_huber}) to the
probability measures $\mu = \Prob_0$, $\nu = \Prob_1$ and the event $A := \{\hat H = 1\}$:
\[
  4\delta \ge \Prob_0(\hat H = 1) + \Prob_1(\hat H = 0) \ge \frac12\exp\big(-\KL(\Prob_0\|\Prob_1)\big),
\]
so $\KL(\Prob_0\|\Prob_1) \ge \log\frac{1}{8\delta} > 0$ (as $\delta < 1/8$). By
Lemma~\ref{lem:mixture_kl}, $\frac{9N\eps^2}{2d} \ge \log\frac{1}{8\delta}$.
\end{proof}

\textbf{Lean remark.} The formalized statement is the intermediate inequality
$\frac12\exp(-9N\eps^2/(2d)) \le \alpha + \beta$ for every measurable set $E$ of histories
with $\Prob_0(E) \le \alpha$ and $\Prob^{\mathrm{test}}_{\theta^{(j)}}(E^c) \le \beta$ for all $j$
(\texttt{COLT83.IsGOptimalDesign.exp\_neg\_le\_of\_mixture}); the passage to the bound on $N$
is done in the proof of Theorem~\ref{thm:lower_adaptive}.

\begin{lemma}[A test built from a PAC algorithm]\label{lem:test_from_alg}
\lean{Learning.Algorithm.thenRepeat, Learning.IdentAlg.testAlg, Learning.LinearBandit.phaseMean, Learning.IsAlgEnvSeq.isAlgEnvSeqUntil_of_thenRepeat, Learning.IsAlgEnvSeq.hasCondDistrib_action_succ_of_thenRepeat, Learning.IsAlgEnvSeq.action_add_ae_eq_of_thenRepeat, Learning.IdentAlg.IsPAC.le_measureReal_action_succ_of_testAlg, Learning.LinearBandit.phaseMean_ae_eq_of_thenRepeat, Learning.LinearBandit.measureReal_inner_add_le_phaseMean_le, Learning.LinearBandit.measureReal_phaseMean_le_inner_sub_le, Learning.LinearBandit.measureReal_lt_phaseMean_le_of_testAlg, Learning.LinearBandit.measureReal_phaseMean_le_le_of_testAlg}\leanok
\uses{def:mixture_test, def:pac, lem:composition, lem:prefix_transfer}
Let $\eps > 0$, $\delta \in (0,1)$ and let $\mathsf{Alg} = (\mathrm{alg},\rho)$ be an
$(\eps,\delta)$-PAC identification algorithm on $\Xs'$ with budget $n$. Let
\[
  n_{\mathrm{est}} := \Big\lceil \frac{8\log(2/\delta)}{\eps^2} \Big\rceil, \qquad N := n + n_{\mathrm{est}} .
\]
Define the test $\mathsf{Test}$ with budget $N$ on $\Xs'$ as follows: for $t < n$ play as
$\mathrm{alg}$; at time $n$ draw $x_n \sim \rho(\hist_n)$ (the recommendation of
$\mathsf{Alg}$); for $n < t < N$ play $x_t := x_n$; let
$\hat v := \frac{1}{n_{\mathrm{est}}}\sum_{s < n_{\mathrm{est}}} y_{n+s}$ and decide
$\hat H := \indic\{\hat v > \eps\}$. Then
\[
  \Prob_0(\hat H = 1) \le \delta, \qquad
  \Prob^{\mathrm{test}}_{\theta^{(j)}}(\hat H = 0) \le 2\delta \ \text{ for every } j ,
\]
and therefore $\Prob_1(\hat H = 0) \le 2\delta$ as well.
\end{lemma}

\begin{proof}
\leanok
\uses{lem:composition, lem:prefix_transfer, lem:noise_representation, lem:noise_azuma, lem:gauss_mean_concentration, def:normalized_instance, def:pac, def:bai_alg}
\emph{Construction.} $\mathsf{Test}$ is the sequential composition (Lemma~\ref{lem:composition})
of $\mathrm{alg}_1 := \mathrm{alg}$ with $T_1 := n$ and, for a history $h$ of length $n$, the
algorithm $\mathrm{alg}_2(h)$ whose initial action distribution is $\rho(h)$ and whose policy
at every later time is deterministic: repeat the first action of the phase. The map
$h \mapsto \mathrm{alg}_2(h)$ is measurable because $\rho$ is a kernel. The decision
$\hat H$ is a deterministic measurable function of $\hist_N$, so $\kappa$ is the
deterministic kernel, and $\Prob^{\mathrm{test}}_\theta$ is the law of $(\hist_N,\hat H)$ under
the composition run against $\nu'_\theta$. Fix $\theta \in \R^d$.

\emph{The second phase for a fixed first history.} Fix a history $h$ of length $n$ and write
$\Prob^{h}$ for the law of the algorithm-environment sequence
$(x_{n+s}, y_{n+s})_{s < n_{\mathrm{est}}}$ for $(\mathrm{alg}_2(h), \nu'_\theta)$. Under
$\Prob^{h}$, $x_n \sim \rho(h)$, $x_{n+s} = x_n$ for all $s$, and by
Lemma~\ref{lem:noise_representation} $y_{n+s} = \ip{x_n}{\theta} + \eta_{n+s}$ with
$(\eta_{n+s})_{s<n_{\mathrm{est}}}$ i.i.d.\ $\Normal(0,1)$ independent of $x_n$. Let
\[
  \bar\eta := \frac{1}{n_{\mathrm{est}}}\sum_{s < n_{\mathrm{est}}}\big(y_{n+s} - \ip{x_{n+s}}{\theta}\big),
  \qquad\text{so that}\qquad \hat v = \ip{x_n}{\theta} + \bar\eta ;
\]
$\bar\eta$ is a measurable function of the second-phase trajectory alone, and
Lemma~\ref{lem:gauss_mean_concentration} gives
$\Prob^{h}(\abs{\bar\eta} \ge \eps/2) \le 2\exp(-n_{\mathrm{est}}\eps^2/8) \le \delta$, because
$n_{\mathrm{est}} \ge 8\log(2/\delta)/\eps^2$.

\emph{Transfer to the composition.} Let $E := \{\abs{\bar\eta} < \eps/2\}$, a measurable set of
second-phase trajectories, and $E_h := E$ for every $h$. The family $(E_h)_h$ does not depend on
$h$, so it is jointly measurable in $(h,\cdot)$, and $\Prob^{h}(E_h) \ge 1 - \delta$ for
\emph{every} $h$. The last assertion of Lemma~\ref{lem:composition}, applied with this family
and $\delta' = 0$, gives $\Prob^{\mathrm{test}}_\theta(\abs{\bar\eta} < \eps/2) \ge 1 - \delta$.
Moreover, by the composition-product structure of Lemma~\ref{lem:composition} ($\hist_n$ has the
law of $\mathrm{alg}$ against $\nu'_\theta$, and given $\hist_n = h$ the first action of the
second phase has law $\rho(h)$), the joint law of $(\hist_n, x_n)$ under the composition is
$(\text{law of } \hist_n)\otimes\rho$, which is exactly the law $\Prob_\theta$ of $(\hist_n,\rec)$
of $\mathsf{Alg}$ against $\nu'_\theta$ (Definition~\ref{def:bai_alg}); so every statement about
$(\hist_n,\rec)$ under $\mathsf{Alg}$ holds for $(\hist_n, x_n)$ under $\mathsf{Test}$.

\emph{Under $H_0$.} With $\theta = 0$, $\hat v = \bar\eta$, so
$\{\hat H = 1\} = \{\bar\eta > \eps\} \subseteq \{\abs{\bar\eta} \ge \eps/2\}$ and
$\Prob_0(\hat H = 1) \le \delta$. (The PAC property is not needed here.)

\emph{Under $\theta^{(j)}$.} Since $v_j \in \Xs'$ and $\norm{v_j} = 1$,
$\max_{x\in\Xs'}\ip{x}{\theta^{(j)}} \ge \ip{v_j}{3\eps v_j} = 3\eps$. By the PAC property
of $\mathsf{Alg}$ transferred to $(\hist_n,x_n)$, with
$\Prob^{\mathrm{test}}_{\theta^{(j)}}$-probability at least $1-\delta$,
$\ip{x_n}{\theta^{(j)}} \ge \max_x\ip{x}{\theta^{(j)}} - \eps \ge 2\eps$. On the
intersection of this event with $\{\abs{\bar\eta} < \eps/2\}$, which has probability at least
$1 - 2\delta$ by the union bound, $\hat v \ge 2\eps - \eps/2 = 3\eps/2 > \eps$, i.e.\ $\hat H = 1$.
Hence $\Prob^{\mathrm{test}}_{\theta^{(j)}}(\hat H = 0) \le 2\delta$.

\emph{Under $H_1$.} $\Prob_1(\hat H = 0) = \sum_j\lambda_j\Prob^{\mathrm{test}}_{\theta^{(j)}}(\hat H = 0)
\le 2\delta\sum_j\lambda_j = 2\delta$.
\end{proof}

\textbf{Lean remark.} The formalization does not use the general composition lemma. The
test's sampling rule is the single LML algorithm \texttt{Learning.Algorithm.thenRepeat}
(\texttt{Learning.IdentAlg.testAlg} when $\rho$ is the recommendation rule of $\mathsf{Alg}$):
its policy is that of $\mathrm{alg}$ at times $t < n$, the kernel $\rho$ applied to the
history of the first $n$ rounds at time $t = n$, and the deterministic policy ``play the action
recorded at index $n$ of the history'' afterwards; the recommendation of $\mathsf{Alg}$ is thus
the action $x_n$. The three facts used about a run of this algorithm are: it is an
algorithm-environment sequence for $\mathrm{alg}$ until time $n-1$
(\texttt{IsAlgEnvSeq.isAlgEnvSeqUntil\_of\_thenRepeat}), $x_n$ has conditional law
$\rho(\hist_n)$ given $\hist_n$ (\texttt{hasCondDistrib\_action\_succ\_of\_thenRepeat}), and
$x_{n+s} = x_n$ almost surely (\texttt{action\_add\_ae\_eq\_of\_thenRepeat}). The PAC guarantee
of $\mathsf{Alg}$ is transferred to $x_n$ by Lemma~\ref{lem:prefix_transfer}
(\texttt{IsPAC.le\_measureReal\_action\_succ\_of\_testAlg}). The concentration of
$\bar\eta$ is Lemma~\ref{lem:noise_azuma} applied to the run of the test itself (so no
conditioning on $\hist_n$ is needed): $\hat v = \ip{x_n}{\theta} + \bar\eta$ almost surely
(\texttt{phaseMean\_ae\_eq\_of\_thenRepeat}, with $\hat v$ the empirical mean
\texttt{Learning.LinearBandit.phaseMean} of the observations of rounds $n, \dots, N-1$), and the
two error bounds are \texttt{measureReal\_lt\_phaseMean\_le\_of\_testAlg} ($\theta = 0$) and
\texttt{measureReal\_phaseMean\_le\_le\_of\_testAlg} (an action of value at least $3\eps$
exists), both for runs on any standard Borel probability space. The decision is the event
$\{\hat v > \eps\}$; no decision kernel is needed because the laws of
Definition~\ref{def:mixture_test} are laws of histories.

\section{The adaptive lower bound}
\label{sec:lower_adaptive_theorem}

\begin{theorem}[Lower bound for adaptive algorithms]\label{thm:lower_adaptive}
\lean{COLT83.le_budget_of_isPAC}\leanok
\uses{def:arm_set, def:pac}
Let $\Xs$ be spanning, $d \ge 2$, $\eps > 0$ and $\delta \in (0, 1/16)$. Every $(\eps,\delta)$-PAC identification
algorithm on $\Xs$ (adaptive or not, with any recommendation kernel) with budget $n$
satisfies
\[
  n \;\ge\; \frac{d\log(1/\delta)}{20000\,\eps^2} .
\]
More precisely the proof gives $n \ge c_{\mathrm{ad}}\, d\log(1/\delta)/\eps^2$ with
$c_{\mathrm{ad}} := \frac{1}{11556 + 5760\sqrt2} \ge \frac{1}{20000}$; the constant
$c_{\mathrm{ad}}$ is used in the rest of the blueprint.
\end{theorem}

\begin{proof}
\leanok
\uses{lem:transform_preserves_pac, def:normalized_instance, lem:baseline_lower, lem:test_from_alg, lem:test_lower}
\emph{Step 0: normalization.} Let $\mathsf{Alg}$ be $(\eps,\delta)$-PAC on $\Xs$ with budget
$n$. By Lemma~\ref{lem:transform_preserves_pac} with $M = d^{-1/2}A_G^{-1/2}$
(Definition~\ref{def:normalized_instance}, which uses that $\Xs$ is spanning), the transformed algorithm $\mathsf{Alg}'$ is
$(\eps,\delta)$-PAC on $\Xs'$ with the same budget $n$. We apply the lemmas of
Sections~\ref{sec:lower_adaptive_baseline} and~\ref{sec:lower_adaptive_test} to $\mathsf{Alg}'$.
(In the formalization these lemmas are stated in the original coordinates and applied to
$\mathsf{Alg}$ directly; see the Lean remark after Lemma~\ref{lem:transform_preserves_pac}.)

\emph{Step 1: three consequences of $\delta < 1/16$.} Write $L := \log(1/\delta) > \log 16 = 4\log 2$.
Then
\begin{enumerate}
  \item[(a)] $\log\frac{1}{4\delta} = L - 2\log 2 \ge L - L/2 = L/2$;
  \item[(b)] $\log\frac{2}{\delta} = L + \log 2 \le L + L/4 = \frac54 L$;
  \item[(c)] $\log\frac{1}{8\delta} = L - 3\log 2 \ge L - \frac34 L = L/4$.
\end{enumerate}

\emph{Step 2: the baseline.} Let $c_0 := (2-\sqrt2)/32$. Lemma~\ref{lem:baseline_lower} and (a)
give $n \ge c_0 L/(2\eps^2)$, i.e.\ $L/\eps^2 \le 2n/c_0$.

\emph{Step 3: the budget of the test.} By (b) and $n \ge 1$,
\[
  n_{\mathrm{est}} = \Big\lceil\frac{8\log(2/\delta)}{\eps^2}\Big\rceil
  \le \frac{8\log(2/\delta)}{\eps^2} + 1 \le \frac{10 L}{\eps^2} + 1
  \le \frac{20}{c_0}\, n + n ,
\]
so the test of Lemma~\ref{lem:test_from_alg} built from $\mathsf{Alg}'$ has budget
$N = n + n_{\mathrm{est}} \le (2 + 20/c_0)\, n$.

\emph{Step 4: the testing lower bound.} That test has errors $\Prob_0(\hat H = 1) \le \delta \le 2\delta$
and $\Prob_1(\hat H = 0) \le 2\delta$ (Lemma~\ref{lem:test_from_alg}), so Lemma~\ref{lem:test_lower}
(applicable since $\delta < 1/16 < 1/8$) and (c) give
\[
  N \ge \frac{2d}{9\eps^2}\log\frac{1}{8\delta} \ge \frac{2d}{9\eps^2}\cdot\frac{L}{4} = \frac{d L}{18\,\eps^2}.
\]

\emph{Step 5: the constant.} Combining Steps~3 and~4,
$n \ge \frac{N}{2 + 20/c_0} \ge \frac{dL}{18(2 + 20/c_0)\,\eps^2}$. Now
$20/c_0 = 640/(2-\sqrt2) = 320(2+\sqrt2) = 640 + 320\sqrt2$, so
$18(2 + 20/c_0) = 18(642 + 320\sqrt2) = 11556 + 5760\sqrt2$, which is $c_{\mathrm{ad}}^{-1}$.
Finally $\sqrt2 \le 1.4143$ gives $5760\sqrt2 \le 8146.4$ and
$11556 + 5760\sqrt2 \le 19702.4 \le 20000$, so $c_{\mathrm{ad}} \ge 1/20000$.
\end{proof}

\begin{remark}[On the constant and the hypotheses]\label{rem:lower_adaptive_constant}
The paper only states $n = \Omega(d\log(1/\delta)/\eps^2)$; the constant
$c_{\mathrm{ad}} \approx 5.1\cdot 10^{-5}$ obtained here is not optimized (the bound
$n_{\mathrm{est}} \le (20/c_0 + 1)n$ of Step~3 is very crude: for the regime of interest,
$n_{\mathrm{est}}$ is a small multiple of $\log(1/\delta)/\eps^2$ while $n$ is of order
$d\log(1/\delta)/\eps^2$). The hypothesis $d \ge 2$ is used only in Lemma~\ref{lem:separated_pair}
(through $1/\sqrt d \le 1/\sqrt2$), the spanning hypothesis only through
Theorem~\ref{thm:kiefer_wolfowitz} in Definition~\ref{def:normalized_instance}, and $\delta < 1/16$ only in Step~1 of the proof of
Theorem~\ref{thm:lower_adaptive}; the baseline lemma holds for $\delta < 1/4$ and the testing
lower bound for $\delta < 1/8$. No upper bound on $\eps$ is needed. The bound applies in
particular to fixed-design algorithms (Definition~\ref{def:fixed_design}), which are
identification algorithms, and it is the $d\log(1/\delta)/\eps^2$ part of the paper's
Theorem~3; the unit-ball case is sharpened in Corollary~\ref{cor:lower_ball_adaptive}.
\end{remark}

\chapter{Properties of the Gaussian width}
\label{chap:width}

This chapter studies the Gaussian width term $\gw(\Xs)$ of Definition~\ref{def:width}
(Section~2.3 of the paper and its Appendices on the Gaussian width properties and on the
multi-task computations). It proves the three parts of Proposition~4 of the paper as
separate statements (Propositions~\ref{prop:width_le_d}, \ref{prop:width_finite}
and~\ref{prop:width_lower_sqrt_dlogd}), the bounds for the unit ball, the hypercubes and the
$m$-sets, and Theorem~5 (Theorem~\ref{thm:width_separation}): a family of finite sets whose
width is polynomially smaller than both generic upper bounds. It also introduces the
\emph{regionalized} width (Definition~\ref{def:regionalized_width}) consumed by the adaptive
algorithm of Chapter~\ref{chap:log_gains}, and the \emph{intrinsic} width
$\gw_{\mathrm{int}}(\Xs)$ of a compact set that does not span $\R^d$
(Definition~\ref{def:width_subspace}). The Gaussian width $\gw(\Xs)$ of
Definition~\ref{def:width} is only defined for spanning action sets
($\Span(\Xs) = \R^d$), and so is the $G$-optimal design $A_G$; every statement below that
involves them carries the spanning hypothesis explicitly. The lemmas on the width for a
matrix (Lemmas~\ref{lem:width_matrix_wd}, \ref{lem:width_homog}, \ref{lem:width_mono}
and~\ref{lem:width_symm}) hold for an arbitrary nonempty compact index set and are used as
such for the regions of a partition.

The inputs are: the $G$-optimal design of the Kiefer--Wolfowitz theorem
(Theorem~\ref{thm:kiefer_wolfowitz}, Chapter~\ref{chap:design}); the expectation of the
maximum of Gaussians (Lemma~\ref{lem:gauss_max_upper}) and Gaussian norm moments
(Lemma~\ref{lem:gauss_norm_moments}) of Chapter~\ref{chap:pre_gaussian}; for the secondary
lower bound $\gw(\Xs) \ge \Omega(\sqrt{d \log d})$, the Sudakov--Fernique inequality of
Chapter~\ref{chap:pre_sudakov}; and for the widths of structured sets, the lower bounds of
Chapter~\ref{chap:log_gains} combined with the upper bound of Theorem~\ref{thm:upper}. All
constants are explicit. Compared with the paper, the lower bounds for structured sets are
stated as explicit inequalities (with admittedly large thresholds on $d$ inherited from the
constants of the lower bounds), and the multi-task computation is organized around the
Helmert matrices with every linear-algebra step isolated.

\section{Upper bounds}
\label{sec:width_upper}

\begin{proposition}[The width is at most $d$]\label{prop:width_le_d}
\uses{def:width, thm:kiefer_wolfowitz}
Let $\Xs \subset \R^d$ be a spanning action set. Then $\gw(\Xs) \le d$. More precisely, the
$G$-optimal design matrix $A_G$ of Theorem~\ref{thm:kiefer_wolfowitz} (which exists since
$\Xs$ is spanning) satisfies $\gwidth{\Xs}{A_G} \le d$.
\end{proposition}

\begin{proof}
\uses{cor:g_optimal_norm_bound, lem:gauss_norm_moments, lem:width_matrix_wd}
By Theorem~\ref{thm:kiefer_wolfowitz}, $A_G \in \designset(\Xs) \cap \PD$, so
$\gw(\Xs) \le \gwidth{\Xs}{A_G}$ by definition of the infimum. For every $g \in \R^d$ and
$x \in \Xs$, the Cauchy--Schwarz inequality and Corollary~\ref{cor:g_optimal_norm_bound}
give $\ip{A_G^{-1/2}x}{g} \le \norm{A_G^{-1/2}x}\,\norm{g} \le \sqrt d\,\norm{g}$, hence
$\sup_{x\in\Xs}\ip{A_G^{-1/2}x}{g} \le \sqrt d\,\norm{g}$. Taking expectations
(Lemma~\ref{lem:width_matrix_wd} for integrability) and using
$\E\norm{g} \le \sqrt{\E\norm{g}^2} = \sqrt d$ (Lemma~\ref{lem:gauss_norm_moments} with
$S = I_d$) yields $\gwidth{\Xs}{A_G} \le \sqrt d \cdot \sqrt d = d$.
\end{proof}

\begin{proposition}[Width of a finite set]\label{prop:width_finite}
\uses{def:width, thm:kiefer_wolfowitz}
Let $\Xs$ be a finite spanning action set with $\abs{\Xs} = m$. Then
$\gw(\Xs) \le \sqrt{2 d \log m}$; more precisely $\gwidth{\Xs}{A_G} \le \sqrt{2d\log m}$
for the $G$-optimal design matrix $A_G$ of Theorem~\ref{thm:kiefer_wolfowitz}.
\end{proposition}

\begin{proof}
\uses{thm:kiefer_wolfowitz, cor:g_optimal_norm_bound, lem:gauss_max_upper, lem:gauss_linear_image}
Since $\Xs$ spans $\R^d$, $m \ge d \ge 1$. If $m = 1$ then $d = 1$, $\Xs = \{x\}$ and
$\gwidth{\Xs}{A} = \E\ip{A^{-1/2}x}{g} = 0 = \sqrt{2d\log 1}$ for every $A \in \PD$.
Assume $m \ge 2$ and let $u_x := A_G^{-1/2}x$ for $x \in \Xs$. The family
$(Z_x)_{x\in\Xs}$, $Z_x := \ip{u_x}{g}$, is a centered Gaussian vector
(Lemma~\ref{lem:gauss_linear_image}) with $\Var Z_x = \norm{u_x}^2 \le d$
(Corollary~\ref{cor:g_optimal_norm_bound}). Lemma~\ref{lem:gauss_max_upper} with
$\sigma^2 = d$ gives $\E\max_{x\in\Xs}Z_x \le \sqrt{d}\sqrt{2\log m}$, i.e.\
$\gwidth{\Xs}{A_G} \le \sqrt{2d\log m}$, and $\gw(\Xs) \le \gwidth{\Xs}{A_G}$ as in
Proposition~\ref{prop:width_le_d}.
\end{proof}

\begin{definition}[Regionalized width]\label{def:regionalized_width}
\uses{def:width_matrix, lem:width_matrix_wd}
Let $\Xs$ be an action set (spanning or not), $A \in \PD$ and let
$\mathcal R = (R_1, \dots, R_K)$, $K \ge 1$, be a partition of
$\Xs$ into nonempty compact subsets ($R_i \cap R_j = \emptyset$ for $i \ne j$ and
$\bigcup_i R_i = \Xs$). The \emph{regionalized Gaussian width of $\Xs$ for $A$ and
$\mathcal R$} is
\[
  \gw(A; \Xs, \mathcal R) := \max_{1 \le i \le K} \gwidth{R_i}{A}
  = \max_{1 \le i \le K}
  \E\Big[\max_{x \in R_i}\ip{x}{A^{-1/2}g}\Big], \qquad g \sim \Normal(0, I_d),
\]
where $\gwidth{R_i}{A}$ is the width for the matrix $A$ of Definition~\ref{def:width_matrix}
with the nonempty compact index set $R_i$; it is well defined and nonnegative by
Lemma~\ref{lem:width_matrix_wd} applied to the index set $R_i$. Since $R_i \subseteq \Xs$ for
every $i$, the supremum over $R_i$ is at most the supremum over $\Xs$ pathwise, so
$0 \le \gw(A;\Xs,\mathcal R) \le \gwidth{\Xs}{A}$, with equality when $K = 1$.
\end{definition}

\textbf{Lean remark.} A partition of a finite set is most conveniently encoded as a
function $\mathrm{region} : \Xs \to \mathrm{Fin}\ K$ that is surjective; the pieces are its
fibers $R_i = \mathrm{region}^{-1}(i)$ (as \texttt{Finset}s, \texttt{Finset.filter}), which
are automatically nonempty, disjoint and cover $\Xs$. For a general compact $\Xs$ the
pieces must be assumed closed (hence compact); only the finite case is used below. The inner
expectation is exactly the width for a matrix of Definition~\ref{def:width_matrix} with the
index set $R_i$, which that definition allows (any nonempty compact index set); hence
Lemmas~\ref{lem:width_matrix_wd}, \ref{lem:width_homog} and~\ref{lem:width_mono} apply to
$\gwidth{R_i}{A}$ verbatim, and Chapter~\ref{chap:log_gains} cites them directly for the
regions.

\begin{proposition}[Regionalized width of a finite set]\label{prop:width_partition}
\uses{def:regionalized_width, thm:kiefer_wolfowitz}
Let $\Xs$ be a finite spanning action set and $\mathcal R = (R_1,\dots,R_K)$ a partition of
$\Xs$ with $\abs{R_i} \le p$ for all $i$, where $p \ge 1$. Then, for the $G$-optimal design
matrix $A_G$ of Theorem~\ref{thm:kiefer_wolfowitz} (which exists since $\Xs$ is spanning),
\[
  \gw(A_G; \Xs, \mathcal R) \le \sqrt{2 d \log p}.
\]
\end{proposition}

\begin{proof}
\uses{cor:g_optimal_norm_bound, lem:gauss_max_upper, lem:gauss_linear_image}
Fix $i$. If $\abs{R_i} = 1$, the expectation $\E\max_{x\in R_i}\ip{A_G^{-1/2}x}{g}$ is the
expectation of a single centered Gaussian, i.e.\ $0 \le \sqrt{2d\log p}$. Otherwise
$2 \le \abs{R_i} \le p$ and, exactly as in the proof of Proposition~\ref{prop:width_finite},
the variables $Z_x = \ip{A_G^{-1/2}x}{g}$, $x \in R_i$, are centered Gaussians with variance
at most $d$ (Corollary~\ref{cor:g_optimal_norm_bound}), so Lemma~\ref{lem:gauss_max_upper}
gives $\E\max_{x\in R_i}Z_x \le \sqrt{2d\log\abs{R_i}} \le \sqrt{2d\log p}$. Take the
maximum over $i$.
\end{proof}

\section{Lower bounds}
\label{sec:width_lower}

\begin{proposition}[Width of the unit ball]\label{prop:width_ball}
\uses{def:width}
The unit ball $\Xs = \ball_d$ is a spanning action set (it is compact and contains the
standard basis $e_1,\dots,e_d$), so its Gaussian width $\gw(\ball_d)$ of
Definition~\ref{def:width} is defined, and
\[
  \frac{d}{\sqrt 3} \le \gw(\ball_d) \le d .
\]
\end{proposition}

\begin{proof}
\uses{prop:width_le_d, lem:designset_basic, lem:gauss_norm_moments, lem:trace_inv_amhm, lem:gauss_linear_image, lem:width_matrix_wd}
The upper bound is Proposition~\ref{prop:width_le_d} ($\ball_d$ being spanning). For the
lower bound, fix
$A \in \designset(\ball_d)\cap\PD$. By Lemma~\ref{lem:designset_basic},
$\tr A \le \sup_{x\in\ball_d}\norm{x}^2 = 1$. For every $v \in \R^d$,
$\sup_{x\in\ball_d}\ip{x}{v} = \norm{v}$ (Cauchy--Schwarz, with equality at
$x = v/\norm{v}$ when $v \ne 0$), so
$\gwidth{\ball_d}{A} = \E\norm{A^{-1/2}g}$. The vector $Y := A^{-1/2}g$ has law
$\Normal(0, A^{-1})$ (Lemma~\ref{lem:gauss_linear_image}), and
Lemma~\ref{lem:gauss_norm_moments} gives $\E\norm{Y} \ge \sqrt{\tr(A^{-1})}/\sqrt 3$.
Finally $\tr(A^{-1}) \ge d^2/\tr(A) \ge d^2$ by Lemma~\ref{lem:trace_inv_amhm}. Hence
$\gwidth{\ball_d}{A} \ge d/\sqrt3$ for every admissible $A$, and the infimum
$\gw(\ball_d)$ is at least $d/\sqrt 3$.
\end{proof}

\begin{lemma}[Well-separated points in the whitened set]\label{lem:whitened_separated_points}
\uses{def:design_matrix, lem:sqrt_props}
Let $d \ge 2$, let $\Xs$ be an action set and $\lambda \in \simplex_{\Xs}$ with
$\Sigma := A(\lambda) \succ 0$ (such a $\lambda$ exists if and only if $\Xs$ is spanning), and let
$\mathcal T := \Sigma^{-1/2}\Xs = \{\Sigma^{-1/2}x : x \in \Xs\}$ and $m := \lfloor d/2 \rfloor$.
There exist $t_1, \dots, t_m \in \mathcal T$ such that $\norm{t_i - t_j} \ge \sqrt{d/2}$
for all $i \ne j$.
\end{lemma}

\begin{proof}
\uses{lem:sqrt_props}
Write $u_x := \Sigma^{-1/2}x$ for $x \in \supp\lambda$. Then
$\sum_{x}\lambda_x u_xu_x^\top = \Sigma^{-1/2}A(\lambda)\Sigma^{-1/2} = I_d$
(Lemma~\ref{lem:sqrt_props}). We construct the points greedily. Pick $t_1 \in \mathcal T$
arbitrarily. Assume $t_1,\dots,t_k$ have been chosen for some $1 \le k < m$, let
$V_k := \Span\{t_1,\dots,t_k\}$ and let $Q_k$ be the orthogonal projection onto
$V_k^\perp$. Then $Q_k$ is symmetric and idempotent, and
$\tr Q_k = \dim V_k^\perp = d - \dim V_k \ge d - k$. Therefore
\[
  \sum_x \lambda_x \norm{Q_k u_x}^2 = \sum_x\lambda_x\, u_x^\top Q_k u_x
  = \tr\Big(Q_k \sum_x\lambda_x u_xu_x^\top\Big) = \tr Q_k \ge d - k ,
\]
so (the $\lambda_x$ being a probability vector) some $x \in \supp\lambda$ has
$\norm{Q_ku_x}^2 \ge d-k$; set $t_{k+1} := u_x \in \mathcal T$. For $i \le k$ we have
$Q_kt_i = 0$ since $t_i \in V_k$, and orthogonal projections do not increase norms, so
\[
  \norm{t_{k+1} - t_i} \ge \norm{Q_k(t_{k+1}-t_i)} = \norm{Q_kt_{k+1}} \ge \sqrt{d-k}
  \ge \sqrt{d - (m-1)} \ge \sqrt{d/2},
\]
where the last step uses $m - 1 \le d/2 - 1$. This constructs $t_1,\dots,t_m$ with the
required pairwise separation.
\end{proof}

\begin{proposition}[Lower bound $\sqrt{d\log d}$; secondary]\label{prop:width_lower_sqrt_dlogd}
\uses{def:width}
Let $\Xs \subset \R^d$ be a spanning action set with $d \ge 4$. Then
\[
  \gw(\Xs) \;\ge\; \frac18\sqrt{d\,\log\lfloor d/2\rfloor},
\]
and consequently, for $d \ge 6$,
\[
  \gw(\Xs) \;\ge\; \frac{1}{8\sqrt 2}\sqrt{d\log d} .
\]
\end{proposition}

\begin{proof}
\uses{lem:whitened_separated_points, thm:sudakov_fernique, lem:gauss_max_lower, lem:width_matrix_wd, lem:gauss_linear_image, lem:exists_posdef_design}
The constant $\frac18$ is $c_{\max}/2$ with $c_{\max} = \frac14$ the constant of
Lemma~\ref{lem:gauss_max_lower}, which holds for every $m \ge 2$ (no threshold); the
condition $d \ge 4$ ensures $m := \lfloor d/2\rfloor \ge 2$.
Fix $\lambda \in \simplex_{\Xs}$ with $\Sigma := A(\lambda) \succ 0$ (it exists since $\Xs$ is
spanning, Lemma~\ref{lem:exists_posdef_design}); it suffices to lower bound
$\gwidth{\Xs}{\Sigma} = \E\sup_{t\in\mathcal T}\ip{t}{g}$, $\mathcal T = \Sigma^{-1/2}\Xs$.
Take $t_1,\dots,t_m \in \mathcal T$ from
Lemma~\ref{lem:whitened_separated_points}. Set $Z_i := \ip{t_i}{g}$; $(Z_1,\dots,Z_m)$ is a
centered Gaussian vector (Lemma~\ref{lem:gauss_linear_image}) with
$\E(Z_i - Z_j)^2 = \norm{t_i-t_j}^2 \ge d/2$ for $i \ne j$. Let $\xi_1,\dots,\xi_m$ be
i.i.d.\ standard Gaussians and $Y_i := \frac{\sqrt d}{2}\xi_i$, so that
$\E(Y_i-Y_j)^2 = \frac{d}{4}\cdot 2 = d/2 \le \E(Z_i-Z_j)^2$ for $i\ne j$ (and $0 = 0$ for
$i = j$). The Sudakov--Fernique inequality (Theorem~\ref{thm:sudakov_fernique}, applied to the
pair $(Y, Z)$ whose increments are ordered as required) gives
$\E\max_i Y_i \le \E\max_i Z_i$. Hence
\[
  \gwidth{\Xs}{\Sigma} \ge \E\max_{i\le m}Z_i \ge \E\max_{i\le m}Y_i
  = \frac{\sqrt d}{2}\,\E\max_{i\le m}\xi_i \ge \frac{\sqrt d}{2}\cdot\frac14\sqrt{\log m}
  = \frac18\sqrt{d\log m},
\]
by Lemma~\ref{lem:gauss_max_lower} (applicable as $m \ge 2$). Taking the infimum over
$\lambda$ proves the first inequality. For the second, $\lfloor d/2\rfloor \ge \sqrt d$ for
$d \ge 6$ (check $d = 6, 7$ directly: $3 \ge \sqrt 7$; for $d\ge 8$,
$\lfloor d/2\rfloor \ge (d-1)/2 \ge \sqrt d$ because $(\sqrt d - 1)^2 \ge 2$), so
$\log\lfloor d/2\rfloor \ge \frac12\log d$ and
$\frac18\sqrt{d\log\lfloor d/2\rfloor} \ge \frac18\sqrt{\tfrac12 d\log d} = \frac{1}{8\sqrt2}\sqrt{d\log d}$.
\end{proof}

\begin{proposition}[Width lower bounds from sample complexity lower bounds]\label{prop:width_lower_from_pac}
\uses{def:pac, def:width}
Let $\Xs$ be a spanning action set, $\delta_0 \in (0,1)$ and $L \ge 0$. Assume that for every
$\eps \in (0,1]$, every $(\eps,\delta_0)$-PAC identification algorithm on $\Xs$ with budget
$T \ge 1$ satisfies $T \ge L/\eps^2$. Then
\[
  \gw(\Xs)^2 \;\ge\; \frac{L}{600} - d\log\frac{2}{\delta_0}.
\]
\end{proposition}

\begin{proof}
\uses{thm:upper, def:fixed_design_algorithm}
Fix $\eps \in (0,1]$. By Theorem~\ref{thm:upper} (which requires $\Xs$ spanning), the
fixed-design algorithm of
Definition~\ref{def:fixed_design_algorithm} with parameters $(\eps, \delta_0)$ is
$(\eps,\delta_0)$-PAC with budget
$T(\eps) = \lceil 600(\gw(\Xs)^2 + d\log(2/\delta_0))/\eps^2\rceil
< 600(\gw(\Xs)^2 + d\log(2/\delta_0))/\eps^2 + 1$, and $T(\eps) \ge 4d \ge 1$. The hypothesis gives
$L/\eps^2 \le T(\eps)$, hence, multiplying by $\eps^2$,
$L < 600(\gw(\Xs)^2 + d\log(2/\delta_0)) + \eps^2$. Letting $\eps \to 0$ gives
$L \le 600(\gw(\Xs)^2 + d\log(2/\delta_0))$, which is the claim.
\end{proof}

\begin{remark}
The outline of this chapter allowed an additional $-1$ on the right-hand side (coming from
the ceiling in the budget at $\eps = 1$); letting $\eps \to 0$ removes it, as shown above.
The hypothesis of Proposition~\ref{prop:width_lower_from_pac} is exactly what the lower
bounds of Section~\ref{sec:log_gains_lower} provide: they show that if
$1 \le T \le L/\eps^2$ then some instance fails with probability at least some $p_0 > \delta_0$,
so an $(\eps,\delta_0)$-PAC algorithm with budget $T \ge 1$ must have $T > L/\eps^2$. The
restriction to budgets $T \ge 1$ is harmless since the algorithm of Theorem~\ref{thm:upper}
has budget at least $4d$.
\end{remark}

\begin{corollary}[Widths of hypercubes and $m$-sets]\label{cor:width_structured}
\uses{def:width, def:msets}
The three action sets below are spanning (see the end of the proof), so their Gaussian
widths are defined.
\begin{enumerate}
  \item[(i)] For $\Xs = \{-1,1\}^d$: $\gw(\Xs)^2 \ge d^2/60000 - 3d$; in particular
        $\gw(\Xs) \ge d/350$ for $d \ge 360000$.
  \item[(ii)] For $\Xs = \{0,1\}^d$: $\gw(\Xs)^2 \ge d^2/240000 - 3d$; in particular
        $\gw(\Xs) \ge d/700$ for $d \ge 1500000$.
  \item[(iii)] For the $m$-sets $\Xs = \{x \in \{0,1\}^d : \sum_i x_i = m\}$
        (Definition~\ref{def:msets}) with $n := d - m + 1 \ge 20m$:
        $\gw(\Xs)^2 \ge mn/(1.5\cdot 10^6) - 3d$; in particular
        $\gw(\Xs) \ge \sqrt{md}/1800$ for $m \ge 10^7$.
\end{enumerate}
Together with Propositions~\ref{prop:width_le_d} and~\ref{prop:width_finite}, this shows
$\gw(\{-1,1\}^d) = \Theta(d)$, $\gw(\{0,1\}^d) = \Theta(d)$ and, for $m$-sets with
$d - m + 1 \ge 20m$, $\sqrt{md} \lesssim \gw(\Xs) \lesssim \sqrt{md\log(ed/m)}$
(using $\log\binom{d}{m} \le m\log(ed/m)$).
\end{corollary}

\begin{proof}
\uses{prop:width_lower_from_pac, thm:lower_hypercube_pm, thm:lower_hypercube_01, thm:lower_msets}
(i) Theorem~\ref{thm:lower_hypercube_pm}: for every $\eps > 0$, every identification
algorithm on $\{-1,1\}^d$ with budget $1 \le T \le d^2/(100\eps^2)$ has an instance on which it
fails with probability at least $1/6$. Take $\delta_0 := 1/7 < 1/6$: an
$(\eps,\delta_0)$-PAC algorithm with budget $T \ge 1$ must have $T > d^2/(100\eps^2)$, so
Proposition~\ref{prop:width_lower_from_pac} applies with $L = d^2/100$ and gives
$\gw(\Xs)^2 \ge d^2/60000 - d\log 14 \ge d^2/60000 - 3d$ ($\log 14 < 2.64$). If
$d \ge 360000$ then $3d \le d^2/60000 - d^2/122500$ (equivalent to
$d \ge 3/(1/60000 - 1/122500) \approx 352800$), so $\gw(\Xs)^2 \ge d^2/122500$ and
$\gw(\Xs) \ge d/350$.

(ii) Theorem~\ref{thm:lower_hypercube_01} gives the failure probability $1/6$ for
$T \le d^2/(400\eps^2)$; with $\delta_0 = 1/7$ and $L = d^2/400$,
$\gw(\Xs)^2 \ge d^2/240000 - 3d$, and for $d \ge 1500000$ this is at least
$d^2/490000$ (equivalent to $d \ge 3/(1/240000 - 1/490000) \approx 1411200$), so
$\gw(\Xs) \ge d/700$.

(iii) Theorem~\ref{thm:lower_msets} gives the failure probability $2/9$ for
$T \le mn/(2500\eps^2)$; with $\delta_0 = 1/5 < 2/9$ and $L = mn/2500$,
$\gw(\Xs)^2 \ge mn/(1.5\cdot10^6) - d\log 10 \ge mn/(1.5\cdot 10^6) - 3d$. Since
$m \le n/20$ we have $d = n + m - 1 \le 21n/20$, i.e.\ $n \ge 20d/21$, so
$mn/(1.5\cdot10^6) \ge md/(1.575\cdot10^6)$. For $m \ge 10^7$,
$3d \le md/(1.575\cdot10^6) - md/(3.15\cdot 10^6)$ (equivalent to $m \ge 9.45\cdot10^6$),
hence $\gw(\Xs)^2 \ge md/(3.15\cdot10^6)$ and $\gw(\Xs) \ge \sqrt{md}/1775 \ge \sqrt{md}/1800$.

The spanning hypotheses (needed by Proposition~\ref{prop:width_lower_from_pac}, i.e.\ by
Theorem~\ref{thm:upper}) hold: $\{-1,1\}^d$ and $\{0,1\}^d$ contain a basis, and the
$m$-sets span $\R^d$ for $1 \le m \le d-1$ (Definition~\ref{def:msets}), which is implied by
$n \ge 20m$.
\end{proof}

\begin{remark}
The thresholds on $d$ and $m$ in Corollary~\ref{cor:width_structured} come from the additive
term $d\log(2/\delta_0)$ in the budget of Theorem~\ref{thm:upper}, which must be dominated
by $L/600$; they are not intrinsic. The unit ball is treated directly in
Proposition~\ref{prop:width_ball}, which is much sharper than what
Corollary~\ref{cor:lower_ball_adaptive} would give through
Proposition~\ref{prop:width_lower_from_pac}.
\end{remark}

\section{Multi-task multi-armed bandits}
\label{sec:width_multitask}

The multi-task bandit action set below does not span $\R^d$, so the Gaussian width
$\gw(\Xs)$ of Definition~\ref{def:width} is not defined for it. As in the paper, one
measures it instead through an isometric identification of its span with a Euclidean space
of the right dimension; to avoid overloading the notation $\gw(\Xs)$, this second notion is
called the \emph{intrinsic width} and written $\gw_{\mathrm{int}}(\Xs)$.

\begin{definition}[Intrinsic width of a set spanning a subspace]\label{def:width_subspace}
\uses{def:width, lem:width_iso}
Let $\Xs \subset \R^d$ be an action set (nonempty compact, not necessarily spanning) and let
$V := \Span(\Xs)$ have dimension $r \ge 1$. Let
$U \in \R^{d\times r}$ be any matrix with orthonormal columns ($U^\top U = I_r$) whose range
is $V$. Then $\Xs_r := U^\top\Xs = \{U^\top x : x\in\Xs\} \subset \R^r$ is compact and spans
$\R^r$ (the map $z \mapsto U^\top z$ restricted to $V$ is a linear isometry onto $\R^r$ with
inverse $u \mapsto Uu$; in particular $x = UU^\top x$ for $x \in \Xs$), so it is a spanning
action set in $\R^r$ to which Definition~\ref{def:width} applies, and we define the
\emph{intrinsic width} of $\Xs$ as
\[
  \gw_{\mathrm{int}}(\Xs) := \gw(\Xs_r) = \gw(U^\top \Xs).
\]
This does not depend on the choice of $U$: if $U'$ is another such matrix then
$M := U^\top U' \in \R^{r\times r}$ is orthogonal ($M^\top M = U'^\top UU^\top U' = U'^\top U' = I_r$
since $UU^\top$ is the orthogonal projection onto $V$ and $U'$ has range $V$) and
$U'^\top x = M^\top U^\top x$ for $x \in V$, so $U'^\top\Xs = M^\top(U^\top\Xs)$ and
$\gw(U'^\top\Xs) = \gw(U^\top\Xs)$ by Lemma~\ref{lem:width_iso} (applied to the spanning set
$U^\top\Xs \subset \R^r$ and the invertible matrix $M^\top$). When $V = \R^d$ (i.e.\ $\Xs$ is
spanning) one has $\gw_{\mathrm{int}}(\Xs) = \gw(\Xs)$, again by Lemma~\ref{lem:width_iso},
$U$ being an orthogonal $d\times d$ matrix.
\end{definition}

\textbf{Lean remark.} The subspace is \texttt{Submodule.span ℝ X}; an orthonormal basis of a
finite-dimensional inner product subspace is provided by \texttt{stdOrthonormalBasis} or
\texttt{OrthonormalBasis} of the subspace; the map $x \mapsto U^\top x$ is the coordinate map
of this basis composed with the orthogonal projection \texttt{orthogonalProjection}, and
$u \mapsto Uu$ is the basis' linear isometry. The identification of action sets in $\R^r$
with \texttt{EuclideanSpace ℝ (Fin r)} then makes Definition~\ref{def:width} applicable.
In Lean, $\gw_{\mathrm{int}}$ is a separate declaration from $\gw$ (the latter takes the
spanning hypothesis as an argument), related by the identity
$\gw_{\mathrm{int}}(\Xs) = \gw(\Xs)$ for spanning $\Xs$.

\begin{definition}[Multi-task bandit action set]\label{def:multitask_set}
\uses{def:arm_set, def:width_subspace}
Let $m \ge 1$ and $d_1,\dots,d_m \ge 2$ be integers, $d := \sum_{j\le m} d_j$,
$d_{1:j} := \sum_{s\le j}d_s$ (with $d_{1:0} = 0$), and let
$B_j := \{d_{1:j-1}+1, \dots, d_{1:j}\} \subseteq [d]$ be the $j$-th \emph{block} of
coordinates, $\abs{B_j} = d_j$. The \emph{multi-task action set} is
\[
  \Xs := \Big\{x \in \{0,1\}^d : \sum_{i\in B_j}x_i = 1 \text{ for every } j \le m\Big\}.
\]
Thus $\abs{\Xs} = \prod_j d_j$; an element of $\Xs$ is determined by the positions
$(k_j)_{j\le m}$ of its ones, $k_j \in B_j$. For $m \ge 2$ the set $\Xs$ does not span $\R^d$
(Lemma~\ref{lem:multitask_basis}); it is nevertheless an action set in the sense of
Definition~\ref{def:arm_set}, and the relevant width is its intrinsic width
$\gw_{\mathrm{int}}(\Xs)$ of Definition~\ref{def:width_subspace}.
We write $\mathbf 1_n$ for the all-ones vector of $\R^n$ and $e_k$ for the $k$-th basis vector.
\end{definition}

\textbf{Lean remark.} The natural index type of the coordinates is the sigma type
$\Sigma_{j : \mathrm{Fin}\ m}\ \mathrm{Fin}\ (d_j)$ (block index, position in the block), with
$\Xs$ the image of $\prod_{j}\mathrm{Fin}\ (d_j)$ under
$(k_j)_j \mapsto \sum_j e_{(j,k_j)}$; the block $B_j$ is then the set of indices with first
component $j$. If one prefers $\mathrm{Fin}\ d$ as index type, use the block map
$b : \mathrm{Fin}\ d \to \mathrm{Fin}\ m$ and $B_j = b^{-1}(j)$; the offsets $d_{1:j}$ are
only needed to translate between the two encodings. All statements below only use the
partition $(B_j)_j$ of the coordinates and the sizes $d_j$.

\begin{definition}[Helmert matrices]\label{def:helmert}
For $n \ge 2$, the \emph{Helmert matrix} $H_n \in \R^{n\times(n-1)}$ has entries, for
$\ell \in [n]$ and $k \in [n-1]$,
\[
  (H_n)_{\ell,k} :=
  \begin{cases}
    \dfrac{1}{\sqrt{k(k+1)}} & \text{if } \ell \le k,\\[8pt]
    -\dfrac{k}{\sqrt{k(k+1)}} & \text{if } \ell = k+1,\\[8pt]
    0 & \text{if } \ell \ge k+2 .
  \end{cases}
\]
Its $k$-th column is $h_k := (1,\dots,1,-k,0,\dots,0)^\top/\sqrt{k(k+1)}$ with $k$ ones.
\end{definition}

\begin{lemma}[Properties of Helmert matrices]\label{lem:helmert_props}
\uses{def:helmert}
For $n\ge 2$: $H_n^\top H_n = I_{n-1}$, $H_n^\top\mathbf 1_n = 0$, and
$H_nH_n^\top = I_n - \frac1n\mathbf 1_n\mathbf 1_n^\top$. Equivalently, the columns of $H_n$
together with $\mathbf 1_n/\sqrt n$ form an orthonormal basis of $\R^n$, and $H_nH_n^\top$ is
the orthogonal projection onto $\mathbf 1_n^\perp = \{v : \sum_\ell v_\ell = 0\}$.
\end{lemma}

\begin{proof}
\uses{def:helmert}
$\norm{h_k}^2 = (k + k^2)/(k(k+1)) = 1$ and $\ip{h_k}{\mathbf 1_n} = (k - k)/\sqrt{k(k+1)} = 0$.
For $k < k'$, the entries of $h_{k'}$ on the first $k+1$ coordinates are all equal to
$1/\sqrt{k'(k'+1)}$ (because $k+1 \le k'$), and $h_k$ vanishes beyond coordinate $k+1$, so
$\ip{h_k}{h_{k'}} = \ip{h_k}{\mathbf 1_n}/\sqrt{k'(k'+1)} = 0$. Hence $H_n^\top H_n = I_{n-1}$,
$H_n^\top\mathbf 1_n = 0$, and $\{h_1,\dots,h_{n-1},\mathbf 1_n/\sqrt n\}$ is an orthonormal
family of $n$ vectors in $\R^n$, i.e.\ an orthonormal basis. The resolution of the identity
for an orthonormal basis, $\sum_k h_kh_k^\top + \frac1n\mathbf 1_n\mathbf 1_n^\top = I_n$,
is the last identity.
\end{proof}

\begin{lemma}[An orthonormal basis of the span of the multi-task set]\label{lem:multitask_basis}
\uses{def:multitask_set, def:helmert}
In the setting of Definition~\ref{def:multitask_set}, define $\mu \in \R^d$ by
$\mu_i := 1/d_j$ for $i \in B_j$, $\kappa := \norm{\mu}^2 = \sum_{j\le m}1/d_j$,
$u_0 := \mu/\sqrt\kappa$, and for $j \le m$ the matrix $Q_j \in \R^{d\times(d_j-1)}$ whose
rows indexed by $B_j$ form $H_{d_j}$ (row $d_{1:j-1}+\ell$ of $Q_j$ is row $\ell$ of
$H_{d_j}$) and whose other rows are zero. Let
\[
  U := [\,u_0\ \ Q_1\ \cdots\ Q_m\,] \in \R^{d\times r}, \qquad r := 1 + \sum_{j\le m}(d_j - 1) = d - m + 1 .
\]
Then $U^\top U = I_r$, $\mathrm{range}(U) = \Span(\Xs)$, and $\dim\Span(\Xs) = d - m + 1$.
\end{lemma}

\begin{proof}
\uses{lem:helmert_props, def:multitask_set}
\emph{Orthonormality.} $\norm{u_0} = 1$. The columns of $Q_j$ and $Q_\ell$ have disjoint
supports for $j \ne \ell$, so $Q_j^\top Q_\ell = 0$, and $Q_j^\top Q_j = H_{d_j}^\top H_{d_j} = I_{d_j-1}$
(Lemma~\ref{lem:helmert_props}). Since $u_0$ is constant on $B_j$,
$u_0^\top Q_j = \frac{1}{d_j\sqrt\kappa}\mathbf 1_{d_j}^\top H_{d_j} = 0$. Hence $U^\top U = I_r$.

\emph{Range.} Let $R_j := \{v \in \R^d : \supp v \subseteq B_j,\ \sum_{i\in B_j}v_i = 0\}$,
a subspace of dimension $d_j - 1$. The columns of $Q_j$ lie in $R_j$
(Lemma~\ref{lem:helmert_props}: $H_{d_j}^\top\mathbf 1 = 0$) and are $d_j - 1$ orthonormal
vectors, so $\mathrm{range}(Q_j) = R_j$, and
$\mathrm{range}(U) = \R u_0 \oplus R_1\oplus\dots\oplus R_m$ (a direct sum, the summands being
pairwise orthogonal). Now $\mu \in \Span\Xs$: $\mu = \frac{1}{\abs{\Xs}}\sum_{x\in\Xs}x$,
since for $i \in B_j$ exactly $\abs{\Xs}/d_j$ elements of $\Xs$ have $x_i = 1$. For
$i, i' \in B_j$, $e_i - e_{i'}$ is the difference of two elements of $\Xs$ that agree
outside $B_j$, hence lies in $\Span\Xs$; these differences span $R_j$. Therefore
$\mathrm{range}(U) \subseteq \Span\Xs$. Conversely, for $x \in \Xs$ and every $j$,
$\sum_{i\in B_j}(x - \mu)_i = 1 - 1 = 0$, so $x - \mu \in R_1\oplus\dots\oplus R_m$ (split
$x - \mu$ according to the blocks) and $x \in \mathrm{range}(U)$. Hence
$\mathrm{range}(U) = \Span\Xs$ and its dimension is the number $r$ of (orthonormal) columns
of $U$.
\end{proof}

\begin{lemma}[Covariance of the uniform design in the intrinsic coordinates]\label{lem:multitask_design_cov}
\uses{def:multitask_set, lem:multitask_basis, def:design_matrix}
Let $\lambda_u$ be the uniform design on $\Xs$ ($\lambda_x = 1/\abs\Xs$ for every
$x \in \Xs$) and $\Sigma := A(\lambda_u) = \frac{1}{\abs\Xs}\sum_{x\in\Xs}xx^\top \in \R^{d\times d}$.
Then, for $i \in B_j$ and $k \in B_\ell$,
\[
  \Sigma_{ik} = \frac{1}{d_j}\indic\{i = k\} \ \text{ if } j = \ell, \qquad
  \Sigma_{ik} = \frac{1}{d_jd_\ell} \ \text{ if } j \ne \ell ,
\]
and, with $U$, $\kappa$ as in Lemma~\ref{lem:multitask_basis},
\[
  \Sigma_r := U^\top\Sigma U = \operatorname{diag}\Big(\kappa,\ \tfrac{1}{d_1}I_{d_1-1},\ \dots,\ \tfrac{1}{d_m}I_{d_m-1}\Big)  \text{ (positive definite, } r\times r) .
\]
Moreover $\Sigma_r$ is the design matrix of the uniform design on $U^\top\Xs$:
$\Sigma_r = \frac{1}{\abs\Xs}\sum_{x\in\Xs}(U^\top x)(U^\top x)^\top \in \designset(U^\top\Xs)$.
\end{lemma}

\begin{proof}
\uses{lem:helmert_props, lem:multitask_basis}
$\Sigma_{ik}$ is the fraction of $x \in \Xs$ with $x_i = x_k = 1$. Within a block, two
distinct coordinates are never both $1$, and $x_i = 1$ for a fraction $1/d_j$ of the
elements; across blocks $j \ne \ell$, the positions of the ones in different blocks are
chosen independently and uniformly (the set $\Xs$ is in bijection with $\prod_j B_j$), so the
fraction is $1/(d_jd_\ell)$. For $v \in R_j$ (support in $B_j$, coordinates summing to $0$):
for $i \in B_j$, $(\Sigma v)_i = \frac{1}{d_j}v_i$, and for $i \in B_\ell$ with
$\ell\ne j$, $(\Sigma v)_i = \frac{1}{d_jd_\ell}\sum_{k\in B_j}v_k = 0$; thus
$\Sigma v = v/d_j$. For $\mu$: for $i \in B_j$,
$(\Sigma\mu)_i = \frac{1}{d_j}\cdot\frac{1}{d_j} + \sum_{\ell\ne j}\sum_{k\in B_\ell}\frac{1}{d_jd_\ell}\cdot\frac1{d_\ell}
= \frac{1}{d_j}\sum_{\ell\le m}\frac{1}{d_\ell} = \kappa\mu_i$, so $\Sigma\mu = \kappa\mu$ and
$\Sigma u_0 = \kappa u_0$. Consequently $u_0^\top\Sigma u_0 = \kappa$,
$u_0^\top\Sigma Q_j = \kappa u_0^\top Q_j = 0$, $Q_j^\top\Sigma Q_\ell = \frac{1}{d_\ell}Q_j^\top Q_\ell
= \frac{1}{d_\ell}\indic\{j=\ell\}I_{d_j-1}$ (Lemma~\ref{lem:multitask_basis}), which is the
block-diagonal form of $U^\top\Sigma U$; it is positive definite since its diagonal entries
are positive. The last claim is the identity $U^\top(\sum_x\lambda_x xx^\top)U =
\sum_x\lambda_x(U^\top x)(U^\top x)^\top$.
\end{proof}

\begin{lemma}[Reduction of the multi-task width to block maxima]\label{lem:multitask_width_reduction}
\uses{def:width_subspace, def:multitask_set, lem:multitask_basis, lem:multitask_design_cov}
In the setting of Definition~\ref{def:multitask_set}, let $g^{(j)} \sim \Normal(0, I_{d_j-1})$,
$j \le m$. Then the intrinsic width of Definition~\ref{def:width_subspace} satisfies
\[
  \gw_{\mathrm{int}}(\Xs) \;\le\; \sum_{j\le m}\sqrt{d_j}\ \E\Big[\max_{k \le d_j}\big(H_{d_j}g^{(j)}\big)_k\Big].
\]
\end{lemma}

\begin{proof}
\uses{def:width, lem:width_matrix_wd, lem:sqrt_props, lem:multitask_design_cov, lem:multitask_basis, lem:gauss_pi}
By Definition~\ref{def:width_subspace} with the matrix $U$ of Lemma~\ref{lem:multitask_basis},
$\gw_{\mathrm{int}}(\Xs) = \gw(U^\top\Xs)$, where $U^\top\Xs$ is a spanning action set in
$\R^r$, $r = d - m + 1$, so that Definition~\ref{def:width} applies to it; and by
Lemma~\ref{lem:multitask_design_cov},
$\Sigma_r \in \designset(U^\top\Xs)\cap\{\text{positive definite } r\times r \text{ matrices}\}$, so
\[
  \gw_{\mathrm{int}}(\Xs) \le \gwidth{U^\top\Xs}{\Sigma_r} = \E\Big[\max_{x\in\Xs}\ip{U^\top x}{\Sigma_r^{-1/2}g_r}\Big],
  \qquad g_r \sim \Normal(0, I_r).
\]
The matrix $\Sigma_r$ is diagonal with positive entries, so
$\Sigma_r^{-1/2} = \operatorname{diag}(\kappa^{-1/2}, \sqrt{d_1}I_{d_1-1},\dots,\sqrt{d_m}I_{d_m-1})$
(this diagonal matrix is positive definite and its square is $\Sigma_r^{-1}$; by uniqueness
of the positive square root, Lemma~\ref{lem:sqrt_props}, it is $\Sigma_r^{-1/2}$). Write
$g_r = (g_0, g^{(1)},\dots,g^{(m)})$ with $g_0 \sim \Normal(0,1)$ and
$g^{(j)}\sim\Normal(0,I_{d_j-1})$ independent (Lemma~\ref{lem:gauss_pi}). For $x \in \Xs$
with ones at positions $k_j \in B_j$, $U^\top x = (\ip{u_0}{x}, Q_1^\top x,\dots,Q_m^\top x)$
with $\ip{u_0}{x} = \ip{\mu}{x}/\sqrt\kappa = \kappa/\sqrt\kappa = \sqrt\kappa$ and
$Q_j^\top x = H_{d_j}^\top e_{k_j}$ ($e_{k_j} \in \R^{d_j}$ with the block-local index).
Hence
\[
  \ip{U^\top x}{\Sigma_r^{-1/2}g_r} = \kappa^{-1/2}\sqrt\kappa\, g_0
  + \sum_{j\le m}\sqrt{d_j}\ip{H_{d_j}^\top e_{k_j}}{g^{(j)}}
  = g_0 + \sum_{j\le m}\sqrt{d_j}\,\big(H_{d_j}g^{(j)}\big)_{k_j}.
\]
As $x$ ranges over $\Xs$, $(k_j)_j$ ranges over all of $\prod_j B_j$ independently across
blocks, so the maximum over $x$ is
$g_0 + \sum_j\sqrt{d_j}\max_{k\le d_j}(H_{d_j}g^{(j)})_k$. Taking expectations and using
$\E g_0 = 0$ gives the claim (each maximum is integrable by Lemma~\ref{lem:width_matrix_wd}).
\end{proof}

\begin{lemma}[Maximum of centered Gaussians]\label{lem:centered_gaussian_max}
\uses{def:helmert, lem:helmert_props}
Let $n \ge 2$, $Z \sim \Normal(0, I_n)$, $\bar Z := \frac1n\sum_{k\le n}Z_k$ and
$g \sim \Normal(0, I_{n-1})$. Then
\[
  \E\Big[\max_{k\le n}(H_ng)_k\Big] = \E\Big[\max_{k\le n}(Z_k - \bar Z)\Big]
  = \E\Big[\max_{k\le n}Z_k\Big] \le \sqrt{2\log n}.
\]
\end{lemma}

\begin{proof}
\uses{lem:gauss_linear_image, lem:gauss_law_of_cov, lem:gauss_max_upper, lem:helmert_props}
By Lemma~\ref{lem:gauss_linear_image}, $H_ng \sim \Normal(0, H_nH_n^\top) = \Normal(0, P)$
with $P := I_n - \frac1n\mathbf 1_n\mathbf 1_n^\top$ (Lemma~\ref{lem:helmert_props}), and
$Z - \bar Z\mathbf 1_n = PZ \sim \Normal(0, PP^\top) = \Normal(0,P)$ since $P$ is a symmetric
projection. Two centered Gaussian vectors with the same covariance have the same law
(Lemma~\ref{lem:gauss_law_of_cov}), so $\E\max_k(H_ng)_k = \E\max_k(Z_k - \bar Z)$. Next,
$\max_k(Z_k - \bar Z) = \max_kZ_k - \bar Z$ and $\E\bar Z = 0$. Finally
$\E\max_kZ_k \le \sqrt{2\log n}$ by Lemma~\ref{lem:gauss_max_upper} with $\sigma^2 = 1$.
\end{proof}

\begin{proposition}[Width of the multi-task set]\label{prop:width_multitask}
\uses{def:multitask_set, def:width_subspace}
In the setting of Definition~\ref{def:multitask_set}, the intrinsic width satisfies
\[
  \gw_{\mathrm{int}}(\Xs) \le \sum_{j\le m}\sqrt{2 d_j\log d_j}.
\]
Equivalently, for the matrix $U$ of Lemma~\ref{lem:multitask_basis}, the spanning action set
$U^\top\Xs \subset \R^{d-m+1}$ has Gaussian width $\gw(U^\top\Xs) \le \sum_{j\le m}\sqrt{2d_j\log d_j}$.
\end{proposition}

\begin{proof}
\uses{lem:multitask_width_reduction, lem:centered_gaussian_max}
Combine Lemma~\ref{lem:multitask_width_reduction} with Lemma~\ref{lem:centered_gaussian_max}
applied to each block ($n = d_j \ge 2$):
$\sqrt{d_j}\,\E\max_k(H_{d_j}g^{(j)})_k \le \sqrt{d_j}\sqrt{2\log d_j}$.
\end{proof}

\begin{theorem}[A finite set with a small width; Theorem~5 of the paper]\label{thm:width_separation}
\uses{def:multitask_set, def:width_subspace, prop:width_le_d, prop:width_finite}
Let $m \ge 2$ and consider the multi-task action set $\Xs$ with $m-1$ blocks of size
$d_j = 2$ ($j < m$) and one block of size $d_m = m^2$, so that $d = m^2 + 2m - 2$ and
$\abs{\Xs} = 2^{m-1}m^2$. Write $r := \dim\Span(\Xs)$ and let $U \in \R^{d\times r}$ be as in
Definition~\ref{def:width_subspace}, so that $U^\top\Xs \subset \R^r$ is a finite spanning
action set with $\abs{U^\top\Xs} = \abs\Xs$ and $\gw(U^\top\Xs) = \gw_{\mathrm{int}}(\Xs)$. Then:
\begin{enumerate}
  \item[(i)] $r = \dim\Span(\Xs) = m^2 + m - 1$, and $d/2 \le r \le d$;
  \item[(ii)] $\gw_{\mathrm{int}}(\Xs) \le 2(m-1)\sqrt{\log 2} + 2m\sqrt{\log m} \le 3\sqrt{d\log d}$;
  \item[(iii)] $\sqrt{d\log\abs{\Xs}} \ge 0.49\, d^{3/4}$, and also
        $\sqrt{r\log\abs\Xs} \ge 0.34\,d^{3/4}$ for the intrinsic dimension $r$.
\end{enumerate}
In particular the generic upper bounds $\gw(U^\top\Xs) \le r$
(Proposition~\ref{prop:width_le_d}) and $\gw(U^\top\Xs) \le \sqrt{2r\log\abs\Xs}$
(Proposition~\ref{prop:width_finite}), both applicable to the spanning set
$U^\top\Xs \subset \R^r$, are of order at least $d^{3/4}$, while the intrinsic width
$\gw_{\mathrm{int}}(\Xs) = \gw(U^\top\Xs)$ is $O(\sqrt{d\log d})$.
\end{theorem}

\begin{proof}
\uses{prop:width_multitask, lem:multitask_basis, prop:width_le_d, prop:width_finite}
(i) Lemma~\ref{lem:multitask_basis}: $\dim\Span\Xs = d - m + 1 = m^2 + m - 1$. It is at most
$d$, and $d - m + 1 \ge d/2$ is equivalent to $d \ge 2m - 2$, which holds.

(ii) Proposition~\ref{prop:width_multitask}:
$\gw_{\mathrm{int}}(\Xs) \le (m-1)\sqrt{2\cdot 2\log 2} + \sqrt{2m^2\log(m^2)} = 2(m-1)\sqrt{\log2} + 2m\sqrt{\log m}$.
Since $m^2 \le d \le 2m^2$ (the second inequality is $2m - 2 \le m^2$), we have
$m \le \sqrt d$ and $\log m \le \frac12\log d$; also $\log 2 \le \frac12\log d$ because
$d \ge 6 \ge 4$. Hence
$\gw_{\mathrm{int}}(\Xs) \le 2m(\sqrt{\log 2} + \sqrt{\log m}) \le 2\sqrt d\cdot 2\sqrt{\tfrac12\log d}
= 2\sqrt2\sqrt{d\log d} \le 3\sqrt{d\log d}$.

(iii) $\log\abs\Xs = (m-1)\log 2 + 2\log m \ge (m-1)\log2 \ge \frac{m}{2}\log 2$ for
$m\ge2$, and $m \ge \sqrt{d/2}$ from $d \le 2m^2$; so
$\log\abs\Xs \ge \frac{\log 2}{2\sqrt2}\sqrt d \ge 0.245\sqrt d$ and
$\sqrt{d\log\abs\Xs} \ge \sqrt{0.245}\,d^{3/4} \ge 0.49\,d^{3/4}$. With $r \ge d/2$,
$\sqrt{r\log\abs\Xs} \ge \sqrt{0.1225}\,d^{3/4} \ge 0.34\,d^{3/4}$.

The ``in particular'' clause: $U^\top\Xs$ is a finite spanning action set in $\R^r$ with
$\abs{U^\top\Xs} = \abs\Xs$ (the map $x \mapsto U^\top x$ is injective on $\Span\Xs$), so
Propositions~\ref{prop:width_le_d} and~\ref{prop:width_finite} apply to it and give
$\gw(U^\top\Xs) \le r$ and $\gw(U^\top\Xs) \le \sqrt{2r\log\abs\Xs}$; the right-hand sides are
at least $d/2$ and $0.34\sqrt2\,d^{3/4}$ respectively by (i) and (iii), whereas
$\gw(U^\top\Xs) = \gw_{\mathrm{int}}(\Xs) \le 3\sqrt{d\log d}$ by (ii).
\end{proof}

\begin{remark}
The paper states Theorem~5 as: there is a finite $\Xs$ with $\dim\Span\Xs = \Theta(d)$,
$\gw(\Xs) \le O(\sqrt{d\log d})$ (the width being understood as the intrinsic width
$\gw_{\mathrm{int}}(\Xs)$, since $\Xs$ is not spanning) and
$\sqrt{d\log\abs\Xs} \ge \Omega(d^{3/4})$; the
explicit form above is what will be formalized. The intuition is that the width of the
multi-task set behaves like $\sum_j\sqrt{d_j\log d_j}$, whereas the cardinality bound only
sees $\log\abs\Xs = \sum_j\log d_j$ and the dimension bound only sees $\sum_j d_j$; blocks of
size $2$ are cheap for the width but expensive for the cardinality bound. Section~\ref{sec:log_gains_lower}
shows a matching lower bound of order $(\sum_j\sqrt{d_j})^2/\eps^2$ on the sample complexity
of adaptive algorithms for this set.
\end{remark}

\chapter{Structured sets with only logarithmic gains from adaptivity}
\label{chap:log_gains}

This chapter formalizes Section~2.4 of the paper. Its first part
(Section~\ref{sec:log_gains_algorithm}) makes the fixed-design algorithm of
Chapter~\ref{chap:upper} adaptive, following Appendix~C of the paper: the finite action set is
partitioned into regions, a fixed design identifies a candidate arm in each region, and
Median Elimination (Chapter~\ref{chap:pre_median}) selects among the candidates. The
resulting sample complexity, Theorem~6 of the paper (Theorem~\ref{thm:log_gains}), is
$O(d\log((\abs\Xs/d)_+/\delta)/\eps^2)$, a $\log d$ improvement over the finite-set bound of
Theorem~\ref{thm:upper} in some regimes. The second part
(Section~\ref{sec:log_gains_lower}) proves, for adaptive and randomized algorithms, the
lower bounds on the $\delta$-independent term of the sample complexity summarized in
Table~1 of the paper: multi-task bandits (Theorem~\ref{thm:lower_multitask}), the two
hypercubes (Theorems~\ref{thm:lower_hypercube_pm} and~\ref{thm:lower_hypercube_01}), the
$m$-sets (Theorem~\ref{thm:lower_msets}) and the unit ball
(Corollary~\ref{cor:lower_ball_adaptive}, from Chapter~\ref{chap:pre_bayes}).

Inputs: the $G$-optimal design and the rounding theorem of Chapter~\ref{chap:design}, the
least-squares estimator and the concentration of Chapters~\ref{chap:upper}
and~\ref{chap:pre_concentration}, Median Elimination (Theorem~\ref{thm:median_elimination},
Corollary~\ref{cor:median_elimination_linear}) and the sequential composition of
Lemma~\ref{lem:composition}; for the lower bounds, Pinsker's inequality and the divergence
decomposition of Chapter~\ref{chap:pre_info}. Deviations from the paper (decided in the
overview): the algorithm only uses the $G$-optimal design (not the width minimizer), the
lower bounds are proved directly for randomized algorithms without Yao's principle (the
uniform average over a finite family of instances is a finite sum of probabilities under
each instance), and the multi-task lower bound treats all block sizes $d_j \ge 2$ at once by
a single Pinsker--Cauchy--Schwarz argument instead of the paper's two cases $d_j \ge 100$ and
$d_j < 100$. All constants are explicit.

\section{An adaptive version of the fixed-design algorithm}
\label{sec:log_gains_algorithm}

Throughout this section $\Xs \subset \R^d$ is a \emph{finite spanning} action set
($\Span(\Xs) = \R^d$, hence $\abs\Xs \ge d$), $\eps \in (0,1]$ and $\delta \in (0,1)$. The
spanning hypothesis is what makes the $G$-optimal design of
Theorem~\ref{thm:kiefer_wolfowitz} available; it is repeated in each statement.

\begin{definition}[Region-based adaptive algorithm]\label{def:region_algorithm}
\uses{def:bai_alg, def:regionalized_width, thm:kiefer_wolfowitz, thm:rounding, def:least_squares, lem:approx_argmax_selector, def:median_elimination, cor:median_elimination_linear, lem:composition}
Let $\Xs$ be a finite spanning action set and let $\mathcal R = (R_1,\dots,R_K)$, $K \ge 1$,
be a partition of $\Xs$ into nonempty subsets. Let $\lambda_G \in \simplex_{\Xs}$ and
$A_G = A(\lambda_G) \succ 0$ be the $G$-optimal design of Theorem~\ref{thm:kiefer_wolfowitz}
(which exists since $\Xs$ is spanning), and let
$w_{\mathcal R} := \gw(A_G;\Xs,\mathcal R)$ (Definition~\ref{def:regionalized_width}). Set
\[
  T_1 := \Big\lceil \frac{640\,\big(w_{\mathcal R}^2 + d\log(4/\delta)\big)}{\eps^2}\Big\rceil ,
\]
which satisfies $T_1 \ge 640\log 4\cdot d \ge 4d$. The \emph{region algorithm} with
parameters $(\mathcal R, \eps, \delta)$ is the sequential composition
(Lemma~\ref{lem:composition}) of the following two phases.
\begin{enumerate}
  \item[(1)] \emph{Fixed design.} Let $x_0,\dots,x_{T_1-1} \in \supp\lambda_G$ be the design
        given by Theorem~\ref{thm:rounding} applied to $\lambda_G$ and $T_1$, so that
        $\Sigma_{T_1} := \sum_{t<T_1}x_tx_t^\top \succeq \frac{T_1}{4}A_G$. Play it, observe
        $y \in \R^{T_1}$, and form the least-squares estimator
        $\hat\theta := \Sigma_{T_1}^{-1}X^\top y$ (Definition~\ref{def:least_squares}). For each
        $i \le K$ let $x^{(i)} := s_i(\hat\theta)$, where $s_i : \R^d \to R_i$ is a measurable
        exact argmax selector for the finite set $R_i$ (Lemma~\ref{lem:approx_argmax_selector}
        with $\eta = 0$): $\ip{x^{(i)}}{\hat\theta} = \max_{x\in R_i}\ip{x}{\hat\theta}$.
  \item[(2)] \emph{Median Elimination.} Run Median Elimination
        (Definition~\ref{def:median_elimination}) with parameters $(\eps/2, \delta/2)$ on the
        $K$-armed bandit induced by the candidates $x^{(1)},\dots,x^{(K)}$
        (Corollary~\ref{cor:median_elimination_linear}: arm $a$ plays $x^{(a)}$); it uses
        exactly $N_{ME}(K,\eps/2,\delta/2)$ rounds and returns an index $\hat a$.
\end{enumerate}
The recommendation is $\rec := x^{(\hat a)}$. The budget is the deterministic number
$T := T_1 + N_{ME}(K, \eps/2, \delta/2)$.
\end{definition}

\textbf{Lean remark.} Phase~(2) is the algorithm $\mathrm{alg}_2(h)$ of
Lemma~\ref{lem:composition}: it depends on the phase-(1) history $h$ only through the finite
list of candidates $(x^{(i)})_{i\le K}$, which is a measurable function of $\hat\theta(h)$,
itself a linear function of $y$. The fixed design of phase~(1) is obtained from the existence
statement of Theorem~\ref{thm:rounding} by choice. The final recommendation kernel is
deterministic: it maps the full history to $x^{(\hat a)}$, where both the candidates and
$\hat a$ are functions of the history.

\begin{lemma}[Width of a region under a scalar Loewner bound]\label{lem:log_gains_region_width_bound}
\uses{def:width_matrix, lem:width_matrix_wd}
Let $R \subset \R^d$ be nonempty and compact (typically a region of a partition of $\Xs$),
$A, B \in \PD$ and $c > 0$ with $B \succeq cA$. Then, in the notation of
Definition~\ref{def:width_matrix} with the index set $R$ (for $g \sim \Normal(0,I_d)$,
$\gwidth{R}{A} = \E[\max_{x\in R}\ip{x}{A^{-1/2}g}]$, well defined by
Lemma~\ref{lem:width_matrix_wd}),
\[
  \gwidth{R}{B} \le \frac{1}{\sqrt c}\,\gwidth{R}{A},
  \qquad\text{i.e.}\qquad
  \E\Big[\max_{x\in R}\ip{x}{B^{-1/2}g}\Big] \le \frac{1}{\sqrt c}\,\E\Big[\max_{x\in R}\ip{x}{A^{-1/2}g}\Big].
\]
\end{lemma}

\begin{proof}
\uses{lem:width_mono, lem:width_homog}
Both lemmas invoked below are stated in Chapter~\ref{chap:model} for an arbitrary nonempty
compact index set, here $R$. Since $cA \in \PD$ and $cA \preceq B$, Lemma~\ref{lem:width_mono}
(monotonicity in the Loewner order, index set $R$) gives $\gwidth{R}{B} \le \gwidth{R}{cA}$,
and Lemma~\ref{lem:width_homog} (homogeneity, index set $R$) gives
$\gwidth{R}{cA} = c^{-1/2}\gwidth{R}{A}$.
\end{proof}

\begin{lemma}[The candidate of the optimal region is good]\label{lem:region_phase1}
\uses{def:region_algorithm}
Let $\Xs$ be a finite spanning action set and $\mathcal R$ a partition of $\Xs$ as in
Definition~\ref{def:region_algorithm}.
Fix $\theta \in \R^d$, let $x_\star \in \argmax_{x\in\Xs}\ip{x}{\theta}$ and let $i_*$ be
the index of the region containing $x_\star$. Under $\nu_\theta$, phase~(1) of the region
algorithm satisfies
\[
  \Prob_\theta\Big( \ip{x^{(i_*)}}{\theta} \ge \max_{x\in\Xs}\ip{x}{\theta} - \frac\eps2 \Big)
  \ge 1 - \frac\delta2 .
\]
\end{lemma}

\begin{proof}
\uses{cor:fixed_design_from_distribution, cor:g_optimal_norm_bound, lem:least_squares_law, lem:log_gains_region_width_bound, cor:sup_bound_symmetric, lem:width_matrix_wd}
Let $\Delta := \hat\theta - \theta$ and $D := \max_{x,x'\in R_{i_*}}\ip{x-x'}{\Delta}$.

\emph{Deterministic step.} Since $x^{(i_*)}$ maximizes $\ip{\cdot}{\hat\theta}$ over
$R_{i_*} \ni x_\star$,
\[
  \ip{x^{(i_*)}}{\theta} = \ip{x^{(i_*)}}{\hat\theta} - \ip{x^{(i_*)}}{\Delta}
  \ge \ip{x_\star}{\hat\theta} - \ip{x^{(i_*)}}{\Delta}
  = \ip{x_\star}{\theta} + \ip{x_\star - x^{(i_*)}}{\Delta} \ge \ip{x_\star}{\theta} - D .
\]
So it suffices to show $\Prob_\theta(D \le \eps/2) \ge 1 - \delta/2$.

\emph{Law of $\Delta$.} By Corollary~\ref{cor:fixed_design_from_distribution} (applied to
$\lambda_G$ and the design of phase~(1)), $\Sigma_{T_1} \succ 0$ and
$\norm{x}^2_{\Sigma_{T_1}^{-1}} \le \frac{4}{T_1}\norm{x}^2_{A_G^{-1}} \le \frac{4d}{T_1}$ for
all $x \in \Xs$ (Corollary~\ref{cor:g_optimal_norm_bound}). By
Lemma~\ref{lem:least_squares_law}, $\Delta \sim \Normal(0, \Sigma_{T_1}^{-1})$.

\emph{Expectation.} Since the expectation of the supremum of the process
$x \mapsto \ip{x}{\Delta}$ over the compact index set $R_{i_*}$ only depends on the law of
$\Delta$ (Lemma~\ref{lem:width_matrix_wd} with the index set $R_{i_*}$),
$\E\max_{x\in R_{i_*}}\ip{x}{\Delta} = \gwidth{R_{i_*}}{\Sigma_{T_1}}$, and
Lemma~\ref{lem:log_gains_region_width_bound} with $B = \Sigma_{T_1} \succeq \frac{T_1}{4}A_G$,
$c = T_1/4$, gives
$\gwidth{R_{i_*}}{\Sigma_{T_1}} \le \frac{2}{\sqrt{T_1}}\gwidth{R_{i_*}}{A_G}
\le \frac{2 w_{\mathcal R}}{\sqrt{T_1}}$, the last step by Definition~\ref{def:regionalized_width}
($w_{\mathcal R} = \max_i\gwidth{R_i}{A_G}$).

\emph{Concentration.} Corollary~\ref{cor:sup_bound_symmetric} for the centered Gaussian
process $x \mapsto \ip{x}{\Delta}$ on the finite index set $R_{i_*}$, with
$\sigma^2 := \max_{x\in R_{i_*}}\norm{x}^2_{\Sigma_{T_1}^{-1}} \le 4d/T_1$ and failure
probability $\delta/2$: $\E D = 2\E\max_{x\in R_{i_*}}\ip{x}{\Delta}$ and, with probability at
least $1 - \delta/2$,
\[
  D \le \E D + 2\sigma\sqrt{2\cG\log(2/\delta)}
  \le \frac{4 w_{\mathcal R}}{\sqrt{T_1}} + \frac{4\sqrt d}{\sqrt{T_1}}\sqrt{2\cG\log(2/\delta)}
  = \frac{4}{\sqrt{T_1}}\Big(w_{\mathcal R} + \sqrt{2\cG\,d\log(2/\delta)}\Big).
\]

\emph{Arithmetic.} Using $(a+b)^2 \le 2a^2 + 2b^2$, $2\cG = \pi^2/2 < 5$ and
$\log(2/\delta) \le \log(4/\delta)$,
\[
  \frac{16}{T_1}\Big(w_{\mathcal R} + \sqrt{2\cG\,d\log(2/\delta)}\Big)^2
  \le \frac{32\,(w_{\mathcal R}^2 + 5\,d\log(4/\delta))}{T_1}
  \le \frac{160\,(w_{\mathcal R}^2 + d\log(4/\delta))}{T_1} \le \frac{\eps^2}{4},
\]
by the choice $T_1 \ge 640(w_{\mathcal R}^2 + d\log(4/\delta))/\eps^2$. Hence $D \le \eps/2$
on the good event.
\end{proof}

\begin{theorem}[Theorem~6 of the paper]\label{thm:log_gains}
\uses{def:region_algorithm, def:pac}
Let $\Xs$ be a finite spanning action set, $\eps\in(0,1]$, $\delta\in(0,1)$.
\begin{enumerate}
  \item[(i)] For every partition $\mathcal R = (R_1,\dots,R_K)$ of $\Xs$, the region
        algorithm of Definition~\ref{def:region_algorithm} is $(\eps,\delta)$-PAC on $\Xs$,
        with budget $T_1 + N_{ME}(K,\eps/2,\delta/2)$.
  \item[(ii)] Let $m := \abs\Xs$, $K := d$, and let $\mathcal R$ be a partition of $\Xs$ into
        $d$ nonempty regions of size at most $\lceil m/d\rceil$ (it exists since $m \ge d$).
        Then the budget of the region algorithm satisfies
        \[
          T \le C_{\mathrm{reg}}\ \frac{d\,\big(\log((m/d)_+) + \log(4/\delta)\big)}{\eps^2},
          \qquad C_{\mathrm{reg}} := 1281 + 6\,C_{ME},
        \]
        where $(u)_+ := \max(1,u)$ and $C_{ME}$ is the constant of Lemma~\ref{lem:me_budget}.
        In particular, for $\delta \le 1/4$,
        $T \le 2C_{\mathrm{reg}}\, d\log\big((m/d)_+/\delta\big)/\eps^2$.
\end{enumerate}
\end{theorem}

\begin{proof}
\uses{lem:region_phase1, cor:median_elimination_linear, lem:composition, lem:me_budget, prop:width_partition}
(i) Fix $\theta$, $x_\star$ and $i_*$ as in Lemma~\ref{lem:region_phase1}. Let $E_1$ be the
event $\{\ip{x^{(i_*)}}{\theta} \ge \max_x\ip{x}{\theta} - \eps/2\}$, which depends on the
phase-(1) history only and has probability at least $1-\delta/2$
(Lemma~\ref{lem:region_phase1}). For every phase-(1) history $h$ (with candidates
$x^{(1)},\dots,x^{(K)}$ determined by $h$), Corollary~\ref{cor:median_elimination_linear}
gives, for the induced $K$-armed Gaussian bandit with means $\ip{x^{(a)}}{\theta}$, that the
output $\hat a$ satisfies $\ip{x^{(\hat a)}}{\theta} \ge \max_{a\le K}\ip{x^{(a)}}{\theta} - \eps/2$
with probability at least $1 - \delta/2$; call $E_2(h)$ this event of the phase-(2)
trajectory. By Lemma~\ref{lem:composition} (with $\delta' = \delta/2$ for the set of
histories $E_1$ and $\delta/2$ for the conditional bound), the composed algorithm satisfies
$\Prob_\theta(E_1 \cap E_2(\hist_{T_1})) \ge 1 - \delta$. On this event,
\[
  \ip{\rec}{\theta} = \ip{x^{(\hat a)}}{\theta} \ge \ip{x^{(i_*)}}{\theta} - \frac\eps2
  \ge \max_{x\in\Xs}\ip{x}{\theta} - \eps ,
\]
i.e.\ $r(\rec,\theta) \le \eps$. Since $\theta$ is arbitrary, the algorithm is
$(\eps,\delta)$-PAC.

(ii) \emph{Phase (1).} By Proposition~\ref{prop:width_partition} with $p = \lceil m/d\rceil$,
$w_{\mathcal R}^2 \le 2d\log\lceil m/d\rceil$. Now $\lceil u\rceil = 1$ for $u \le 1$ and
$\lceil u\rceil \le 2u$ for $u \ge 1$, so $\log\lceil m/d\rceil \le \log((m/d)_+) + \log 2$.
Hence, using $2\log 2 = \log 4 \le \log(4/\delta)$,
\[
  T_1 \le \frac{640\big(2d\log((m/d)_+) + 2d\log2 + d\log(4/\delta)\big)}{\eps^2} + 1
  \le \frac{1280\,d\big(\log((m/d)_+) + \log(4/\delta)\big)}{\eps^2} + 1 ,
\]
and $1 \le d\log(4/\delta)/\eps^2$ (as $\eps \le 1$, $\log 4 > 1$), so
$T_1 \le 1281\,d(\log((m/d)_+) + \log(4/\delta))/\eps^2$.

\emph{Phase (2).} Lemma~\ref{lem:me_budget} with $K = d$ and parameters $(\eps/2,\delta/2)$
(admissible: $\eps/2 \le 1$, $\delta/2 < 1$):
$N_{ME}(d,\eps/2,\delta/2) \le C_{ME}\,d\log(6/\delta)/(\eps/2)^2 = 4C_{ME}\,d\log(6/\delta)/\eps^2$.
Since $\log(6/\delta) = \log(4/\delta) + \log(3/2)$ and $\log(3/2) < 0.41 < 0.3\log 4 \le 0.3\log(4/\delta)$,
we get $\log(6/\delta) \le 1.3\log(4/\delta)$ and
$N_{ME}(d,\eps/2,\delta/2) \le 5.2\,C_{ME}\,d\log(4/\delta)/\eps^2 \le 6C_{ME}\,d\log(4/\delta)/\eps^2$.

\emph{Total.} $T = T_1 + N_{ME} \le (1281 + 6C_{ME})\,d(\log((m/d)_+) + \log(4/\delta))/\eps^2$.
For $\delta \le 1/4$, $\log(4/\delta) \le 2\log(1/\delta)$ (equivalent to $\delta \le 1/4$),
so $\log((m/d)_+) + \log(4/\delta) \le 2\log((m/d)_+/\delta)$.
\end{proof}

\begin{remark}[On the form of the bound]
The paper states the budget as $O(d\log((\abs\Xs/d)_+/\delta)/\eps^2)$. That form cannot
hold uniformly in $\delta \in (0,1)$ with a fixed constant (the right-hand side tends to
$d\log((\abs\Xs/d)_+)/\eps^2$, possibly $0$, as $\delta \to 1$, while any algorithm with the
guarantee needs a positive budget), which is why the statement carries the additive
$\log(4/\delta)$ and the clean form is given for $\delta \le 1/4$. When $\abs\Xs \le d$
(only possible with $\abs\Xs = d$, a basis), $\log((m/d)_+) = 0$ and the bound is
$O(d\log(1/\delta)/\eps^2)$, matching the adaptive lower bound
Theorem~\ref{thm:lower_adaptive}. The improvement over Theorem~\ref{thm:upper} (whose
finite-set form is $O(d\log(\abs\Xs/\delta)/\eps^2)$ by Proposition~\ref{prop:width_finite})
is the replacement of $\log\abs\Xs$ by $\log(\abs\Xs/d)$, e.g.\ a factor $\log d$ when
$\abs\Xs = d^{1+o(1)}$. For the multi-armed bandit ($\Xs = \{e_1,\dots,e_d\}$) this is the
classical $\log d$ gain of Median Elimination~\cite{even2002pac}.
\end{remark}

\section{Adaptive lower bounds for structured sets}
\label{sec:log_gains_lower}

All statements of this section concern an arbitrary identification algorithm
$\mathsf{Alg}$ with budget $T \ge 1$ on the action set at hand (Definition~\ref{def:bai_alg});
both the sampling rule and the recommendation may be randomized. The action sets of this
section are nonempty compact sets that need not be spanning (the multi-task set of
Definition~\ref{def:multitask_set} is not): Definitions~\ref{def:env}, \ref{def:bai_alg},
\ref{def:regret} and~\ref{def:pac} and the information-theoretic tools
(Corollary~\ref{cor:divergence_decomposition_gaussian},
Lemma~\ref{lem:kl_data_processing_recommendation}) are stated for arbitrary action sets in
the sense of Definition~\ref{def:arm_set}, so nothing below requires spanning. Lower bounds are obtained
by exhibiting a finite family of reward vectors $(\theta_I)_{I\in\mathcal I}$ and bounding the
\emph{uniform average} $\frac{1}{\abs{\mathcal I}}\sum_{I\in\mathcal I}\E_{\theta_I}[\cdot]$ of
expectations under the individual instances. No minimax (Yao) principle is needed: a bound
on the average implies the existence of an instance realizing it.

\textbf{Lean remark.} A family of instances is a function $\theta : \mathcal I \to \R^d$ on a
\texttt{Fintype} $\mathcal I$ (in most cases $\mathcal I = \Xs$ itself, or a type of
subsets), and the uniform average is
\texttt{(∑ I, E[θ I] f) / Fintype.card ℐ}. The conclusion ``some $I$ has
$\Prob_{\theta_I}(\cdot) \ge p$'' follows from \texttt{Finset.exists\_le\_of\_sum\_le}
(or \texttt{Finset.exists\_lt\_of\_sum\_lt}) applied to the average.

\begin{lemma}[Two-instance Pinsker inequality for identification algorithms]\label{lem:two_point_pinsker}
\uses{def:bai_alg, def:arm_set}
Let $\Xs$ be an arbitrary action set (nonempty compact, not necessarily spanning), let
$\mathsf{Alg}$ be an identification algorithm with budget $T \ge 1$ on $\Xs$, and
let $\theta, \theta' \in \R^d$. Define
\[
  \KL(\theta,\theta') := \frac12\sum_{t<T}\E_\theta\big[\ip{x_t}{\theta - \theta'}^2\big].
\]
Then for every measurable event $E$ of $(\hist_T, \rec)$,
\[
  \Prob_{\theta'}(E) \le \Prob_\theta(E) + \sqrt{\KL(\theta,\theta')/2}
  \qquad\text{and}\qquad
  \Prob_{\theta}(E) \le \Prob_{\theta'}(E) + \sqrt{\KL(\theta,\theta')/2}.
\]
\end{lemma}

\begin{proof}
\uses{lem:pinsker, cor:divergence_decomposition_gaussian, lem:kl_data_processing_recommendation}
By Corollary~\ref{cor:divergence_decomposition_gaussian} (stated for an arbitrary action
set: it only uses that $\sup_{x\in\Xs}\norm x < \infty$), the Kullback--Leibler divergence
between the laws of $\hist_T$ under $\nu_\theta$ and $\nu_{\theta'}$ equals
$\KL(\theta,\theta')$, and by Lemma~\ref{lem:kl_data_processing_recommendation} the same
holds for the laws $\Prob_\theta$, $\Prob_{\theta'}$ of $(\hist_T,\rec)$ (the recommendation
kernel is common to both). Pinsker's inequality (Lemma~\ref{lem:pinsker}) applied to
$\mu = \Prob_\theta$, $\nu = \Prob_{\theta'}$ and the set $E$ gives the second inequality;
applied to the complement $E^c$ it gives
$\Prob_\theta(E^c) - \Prob_{\theta'}(E^c) \le \sqrt{\KL/2}$, i.e.\ the first one.
\end{proof}

\begin{lemma}[Three-point inequality]\label{lem:log_gains_three_point}
\uses{def:bai_alg}
Let $\mathsf{Alg}$ be an identification algorithm with budget $T$, let
$\theta^0, \theta^+, \theta^- \in \R^d$ and $\kappa \ge 0$ with
$\KL(\theta^0,\theta^+) \le \kappa$ and $\KL(\theta^0,\theta^-) \le \kappa$ (notation of
Lemma~\ref{lem:two_point_pinsker}). Then for every measurable event $E$ of $(\hist_T,\rec)$,
\[
  \Prob_{\theta^+}(E) + \Prob_{\theta^-}(E^c) \ge 1 - 2\sqrt{\kappa/2} .
\]
\end{lemma}

\begin{proof}
\uses{lem:two_point_pinsker}
Lemma~\ref{lem:two_point_pinsker} with $(\theta,\theta') = (\theta^0,\theta^+)$ and the event
$E^c$ gives $\Prob_{\theta^+}(E^c) \le \Prob_{\theta^0}(E^c) + \sqrt{\kappa/2}$, i.e.\
$\Prob_{\theta^+}(E) \ge \Prob_{\theta^0}(E) - \sqrt{\kappa/2}$; with
$(\theta^0,\theta^-)$ and the event $E$ it gives
$\Prob_{\theta^-}(E^c) \ge \Prob_{\theta^0}(E^c) - \sqrt{\kappa/2}$. Sum and use
$\Prob_{\theta^0}(E) + \Prob_{\theta^0}(E^c) = 1$.
\end{proof}

\begin{lemma}[From an average regret bound to a failure probability]\label{lem:fail_prob_from_expected_value}
\uses{def:bai_alg, def:regret}
Let $\mathsf{Alg}$ be an identification algorithm on a compact $\Xs$, let $\eps_* > 0$ and let
$(\theta_I)_{I\in\mathcal I}$ be a finite family of reward vectors with
$Z(\theta_I) = \max_{x,y\in\Xs}\ip{x-y}{\theta_I} = 10\eps_*$ for every $I$. If
\[
  \frac{1}{\abs{\mathcal I}}\sum_{I\in\mathcal I}\E_{\theta_I}\big[r(\rec,\theta_I)\big] \ge b\,\eps_*
  \qquad\text{for some } b > 1,
\]
then there exists $I \in \mathcal I$ with
$\Prob_{\theta_I}\big(r(\rec,\theta_I) > \eps_*\big) \ge (b-1)/9$. Equivalently, writing
$G(\rec,\theta) := \ip{\rec}{\theta} - \min_{x\in\Xs}\ip{x}{\theta} = Z(\theta) - r(\rec,\theta)$
for the gap to the worst arm: if
$\frac{1}{\abs{\mathcal I}}\sum_I\E_{\theta_I}[G(\rec,\theta_I)] \le a\,\eps_*$ with $a < 9$,
then some $I$ has $\Prob_{\theta_I}(r(\rec,\theta_I) > \eps_*) \ge (9-a)/9$. In particular,
when $\min_x\ip{x}{\theta_I} = 0$ for all $I$, $G(\rec,\theta_I) = \ip{\rec}{\theta_I}$.
\end{lemma}

\begin{proof}
\uses{def:regret}
For each $I$, $0 \le r(\rec,\theta_I) \le Z(\theta_I) = 10\eps_*$ pathwise, so
$r \le \eps_*\indic\{r \le \eps_*\} + 10\eps_*\indic\{r > \eps_*\} \le \eps_* + 9\eps_*\indic\{r > \eps_*\}$
and $\E_{\theta_I}[r] \le \eps_* + 9\eps_*\,p_I$ with $p_I := \Prob_{\theta_I}(r > \eps_*)$.
Averaging over $I$: $b\eps_* \le \eps_* + 9\eps_*\,\frac{1}{\abs{\mathcal I}}\sum_Ip_I$, so the
average of the $p_I$ is at least $(b-1)/9$ and some $p_I$ is at least that. The second
formulation is the first with $b = 10 - a$, since $\E[G] = 10\eps_* - \E[r]$.
\end{proof}

\subsection{Multi-task bandits}
\label{sec:log_gains_lower_multitask}

\begin{definition}[Hard instances for the multi-task set]\label{def:multitask_instances}
\uses{def:multitask_set, def:regret}
Let $\Xs$ be the multi-task action set of Definition~\ref{def:multitask_set} (blocks
$B_1,\dots,B_m$ of sizes $d_j \ge 2$) and $\eps > 0$. Let
\[
  S_d := \sum_{s\le m}\sqrt{d_s}, \qquad \eps_j := \frac{\eps\sqrt{d_j}}{S_d}\ (j\le m),
  \qquad\text{so that}\qquad \sum_{j\le m}\eps_j = \eps .
\]
Let $\tilde{\Xs} := \{\tilde x \in \{0,1\}^d : \sum_{i\in B_j}\tilde x_i \le 1\ \forall j\} \supseteq \Xs$.
For $\tilde x \in \tilde{\Xs}$ the reward vector $\theta_{\tilde x} \in \R^d$ is
$\theta_{\tilde x}[i] := 10\,\eps_j\,\tilde x[i]$ for $i \in B_j$. For $j \le m$ let
$\Xs^{(j)} := \{x \in \tilde{\Xs} : \sum_{i\in B_\ell}x_i = 1 \text{ for } \ell\ne j,\ \sum_{i\in B_j}x_i = 0\}$
and, for $x \in \Xs^{(j)}$ and $i \in B_j$, $x^{+i} := x + e_i \in \Xs$. The map
$(x, i) \mapsto x^{+i}$ is a bijection from $\Xs^{(j)}\times B_j$ onto $\Xs$, and
$\abs{\Xs^{(j)}} = \prod_{s\ne j}d_s$.

For $x \in \Xs$ and $\tilde x \in \tilde{\Xs}$,
$\ip{x}{\theta_{\tilde x}} = \sum_{j}10\eps_j\sum_{i\in B_j}x_i\tilde x_i \in [0, 10\eps]$; for
$\tilde x \in \Xs$, $\max_{x\in\Xs}\ip{x}{\theta_{\tilde x}} = 10\eps$ (attained at
$x = \tilde x$) and $\min_{x\in\Xs}\ip{x}{\theta_{\tilde x}} = 0$ (choose in every block a
position different from that of $\tilde x$, possible since $d_j \ge 2$); hence
$Z(\theta_{\tilde x}) = 10\eps$.
\end{definition}

\textbf{Lean remark.} The multi-task set does not span $\R^d$. This is allowed: an action
set in the sense of Definition~\ref{def:arm_set} is any nonempty compact set, spanning is a
separate hypothesis, and Definitions~\ref{def:env}, \ref{def:bai_alg}, \ref{def:regret}
and~\ref{def:pac} as well as Lemma~\ref{lem:two_point_pinsker} apply to $\Xs$ as they stand,
so the lower bound below needs no transport. (Should one prefer to work on a spanning set,
Lemma~\ref{lem:log_gains_subspace_transport} identifies identification algorithms on $\Xs$
with those on $U^\top\Xs \subset \R^{d-m+1}$, $U$ as in Definition~\ref{def:width_subspace},
through $\theta \mapsto U^\top\theta$; this is what the non-adaptive upper bound of
Remark~\ref{rem:table} uses.)
In the family $(\theta_{\tilde x})$, the coordinate $i \in B_j$ of $\rec$ is written
$\rec_i \in \{0,1\}$, and $\rec$ has exactly one nonzero coordinate in each block.

\begin{lemma}[Average value of a block]\label{lem:multitask_block_average}
\uses{def:multitask_instances, def:bai_alg}
Let $\mathsf{Alg}$ be an identification algorithm with budget $T$ on the multi-task set
$\Xs$. Fix $j \le m$ and $x \in \Xs^{(j)}$. For $i \in B_j$ let
$p_i := \Prob_{\theta_x}(\rec_i = 1)$ and $N_i := \sum_{t<T}x_t[i]$. Then
\begin{enumerate}
  \item[(i)] $\sum_{i\in B_j}p_i = 1$ and $\sum_{i\in B_j}\E_{\theta_x}[N_i] = T$;
  \item[(ii)] $\KL(\theta_x, \theta_{x^{+i}}) = 50\,\eps_j^2\,\E_{\theta_x}[N_i]$ for every
        $i \in B_j$ (notation of Lemma~\ref{lem:two_point_pinsker});
  \item[(iii)] $\Prob_{\theta_{x^{+i}}}(\rec_i = 1) \le p_i + 5\eps_j\sqrt{\E_{\theta_x}[N_i]}$
        for every $i \in B_j$;
  \item[(iv)] the block-$j$ value, averaged over the $d_j$ instances $x^{+i}$, satisfies
        \[
          \frac{1}{d_j}\sum_{i\in B_j}\E_{\theta_{x^{+i}}}\Big[\sum_{i'\in B_j}\rec_{i'}\,\theta_{x^{+i}}[i']\Big]
          \le \frac{10\,\eps_j}{d_j}\Big(1 + 5\,\eps_j\sqrt{d_j\,T}\Big).
        \]
\end{enumerate}
\end{lemma}

\begin{proof}
\uses{lem:two_point_pinsker, def:multitask_instances}
(i) Both $\rec$ and every $x_t$ belong to $\Xs$, so they have exactly one coordinate equal to
$1$ in $B_j$: $\sum_{i\in B_j}\rec_i = 1$ pathwise, and $\sum_{i\in B_j}N_i = T$ pathwise;
take expectations.

(ii) $\theta_x - \theta_{x^{+i}} = -10\eps_j e_i$, so
$\ip{x_t}{\theta_x - \theta_{x^{+i}}}^2 = 100\eps_j^2x_t[i]^2 = 100\eps_j^2x_t[i]$ (as
$x_t[i]\in\{0,1\}$), and $\KL(\theta_x,\theta_{x^{+i}}) = \frac12\sum_t\E_{\theta_x}[100\eps_j^2x_t[i]]
= 50\eps_j^2\E_{\theta_x}[N_i]$.

(iii) Lemma~\ref{lem:two_point_pinsker} with $\theta = \theta_x$, $\theta' = \theta_{x^{+i}}$
and $E = \{\rec_i = 1\}$: $\Prob_{\theta_{x^{+i}}}(E) \le p_i + \sqrt{25\eps_j^2\E_{\theta_x}[N_i]}$.

(iv) In block $j$ the vector $\theta_{x^{+i}}$ has the single nonzero coordinate
$\theta_{x^{+i}}[i] = 10\eps_j$, so $\sum_{i'\in B_j}\rec_{i'}\theta_{x^{+i}}[i'] = 10\eps_j\rec_i$
and its expectation is $10\eps_j\Prob_{\theta_{x^{+i}}}(\rec_i = 1)$. Summing (iii) over
$i \in B_j$ and using (i) and the Cauchy--Schwarz inequality,
\[
  \sum_{i\in B_j}\Prob_{\theta_{x^{+i}}}(\rec_i = 1) \le 1 + 5\eps_j\sum_{i\in B_j}\sqrt{\E_{\theta_x}[N_i]}
  \le 1 + 5\eps_j\sqrt{d_j\sum_{i\in B_j}\E_{\theta_x}[N_i]} = 1 + 5\eps_j\sqrt{d_jT}.
\]
Multiply by $10\eps_j/d_j$.
\end{proof}

\begin{theorem}[Lower bound for multi-task bandits]\label{thm:lower_multitask}
\uses{def:multitask_instances, def:pac}
Let $\Xs$ be the multi-task action set, $\eps > 0$, and let $\mathsf{Alg}$ be an
identification algorithm on $\Xs$ with budget $T \ge 1$ satisfying
\[
  T \le \frac{S_d^2}{20000\,\eps^2}, \qquad S_d = \sum_{j\le m}\sqrt{d_j}.
\]
Then $\frac{1}{\abs\Xs}\sum_{\tilde x\in\Xs}\E_{\theta_{\tilde x}}\ip{\rec}{\theta_{\tilde x}} \le 5.36\,\eps$,
and there exists $\tilde x \in \Xs$ with
$\Prob_{\theta_{\tilde x}}\big(r(\rec,\theta_{\tilde x}) > \eps\big) \ge 0.4$. Consequently an
$(\eps,\delta)$-PAC algorithm on $\Xs$ with budget $T \ge 1$ and $\delta < 0.4$ must have
$T > S_d^2/(20000\,\eps^2)$.
\end{theorem}

\begin{proof}
\uses{lem:multitask_block_average, lem:fail_prob_from_expected_value, def:multitask_instances}
For $\tilde x \in \Xs$, $\ip{\rec}{\theta_{\tilde x}} = \sum_{j\le m}V_j(\tilde x)$ with
$V_j(\tilde x) := \sum_{i\in B_j}\rec_i\theta_{\tilde x}[i]$. Fix $j$. Using the bijection
$\Xs^{(j)}\times B_j \to \Xs$, $(x,i)\mapsto x^{+i}$, and $\abs\Xs = \abs{\Xs^{(j)}}\,d_j$,
\[
  \frac{1}{\abs\Xs}\sum_{\tilde x\in\Xs}\E_{\theta_{\tilde x}}[V_j(\tilde x)]
  = \frac{1}{\abs{\Xs^{(j)}}}\sum_{x\in\Xs^{(j)}}\frac{1}{d_j}\sum_{i\in B_j}\E_{\theta_{x^{+i}}}[V_j(x^{+i})]
  \le \frac{10\eps_j}{d_j}\Big(1 + 5\eps_j\sqrt{d_jT}\Big)
\]
by Lemma~\ref{lem:multitask_block_average}(iv). With the budget assumption,
\[
  \eps_j\sqrt{d_jT} = \frac{\eps\sqrt{d_j}}{S_d}\sqrt{d_j}\sqrt T
  \le \frac{\eps\,d_j}{S_d}\cdot\frac{S_d}{\sqrt{20000}\,\eps} = \frac{d_j}{\sqrt{20000}},
\]
so, using $d_j \ge 2$ and $50/\sqrt{20000} < 0.36$,
\[
  \frac{10\eps_j}{d_j}\Big(1 + 5\eps_j\sqrt{d_jT}\Big) \le \frac{10\eps_j}{d_j} + \frac{50\,\eps_j}{\sqrt{20000}}
  \le 5\eps_j + 0.36\,\eps_j = 5.36\,\eps_j .
\]
Summing over $j$ and using $\sum_j\eps_j = \eps$ gives the bound $5.36\,\eps$ on the
average value. Since $\min_x\ip{x}{\theta_{\tilde x}} = 0$ and $Z(\theta_{\tilde x}) = 10\eps$
for $\tilde x \in \Xs$ (Definition~\ref{def:multitask_instances}),
Lemma~\ref{lem:fail_prob_from_expected_value} (gap form, $\eps_* = \eps$, $a = 5.36$) gives
an instance with failure probability at least $(9 - 5.36)/9 = 0.4044 \ge 0.4$. The PAC
consequence is immediate: an $(\eps,\delta)$-PAC algorithm has failure probability at most
$\delta < 0.4$ on every instance.
\end{proof}

\begin{remark}
The paper proves this bound (with the average value $7\eps$ and failure probability $0.1$)
by splitting into $d_j \ge 100$ and $d_j < 100$ and using deterministic algorithms plus Yao's
principle. The argument above, a single Pinsker plus Cauchy--Schwarz step per block, handles
all $d_j \ge 2$ uniformly and directly for randomized algorithms; the equality
$\eps_j\sqrt{d_jT} = d_j/\sqrt{20000}$ at the maximal budget makes the dependence on $d_j$
cancel exactly. Since $d \le S_d^2 \le md$, the lower bound is at least $d/(20000\eps^2)$
and at most $md/(20000\eps^2)$; for $m = 1$ it is the classical $\Omega(K/\eps^2)$ bound for
$K = d_1$ arms.
\end{remark}

\subsection{Hypercubes}
\label{sec:log_gains_lower_hypercubes}

\begin{theorem}[Lower bound for the hypercube $\{-1,1\}^d$]\label{thm:lower_hypercube_pm}
\uses{def:bai_alg, def:pac, def:regret}
Let $\Xs = \{-1,1\}^d$ and $\eps > 0$. For $s \in \{-1,1\}^d$ let $\theta_s := \frac{5\eps}{d}s$.
Let $\mathsf{Alg}$ be an identification algorithm on $\Xs$ with budget $T \ge 1$ satisfying
$T \le d^2/(100\,\eps^2)$. Then
$\frac{1}{2^d}\sum_{s}\E_{\theta_s}[r(\rec,\theta_s)] \ge 2.5\,\eps$ and there exists
$s \in \{-1,1\}^d$ with $\Prob_{\theta_s}(r(\rec,\theta_s) > \eps) \ge 1/6$. Consequently an
$(\eps,\delta)$-PAC algorithm on $\{-1,1\}^d$ with budget $T \ge 1$ and $\delta < 1/6$ has
$T > d^2/(100\,\eps^2)$.
\end{theorem}

\begin{proof}
\uses{lem:log_gains_three_point, lem:fail_prob_from_expected_value, lem:two_point_pinsker}
\emph{Regret as a Hamming distance.} $\max_{x\in\Xs}\ip{x}{\theta_s} = 5\eps$ (at $x = s$),
$\min_x\ip{x}{\theta_s} = -5\eps$ (at $x = -s$), so $Z(\theta_s) = 10\eps$, and for
$\rec\in\{-1,1\}^d$,
\[
  r(\rec,\theta_s) = \frac{5\eps}{d}\sum_{j\le d}(1 - \rec_js_j) = \frac{10\eps}{d}\sum_{j\le d}\indic\{\rec_j \ne s_j\}.
\]

\emph{One coordinate.} Fix $j \le d$ and $s_{-j} \in \{-1,1\}^{d-1}$ (the coordinates other
than $j$); let $s^+, s^-$ be the two completions (with $s_j = +1$, resp.\ $-1$) and let
$\theta^0 \in \R^d$ be $\theta_{s^\pm}$ with coordinate $j$ set to $0$ (the same vector for
both). Then $\theta^0 - \theta_{s^\pm} = \mp\frac{5\eps}{d}e_j$ and, since $x_t[j]^2 = 1$,
$\KL(\theta^0,\theta_{s^\pm}) = \frac12\sum_{t<T}\E_{\theta^0}\big[\tfrac{25\eps^2}{d^2}x_t[j]^2\big]
= \frac{25\eps^2T}{2d^2} =: \kappa$. Lemma~\ref{lem:log_gains_three_point} with the event
$E := \{\rec_j = -1\}$ gives
\[
  \Prob_{\theta_{s^+}}(\rec_j \ne s^+_j) + \Prob_{\theta_{s^-}}(\rec_j \ne s^-_j)
  = \Prob_{\theta_{s^+}}(E) + \Prob_{\theta_{s^-}}(E^c) \ge 1 - 2\sqrt{\kappa/2}
  = 1 - \frac{5\eps\sqrt T}{d} \ge \frac12 ,
\]
using $T \le d^2/(100\eps^2)$. Averaging over the $2^{d-1}$ choices of $s_{-j}$:
$\frac{1}{2^d}\sum_{s}\Prob_{\theta_s}(\rec_j \ne s_j) \ge \frac14$.

\emph{Average regret.} Summing over $j$,
$\frac{1}{2^d}\sum_s\E_{\theta_s}[r(\rec,\theta_s)] = \frac{10\eps}{d}\sum_{j\le d}\frac{1}{2^d}\sum_s\Prob_{\theta_s}(\rec_j\ne s_j)
\ge \frac{10\eps}{d}\cdot d\cdot\frac14 = 2.5\,\eps$.
Lemma~\ref{lem:fail_prob_from_expected_value} with $\eps_* = \eps$, $b = 2.5$ gives an
instance with failure probability at least $1.5/9 = 1/6$.
\end{proof}

\begin{theorem}[Lower bound for the hypercube $\{0,1\}^d$]\label{thm:lower_hypercube_01}
\uses{def:bai_alg, def:pac, def:regret}
Let $\Xs = \{0,1\}^d$ and $\eps > 0$. For $s \in \{-1,1\}^d$ let $\theta_s := \frac{10\eps}{d}s$
and $\sigma(s) := \indic\{s = 1\} \in \{0,1\}^d$ (coordinatewise). Let $\mathsf{Alg}$ be an
identification algorithm on $\Xs$ with budget $T \ge 1$ satisfying $T \le d^2/(400\,\eps^2)$. Then
$\frac{1}{2^d}\sum_s\E_{\theta_s}[r(\rec,\theta_s)] \ge 2.5\,\eps$ and there exists $s$ with
$\Prob_{\theta_s}(r(\rec,\theta_s) > \eps) \ge 1/6$. Consequently an $(\eps,\delta)$-PAC
algorithm on $\{0,1\}^d$ with budget $T \ge 1$ and $\delta < 1/6$ has $T > d^2/(400\,\eps^2)$.
\end{theorem}

\begin{proof}
\uses{lem:log_gains_three_point, lem:fail_prob_from_expected_value, lem:two_point_pinsker}
\emph{Regret.} $\max_{x\in\{0,1\}^d}\ip{x}{\theta_s} = \frac{10\eps}{d}\abs{\{j : s_j = 1\}}$
(at $x = \sigma(s)$) and $\min_x\ip{x}{\theta_s} = -\frac{10\eps}{d}\abs{\{j : s_j = -1\}}$, so
$Z(\theta_s) = 10\eps$, and for $\rec \in \{0,1\}^d$,
\[
  r(\rec,\theta_s) = \frac{10\eps}{d}\sum_{j\le d}\big(\sigma(s)_j - \rec_js_j\big)
  = \frac{10\eps}{d}\sum_{j\le d}\indic\{\rec_j \ne \sigma(s)_j\},
\]
since $\sigma_j - \rec_js_j$ equals $1 - \rec_j$ when $s_j = 1$ and $\rec_j$ when $s_j = -1$.

\emph{One coordinate.} Fix $j$ and $s_{-j}$, with completions $s^\pm$ and the common
alternative $\theta^0$ (coordinate $j$ zeroed). Now
$\KL(\theta^0,\theta_{s^\pm}) = \frac12\sum_t\E_{\theta^0}\big[\tfrac{100\eps^2}{d^2}x_t[j]^2\big]
= \frac{50\eps^2}{d^2}\E_{\theta^0}[N_j] \le \frac{50\eps^2T}{d^2} =: \kappa$, where
$N_j := \sum_{t<T}x_t[j] \le T$ (as $x_t[j]\in\{0,1\}$).
Lemma~\ref{lem:log_gains_three_point} with $E := \{\rec_j = 0\}$ gives
\[
  \Prob_{\theta_{s^+}}(\rec_j \ne \sigma(s^+)_j) + \Prob_{\theta_{s^-}}(\rec_j \ne \sigma(s^-)_j)
  \ge 1 - 2\sqrt{\kappa/2} = 1 - \frac{10\eps\sqrt T}{d} \ge \frac12 ,
\]
using $T \le d^2/(400\eps^2)$. The rest is identical to the proof of
Theorem~\ref{thm:lower_hypercube_pm}: averaging over $s_{-j}$ and summing over $j$ gives the
average regret bound $2.5\eps$, and Lemma~\ref{lem:fail_prob_from_expected_value} with
$b = 2.5$ gives the failure probability $1/6$.
\end{proof}

\begin{remark}
The paper's Appendix on hypercubes only sketches these proofs (``analogous to the $d_j = 2$
case of the multi-task argument, with the expected gap to the worst arm instead of the
expected value''). The constant differs between the two hypercubes: for $\{-1,1\}^d$ the
instances are at distance $5\eps/d$ from the zeroed alternative in coordinate $j$, for
$\{0,1\}^d$ at distance $10\eps/d$ (the range of $\ip{x}{\theta}$ over the cube must be
$10\eps$ in both cases), so the divergence is four times larger and the admissible budget
four times smaller. Both give $H_2 \ge \Omega(d^2/\eps^2)$, matching the upper bound
$O((d\log(1/\delta) + d^2)/\eps^2)$ of Theorem~\ref{thm:upper} since $\gw \le d$
(Proposition~\ref{prop:width_le_d}).
\end{remark}

\subsection{\texorpdfstring{$m$}{m}-sets}
\label{sec:log_gains_lower_msets}

\begin{definition}[$m$-sets and their hard instances]\label{def:msets}
\uses{def:arm_set, def:regret}
Let $1 \le m \le d-1$. The \emph{$m$-set} action set is
$\Xs := \{x \in \{0,1\}^d : \sum_{i\le d}x_i = m\}$, with $\abs\Xs = \binom dm$; it spans
$\R^d$ (it contains $x_S$ for every $m$-subset $S$; $x_S - x_{S'}$ with
$S \triangle S' = \{i,i'\}$ gives $e_i - e_{i'}$, and $\sum_S x_S$ is a positive multiple of
$\mathbf 1$; together these span $\R^d$). We write $n := d - m + 1$ and, for $\eps > 0$,
$\Delta := 10\eps/m$. For every $S \subseteq [d]$ (of any size) let
$\theta^{(S)} := \Delta\,\indic_S \in \R^d$ ($\Delta$ on $S$, $0$ elsewhere). For a
recommendation $\rec \in \Xs$ let $\hat S := \supp\rec$, so $\abs{\hat S} = m$ and
$\ip{\rec}{\theta^{(S)}} = \Delta\abs{\hat S\cap S}$. If $\abs S = m$ and $d \ge 2m$, then
$\max_{x\in\Xs}\ip{x}{\theta^{(S)}} = m\Delta = 10\eps$, $\min_{x\in\Xs}\ip{x}{\theta^{(S)}} = 0$
and $Z(\theta^{(S)}) = 10\eps$.
\end{definition}

\begin{lemma}[Equivalent sampling of $(B, i)$]\label{lem:msets_equivalent_sampling}
\uses{def:msets}
For every function $F$ on pairs $(B, i)$ with $B \subseteq [d]$, $\abs B = m-1$, $i \notin B$,
\[
  \frac{1}{\binom dm}\sum_{\abs S = m}\frac1m\sum_{i\in S}F(S\setminus\{i\}, i)
  = \frac{1}{\binom{d}{m-1}}\sum_{\abs B = m-1}\frac1n\sum_{i\notin B}F(B, i).
\]
In words: drawing $S$ uniformly among $m$-subsets and then $i$ uniformly in $S$ gives the
same law for $(S\setminus\{i\}, i)$ as drawing $B$ uniformly among $(m-1)$-subsets and then
$i$ uniformly in $[d]\setminus B$.
\end{lemma}

\begin{proof}
\uses{def:msets}
The map $(S, i) \mapsto (S\setminus\{i\}, i)$ is a bijection from
$\{(S,i) : \abs S = m,\ i \in S\}$ onto $\{(B,i) : \abs B = m-1,\ i\notin B\}$ (inverse
$(B,i)\mapsto(B\cup\{i\},i)$), so both sides are sums of $F$ over the same set of pairs,
with the constant weights $1/(m\binom dm)$ and $1/(n\binom{d}{m-1})$ respectively. These
are equal: $n\binom{d}{m-1} = (d-m+1)\frac{d!}{(m-1)!\,(d-m+1)!} = \frac{d!}{(m-1)!\,(d-m)!} = m\binom dm$.
\end{proof}

\begin{lemma}[Good coordinates under an alternative instance]\label{lem:msets_good_set}
\uses{def:msets, def:bai_alg}
Let $\mathsf{Alg}$ be an identification algorithm with budget $T$ on the $m$-sets and let
$B \subseteq [d]$ with $\abs B = m-1$. For $i \in [d]$ let $N_i := \sum_{t<T}x_t[i]$ and
$A_i := \{i \in \hat S\}$. Then
$\sum_{i\notin B}\E_{\theta^{(B)}}[N_i] \le mT$ and $\sum_{i\notin B}\Prob_{\theta^{(B)}}(A_i) \le m$,
and the set
\[
  G_B := \Big\{i \in [d]\setminus B : \E_{\theta^{(B)}}[N_i] \le \frac{4mT}{n}
  \ \text{ and }\ \Prob_{\theta^{(B)}}(A_i) \le \frac{4m}{n}\Big\}
\]
satisfies $\abs{G_B} \ge n/2$.
\end{lemma}

\begin{proof}
\uses{def:msets}
Every $x_t$ and $\rec$ have exactly $m$ ones, so $\sum_{i\le d}N_i = mT$ and
$\sum_{i\le d}\indic_{A_i} = m$ pathwise; restrict to $i \notin B$ and take expectations.
Let $G^1 := \{i\notin B : \E N_i > 4mT/n\}$; then $\abs{G^1}\cdot 4mT/n < \sum_{i\in G^1}\E N_i \le mT$
unless $G^1 = \emptyset$, so $\abs{G^1} < n/4$. Likewise
$G^2 := \{i \notin B : \Prob(A_i) > 4m/n\}$ has $\abs{G^2} < n/4$. Since
$\abs{[d]\setminus B} = n$, $\abs{G_B} = n - \abs{G^1\cup G^2} > n/2$.
\end{proof}

\begin{lemma}[Pinsker step for $m$-sets]\label{lem:msets_pinsker_step}
\uses{def:msets, lem:msets_good_set}
In the setting of Lemma~\ref{lem:msets_good_set}, for every $i \in G_B$,
\[
  \KL(\theta^{(B)}, \theta^{(B\cup\{i\})}) \le \frac{2\Delta^2 mT}{n}
  \qquad\text{and}\qquad
  \Prob_{\theta^{(B\cup\{i\})}}(i \in \hat S) \le \frac{4m}{n} + \Delta\sqrt{\frac{mT}{n}} .
\]
\end{lemma}

\begin{proof}
\uses{lem:two_point_pinsker, lem:msets_good_set}
$\theta^{(B)} - \theta^{(B\cup\{i\})} = -\Delta e_i$, so
$\KL(\theta^{(B)},\theta^{(B\cup\{i\})}) = \frac12\sum_t\E_{\theta^{(B)}}[\Delta^2x_t[i]^2]
= \frac{\Delta^2}{2}\E_{\theta^{(B)}}[N_i] \le \frac{\Delta^2}{2}\cdot\frac{4mT}{n}$ by
Lemma~\ref{lem:msets_good_set}. Lemma~\ref{lem:two_point_pinsker} with $E = A_i$ gives
$\Prob_{\theta^{(B\cup\{i\})}}(A_i) \le \Prob_{\theta^{(B)}}(A_i) + \sqrt{\KL/2}
\le \frac{4m}{n} + \sqrt{\Delta^2mT/n}$.
\end{proof}

\begin{theorem}[Lower bound for $m$-sets]\label{thm:lower_msets}
\uses{def:msets, def:pac}
Let $\Xs$ be the $m$-sets with $n = d - m + 1 \ge 20m$, let $\eps > 0$ and let
$\mathsf{Alg}$ be an identification algorithm on $\Xs$ with budget $T \ge 1$ satisfying
\[
  T \le \frac{m\,n}{2500\,\eps^2}.
\]
Then $\frac{1}{\binom dm}\sum_{\abs S = m}\E_{\theta^{(S)}}\ip{\rec}{\theta^{(S)}} \le 7\eps$ and
there exists an $m$-subset $S$ with $\Prob_{\theta^{(S)}}(r(\rec,\theta^{(S)}) > \eps) \ge 2/9$.
Consequently an $(\eps,\delta)$-PAC algorithm on $\Xs$ with budget $T \ge 1$ and $\delta < 2/9$ has
$T > mn/(2500\,\eps^2)$.
\end{theorem}

\begin{proof}
\uses{lem:msets_good_set, lem:msets_pinsker_step, lem:msets_equivalent_sampling, lem:fail_prob_from_expected_value, def:msets}
Note that $n \ge 20m$ implies $d = n + m - 1 \ge 21m - 1 \ge 2m$, so
Definition~\ref{def:msets} gives $Z(\theta^{(S)}) = 10\eps$ and
$\min_x\ip{x}{\theta^{(S)}} = 0$ for every $m$-subset $S$.

\emph{Constants.} $4m/n \le 0.2$, and
$\Delta\sqrt{mT/n} \le \frac{10\eps}{m}\sqrt{\frac{m}{n}\cdot\frac{mn}{2500\eps^2}} = \frac{10\eps}{m}\cdot\frac{m}{50\eps} = 0.2$.
Hence, by Lemma~\ref{lem:msets_pinsker_step}, $\Prob_{\theta^{(B\cup\{i\})}}(i\in\hat S) \le 0.4$
for every $(m-1)$-subset $B$ and every $i \in G_B$.

\emph{Average over $i \notin B$.} For fixed $B$, bounding the terms with $i \in G_B$ by
$0.4$ and the others by $1$, and using $\abs{G_B} \ge n/2$ (Lemma~\ref{lem:msets_good_set}),
\[
  \frac1n\sum_{i\notin B}\Prob_{\theta^{(B\cup\{i\})}}(i \in \hat S)
  \le \frac1n\Big(0.4\,\abs{G_B} + (n - \abs{G_B})\Big) \le \frac1n\Big(0.4\cdot\frac n2 + \frac n2\Big) = 0.7 .
\]

\emph{Average value.} For an $m$-subset $S$,
$\E_{\theta^{(S)}}\ip{\rec}{\theta^{(S)}} = \Delta\,\E_{\theta^{(S)}}\abs{\hat S\cap S}
= \Delta\sum_{i\in S}\Prob_{\theta^{(S)}}(i\in\hat S) = 10\eps\cdot\frac1m\sum_{i\in S}\Prob_{\theta^{(S)}}(i\in\hat S)$.
Averaging over $S$ and applying Lemma~\ref{lem:msets_equivalent_sampling} to
$F(B,i) := \Prob_{\theta^{(B\cup\{i\})}}(i\in\hat S)$ (note $\theta^{(S)} = \theta^{(B\cup\{i\})}$ for
$B = S\setminus\{i\}$),
\[
  \frac{1}{\binom dm}\sum_{\abs S = m}\E_{\theta^{(S)}}\ip{\rec}{\theta^{(S)}}
  = 10\eps\cdot\frac{1}{\binom{d}{m-1}}\sum_{\abs B = m-1}\frac1n\sum_{i\notin B}\Prob_{\theta^{(B\cup\{i\})}}(i\in\hat S)
  \le 10\eps\cdot 0.7 = 7\eps .
\]

\emph{Conclusion.} Lemma~\ref{lem:fail_prob_from_expected_value} (gap form, $a = 7$) gives
an $m$-subset $S$ with $\Prob_{\theta^{(S)}}(r > \eps) \ge 2/9$.
\end{proof}

\textbf{Lean remark.} The instances are indexed by \texttt{Finset (Fin d)} of card $m$
(\texttt{Finset.powersetCard m Finset.univ}) and the alternatives by subsets of card $m-1$;
$\hat S$ is the support of $\rec$ viewed as a \texttt{Finset}. Lemma~\ref{lem:msets_equivalent_sampling}
is a reindexing of a double sum along the bijection $(S,i)\mapsto(S\setminus\{i\},i)$
together with the identity $n\binom{d}{m-1} = m\binom dm$ (\texttt{Nat.choose\_succ\_right\_eq}
or \texttt{Nat.succ\_mul\_choose\_eq}). Lemma~\ref{lem:msets_good_set} is a counting
(Markov) argument on a finite sum.

\subsection{The unit ball}
\label{sec:log_gains_lower_ball}

\begin{corollary}[Adaptive lower bound for the unit ball]\label{cor:lower_ball_adaptive}
\uses{def:pac, cor:ball_pac_lower}
Let $d \ge 2$, $\eps > 0$ and $\delta \le 1/500$. Every $(\eps,\delta)$-PAC identification
algorithm on $\ball_d$ with budget $T$ satisfies $T \ge d^2/(1000\,\eps^2)$.
\end{corollary}

\begin{proof}
\uses{cor:ball_pac_lower}
This is Corollary~\ref{cor:ball_pac_lower} of Chapter~\ref{chap:pre_bayes}, which replaces
the citation of~\cite{chen2024assouad} in the paper's Appendix on the unit ball: the
Gaussian-prior argument (Theorem~\ref{thm:ball_simple_regret_lower}(ii)) gives an instance
$\theta$ with $\E_\theta[r(\rec,\theta)] \ge 3d/(40\sqrt T)$ and
$\norm\theta \le 10d/\sqrt T$, and the PAC property bounds
$\E_\theta[r] \le \eps + 2\norm\theta\delta \le \eps + 20\delta d/\sqrt T$; for
$\delta \le 1/500$ this gives $(3/40 - 1/25)\,d/\sqrt T \le \eps$, i.e.\
$T \ge (7/200)^2d^2/\eps^2 = 49d^2/(40000\eps^2) \ge d^2/(1000\eps^2)$.
\end{proof}

\section{Summary of the sample complexities}
\label{sec:log_gains_table}

The non-adaptive upper bound of Theorem~\ref{thm:upper} requires a spanning action set. For
the multi-task set, which is not spanning, it is applied to the spanning set
$U^\top\Xs \subset \R^{d-m+1}$ of Definition~\ref{def:width_subspace} and transported back
by the following lemma, the analogue for a subspace of the transformation lemma
\ref{lem:transform_preserves_pac} (which handles invertible maps of $\R^d$).

\begin{lemma}[Transport of identification algorithms along the span]\label{lem:log_gains_subspace_transport}
\uses{def:arm_set, def:bai_alg, def:pac, def:width_subspace}
Let $\Xs \subset \R^d$ be an action set with $V := \Span(\Xs)$ of dimension $r \ge 1$, let
$U \in \R^{d\times r}$ have orthonormal columns with range $V$, and let
$\Xs_r := U^\top\Xs \subset \R^r$, a spanning action set in $\R^r$
(Definition~\ref{def:width_subspace}). The map $\phi : \Xs_r \to \Xs$, $u \mapsto Uu$, is a
homeomorphism with inverse $x \mapsto U^\top x$, and for all $u \in \R^r$, $\theta \in \R^d$,
\[
  \ip{Uu}{\theta} = \ip{u}{U^\top\theta}, \qquad\text{in particular}\qquad
  \ip{U^\top x}{U^\top\theta} = \ip{x}{\theta} \ \text{ for } x \in \Xs,\ \theta \in V .
\]
For every identification algorithm $\mathsf{Alg}_r = (\mathrm{alg}_r,\rho_r)$ on $\Xs_r$ with
budget $T$ there is an identification algorithm $\mathsf{Alg} = (\mathrm{alg},\rho)$ on
$\Xs$ with budget $T$ (the \emph{transported algorithm}: play $Uu$ whenever $\mathsf{Alg}_r$
plays $u$, feed it the same observations, recommend $U\rec_r$) such that, for every
$\theta \in \R^d$ with $\theta_r := U^\top\theta$, the law of $(\hist_T,\rec)$ under
$\nu_\theta$ is the image of the law of $(\hist_{T,r},\rec_r)$ under $\nu_{\theta_r}$ by
$((u_t,y_t)_{t<T},\rec_r) \mapsto ((Uu_t,y_t)_{t<T},U\rec_r)$, and the simple regrets agree
pathwise: $r_{\Xs}(U\rec_r,\theta) = r_{\Xs_r}(\rec_r,\theta_r)$. Consequently
$\mathsf{Alg}$ is $(\eps,\delta)$-PAC on $\Xs$ if and only if $\mathsf{Alg}_r$ is
$(\eps,\delta)$-PAC on $\Xs_r$.
\end{lemma}

\begin{proof}
\uses{def:env, def:regret, lem:transform_preserves_pac}
The identities: $\ip{Uu}{\theta} = u^\top U^\top\theta$, and for $\theta \in V$,
$UU^\top\theta = \theta$ ($UU^\top$ is the orthogonal projection onto $V$), so
$\ip{U^\top x}{U^\top\theta} = \ip{x}{UU^\top\theta} = \ip{x}{\theta}$. The map $\phi$ is
continuous with continuous inverse $U^\top$ (as $U^\top Uu = u$ and $UU^\top x = x$ for
$x \in \Xs \subset V$), hence a measurable bijection $\Xs_r \to \Xs$ with measurable
inverse. The reward kernel of $\nu_\theta$ at $Uu$ is $\Normal(\ip{Uu}{\theta},1)
= \Normal(\ip{u}{\theta_r},1)$, the reward kernel of $\nu_{\theta_r}$ at $u$: the
environments correspond under $\phi$. The transported algorithm is obtained by mapping the
policy kernels of $\mathrm{alg}_r$ through $\phi$ (and the histories back through
$\phi^{-1}$), exactly as in Lemma~\ref{lem:transform_preserves_pac}, whose proof only uses
that $x \mapsto Mx$ is a measurable bijection of action sets under which the environments
correspond; the same argument (trajectory measures are characterized by their kernels,
Ionescu--Tulcea) gives the image-law statement. For the regrets,
$\max_{x\in\Xs}\ip{x}{\theta} = \max_{u\in\Xs_r}\ip{Uu}{\theta} = \max_{u\in\Xs_r}\ip{u}{\theta_r}$
($\phi$ is onto $\Xs$) and $\ip{U\rec_r}{\theta} = \ip{\rec_r}{\theta_r}$, so the regrets
agree. Finally, $\theta \mapsto U^\top\theta$ is onto $\R^r$ (with right inverse
$\theta_r \mapsto U\theta_r$), so ``for every $\theta$'' on $\Xs$ and ``for every $\theta_r$''
on $\Xs_r$ are the same condition, and the PAC equivalence follows from the equality of the
probabilities of the events $\{r \le \eps\}$.
\end{proof}

\textbf{Lean remark.} The lemma is a special case of a general ``change of action type''
statement for LML algorithms along a measurable bijection $\phi$ with measurable inverse
between action types, such that the environments correspond
($\kappa_\theta \circ \phi = \kappa_{\theta_r}$); Lemma~\ref{lem:transform_preserves_pac} and
Lemma~\ref{lem:separation_block_embedding} are two more instances. Only the direction
``from $\Xs_r$ to $\Xs$'' is used (to transport the upper bound of Theorem~\ref{thm:upper});
the converse direction uses the identity $\ip{U^\top x}{U^\top\theta} = \ip{x}{\theta}$ for
$\theta \in V$ together with the fact that the environment only sees $\ip{x}{\theta}$,
$x \in \Xs$, i.e.\ $\theta$ only through its projection $UU^\top\theta \in V$.

\begin{remark}[Table~1 of the paper]\label{rem:table}
For the structured sets of this chapter, the minimax sample complexity has the form
$H_1\log(1/\delta) + H_2$ with $H_1 = \Theta(d/\eps^2)$ (upper bound: Theorem~\ref{thm:upper};
lower bound: Theorem~\ref{thm:lower_adaptive}). The chapter establishes the following
bounds on $H_2$; the non-adaptive upper bounds come from Theorem~\ref{thm:upper} and the
width bounds of Chapter~\ref{chap:width}, the adaptive lower bounds from
Section~\ref{sec:log_gains_lower}, so that on each of these sets adaptivity can only improve
the sample complexity of the fixed-design algorithm by logarithmic factors.
\begin{itemize}
  \item Unit ball $\ball_d$: non-adaptive upper bound $O((d\log(1/\delta) + d^2)/\eps^2)$
        (Proposition~\ref{prop:width_ball}, $\gw(\ball_d) \le d$); adaptive lower bound
        $T \ge d^2/(1000\eps^2)$ for $\delta \le 1/500$ (Corollary~\ref{cor:lower_ball_adaptive}).
  \item Hypercubes $\{-1,1\}^d$ and $\{0,1\}^d$: non-adaptive upper bound
        $O((d\log(1/\delta) + d^2)/\eps^2)$ (Proposition~\ref{prop:width_le_d}); adaptive lower
        bounds $T > d^2/(100\eps^2)$, resp.\ $T > d^2/(400\eps^2)$, for $\delta < 1/6$
        (Theorems~\ref{thm:lower_hypercube_pm} and~\ref{thm:lower_hypercube_01}).
  \item $m$-sets with $d - m + 1 \ge 20m$ (the paper writes $m \le d/21$): non-adaptive
        upper bound $O((d\log(1/\delta) + md\log(ed/m))/\eps^2)$
        (Proposition~\ref{prop:width_finite} with $\log\binom dm \le m\log(ed/m)$); adaptive
        lower bound $T > m(d-m+1)/(2500\eps^2)$ for $\delta < 2/9$ (Theorem~\ref{thm:lower_msets}).
        The gap is the factor $\log(d/m)$.
  \item Multi-task bandits with block sizes $d_1,\dots,d_m \ge 2$: non-adaptive upper bound
        $O\big((d\log(1/\delta) + (\sum_j\sqrt{d_j\log d_j})^2)/\eps^2\big)$. Here the
        multi-task set $\Xs$ is not spanning, so Theorem~\ref{thm:upper} is applied to the
        spanning action set $\Xs_r := U^\top\Xs \subset \R^{r}$, $r = d - m + 1 \le d$, with
        $U$ the orthonormal basis of $\Span\Xs$ of Lemma~\ref{lem:multitask_basis}: it gives
        a fixed-design $(\eps,\delta)$-PAC algorithm on $\Xs_r$ with budget
        $\lceil 600(\gw(\Xs_r)^2 + r\log(2/\delta))/\eps^2\rceil$, where
        $\gw(\Xs_r) = \gw_{\mathrm{int}}(\Xs) \le \sum_j\sqrt{2d_j\log d_j}$
        (Definition~\ref{def:width_subspace}, Proposition~\ref{prop:width_multitask}); this
        algorithm is transported to $\Xs$ by Lemma~\ref{lem:log_gains_subspace_transport}
        (play $Uu$ instead of $u$; $\ip{Uu}{\theta} = \ip{u}{U^\top\theta}$), and the transported
        algorithm is again a deterministic fixed design on $\Xs$ with a deterministic
        recommendation. Adaptive lower bound
        $T > (\sum_j\sqrt{d_j})^2/(20000\eps^2)$ for $\delta < 0.4$
        (Theorem~\ref{thm:lower_multitask}). The gap is at most the factor $\max_j\log d_j$.
\end{itemize}
For the hypercubes and the $m$-sets, the adaptive lower bounds are also what yields the
width lower bounds of Corollary~\ref{cor:width_structured}. The region algorithm of
Theorem~\ref{thm:log_gains} gives, for the multi-armed bandit $\Xs = \{e_1,\dots,e_d\}$, the
bound $O(d\log(1/\delta)/\eps^2)$ of Median Elimination, and more generally saves a factor
$\log d$ over the finite-set bound whenever $\abs\Xs = d^{O(1)}$ and $\delta$ is constant.
\end{remark}

\chapter{Estimation of the \texorpdfstring{$\ell_2$}{l2}-norm of the reward vector}
\label{chap:norm}

This chapter formalizes Section~3 and Appendix~D of the paper (the four subsections on
$\ell_2$-norm estimation): on the action set $\Xs = \ball_d$ one wants to estimate
$r := \norm{\theta}$ up to an additive error $\eps$ with probability at least $1-\delta$,
using $O(d\log(1/\delta)/\eps^2)$ observations. The target is Theorem~8 of the paper
(Theorem~\ref{thm:norm_estimation}). The procedure has three ingredients: a statistic
built from random Rademacher directions, each sampled repeatedly, which estimates $r^2$ with
an additive error proportional to $r\eps$ once a constant-factor estimate $r_0$ of $r$ is
available (Section~\ref{sec:norm_additive}); an adaptive multi-scale test which produces
such an $r_0$, or detects that $r$ is either at most $\eps$ or at least $\sqrt d$
(Section~\ref{sec:norm_test}); and a direct estimator for the large-norm regime
$r \ge \sqrt d$ (Section~\ref{sec:norm_large}). The three are composed into a
meta-algorithm in Section~\ref{sec:norm_meta}. The chapter is fully formalized.

From Chapter~\ref{chap:pre_subexp} we use the sub-exponential calculus: the definition
\ref{def:subexp}, the tail bound for averages (Corollary~\ref{cor:subexp_average_tail}), the
sub-exponential parameters of squared Gaussians (Lemma~\ref{lem:gaussian_square_subexp}), of
chi-square variables (Lemma~\ref{lem:chi_square_subexp}) and of squared Rademacher linear
forms (Lemma~\ref{lem:rademacher_square_subexp}). From the model chapter we use the seed
representation of the runs of a seeded algorithm (Lemma~\ref{lem:seeded}): the internal
randomness of the algorithm (the Rademacher directions) and the Gaussian noises are the
coordinates of an i.i.d.\ sequence, and every statistic of the run is a measurable
function of finitely many of these coordinates. Compared with the paper, the constants are
all explicit (no attempt was made to optimize them), the Hanson--Wright inequality is
replaced by Lemma~\ref{lem:rademacher_square_subexp}, the conditional argument for the noise
term is a Fubini argument over the directions, the large-norm regime uses a deterministic
basis design instead of random Rademacher rows, which removes the random-matrix
concentration step of the paper, and the sequential composition of the phases is replaced
by a deterministic round schedule on the seed space. Throughout the chapter $\eps \in (0,1]$, $\delta \in (0,1)$, $d \ge 1$,
$\Xs = \ball_d$ and $r := \norm{\theta}$.

\section{Estimation algorithms and the Rademacher statistic}
\label{sec:norm_setup}

\begin{definition}[Norm estimation algorithm]\label{def:norm_estimator}
\lean{Maiti2026Power.IsAccurateNormEst}\leanok
\uses{def:env, def:bai_alg}
An \emph{estimation algorithm with budget $T$} on $\ball_d$ is a pair $(\mathrm{alg}, \rho)$
of an LML algorithm with actions in $\ball_d$ and observations in $\R$, and a Markov kernel
$\rho$ from histories of length $T$ to $[0,\infty)$. Run against $\nu_\theta$
(Definition~\ref{def:env}) it produces $\hist_T$ and the estimate $\hat r \sim \rho(\hist_T)$;
we write $\Prob_\theta$ for the law of $(\hist_T, \hat r)$. It is
\emph{$(\eps,\delta)$-accurate} if for every $\theta \in \R^d$,
\[
  \Prob_\theta\big( \abs{\hat r - \norm{\theta}} \le \eps \big) \ge 1 - \delta .
\]
\end{definition}

\textbf{Lean remark.} This is Definition~\ref{def:bai_alg} with the recommendation type
$\Xs$ replaced by $[0,\infty)$ (or simply $\R$): the LML notion of an algorithm with a
recommendation is polymorphic in the recommendation type, so no new framework is needed.
All the estimators below are deterministic functions of the history, so $\rho$ is
\texttt{Kernel.deterministic}.

\begin{definition}[Rademacher direction]\label{def:rademacher_direction}
\lean{Maiti2026Power.radDir, Maiti2026Power.radDir_mem_unitBall, Maiti2026Power.measurable_radDir}\leanok
\uses{def:arm_set}
A \emph{Rademacher direction} is the random vector $x := \sigma/\sqrt d \in \R^d$ where
$\sigma = (\sigma_1, \dots, \sigma_d)$ has i.i.d.\ coordinates uniform on $\{-1, +1\}$. We
write $\mathrm{Rad}_d$ for its law (the image of the uniform measure on $\{-1,+1\}^d$ by
$\sigma \mapsto \sigma/\sqrt d$). Since $\norm{x} = 1$, $x \in \sphere^{d-1} \subset \ball_d$.
\end{definition}

\textbf{Lean remark.} The sign vector is encoded by a Boolean vector $u \in \{0,1\}^d$ with
the uniform law (\texttt{uniformBoolVec}, the seed of the meta-algorithm), through
$\sigma_i := \mathrm{signOf}(u_i) \in \{-1,+1\}$; $\mathrm{radDir}(u) := \sigma/\sqrt d$.

\begin{lemma}[Inner products with a Rademacher direction]\label{lem:direction_second_moment}
\lean{Maiti2026Power.inner_radDir, Maiti2026Power.card_mul_inner_radDir_sq, Maiti2026Power.norm_radDir_le_one}\leanok
\uses{def:rademacher_direction}
Let $x = \sigma/\sqrt d$ be a Rademacher direction and $\theta \in \R^d$. Then
$\ip{x}{\theta} = \ip{\sigma}{\theta}/\sqrt d$ and $d\,\ip{x}{\theta}^2 = \ip{\sigma}{\theta}^2$;
in particular $\E\big[\ip{x}{\theta}^2\big] = \norm{\theta}^2/d$.
\end{lemma}

\begin{proof}
\leanok
\uses{lem:rademacher_fourth_moment}
$\E[\sigma_i\sigma_j] = \indic\{i = j\}$ by independence and $\E\sigma_i = 0$,
$\sigma_i^2 = 1$, so $\E[\sigma\sigma^\top] = I_d$ and $\E[xx^\top] = I_d/d$. Then
$\E\ip{x}{\theta}^2 = \theta^\top\E[xx^\top]\theta = \norm{\theta}^2/d$ (this is also the
second-moment identity of Lemma~\ref{lem:rademacher_fourth_moment} divided by $d$). The last
identity is $\ip{x}{\theta} = \ip{\sigma}{\theta}/\sqrt d$.
\end{proof}

\begin{definition}[Rademacher block sampling and the squared statistic]\label{def:squared_statistic}
\lean{Maiti2026Power.rbMeasure, Maiti2026Power.rbMean, Maiti2026Power.rbStat, Maiti2026Power.rbX, Maiti2026Power.rbW, Maiti2026Power.rbStat_sub_sq_eq}\leanok
\uses{def:rademacher_direction, def:env, def:regret}
Let $s \ge 1$ and $K \ge 1$ be integers. The \emph{Rademacher block sampling rule}
$\mathrm{RB}(s, K)$ is the sampling rule with actions in $\ball_d$ and budget $sK$ which, for
each $k < K$, draws at round $ks$ a fresh Rademacher direction $x^{(k)} = \sigma^{(k)}/\sqrt d$
and plays the same action $x^{(k)}$ at rounds $ks + \ell$, $0 \le \ell < s$. Its
\emph{block space} is the product space
\[
  \Omega_{s,K} := \big(\{-1,+1\}^d\big)^K \times \big(\R^s\big)^K
  \quad\text{with the product law}\quad
  \Prob_{s,K} := \mathrm{Unif}\big(\{-1,+1\}^d\big)^{\otimes K} \otimes
    \big(\Normal(0,1)^{\otimes s}\big)^{\otimes K}
\]
(the directions of the $K$ blocks and the noises of the $sK$ rounds). For
$(\sigma, \eta) \in \Omega_{s,K}$ and a reward vector $\theta$ with $r = \norm{\theta}$, define
for $k < K$
\[
  x^{(k)} := \frac{\sigma^{(k)}}{\sqrt d}, \qquad
  \mu_k := \ip{x^{(k)}}{\theta}, \qquad
  \bar y_k := \mu_k + \frac1s \sum_{\ell < s} \eta_{k,\ell}, \qquad
  Z_k := d\Big( \bar y_k^2 - \frac1s \Big), \qquad
  \bar Z := \frac1K \sum_{k<K} Z_k ,
\]
\[
  X_k := d\mu_k^2 - r^2, \qquad
  W_k := Z_k - d\mu_k^2 .
\]
Then, as an algebraic identity,
\[
  \bar Z - r^2 = \frac1K \sum_{k<K} X_k + \frac1K \sum_{k<K} W_k .
\]
When $\mathrm{RB}(s,K)$ is run against $\nu_\theta$, the observations of block $k$ are
$y_{ks+\ell} = \mu_k + \eta_{k,\ell}$ (Lemma~\ref{lem:seeded}), so that $\bar y_k$ is the
average of the observations of block $k$ and $\bar Z$ is a statistic of the history.
\end{definition}

\textbf{Lean remark.} $\mathrm{RB}(s,K)$ is not a standalone algorithm in Lean: it appears as
the \emph{windows} of the meta-algorithm of Definition~\ref{def:meta_algorithm}, whose seeds
at the block starts are the sign vectors $\sigma^{(k)}$ and whose noises are the
$\eta_{k,\ell}$ (Lemma~\ref{lem:statistic_structure}). All the probabilistic statements of
this chapter about $\bar Z$ are statements about the block space
$(\Omega_{s,K}, \Prob_{s,K})$ (\texttt{rbMeasure}), and are transported to the run through the
window projections of Lemma~\ref{lem:statistic_structure}. The identity
$\bar Z - r^2 = \frac1K\sum X_k + \frac1K\sum W_k$ is a rearrangement of
$Z_k = X_k + W_k + r^2$.

\textbf{Lean remark.} The directions are the algorithm's internal randomness: in LML the
policy of $\mathrm{RB}(s,K)$ is the kernel which, at rounds $t$ with $s \mid t$, ignores the
history and returns the measure $\mathrm{Rad}_d$ (a finitely supported measure on the
subtype $\ball_d$), and at the other rounds returns the Dirac mass at $x_{t-1}$. The
quantities $\bar y_k$, $Z_k$, $\bar Z$ are measurable functions of the history of length
$sK$. The identity $\bar Z - r^2 = \frac1K\sum X_k + \frac1K\sum W_k$ is a rearrangement of
$Z_k = X_k + W_k + r^2$.

\begin{lemma}[Window projections of the seed space]\label{lem:statistic_structure}
\lean{Maiti2026Power.nsMeasure, Maiti2026Power.rbProj, Maiti2026Power.nsMeasure_map_rbProj, Maiti2026Power.indepFun_prefix_rbProj, Maiti2026Power.lnProj, Maiti2026Power.nsMeasure_map_lnProj, Maiti2026Power.indepFun_prefix_lnProj, ProbabilityTheory.indepFun_prefix_window, MeasureTheory.Measure.infinitePi_map_eval_comp, MeasureTheory.Measure.infinitePi_map_eval_comp', Maiti2026Power.shiftSeq, Maiti2026Power.nsMeasure_map_shiftSeq, Maiti2026Power.winStat, Maiti2026Power.winStat_eq_rbStat}\leanok
\uses{def:squared_statistic, lem:seeded}
Let $\Omega_{\mathrm{seed}} := \big(\{-1,+1\}^d \times \R\big)^{\N}$ with the product law
$\Prob_{\mathrm{seed}} := \big(\mathrm{Unif}(\{-1,+1\}^d) \otimes \Normal(0,1)\big)^{\otimes\N}$
(the seed space of Lemma~\ref{lem:seeded} for seeds which are uniform sign vectors and
standard Gaussian noises), and write $\omega_t = (\sigma_t, \eta_t)$ for its coordinates.
\begin{enumerate}
  \item[(i)] For $a \in \N$, $s, K \ge 1$, the \emph{Rademacher window projection}
    $\pi_{a,s,K}(\omega) := \big( (\sigma_{a+ks})_{k<K}, (\eta_{a+ks+\ell})_{k<K,\ell<s} \big)
    \in \Omega_{s,K}$ has law $\Prob_{s,K}$ (Definition~\ref{def:squared_statistic}), and it
    is independent of the prefix $(\omega_t)_{t<T}$ whenever $T \le a$.
  \item[(ii)] For $a, n \in \N$, the \emph{basis window projection}
    $\pi'_{a,n}(\omega) := (\eta_{a+in+\ell})_{i<d,\ell<n} \in (\R^n)^d$ has law
    $\Normal(0,1)^{\otimes dn}$, and it is independent of the prefix $(\omega_t)_{t<T}$
    whenever $T \le a$.
  \item[(iii)] For an observation sequence $y : \N \to \R$, let
    $\bar Z_{a,s,K}(y) := \frac1K\sum_{k<K} d\big( (\frac1s\sum_{\ell<s} y_{a+ks+\ell})^2 - \frac1s \big)$
    be the squared statistic of the window of rounds $[a, a + sK)$. If
    $y_{a+ks+\ell} = \ip{\sigma_{a+ks}/\sqrt d}{\theta} + \eta_{a+ks+\ell}$ for all $k < K$,
    $\ell < s$, then $\bar Z_{a,s,K}(y) = \bar Z\big(\pi_{a,s,K}(\omega)\big)$, the statistic of
    Definition~\ref{def:squared_statistic} evaluated at the window projection.
  \item[(iv)] For $a \in \N$, the shifted sequence $(\omega_{a+t})_{t\in\N}$ has law
    $\Prob_{\mathrm{seed}}$.
\end{enumerate}
\end{lemma}

\begin{proof}
\leanok
\uses{lem:seeded, def:squared_statistic}
The coordinates of an i.i.d.\ sequence indexed by an injective family of indices are
i.i.d.\ with the same law: the image of $\mu^{\otimes\N}$ by
$\omega \mapsto (\omega_{f(p)})_{p \in J}$ is $\mu^{\otimes J}$ for injective $f : J \to \N$
(\texttt{Measure.infinitePi\_map\_eval\_comp}, from the restriction and reindexing lemmas of
infinite product measures; for an infinite index set, as in (iv), the image measure is
identified through its finite-dimensional cylinders, \texttt{Measure.eq\_infinitePi}). The window projections are of this form (the indices
$a + ks + \ell$, resp.\ $a + in + \ell$, are pairwise distinct), which gives the laws after
splitting the pairs $(\sigma, \eta)$ and currying the index $(k, \ell)$. The prefix and a
window with $T \le a$ use disjoint sets of coordinates, and disjoint families of
coordinates of an i.i.d.\ sequence are independent (\texttt{iIndepFun\_infinitePi} and
\texttt{indepFun\_finset}). Part (iii) is a rewriting: the block average of the
observations is $\mu_k + \frac1s\sum_\ell \eta_{k,\ell}$.
\end{proof}

\textbf{Lean remark.} The lemma replaces the sequential composition and the conditional
independence arguments of the paper: on the seed space, every event of the multi-scale
phase (rounds $t < T_1$) is a function of the prefix $(\omega_t)_{t<T_1}$, and every
statistic of the second phase is a function of a window projection with $a = T_1$; the two
are independent by (i)--(ii).

\begin{proof}
\uses{lem:noise_representation, lem:gauss_pi, lem:gauss_linear_image}
By Lemma~\ref{lem:noise_representation} the $\eta_t$, $t < sK$, are i.i.d.\ $\Normal(0,1)$
and $\eta_t$ is independent of $(\hist_t, x_t)$. Since $x_{ks+\ell} = x^{(k)}$ for
$\ell < s$, averaging the observations of block $k$ gives $\bar y_k = \mu_k + \zeta_k$.
The direction $x^{(k)}$ is drawn at round $ks$ from the constant kernel $\mathrm{Rad}_d$, so
it is independent of $\hist_{ks}$, hence of $(x^{(j)}, \zeta_j)_{j<k}$; and
$(\eta_{ks+\ell})_{\ell<s}$ is independent of $(\hist_{ks}, x^{(k)})$. An induction on $k$
shows that the law of $\big((x^{(j)},\zeta_j)\big)_{j \le k}$ is the product law: at each
step the conditional law of $(x^{(k)}, (\eta_{ks+\ell})_{\ell<s})$ given the past is the
constant kernel $\mathrm{Rad}_d \otimes \Normal(0,1)^{\otimes s}$ (Lemma~\ref{lem:gauss_pi}).
Finally $\zeta_k = \frac1s\mathbf 1^\top(\eta_{ks+\ell})_{\ell<s} \sim \Normal(0, 1/s)$ by
Lemma~\ref{lem:gauss_linear_image}.
\end{proof}

\textbf{Lean remark.} The proof is an induction on the blocks using the step decomposition
of the trajectory measure (LML \texttt{IsAlgEnvSeq.hasCondDistrib\_step}): since the policy
kernel at round $ks$ does not depend on the history, the conditional law of the block
$k$ given $\hist_{ks}$ is a constant kernel, which is exactly independence from the past.
Alternatively one can view $\mathrm{RB}(s,K)$ as the $K$-fold sequential composition
(Lemma~\ref{lem:composition}) of the one-block rule $\mathrm{RB}(s,1)$.

\section{Additive estimation given a constant-factor estimate}
\label{sec:norm_additive}

Throughout this section $\mathrm{RB}(s,K)$ is run against $\nu_\theta$ and the notation of
Definition~\ref{def:squared_statistic} is in force. The decomposition
$\bar Z - r^2 = \frac1K\sum_k X_k + \frac1K\sum_k W_k$ separates the fluctuations of the
directions (Term~1, the $X_k$) from the fluctuations of the noise (Term~2, the $W_k$).

\begin{lemma}[Term 1 is sub-exponential]\label{lem:term1_subexp}
\lean{Maiti2026Power.hasSubexponentialMGF_rbX, Maiti2026Power.iIndepFun_rbX, Maiti2026Power.rbX_eq_zero_of_eq_zero}\leanok
\uses{def:squared_statistic, lem:statistic_structure, def:subexp}
On the block space $(\Omega_{s,K}, \Prob_{s,K})$, the variables $X_k = d\mu_k^2 - r^2$,
$k < K$, are independent and each has $\mathrm{HasSubexponentialMGF}\ X_k\ (16 r^4)\ (4r^2)$.
If $r = 0$ then $X_k = 0$ for all $k$.
\end{lemma}

\begin{proof}
\leanok
\uses{lem:rademacher_square_subexp, lem:direction_second_moment, def:squared_statistic}
By Lemma~\ref{lem:direction_second_moment}, $d\mu_k^2 = \ip{\sigma^{(k)}}{\theta}^2$ where
$\sigma^{(k)}$ is a uniform sign vector, so $X_k = \ip{\sigma^{(k)}}{\theta}^2 - \norm{\theta}^2$,
which has the sub-exponential parameters $(16r^4, 4r^2)$ by
Lemma~\ref{lem:rademacher_square_subexp}. The $\sigma^{(k)}$ are i.i.d.\ by
Lemma~\ref{lem:statistic_structure}, hence so are the $X_k$. When $\theta = 0$,
$\mu_k = 0 = r$.
\end{proof}

\begin{lemma}[Tail bound for Term 1]\label{lem:term1_bound}
\lean{Maiti2026Power.rbMeasure_real_avg_rbX_gt_le}\leanok
\uses{lem:term1_subexp}
Assume $r > 0$. For every $t \ge 0$,
\[
  \Prob_{s,K}\Big( \Big| \frac1K\sum_{k<K} X_k \Big| > t \Big)
  \le 2\exp\Big( -K \min\Big( \frac{t^2}{32 r^4}, \frac{t}{8 r^2} \Big) \Big).
\]
If $r = 0$ the left-hand side is $0$.
\end{lemma}

\begin{proof}
\leanok
\uses{cor:subexp_average_tail, lem:term1_subexp}
Corollary~\ref{cor:subexp_average_tail} applied to the independent variables $X_k$ with the
common parameters $(V, b) = (16r^4, 4r^2)$ (Lemma~\ref{lem:term1_subexp}) gives the bound
$2\exp(-K\min(t^2/(2V), t/(2b)))$, and $2V = 32r^4$, $2b = 8r^2$. The case $r = 0$ is
$X_k = 0$.
\end{proof}

\begin{lemma}[Term 2 is conditionally sub-exponential]\label{lem:term2_conditional}
\lean{Maiti2026Power.rbNoise, Maiti2026Power.hasLaw_sqrt_mul_rbMean, Maiti2026Power.hasSubexponentialMGF_rbW, Maiti2026Power.iIndepFun_rbW}\leanok
\uses{def:squared_statistic, lem:statistic_structure, def:subexp}
Fix deterministic sign vectors $\sigma^{(0)}, \dots, \sigma^{(K-1)}$, let
$x^{(k)} := \sigma^{(k)}/\sqrt d$, $\mu_k := \ip{x^{(k)}}{\theta}$, and let the noises
$(\eta_{k,\ell})_{k<K,\ell<s}$ be i.i.d.\ $\Normal(0,1)$ (the noise marginal
$\Prob^{\eta}_{s,K} := (\Normal(0,1)^{\otimes s})^{\otimes K}$ of the block space). Define
$\bar y_k$, $Z_k$ and $W_k := Z_k - d\mu_k^2$ as in Definition~\ref{def:squared_statistic}.
Then $\sqrt d\,\bar y_k \sim \Normal(\sqrt d\,\mu_k, d/s)$, the $W_k$ are independent and
$\mathrm{HasSubexponentialMGF}\ W_k\ V_k\ b$ with
\[
  V_k := 8\Big( \frac{d^2}{s^2} + \mu_k^2\,\frac{d^2}{s} \Big), \qquad b := \frac{4d}{s} .
\]
\end{lemma}

\begin{proof}
\leanok
\uses{lem:gaussian_square_subexp, lem:gauss_linear_image}
$\sqrt d\,\bar y_k \sim \Normal(\sqrt d\,\mu_k, d/s)$ (Lemma~\ref{lem:gauss_linear_image}),
and
\[
  W_k = d\bar y_k^2 - \frac ds - d\mu_k^2
      = \big(\sqrt d\,\bar y_k\big)^2 - \Big( (\sqrt d\,\mu_k)^2 + \frac ds \Big),
\]
which is $X^2 - (\mu^2 + \sigma^2)$ for $X \sim \Normal(\mu,\sigma^2)$ with
$\mu = \sqrt d\,\mu_k$, $\sigma^2 = d/s$. Lemma~\ref{lem:gaussian_square_subexp} gives the
parameters $(8(\sigma^4 + \mu^2\sigma^2), 4\sigma^2) = (8(d^2/s^2 + \mu_k^2 d^2/s), 4d/s)$.
Independence of the $W_k$ follows from that of the $\zeta_k$, each $W_k$ being a function
of $\zeta_k$ alone.
\end{proof}

\begin{lemma}[Tail bound for Term 2]\label{lem:term2_bound}
\lean{Maiti2026Power.rbVbar, Maiti2026Power.sum_rbVbar_eq, Maiti2026Power.rbNoise_real_avg_rbW_ge_le, Maiti2026Power.rbMeasure_real_avg_rbW_gt_le}\leanok
\uses{lem:term2_conditional, lem:statistic_structure}
On the block space, let
\[
  \bar V := \frac1K\sum_{k<K} V_k
  = 8\Big( \frac{d^2}{s^2} + \frac{d^2}{s}\cdot\frac1K\sum_{k<K}\mu_k^2 \Big),
  \qquad b := \frac{4d}{s} ,
\]
a function of the directions. Then for every $\tau \in \R$ and $t \ge 0$,
\[
  \Prob_{s,K}\Big( \Big| \frac1K\sum_{k<K} W_k \Big| > t \Big)
  \le \Prob_{s,K}(\bar V > \tau)
    + 2\exp\Big( -K\min\Big( \frac{t^2}{2\tau}, \frac{t}{2b} \Big) \Big).
\]
\end{lemma}

\begin{proof}
\leanok
\uses{lem:term2_conditional, lem:statistic_structure, cor:subexp_average_tail}
By Lemma~\ref{lem:statistic_structure} the joint law of the directions and the noise
averages is the product $\mathrm{Rad}_d^{\otimes K} \otimes \Normal(0,1/s)^{\otimes K}$, and
$\frac1K\sum_k W_k$ is a measurable function $F(x^{(0:K)}, \zeta_{0:K})$. Split the event on
$\{\bar V \le \tau\}$, which depends on the directions only, and apply Fubini:
\[
  \Prob_\theta\Big(\Big|\frac1K\sum_k W_k\Big| > t,\ \bar V \le \tau\Big)
  = \int \indic\{\bar V(x) \le \tau\}\,
    \Prob_\zeta\Big( \Big| \frac1K\sum_k W_k(x, \zeta) \Big| > t \Big)\,
    \mathrm{Rad}_d^{\otimes K}(dx) .
\]
For fixed directions $x$, Lemma~\ref{lem:term2_conditional} and
Corollary~\ref{cor:subexp_average_tail} bound the inner probability by
$2\exp(-K\min(t^2/(2\bar V(x)), t/(2b)))$, which is at most
$2\exp(-K\min(t^2/(2\tau), t/(2b)))$ when $\bar V(x) \le \tau$. Adding
$\Prob_\theta(\bar V > \tau)$ for the complementary event gives the claim.
\end{proof}

\textbf{Lean remark.} The block space is the product
$\mathrm{Unif}(\{-1,+1\}^d)^{\otimes K} \otimes (\Normal(0,1)^{\otimes s})^{\otimes K}$ and the
bound is proved by Fubini (\texttt{MeasureTheory.Measure.prod}) over the directions, with
the inner bound of Lemma~\ref{lem:term2_conditional} for each fixed direction vector. No
conditional expectation is needed.

\begin{lemma}[Bound on the conditional variance proxy]\label{lem:vbar_bound}
\lean{Maiti2026Power.rbMeasure_real_rbVbar_gt_le}\leanok
\uses{lem:term2_bound, lem:term1_bound}
With $\bar V$ as in Lemma~\ref{lem:term2_bound} and
\[
  \tau := 8\Big( \frac{d^2}{s^2} + \frac{3 d r^2}{2 s} \Big),
  \qquad\text{one has}\qquad
  \Prob_{s,K}(\bar V > \tau) \le 2\exp(-K/128) .
\]
\end{lemma}

\begin{proof}
\leanok
\uses{lem:term1_bound, def:squared_statistic}
Since $\frac1K\sum_k d\mu_k^2 = r^2 + \frac1K\sum_k X_k$, on the event
$\{|\frac1K\sum_k X_k| \le r^2/2\}$ we have $\frac1K\sum_k\mu_k^2 \le 3r^2/(2d)$ and hence
$\bar V \le 8(d^2/s^2 + (d^2/s)\cdot 3r^2/(2d)) = \tau$. If $r = 0$ this event is sure and
$\bar V = 8d^2/s^2 \le \tau$. If $r > 0$, Lemma~\ref{lem:term1_bound} with $t = r^2/2$ gives
$\min(t^2/(32r^4), t/(8r^2)) = \min(1/128, 1/16) = 1/128$, so the complementary event has
probability at most $2\exp(-K/128)$.
\end{proof}

\begin{definition}[Additive estimation algorithm; Algorithm~1]\label{def:additive_algorithm}
\lean{Maiti2026Power.addS, Maiti2026Power.addK}\leanok
\uses{def:squared_statistic, def:norm_estimator}
Let $\eps \in (0,1]$, $\delta \in (0,1)$ and $r_0 > 0$ (a scale parameter), and set
\[
  c_0 := 4, \qquad c_1 := 4096, \qquad
  s(r_0) := \Big\lceil \frac{c_0 d}{r_0^2} \Big\rceil, \qquad
  K(r_0) := \Big\lceil \frac{c_1 r_0^2 \log(4/\delta)}{\eps^2} \Big\rceil .
\]
\emph{Algorithm~1} with parameters $(\eps,\delta,r_0)$, written
$\mathrm{AddEst}(\eps,\delta,r_0)$, is the estimation algorithm
(Definition~\ref{def:norm_estimator}) with budget $s(r_0)K(r_0)$ whose sampling rule is the
Rademacher block sampling rule $\mathrm{RB}(s(r_0),K(r_0))$ of
Definition~\ref{def:squared_statistic} and whose estimate is the deterministic function of
the history
\[
  \hat r := \sqrt{\max(\bar Z, 0)} ,
\]
where $\bar Z$ is the squared statistic of Definition~\ref{def:squared_statistic}. When
$\eps$ and $\delta$ are fixed we write $s := s(r_0)$ and $K := K(r_0)$.
\end{definition}

\textbf{Lean remark.} Only the parameters $s(r_0), K(r_0)$ (natural numbers computed by
\texttt{Nat.ceil}) are defined on their own; the algorithm itself is the second phase of the
meta-algorithm of Definition~\ref{def:meta_algorithm} for the data-dependent scale
$r_0 = t_{j_*}$, and its estimate is a case of \texttt{NormEstParam.estimate}. The guarantee
below is stated on the block space.

\textbf{Lean remark.} The algorithm is a pair (LML algorithm, deterministic recommendation
kernel); its parameters $s, K$ are natural numbers computed from $(\eps,\delta,r_0)$ by
\texttt{Nat.ceil}. The meta-algorithm of Definition~\ref{def:meta_algorithm} instantiates it
with a data-dependent $r_0$ taking finitely many values.

\begin{theorem}[Additive estimation with a constant-factor estimate; Algorithm~1]\label{thm:additive_estimation}
\lean{Maiti2026Power.rbMeasure_real_addEst_gt_le, Maiti2026Power.rbMeasure_real_sqrt_addEst_gt_le, Maiti2026Power.NormEstParam.s'_mul_K'_le}\leanok
\uses{def:additive_algorithm, def:norm_estimator, def:multiscale_algorithm}
Let $\eps \in (0,1]$, $\delta \in (0,1)$, $r_0 \ge \eps$, and let $s = s(r_0)$, $K = K(r_0)$
be the parameters of Algorithm~1 (Definition~\ref{def:additive_algorithm}, with
$c_0 = 4$, $c_1 = 4096$). Then for every
$\theta$ with $r/2 < r_0 \le 2r$, on the block space $(\Omega_{s,K}, \Prob_{s,K})$,
\[
  \Prob_{s,K}\Big( \abs{\bar Z - r^2} \le \frac{r\eps}{2} \Big) \ge 1 - \frac{3\delta}{8},
  \qquad\text{hence}\qquad
  \Prob_{s,K}\Big( \abs{\hat r - r} \le \frac{\eps}{2} \Big) \ge 1 - \frac{3\delta}{8}
  \ge 1 - \frac{\delta}{2} .
\]
Moreover, for the scales $r_0 = t_j = 2^j\eps$ with $t_j < 2\sqrt d$ (the scales $j < J$ of
Definition~\ref{def:multiscale_algorithm}), the budget satisfies
\[
  sK \le C\,\frac{d\log(4/\delta)}{\eps^2}, \qquad C := 32776 .
\]
\end{theorem}

\begin{proof}
\leanok
\uses{lem:term1_bound, lem:term2_bound, lem:vbar_bound, def:squared_statistic, def:additive_algorithm, def:multiscale_algorithm}
Write $L := \log(4/\delta) \ge \log 4 > 1$ and note $\log(16/\delta) = \log 4 + L \le 2L$.
Since $\eps \le r_0 \le 2r$ we have $r > 0$; the hypotheses also give $r_0/2 < r < 2r_0$. By
Definition~\ref{def:squared_statistic}, $\{|\bar Z - r^2| > r\eps/2\}$ is contained in the
union of $E_1 := \{|\frac1K\sum X_k| > r\eps/4\}$ and $E_2 := \{|\frac1K\sum W_k| > r\eps/4\}$.

\emph{Term 1.} Lemma~\ref{lem:term1_bound} with $t = r\eps/4$: $t^2/(32r^4) = \eps^2/(512r^2)$
and $t/(8r^2) = \eps/(32r)$. Since $\eps \le 2r \le 16r$, the minimum is $\eps^2/(512r^2)$,
so $\Prob(E_1) \le 2\exp(-K\eps^2/(512r^2))$. This is at most $\delta/8$ as soon as
$K \ge 512r^2\log(16/\delta)/\eps^2$, which holds because $r^2 < 4r_0^2$ and
$\log(16/\delta) \le 2L$ give $512r^2\log(16/\delta)/\eps^2 \le 4096\,r_0^2L/\eps^2 \le K$.

\emph{Variance proxy.} Lemma~\ref{lem:vbar_bound}: $\Prob(\bar V > \tau) \le 2e^{-K/128} \le \delta/8$
because $K \ge 4096 r_0^2L/\eps^2 \ge 4096 L \ge 128\log(16/\delta)$ (using $r_0 \ge \eps$).

\emph{Term 2.} Since $s \ge 4d/r_0^2$ we have $d/s \le r_0^2/4$, hence $b = 4d/s \le r_0^2$ and,
with $r^2 < 4r_0^2$,
\[
  \tau = 8\Big( \frac{d^2}{s^2} + \frac{3dr^2}{2s} \Big)
  \le 8\Big( \frac{r_0^4}{16} + \frac{3 r^2 r_0^2}{8} \Big)
  \le 8\Big( \frac{r_0^4}{16} + \frac{3 r_0^4}{2} \Big) = \frac{25}{2}\,r_0^4 .
\]
With $t = r\eps/4$ and $r > r_0/2$:
$t^2/(2\tau) \ge (r^2\eps^2/16)/(25r_0^4) = r^2\eps^2/(400r_0^4) > \eps^2/(1600 r_0^2)$ and
$t/(2b) \ge r\eps/(8r_0^2) > \eps/(16r_0)$; since $\eps \le r_0$,
$\eps^2/(1600r_0^2) \le \eps/(16r_0)$, so the minimum in Lemma~\ref{lem:term2_bound} is at
least $\eps^2/(1600r_0^2)$ and
$\Prob(E_2) \le \Prob(\bar V > \tau) + 2\exp(-K\eps^2/(1600r_0^2))$. The last term is at
most $\delta/8$ as soon as $K \ge 1600r_0^2\log(16/\delta)/\eps^2$, which holds since
$1600r_0^2\log(16/\delta)/\eps^2 \le 3200\,r_0^2L/\eps^2 \le K$.

\emph{Union bound.} $\Prob(|\bar Z - r^2| > r\eps/2) \le \delta/8 + \delta/8 + \delta/8 = 3\delta/8$.

\emph{From $r^2$ to $r$.} On the event $|\bar Z - r^2| \le r\eps/2$, since $r > r_0/2 \ge \eps/2$,
$\bar Z \ge r(r - \eps/2) > 0$, so $\hat r = \sqrt{\bar Z}$ and
\[
  \abs{\hat r - r} = \frac{\abs{\bar Z - r^2}}{\sqrt{\bar Z} + r} \le \frac{r\eps/2}{r} = \frac{\eps}{2}.
\]

\emph{Budget.} For $r_0 = t_j < 2\sqrt d$, $4d/r_0^2 > 1$, so
$s \le 4d/r_0^2 + 1 \le 8d/r_0^2 = (8d/\eps^2)\,4^{-j}$; and
$K \le 4096\,r_0^2L/\eps^2 + 1 = 4096\cdot 4^j L + 1$. Hence
$sK \le (8d/\eps^2)(4096\,L + 4^{-j}) \le (8d/\eps^2)(4096\,L + 1) \le 8\cdot 4097\, dL/\eps^2
= 32776\, dL/\eps^2$, using $L \ge 1$.
\end{proof}

\begin{remark}[On the constants of Algorithm 1]\label{rem:norm_additive_constants}
The paper takes $c_0, c_1$ ``large enough''; the values $c_0 = 4$, $c_1 = 4096$ are what the
computation above produces with the sub-exponential parameters $(16r^4, 4r^2)$ of
Lemma~\ref{lem:rademacher_square_subexp} and the threshold $t = r\eps/4$ of the paper. The
proof shows that $|\hat r - r| \le \eps/2$, twice as much as needed; using the threshold
$t = r\eps/2$ instead would divide $c_1$ (and $C$) by four. The three failure events are
bounded by $\delta/8$ each, so the total failure probability is $3\delta/8$ rather than the
$\delta/2$ claimed by the paper; the extra room is not used.
\end{remark}

\section{A constant-factor estimate by a multi-scale test}
\label{sec:norm_test}

\begin{definition}[Scale test]\label{def:scale_test}
\lean{Maiti2026Power.testS, Maiti2026Power.testK}\leanok
\uses{def:squared_statistic}
For a scale $t > 0$ and a confidence level $\delta' \in (0,1)$, the test
$\mathrm{Test}(t,\delta')$ is the algorithm $\mathrm{RB}(s,K)$ with
\[
  s := \Big\lceil \frac{c_0 d}{t^2} \Big\rceil, \qquad
  K := \big\lceil c_1' \log(1/\delta') \big\rceil, \qquad
  c_0 := 4, \quad c_1' := 5120 ,
\]
followed by the binary output $\hat H := H_1$ if $U \ge \tfrac32 t^2$ and $\hat H := H_0$
otherwise, where $U := \bar Z$ is the statistic of Definition~\ref{def:squared_statistic}.
The hypothesis $H_0$ is ``$r \le t$'' and $H_1$ is ``$r \ge 2t$''; when $t < r < 2t$ neither
output counts as an error. Its budget is $sK$.
\end{definition}

\textbf{Lean remark.} As for Algorithm~1, only the parameters are defined on their own; the
tests are the windows of the first phase of the meta-algorithm
(Definition~\ref{def:multiscale_algorithm}), and the error bounds below are stated on the
block space.

\begin{lemma}[Error of the test under $H_0$]\label{lem:test_error_h0}
\lean{Maiti2026Power.testS_facts, Maiti2026Power.testK_facts, Maiti2026Power.six_mul_exp_neg_le, Maiti2026Power.rbMeasure_real_test_h0_le}\leanok
\uses{def:scale_test}
Let $\delta' \in (0,1/4]$ and $\theta$ with $r \le t$. Then, on the block space of
$\mathrm{RB}(s,K)$, $\Prob_{s,K}(U \ge \tfrac32 t^2) \le 3\delta'/4$: the test returns $H_1$
with probability at most $3\delta'/4$.
\end{lemma}

\begin{proof}
\leanok
\uses{lem:term1_bound, lem:term2_bound, lem:vbar_bound, def:squared_statistic}
Write $T_1 := \frac1K\sum_k X_k$, $T_2 := \frac1K\sum_k W_k$, so $U = r^2 + T_1 + T_2$. If
$|T_1| \le t^2/8$ and $|T_2| \le t^2/8$ then $U \le r^2 + t^2/4 \le 5t^2/4 < 3t^2/2$ and the
output is $H_0$. Hence the error event is contained in
$\{|T_1| > t^2/8\} \cup \{|T_2| > t^2/8\}$.

\emph{Term 1.} If $r = 0$ then $T_1 = 0$. Otherwise Lemma~\ref{lem:term1_bound} with
$a_0 = t^2/8$ gives $a_0^2/(32r^4) = t^4/(2048 r^4) \ge 1/2048$ and
$a_0/(8r^2) = t^2/(64r^2) \ge 1/64$ since $r \le t$; so
$\Prob(|T_1| > t^2/8) \le 2e^{-K/2048}$.

\emph{Term 2.} Since $s \ge 4d/t^2$: $d/s \le t^2/4$, $b = 4d/s \le t^2$ and
$\tau = 8(d^2/s^2 + 3dr^2/(2s)) \le 8(t^4/16 + 3t^4/8) = 7t^4/2$ (using $r \le t$). With
$\nu_0 = t^2/8$: $\nu_0^2/(2\tau) \ge (t^4/64)/(7t^4) = 1/448$ and
$\nu_0/(2b) \ge (t^2/8)/(2t^2) = 1/16$. Lemmas~\ref{lem:term2_bound} and~\ref{lem:vbar_bound}
give $\Prob(|T_2| > t^2/8) \le 2e^{-K/128} + 2e^{-K/448}$.

\emph{Total.} The error probability is at most
$2e^{-K/2048} + 2e^{-K/128} + 2e^{-K/448} \le 6e^{-K/2048}$. Since $\delta' \le 1/4$,
$\log 8 = \tfrac32\log 4 \le \tfrac32\log(1/\delta')$, hence
$2048\log(8/\delta') \le 2048\cdot\tfrac52\log(1/\delta') = 5120\log(1/\delta') \le K$, and
$6e^{-K/2048} \le 6\delta'/8 = 3\delta'/4$.
\end{proof}

\begin{lemma}[Error of the test under $H_1$]\label{lem:test_error_h1}
\lean{Maiti2026Power.rbMeasure_real_test_h1_le}\leanok
\uses{def:scale_test}
Let $\delta' \in (0,1/4]$ and $\theta$ with $r \ge 2t$. Then
$\Prob_{s,K}(U < \tfrac32 t^2) \le 3\delta'/4$: the test returns $H_0$ with probability at
most $3\delta'/4$.
\end{lemma}

\begin{proof}
\leanok
\uses{lem:term1_bound, lem:term2_bound, lem:vbar_bound, def:squared_statistic}
With $T_1, T_2$ as above, if $|T_1| \le r^2/8$ and $|T_2| \le r^2/8$ then
$U \ge r^2 - r^2/4 = 3r^2/4 \ge 3t^2$ (as $r^2 \ge 4t^2$) and the output is $H_1$. Hence the
error event is contained in $\{|T_1| > r^2/8\} \cup \{|T_2| > r^2/8\}$; here $r > 0$.

\emph{Term 1.} Lemma~\ref{lem:term1_bound} with $a_1 = r^2/8$:
$\min(a_1^2/(32r^4), a_1/(8r^2)) = \min(1/2048, 1/64) = 1/2048$, so
$\Prob(|T_1| > r^2/8) \le 2e^{-K/2048}$.

\emph{Term 2.} As before $d/s \le t^2/4 \le r^2/16$ and $b \le t^2 \le r^2/4$, and
$\tau \le 8(t^4/16 + 3t^2r^2/8) \le 8(r^4/256 + 3r^4/32) = 25r^4/32$. With $\nu_1 = r^2/8$:
$\nu_1^2/(2\tau) \ge (r^4/64)/(25r^4/16) = 1/100$ and $\nu_1/(2b) \ge (r^2/8)/(r^2/2) = 1/4$.
Lemmas~\ref{lem:term2_bound} and~\ref{lem:vbar_bound} give
$\Prob(|T_2| > r^2/8) \le 2e^{-K/128} + 2e^{-K/100}$.

\emph{Total.} At most $6e^{-K/2048} \le 3\delta'/4$ exactly as in
Lemma~\ref{lem:test_error_h0}.
\end{proof}

\begin{definition}[Multi-scale algorithm; Algorithm~2]\label{def:multiscale_algorithm}
\lean{Maiti2026Power.nsScale, Maiti2026Power.nsDelta, Maiti2026Power.nsS, Maiti2026Power.nsK, Maiti2026Power.nsJ, Maiti2026Power.nsStart, Maiti2026Power.one_le_nsJ, Maiti2026Power.nsScale_nsJ_lt, Maiti2026Power.NormEstParam.testStat, Maiti2026Power.NormEstParam.testH1, Maiti2026Power.NormEstParam.jStar}\leanok
\uses{def:scale_test}
Let $\eps \in (0,1]$, $\delta \in (0,1)$ and set, for $j \in \N$,
\[
  t_j := 2^j\eps, \qquad \delta_j := \frac{\delta}{2^{j+2}}, \qquad
  s_j := \Big\lceil \frac{4d}{t_j^2} \Big\rceil, \qquad
  K_j := \big\lceil 5120\log(1/\delta_j) \big\rceil, \qquad
  J := \min\{ j \in \N : t_j \ge 2\sqrt d \} .
\]
Since $\eps \le 1 < 2\sqrt d$, $J \ge 1$; and $t_J < 4\sqrt d$ by minimality. The
\emph{multi-scale algorithm} runs the tests $\mathrm{Test}(t_j, \delta_j)$, $j = 0, \dots, J-1$,
one after the other: test $j$ occupies the rounds $[\mathrm{start}_j, \mathrm{start}_{j+1})$
with $\mathrm{start}_j := \sum_{i<j} s_iK_i$, and its budget is the deterministic number
$N_{\mathrm{ms}} := \mathrm{start}_J = \sum_{j<J} s_jK_j$. Its output is the index
\[
  j_* := \min\big( \{ j < J : \mathrm{Test}(t_j,\delta_j) \text{ returns } H_0 \} \cup \{J\} \big),
\]
i.e.\ the first scale at which the test returns $H_0$, or $J$ if all $J$ tests return $H_1$,
and the constant-factor estimate is $r_0 := t_{j_*}$. It satisfies
$\eps \le r_0 \le t_J < 4\sqrt d$ and $r_0 \in \{t_0, \dots, t_J\}$.
\end{definition}

\textbf{Lean remark.} The paper stops the loop at $j_*$ and idles afterwards; since the
budget is deterministic anyway, the formalization runs all $J$ tests (the tests after $j_*$
are simply not used by the output), which makes the sampling rule of the first phase
non-adaptive. The output $j_*$ is a measurable function of the observations of the rounds
$t < N_{\mathrm{ms}}$ (\texttt{NormEstParam.jStar}, defined by \texttt{Nat.find} from the
test outcomes \texttt{testH1}), with finitely many values; the multi-scale algorithm is the
first phase of the meta-algorithm of Definition~\ref{def:meta_algorithm}.

\textbf{Lean remark.} The loop is a sequential composition of $J$ fixed-length phases
(Lemma~\ref{lem:composition}) whose parameters depend on the past: phase $j$ is
$\mathrm{Test}(t_j,\delta_j)$ if no earlier test returned $H_0$ (a measurable function of the
history, computed from the statistics of the earlier blocks) and the constant algorithm
playing $0$ otherwise. The data-dependent output $r_0$ is a measurable function of the full
history of length $N_{\mathrm{ms}}$ with finitely many values. The algorithm is adaptive, but
in a very restricted sense: only the choice between ``continue'' and ``idle'' depends on the
data.

\begin{lemma}[Good event of the multi-scale algorithm]\label{lem:norm_multiscale_good_event}
\lean{Maiti2026Power.NormEstParam.testOK, Maiti2026Power.NormEstParam.badTest, Maiti2026Power.NormEstParam.nsMeasure_real_badTest_le, Maiti2026Power.NormEstParam.bad, Maiti2026Power.NormEstParam.nsMeasure_real_bad_le, Maiti2026Power.sum_nsDelta_le}\leanok
\uses{def:multiscale_algorithm, lem:test_error_h0, lem:test_error_h1, lem:statistic_structure}
Run the multi-scale algorithm against $\nu_\theta$ on the seed space of
Lemma~\ref{lem:statistic_structure}. For $j < J$ let $E_j$ be the event that the outcome of
test $j$ is wrong: it returns $H_1$ while $r \le t_j$, or $H_0$ while $r \ge 2t_j$.
Then $\Prob_{\mathrm{seed}}(E_j) \le 3\delta_j/4$ and
$\Prob_{\mathrm{seed}}\big( \bigcup_{j<J} E_j \big) \le 3\delta/8$.
\end{lemma}

\begin{proof}
\leanok
\uses{lem:statistic_structure, lem:norm_meta_plan, lem:test_error_h0, lem:test_error_h1, def:multiscale_algorithm}
By Lemma~\ref{lem:norm_meta_plan}, the observations of the window of test $j$ are
$y_{\mathrm{start}_j + ks + \ell} = \ip{\sigma_{\mathrm{start}_j+ks}/\sqrt d}{\theta}
+ \eta_{\mathrm{start}_j+ks+\ell}$, so by Lemma~\ref{lem:statistic_structure}(iii) the
statistic $U_j$ of test $j$ is $\bar Z$ evaluated at the window projection
$\pi_{\mathrm{start}_j, s_j, K_j}$, whose law is the block law $\Prob_{s_j,K_j}$
(Lemma~\ref{lem:statistic_structure}(i)). Hence $E_j$ is the preimage under this
projection of the error event of the test on the block space, and its probability is at
most $3\delta_j/4$ by Lemma~\ref{lem:test_error_h0} (if $r \le t_j$) or
Lemma~\ref{lem:test_error_h1} (if $r \ge 2t_j$); when $t_j < r < 2t_j$, $E_j = \emptyset$.
Here $\delta_j \le \delta/4 < 1/4$. The union bound gives
$\sum_{j<J} 3\delta_j/4 \le \frac{3\delta}{4}\sum_{j\ge0}2^{-j-2} = \frac{3\delta}{8}$.
\end{proof}

\begin{theorem}[Guarantee of the multi-scale algorithm]\label{thm:multiscale_guarantee}
\lean{Maiti2026Power.NormEstParam.jStar_bounds}\leanok
\uses{def:multiscale_algorithm, lem:norm_multiscale_good_event}
Let $E_j$, $j < J$, be the events of Lemma~\ref{lem:norm_multiscale_good_event} and let $G$
be the complement of $\bigcup_{j<J}E_j$ (all tests are correct), so
$\Prob_{\mathrm{seed}}(G) \ge 1 - 3\delta/8$ for every $\theta$. On $G$ the output $j_*$
satisfies one of the following:
\begin{enumerate}
  \item[(a)] $j_* = 0$ (that is, $r_0 = \eps$) and $r < 2\eps$;
  \item[(b)] $0 < j_* < J$, $r/2 < t_{j_*} \le 2r$, and $\eps \le t_{j_*} < 2\sqrt d$;
  \item[(c)] $j_* = J$ (that is, $r_0 = t_J \ge 2\sqrt d$) and $r \ge \sqrt d$.
\end{enumerate}
\end{theorem}

\begin{proof}
\leanok
\uses{lem:norm_multiscale_good_event, def:multiscale_algorithm}
The probability bound is Lemma~\ref{lem:norm_multiscale_good_event}. Fix an outcome in $G$:
every test $j < J$ returns $H_1$ if $r \ge 2t_j$ and $H_0$ if $r \le t_j$. Recall
$t_{j+1} = 2t_j$, $J \ge 1$ and $t_J \ge 2\sqrt d$.

If $j_* = 0$: test $0$ returned $H_0$, which is correct only if $r < 2t_0 = 2\eps$.

If $j_* > 0$: test $j_* - 1$ returned $H_1$, which is correct only if $r > t_{j_*-1} = t_{j_*}/2$,
so $t_{j_*} < 2r$. If moreover $j_* < J$, test $j_*$ returned $H_0$, so $r < 2t_{j_*}$, i.e.\
$r/2 < t_{j_*}$; and $\eps \le t_{j_*} < 2\sqrt d$ by the definition of $J$: this is (b). If
$j_* = J$ then $\sqrt d \le t_J/2 = t_{J-1} < r$: this is (c).
\end{proof}

\begin{lemma}[Budget of the multi-scale algorithm]\label{lem:multiscale_samples}
\lean{Maiti2026Power.sum_pow_le, Maiti2026Power.sum_mul_pow_le, Maiti2026Power.NormEstParam.s_le, Maiti2026Power.NormEstParam.K_le, Maiti2026Power.NormEstParam.T₁_le}\leanok
\uses{def:multiscale_algorithm}
With $L := \log(4/\delta)$, the budget of the multi-scale algorithm satisfies
\[
  N_{\mathrm{ms}} = \sum_{j<J} s_jK_j
  \le \frac{8d}{\eps^2}\Big( 6828\,L + 1578 \Big)
  \le 63736\,\frac{d\log(4/\delta)}{\eps^2}
  \le C'\,\frac{d\log(4/\delta)}{\eps^2}, \qquad C' := 64000 .
\]
\end{lemma}

\begin{proof}
\leanok
\uses{def:multiscale_algorithm}
For $j < J$, $t_j < 2\sqrt d$, so $4d/t_j^2 > 1$ and
$s_j \le 4d/t_j^2 + 1 \le 8d/t_j^2 = 8d\,4^{-j}/\eps^2$. Next,
$\log(1/\delta_j) = L + j\log 2$, so
$K_j \le 5120L + 5120\,j\log 2 + 1 \le 5121\,L + 3549\,j$ using $L \ge \log 4 > 1$ and
$5120\log 2 \le 3549$. With $\sum_{j\ge0}4^{-j} = 4/3$ and $\sum_{j\ge0}j4^{-j} = 4/9$,
\[
  N_{\mathrm{ms}} \le \frac{8d}{\eps^2}\sum_{j\ge0}4^{-j}(5121L + 3549j)
  = \frac{8d}{\eps^2}\Big( \frac{4}{3}\,5121\,L + \frac49\,3549 \Big)
  \le \frac{8d}{\eps^2}\big( 6828\,L + 1578 \big).
\]
Finally $1578 \le 1139\,L$ (as $L \ge \log 4 \ge 1.386$), so
$N_{\mathrm{ms}} \le 8\cdot 7967\,dL/\eps^2 \le 64000\,dL/\eps^2$.
\end{proof}

\section{The large-norm regime}
\label{sec:norm_large}

\begin{definition}[Large-norm estimator; Algorithm~3]\label{def:large_norm_algorithm}
\lean{Maiti2026Power.lnN, Maiti2026Power.lnNoise, Maiti2026Power.lnDelta, Maiti2026Power.lnR, Maiti2026Power.lnEst}\leanok
\uses{def:fixed_design, def:norm_estimator}
Let $\eps \in (0,1]$, $\delta \in (0,1)$ and $n := \lceil 48\log(4/\delta)/\eps^2 \rceil$. The
\emph{large-norm estimator} is the fixed-design estimation algorithm with budget $dn$ which
plays each basis vector $e_i \in \ball_d$, $i = 1,\dots,d$, exactly $n$ times (in any fixed
order), computes the empirical means $\hat\theta_i := \frac1n\sum_{t : x_t = e_i} y_t$, and
outputs
\[
  \hat R := \norm{\hat\theta}^2 - \frac dn, \qquad
  \hat r := \sqrt{\max(\hat R, 0)} .
\]
On the \emph{noise space} $(\R^n)^d$ with the law $\Normal(0,1)^{\otimes dn}$ of the noises
$\eta = (\eta_{i,\ell})_{i<d,\ell<n}$ of these rounds, $\hat\theta = \theta + \Delta$ with
$\Delta_i := \frac1n\sum_{\ell<n}\eta_{i,\ell}$, and $\hat R$, $\hat r$ are functions of
$(\theta, \eta)$.
\end{definition}

\textbf{Lean remark.} The basis design is the second phase of the meta-algorithm of
Definition~\ref{def:meta_algorithm} when $j_* = J$ (rounds $T_1 + in + \ell$ play $e_i$);
the guarantee below is stated on the noise space, and transported through the basis window
projection of Lemma~\ref{lem:statistic_structure}(ii).

\begin{lemma}[Decomposition of the large-norm estimator]\label{lem:large_norm_decomposition}
\lean{Maiti2026Power.lnR_sub_eq, Maiti2026Power.hasLaw_two_mul_inner_lnDelta, Maiti2026Power.lnG, Maiti2026Power.hasLaw_lnG, Maiti2026Power.norm_lnDelta_sq_sub_eq, Maiti2026Power.hasSubexponentialMGF_lnQ}\leanok
\uses{def:large_norm_algorithm, def:subexp}
On the noise space, $\Delta \sim \Normal(0, I_d/n)$ and
\[
  \hat R - r^2 = L + Q, \qquad
  L := 2\ip{\theta}{\Delta} \sim \Normal\big(0, 4r^2/n\big), \qquad
  Q := \norm{\Delta}^2 - \frac dn = \frac{\norm{g}^2 - d}{n},
\]
where $g := \sqrt n\,\Delta \sim \Normal(0, I_d)$; $Q$ has
$\mathrm{HasSubexponentialMGF}\ Q\ (8d/n^2)\ (4/n)$.
\end{lemma}

\begin{proof}
\leanok
\uses{lem:gauss_pi, lem:gauss_linear_image, lem:chi_square_subexp, lem:subexp_const_mul}
The coordinate $\hat\theta_i$ is $\theta_i$ plus the average of the $n$ noise variables of
the rounds where $e_i$ is played. These $d$ averages are independent
$\Normal(0,1/n)$ (a sum of independent Gaussians is Gaussian,
Lemmas~\ref{lem:gauss_pi} and~\ref{lem:gauss_linear_image}), so
$\Delta \sim \Normal(0, I_d/n)$. Expanding
$\norm{\theta + \Delta}^2 = r^2 + 2\ip{\theta}{\Delta} + \norm{\Delta}^2$ gives the
decomposition; $L$ is a linear image of $\Delta$ with variance $4\theta^\top(I/n)\theta$
(Lemma~\ref{lem:gauss_linear_image}). Finally $\norm{g}^2 - d$ has parameters $(8d, 4)$ by
Lemma~\ref{lem:chi_square_subexp}, and scaling by $1/n$ gives $(8d/n^2, 4/n)$ by
Lemma~\ref{lem:subexp_const_mul}.
\end{proof}

\begin{theorem}[Guarantee in the large-norm regime]\label{thm:large_norm_guarantee}
\lean{Maiti2026Power.lnNoise_real_lnEst_gt_le, Maiti2026Power.NormEstParam.n_le}\leanok
\uses{def:large_norm_algorithm}
Let $\eps \in (0,1]$, $\delta \in (0,1)$ and $\theta$ with $r \ge \sqrt d$. Then, on the noise
space,
\[
  \Prob\big( \abs{\hat R - r^2} \ge \eps r \big) \le \frac{\delta}{2},
  \qquad\text{hence}\qquad
  \Prob\big( \abs{\hat r - r} \le \eps \big) \ge 1 - \frac{\delta}{2},
\]
and the budget satisfies $n \le 48\log(4/\delta)/\eps^2 + 1 \le 49\log(4/\delta)/\eps^2$,
so $dn \le 49\,d\log(4/\delta)/\eps^2$.
\end{theorem}

\begin{proof}
\leanok
\uses{lem:large_norm_decomposition, lem:gauss_tail, lem:subexp_tail, lem:gauss_law_of_cov}
Write $L_\delta := \log(4/\delta) > 1$; then $n \ge 48L_\delta/\eps^2 \ge 48$ and
$n\eps^2 \ge 48L_\delta$. By Lemma~\ref{lem:large_norm_decomposition},
$\{|\hat R - r^2| \ge \eps r\} \subseteq \{|L| \ge \eps r/2\} \cup \{|Q| \ge \eps r/2\}$.

\emph{Gaussian term.} $L \sim \Normal(0, 4r^2/n)$ and $-L$ has the same law
(Lemma~\ref{lem:gauss_law_of_cov}), so by Lemma~\ref{lem:gauss_tail} applied to $L$ and $-L$,
\[
  \Prob(|L| \ge \eps r/2) \le 2\exp\Big( -\frac{(\eps r/2)^2}{2\cdot 4r^2/n} \Big)
  = 2\exp\Big( -\frac{n\eps^2}{32} \Big)
  \le 2\exp\Big( -\frac{3}{2}L_\delta \Big) = 2\Big(\frac{\delta}{4}\Big)^{3/2}
  = \frac{\delta^{3/2}}{4} \le \frac{\delta}{4}.
\]

\emph{Chi-square term.} Lemma~\ref{lem:subexp_tail} (both tails) for $Q$ with parameters
$(V,b) = (8d/n^2, 4/n)$ and $u := \eps r/2$:
$u^2/(2V) = \eps^2r^2n^2/(64d) \ge n^2\eps^2/64$ (as $r^2 \ge d$), and
$n^2\eps^2/64 = n(n\eps^2)/64 \ge 48\cdot 48 L_\delta/64 = 36L_\delta$; and
$u/(2b) = \eps r n/16 \ge \eps n/16 \ge 48L_\delta/(16\eps) \ge 3L_\delta$ (as
$r \ge \sqrt d \ge 1$ and $\eps \le 1$). Hence
$\Prob(|Q| \ge \eps r/2) \le 2\exp(-3L_\delta) = 2(\delta/4)^3 \le \delta/4$.

\emph{Union bound.} $\Prob(|\hat R - r^2| \ge \eps r) \le \delta/2$.

\emph{From $r^2$ to $r$.} On the event $|\hat R - r^2| < \eps r$, since $\eps \le 1 \le r$,
$\hat R > r(r - \eps) \ge 0$, so $\hat r = \sqrt{\hat R}$ and
$|\hat r - r| = |\hat R - r^2|/(\sqrt{\hat R} + r) < \eps r/r = \eps$.

\emph{Budget.} $n \le 48L_\delta/\eps^2 + 1$ and $1 \le L_\delta/\eps^2$.
\end{proof}

\begin{remark}[On the deterministic design and the constant $48$]\label{rem:norm_large_design}
The paper draws $n$ random Rademacher rows and uses a random-matrix concentration bound
(Vershynin \cite{vershynin2018high}, Section~4.7) to control the least-squares covariance
$(X^\top X)^{-1}$, together with the Gaussian Hanson--Wright inequality. With the
deterministic basis design the covariance of $\hat\theta$ is exactly $I_d/n$, so only the
one-dimensional Gaussian tail and the chi-square tail are needed, and no invertibility
event has to be handled. The value $n = \lceil 32\log(4/\delta)/\eps^2\rceil$ (which is what
one would first write down, and what the blueprint outline suggested) is not sufficient
for the stated failure probability $\delta/2$: it makes the Gaussian term alone equal to
$2\exp(-\log(4/\delta)) = \delta/2$, leaving nothing for the chi-square term. With $48$ the
two terms are $\delta^{3/2}/4$ and $\delta^3/32$, whose sum is at most $\delta/2$ for every
$\delta \in (0,1)$. The paper also assumes $\eps \le r/2$
for the final step; with $\eps \le 1 \le \sqrt d \le r$ the weaker $\eps \le r$ suffices.
\end{remark}

\section{The meta-algorithm}
\label{sec:norm_meta}

\begin{definition}[Meta-algorithm for norm estimation; Algorithm~4]\label{def:meta_algorithm}
\lean{Maiti2026Power.NormEstParam, Maiti2026Power.nsN2, Maiti2026Power.nsBudget, Maiti2026Power.NormEstParam.T, Maiti2026Power.NormEstParam.phaseOf, Maiti2026Power.NormEstParam.rbBlock, Maiti2026Power.NormEstParam.basisRound, Maiti2026Power.NormEstParam.planAction, Maiti2026Power.NormEstParam.nextAction, Maiti2026Power.NormEstParam.alg, Maiti2026Power.NormEstParam.addStat, Maiti2026Power.NormEstParam.lnRStat, Maiti2026Power.NormEstParam.estimate, Maiti2026Power.NormEstParam.output}\leanok
\uses{def:multiscale_algorithm, def:additive_algorithm, def:large_norm_algorithm, lem:seeded, def:norm_estimator}
Let $\eps \in (0,1]$, $\delta \in (0,1)$, let $T_1 := N_{\mathrm{ms}}$ be the budget of the
multi-scale algorithm and
\[
  N_2 := \max\big( dn,\ \max\{ s(t_j)K(t_j) : j < J \} \big),
\]
where $s(t) = \lceil 4d/t^2\rceil$, $K(t) = \lceil 4096\,t^2\log(4/\delta)/\eps^2\rceil$ are the
parameters of Algorithm~1 with $r_0 = t$ (Definition~\ref{def:additive_algorithm}) and $n$
is that of Definition~\ref{def:large_norm_algorithm}. The \emph{meta-algorithm}
$\mathrm{NormEst}(\eps,\delta)$ is the seeded algorithm (Lemma~\ref{lem:seeded}) with budget
$N_{\mathrm{norm}}(\eps,\delta) := T_1 + N_2$ and seeds uniform on $\{-1,+1\}^d$, defined by the
following \emph{round schedule}. The first phase is the multi-scale algorithm of
Definition~\ref{def:multiscale_algorithm}: for $t < T_1$, if $t = \mathrm{start}_j + ks + \ell$
with $j < J$, $k < K_j$, $\ell < s_j$, the round $t$ belongs to the \emph{Rademacher block}
starting at $b := \mathrm{start}_j + ks$, and the action is $x_t := \sigma_b/\sqrt d$ where
$\sigma_b$ is the seed of round $b$ (the action of round $b$, repeated). The second phase,
rounds $T_1 \le t < T_1 + N_2$, depends on the output $j_*$ of the first phase:
\begin{enumerate}
  \item[(a)] if $j_* = 0$: the action is $0$ and the estimate is $\hat r := \eps$;
  \item[(b)] if $j_* = J$: the rounds $T_1 + in + \ell$ ($i < d$, $\ell < n$) are
        \emph{basis rounds} playing $e_i$, and the estimate is the large-norm estimate
        $\hat r = \sqrt{\max(\hat R, 0)}$ of Definition~\ref{def:large_norm_algorithm} computed
        from their observations;
  \item[(c)] if $0 < j_* < J$: the rounds $T_1 + ks + \ell$ ($k < K(t_{j_*})$,
        $\ell < s(t_{j_*})$) form the Rademacher blocks of Algorithm~1 with $r_0 = t_{j_*}$
        (Definition~\ref{def:additive_algorithm}), and the estimate is
        $\hat r = \sqrt{\max(\bar Z, 0)}$ computed from their observations.
\end{enumerate}
The remaining rounds of the second phase play the action $0$ (their observations are not
used). The \emph{plan} $\mathrm{plan}(j_*, t, \sigma)$ is the action of round $t$ as a
function of $j_*$ and of the seeds $\sigma$ at the block starts; as a seeded algorithm, the
action of round $t$ is computed from the past actions and observations (which determine
$j_*$ once $t \ge T_1$, and the action of the current block start) and the fresh seed of
round $t$.
\end{definition}

\textbf{Lean remark.} The parameters are gathered in a structure \texttt{NormEstParam}
(\texttt{ε}, \texttt{δ} and their ranges). The schedule is encoded by two functions of
$(j_*, t)$: \texttt{rbBlock} returns the start of the Rademacher block of round $t$ (if any)
and \texttt{basisRound} the basis vector of round $t$ (if any); \texttt{planAction} is the
plan and \texttt{nextAction} its computation from the past, so that the seeded algorithm
\texttt{alg} has \texttt{next}$_n(h, u) = \mathrm{nextAction}(n, \mathrm{past}(h), j_*(h), u)$,
where $h$ is the history of the first $n$ rounds.
The estimate is the measurable function \texttt{estimate} of the observation sequence
(\texttt{output} on histories of length $N_{\mathrm{norm}}$), and the identification
algorithm is the fixed-budget algorithm with this deterministic output. The paper composes
the phases sequentially; here the composition is replaced by the deterministic schedule,
and the data-dependence is entirely in $j_*$.

\begin{lemma}[The run of the meta-algorithm follows its plan]\label{lem:norm_meta_plan}
\lean{Maiti2026Power.NormEstParam.seedAction_eq_planAction, Maiti2026Power.NormEstParam.yω_eq, Maiti2026Power.NormEstParam.testStat_eq_rbStat, Maiti2026Power.NormEstParam.addStat_eq_rbStat, Maiti2026Power.NormEstParam.lnRStat_eq_lnR, Maiti2026Power.NormEstParam.estimate_congr, Maiti2026Power.NormEstParam.estimate_yOf_seedFinHist, Maiti2026Power.NormEstParam.jsω_prefixExtend}\leanok
\uses{def:meta_algorithm, lem:seeded, lem:statistic_structure}
Run $\mathrm{NormEst}(\eps,\delta)$ against $\nu_\theta$ on the seed space
$(\Omega_{\mathrm{seed}}, \Prob_{\mathrm{seed}})$ of Lemma~\ref{lem:statistic_structure}, with
seeds $\sigma_t$ and noises $\eta_t$, and let $j_* = j_*(\omega)$ be the output of the first
phase. Then for every round $t$ the action is the plan action,
$x_t = \mathrm{plan}(j_*, t, \sigma)$, and $y_t = \ip{x_t}{\theta} + \eta_t$. Consequently:
\begin{enumerate}
  \item[(i)] for $j < J$, the statistic $U_j$ of test $j$ is $\bar Z$ evaluated at the window
    projection $\pi_{\mathrm{start}_j, s_j, K_j}(\omega)$;
  \item[(ii)] on $\{0 < j_* < J\}$, the statistic $\bar Z$ of the second phase is $\bar Z$
    evaluated at $\pi_{T_1, s(t_{j_*}), K(t_{j_*})}(\omega)$;
  \item[(iii)] on $\{j_* = J\}$, the statistic $\hat R$ of the second phase is the large-norm
    statistic evaluated at the basis window projection $\pi'_{T_1, n}(\omega)$;
  \item[(iv)] $j_*$ is a function of the prefix $(\omega_t)_{t<T_1}$, and the estimate is a
    function of the observations of the rounds $t < N_{\mathrm{norm}}$ only (so the estimate
    of the run is the estimate computed from the history $\hist_{N_{\mathrm{norm}}}$).
\end{enumerate}
\end{lemma}

\begin{proof}
\leanok
\uses{def:meta_algorithm, lem:seeded, lem:statistic_structure}
Strong induction on $t$. For $t < T_1$ the schedule does not depend on $j_*$, and the
action computed from the past is the plan action because the action of the block start
$b \le t$ is, by the induction hypothesis, $\sigma_b/\sqrt d$ (or $t = b$ and the fresh seed is
used). For $t \ge T_1$, the value of $j_*$ computed from the past observations
$(y_r)_{r<T_1}$ is the true $j_*$ (the first phase is complete), and the same argument
applies. The observation identity is the definition of the noise environment.
Parts (i)--(iii) follow from Lemma~\ref{lem:statistic_structure}(iii) and the
corresponding identities for the basis rounds ($\ip{e_i}{\theta} = \theta_i$); part (iv) from
the fact that $j_*$ only reads the rounds $t < T_1$, each of which only depends on the
coordinates $\omega_r$, $r \le t$, and that each statistic of the second phase only reads
rounds $t < T_1 + N_2$.
\end{proof}

\textbf{Lean remark.} The second phase is ``$\mathrm{alg}_2(h)$'' in the notation of
Lemma~\ref{lem:composition}, a measurable function of the first-phase history $h$ through
$r_0(h) \in \{t_0,\dots,t_J\}$; since $r_0$ takes finitely many values, measurability is a
finite case distinction. The estimate $\hat r$ is a deterministic measurable function of the
whole history with values in $[0,\infty)$, i.e.\ a recommendation kernel in the sense of
Definition~\ref{def:norm_estimator}. Padding with a constant action keeps the budget
deterministic, as required by Definition~\ref{def:bai_alg}; the padded rounds are simply
not used by the estimate.

\begin{theorem}[Norm estimation; Theorem~8 of the paper]\label{thm:norm_estimation}
\lean{Maiti2026Power.NormEstParam.fail2, Maiti2026Power.NormEstParam.fail_subset, Maiti2026Power.NormEstParam.nsMeasure_real_jsω_inter_preimage_eq_mul, Maiti2026Power.NormEstParam.nsMeasure_real_fail2_le, Maiti2026Power.NormEstParam.nsMeasure_real_estimate_fail_le, Maiti2026Power.NormEstParam.one_sub_le_seedMeasure_real_output, Maiti2026Power.NormEstParam.N₂_le, Maiti2026Power.NormEstParam.T_le}\leanok
\uses{def:meta_algorithm, def:norm_estimator, lem:seeded}
Let $\eps \in (0,1]$ and $\delta \in (0,1)$. The meta-algorithm $\mathrm{NormEst}(\eps,\delta)$
is an $(\eps,\delta)$-accurate estimation algorithm on $\ball_d$: for every
$\theta \in \R^d$, $\Prob_\theta(|\hat r - \norm{\theta}| \le \eps) \ge 1 - 7\delta/8$. Its
budget satisfies
\[
  N_{\mathrm{norm}}(\eps,\delta) \le C_{\mathrm{norm}}\,\frac{d\log(4/\delta)}{\eps^2},
  \qquad C_{\mathrm{norm}} := 10^5 .
\]
\end{theorem}

\begin{corollary}[Theorem~8 of the paper, existential form]\label{cor:norm_estimation_exists}
\lean{Maiti2026Power.exists_isAccurateNormEst}\leanok
\uses{thm:norm_estimation, def:norm_estimator}
Let $\eps \in (0,1]$ and $\delta \in (0,1)$. There exist a budget
$T \le C_{\mathrm{norm}}\, d\log(4/\delta)/\eps^2$ and an $(\eps,\delta)$-accurate estimation
algorithm on $\ball_d$ with budget $T$.
\end{corollary}
\begin{proof}
\leanok
\uses{thm:norm_estimation, lem:seeded}
Take $\mathrm{NormEst}(\eps,\delta)$ of Theorem~\ref{thm:norm_estimation}; it is
$(\eps, 7\delta/8)$-accurate, hence $(\eps,\delta)$-accurate. (For $d = 0$ every reward
vector is $0$ and the constant estimate $\eps$ with budget $0$ is accurate.)
\end{proof}

\begin{proof}
\leanok
\uses{thm:multiscale_guarantee, thm:additive_estimation, thm:large_norm_guarantee, lem:multiscale_samples, lem:norm_meta_plan, lem:statistic_structure, lem:norm_multiscale_good_event, lem:seeded}
\emph{Correctness.} By Lemma~\ref{lem:seeded} it suffices to prove the bound on the seed
space, where, by Lemma~\ref{lem:norm_meta_plan}(iv), the estimate of the run is the
estimate $\hat r(\omega)$ computed from the observation sequence. Let $G$ be the good event
of Theorem~\ref{thm:multiscale_guarantee}, $\Prob_{\mathrm{seed}}(G^c) \le 3\delta/8$. For
$1 \le j \le J$ let $F_j$ be the failure event of the second phase at outcome $j$: for
$j = J$, $F_J := \{ \eps < |\hat r' - r| \}$ where $\hat r'$ is the large-norm estimate
evaluated at the basis window projection $\pi'_{T_1,n}$; for $j < J$,
$F_j := \{ \eps/2 < |\hat r_j - r| \}$ where $\hat r_j$ is the estimate of Algorithm~1 with
$r_0 = t_j$ evaluated at $\pi_{T_1, s(t_j), K(t_j)}$. By Lemma~\ref{lem:norm_meta_plan}(ii)--(iii),
\[
  \{ \eps < |\hat r - r| \} \subseteq G^c \cup \bigcup_{1 \le j \le J}
    \big( G^c{}^c \cap \{j_* = j\} \cap F_j \big):
\]
on $G \cap \{j_* = 0\}$, $r < 2\eps$ by Theorem~\ref{thm:multiscale_guarantee}(a) and
$\hat r = \eps$, so $|\hat r - r| \le \eps$; on $\{j_* = j\}$ with $j \ge 1$ the estimate is
$\hat r_j$ or $\hat r'$. For each $j \ge 1$: if $G \cap \{j_* = j\} = \emptyset$ the term
vanishes; otherwise Theorem~\ref{thm:multiscale_guarantee}(b)--(c) gives the hypotheses of
Theorem~\ref{thm:additive_estimation} ($\eps \le t_j$, $r/2 < t_j \le 2r$) or of
Theorem~\ref{thm:large_norm_guarantee} ($r \ge \sqrt d$), and
$\Prob_{\mathrm{seed}}(F_j) \le 3\delta/8 \le \delta/2$ (resp.\ $\le \delta/2$) by
Lemma~\ref{lem:statistic_structure}(i)--(ii) (the law of the window projection is the
block, resp.\ noise, law). Since $\{j_* = j\}$ is a function of the prefix
$(\omega_t)_{t<T_1}$ (Lemma~\ref{lem:norm_meta_plan}(iv)) and $F_j$ is a function of a window
projection with $a = T_1$, they are independent (Lemma~\ref{lem:statistic_structure}), so
$\Prob_{\mathrm{seed}}(\{j_* = j\} \cap F_j) = \Prob_{\mathrm{seed}}(j_* = j)\,\Prob_{\mathrm{seed}}(F_j)
\le \Prob_{\mathrm{seed}}(j_* = j)\,\delta/2$. Summing over $j$ (the events $\{j_* = j\}$ are
disjoint, so their probabilities add up to at most $1$),
$\Prob_{\mathrm{seed}}(\eps < |\hat r - r|) \le 3\delta/8 + \delta/2 = 7\delta/8$.

\emph{Budget.} Lemma~\ref{lem:multiscale_samples} gives $N_{\mathrm{ms}} \le 63736\,dL/\eps^2$
with $L = \log(4/\delta)$. For $j < J$, $t_j \in [\eps, 2\sqrt d)$, so
Theorem~\ref{thm:additive_estimation} (budget clause with $r_0 = t_j$) gives
$s(t_j)K(t_j) \le 32776\,dL/\eps^2$, and Theorem~\ref{thm:large_norm_guarantee} gives
$dn \le 49\,dL/\eps^2$; hence $N_2 \le 32776\,dL/\eps^2$ and
$N_{\mathrm{norm}} \le (63736 + 32776)\,dL/\eps^2 \le 10^5\,dL/\eps^2$.
\end{proof}

\begin{remark}[Reading of the constants]\label{rem:norm_constants}
The budget is $O(d\log(1/\delta)/\eps^2)$ as in the paper: for $\delta \le 1/4$,
$\log(4/\delta) \le 2\log(1/\delta)$. The constant $C_{\mathrm{norm}} = 10^5$ is dominated by the
multi-scale phase ($63736$), itself driven by the test constant $c_1' = 5120$; the
constants are what the union bounds produce and were not optimized. In the paper the
sample complexity of the multi-scale phase is written $\sum_j s_jt_j$, a typo for
$\sum_j s_jK_j$. The paper's failure probabilities are $\delta/2$ for the multi-scale phase
and $\delta/2$ for the second phase; the sharper $3\delta/8$ of
Lemma~\ref{lem:norm_multiscale_good_event} is the paper's own computation
$\sum_j 3\delta_j/4$.
\end{remark}

\chapter{A polynomial separation between adaptive and non-adaptive algorithms}
\label{chap:separation}

This chapter formalizes Section~2.5 and Appendix~I of the paper: an action set
$\Xs \subset \R^{kd}$, a union of $k$ unit balls living in $k$ disjoint blocks of $d$
coordinates, on which every non-adaptive algorithm needs $\Omega(kd^2/\eps^2)$ samples while
an adaptive algorithm succeeds with $O(kd\log(k/\delta)/\eps^2 + d^2/\eps^2)$ samples. The
target is Theorem~7 of the paper (Theorem~\ref{thm:separation}). The adaptive algorithm
first runs the norm estimation procedure of Chapter~\ref{chap:norm} on each block to locate
the block containing a near-optimal arm, then identifies a good arm in that block with the
basis design of Lemma~\ref{lem:ball_identification}; as in Chapter~\ref{chap:norm}, the
phases are composed by a deterministic round schedule of a seeded algorithm
(Lemma~\ref{lem:seeded}), and the analysis takes place on the seed space. The chapter is
fully formalized. The non-adaptive lower bound deviates from the paper: instead of
the cited adaptive lower bound of Chen et al.\ \cite{chen2024assouad} we project a fixed
design on the block receiving the fewest points (Lemma~\ref{lem:block_restriction}, an
instance of the transport Lemma~\ref{lem:fixed_design_transport}) and apply the ingredients
of our own non-adaptive lower bound (the Bayesian step,
Lemma~\ref{lem:lower_nonadaptive_bayes_step}) with the width of the unit ball for a matrix
(Proposition~\ref{prop:width_ball}, $\gwidth{\ball_d}{A} \ge \sqrt{2/\pi}\,d/\sqrt{\tr A}$). The class of
non-adaptive algorithms is therefore that of Definition~\ref{def:fixed_design}: deterministic
fixed designs with an arbitrary (possibly randomized) recommendation kernel, whereas the
paper's Theorem~7 speaks of ``possibly randomized non-adaptive'' algorithms
(Remark~\ref{rem:separation_gap}). The block-ball set is spanning
(Lemma~\ref{lem:block_ball_basic}), which is the hypothesis needed wherever
Theorem~\ref{thm:lower_nonadaptive}, Theorem~\ref{thm:lower_adaptive} or the Gaussian width
(Definition~\ref{def:width}, Proposition~\ref{prop:width_ball} on the spanning set $\ball_d$)
is used below. Throughout,
$k, d \ge 1$ are integers with $kd \ge 2$; the comparison of the two bounds uses
$d \le k \le d^2$.

\section{The block-ball action set}
\label{sec:separation_set}

\begin{definition}[Block-ball action set]\label{def:block_ball_set}
\lean{Maiti2026Power.blockBallSet}\leanok
\uses{def:arm_set}
For $i \in \{1,\dots,k\}$ let $B_i := \{(i-1)d + 1, \dots, id\}$ be the $i$-th block of
coordinates of $\R^{kd}$, let $\iota_i : \R^d \to \R^{kd}$ be the isometric embedding which
places a vector on the coordinates of $B_i$ and $0$ elsewhere, and let
$\pi_i : \R^{kd} \to \R^d$ be the projection on the coordinates of $B_i$, so that
$\pi_i\circ\iota_i = \mathrm{id}$, $\ip{\iota_i u}{v} = \ip{u}{\pi_i v}$ and
$\norm{\iota_i u} = \norm u$. Define
\[
  \Xs_i := \{ x \in \R^{kd} : \supp(x) \subseteq B_i,\ \norm{x} \le 1 \} = \iota_i(\ball_d),
  \qquad
  \Xs := \bigcup_{i=1}^k \Xs_i .
\]
For $\theta \in \R^{kd}$ we write $\theta^{(i)} := \pi_i\theta \in \R^d$ for its $i$-th block,
so $\theta = \sum_i\iota_i\theta^{(i)}$ and $\norm{\theta}^2 = \sum_i\norm{\theta^{(i)}}^2$.
\end{definition}

\textbf{Lean remark.} With $\R^{kd}$ encoded as \texttt{EuclideanSpace ℝ (Fin k × Fin d)}
the blocks are $\{i\}\times\mathrm{Fin}\ d$ (0-indexed) and $\iota_i$, $\pi_i$ are the
obvious linear maps; $\Xs_i \cap \Xs_j = \{0\}$ for $i \ne j$, so an element of $\Xs$
determines its block unless it is $0$. Where a block index is needed for a point of $\Xs$
we use the smallest one.

\begin{lemma}[Basic properties of the block-ball set]\label{lem:block_ball_basic}
\lean{Maiti2026Power.blockEmb, Maiti2026Power.blockProj, Maiti2026Power.inner_blockEmb_left, Maiti2026Power.norm_blockEmb, Maiti2026Power.norm_blockProj_le, Maiti2026Power.blockProj_blockEmb_same, Maiti2026Power.blockProj_blockEmb_of_ne, Maiti2026Power.exists_blockProj_eq_zero, Maiti2026Power.blockBallSet_eq_iUnion, Maiti2026Power.isCompact_blockBallSet, Maiti2026Power.blockBallSet_subset_unitBall, Maiti2026Power.span_blockBallSet, Maiti2026Power.supportFn_blockBallSet, Maiti2026Power.supportFn_blockBallSet_blockEmb, Maiti2026Power.simpleRegret_blockBallSet_blockEmb, Maiti2026Power.supportFn_unitBall, Maiti2026Power.simpleRegret_unitBall}\leanok
\uses{def:block_ball_set, def:regret}
$\Xs = \bigcup_i\iota_i(\ball_d)$ is a compact subset of $\ball_{kd}$ with $\Span(\Xs) = \R^{kd}$,
hence an action set in the sense of Definition~\ref{def:arm_set}. Every $x \in \Xs$ satisfies
$\pi_j x = 0$ for all $j$ but one. For every $\theta \in \R^{kd}$,
\[
  \max_{x \in \Xs}\ip{x}{\theta} = \max_{1\le i\le k}\norm{\theta^{(i)}},
\]
and for $\vartheta \in \R^d$ and $\rec \in \R^{kd}$,
$\max_{x\in\Xs}\ip{x}{\iota_i\vartheta} = \norm{\vartheta} = \max_{u \in \ball_d}\ip{u}{\vartheta}$
and
\[
  r_{\Xs}(\rec, \iota_i\vartheta) = \norm{\vartheta} - \ip{\pi_i\rec}{\vartheta}
  = r_{\ball_d}(\pi_i\rec, \vartheta) .
\]
\end{lemma}

\begin{proof}
\leanok
\uses{def:block_ball_set}
Each $\iota_i(\ball_d)$ is the continuous image of a compact set, and a finite union
of compact sets is compact; $x \in \Xs$ lies in $\iota_i(\ball_d)$ for its block $i$, with
$x = \iota_i\pi_i x$. $\Xs$ contains every basis vector $e_m$ ($m \in B_i$ for some
$i$), so it spans $\R^{kd}$; and $\norm{\iota_i u} = \norm u \le 1$ gives $\Xs \subseteq \ball_{kd}$.
For $x = \iota_i u$ with $u \in \ball_d$, $\ip{x}{\theta} = \ip{u}{\theta^{(i)}} \le \norm{\theta^{(i)}}$
by Cauchy--Schwarz, with equality for $u = \theta^{(i)}/\norm{\theta^{(i)}}$ (any $u$ if
$\theta^{(i)} = 0$). Taking the maximum over $i$ gives the second identity; for
$\theta = \iota_i\vartheta$ the blocks are $\theta^{(i)} = \vartheta$ and $\theta^{(j)} = 0$
for $j \ne i$, and the same Cauchy--Schwarz argument on $\ball_d$ gives
$\max_{u\in\ball_d}\ip{u}{\vartheta} = \norm\vartheta$. The regret formula follows from
Definition~\ref{def:regret} and $\ip{\rec}{\iota_i\vartheta} = \ip{\pi_i\rec}{\vartheta}$.
\end{proof}

\section{Lower bound for non-adaptive algorithms}
\label{sec:separation_lower}

\begin{lemma}[Restriction of a fixed design to a block]\label{lem:block_restriction}
\lean{Maiti2026Power.blockProjBall, Maiti2026Power.blockDesign, Maiti2026Power.blockRounds, Maiti2026Power.isPAC_fixedDesignTransport_blockProjBall, Maiti2026Power.trace_sum_outerSelf_blockDesign_le, Maiti2026Power.pairwiseDisjoint_blockRounds, Maiti2026Power.sum_card_blockRounds_le, Maiti2026Power.exists_mul_card_blockRounds_le}\leanok
\uses{def:fixed_design, def:block_ball_set, lem:block_ball_basic, def:pac, lem:fixed_design_transport}
Let $x_0,\dots,x_{T-1} \in \Xs$ be a fixed design with recommendation kernel $\rho$ (from
$\R^T$ to $\Xs$) which is $(\eps,\delta)$-PAC on $\Xs$, and fix a block $i$. The projected
design $x'_t := \pi_i x_t \in \ball_d$ ($t < T$), with the recommendation kernel
$(\pi_i)_*\rho$ (draw $\rec \sim \rho(y)$, the history being relabelled with the actions $x_t$,
and output $\pi_i\rec \in \ball_d$), is a fixed-design algorithm with budget $T$ on $\ball_d$
which is $(\eps,\delta)$-PAC on $\ball_d$. Its design matrix
$\Sigma'_T := \sum_{t<T}(\pi_ix_t)(\pi_ix_t)^\top$ satisfies $\tr\Sigma'_T \le T_i$, where
$T_i := \abs{\{t < T : \pi_i x_t \ne 0\}}$ is the number of rounds played in block $i$; the sets
$\{t < T : \pi_ix_t \ne 0\}$ are pairwise disjoint, so $\sum_i T_i \le T$ and some block $i^*$
has $kT_{i^*} \le T$.
\end{lemma}

\begin{proof}
\leanok
\uses{lem:fixed_design_transport, lem:block_ball_basic}
This is Lemma~\ref{lem:fixed_design_transport} with $L = \iota_i$ and $\pi = \pi_i$: by
Lemma~\ref{lem:block_ball_basic}, $\pi_ix_t \in \ball_d$ (as $\norm{\pi_i x} \le \norm x \le 1$),
$\ip{\pi_ix_t}{\vartheta} = \ip{x_t}{\iota_i\vartheta}$ (adjointness) and
$r_{\ball_d}(\pi_i\rec,\vartheta) = r_{\Xs}(\rec,\iota_i\vartheta)$ for every $\vartheta \in \R^d$.
The trace bound is $\tr\Sigma'_T = \sum_{t<T}\norm{\pi_ix_t}^2 \le T_i$ since each term is at
most $1$ and vanishes outside the rounds played in block $i$. The sets of rounds are disjoint
because each $x_t$ lies in a single block (Lemma~\ref{lem:block_ball_basic}), so
$\sum_i T_i \le T$, and the pigeonhole principle gives $i^*$ with $kT_{i^*} \le T$.
\end{proof}

\textbf{Lean remark.} The paper (and an earlier version of this blueprint) restricts the
design to the rounds $t$ with $x_t$ in block $i$; keeping all $T$ rounds with the projected
actions $\pi_ix_t$ (which are $0$ outside the block and contribute pure noise) gives a
fixed-design algorithm with the same budget $T$, which is exactly the situation of
Lemma~\ref{lem:fixed_design_transport}, and the number $T_i$ of rounds in the block enters
only through the trace of the design matrix. The composition of $\rho$ with the relabelling of
the history and with $\pi_i$ is where the randomness allowed in Definition~\ref{def:bai_alg}
is used.

\begin{theorem}[Non-adaptive lower bound for the block-ball set]\label{thm:separation_lower_nonadaptive}
\lean{Maiti2026Power.blockBallSet_le_budget_of_isFixedDesign_of_isPAC, Maiti2026Power.blockBallSet_le_budget_of_isPAC}\leanok
\uses{def:block_ball_set, lem:block_ball_basic, def:fixed_design, def:pac}
Let $\eps \in (0,1]$. Recall that $\Xs$ is a spanning action set in $\R^{kd}$
(Lemma~\ref{lem:block_ball_basic}).
\begin{enumerate}
  \item[(i)] If a fixed-design algorithm (Definition~\ref{def:fixed_design}: a deterministic
        design $x_0,\dots,x_{T-1} \in \Xs$ with an arbitrary, possibly randomized,
        recommendation kernel; budget $T \ge 0$) is $(\eps,\delta)$-PAC on $\Xs$ with
        $\delta \le 1/10$, then
        \[
          T \ge \frac{k d^2}{210\,\eps^2} .
        \]
  \item[(ii)] If $kd \ge 2$ and any identification algorithm (adaptive or not) with budget $T$
        is $(\eps,\delta)$-PAC on $\Xs$ with $\delta < 1/16$, then
        $T \ge kd\log(1/\delta)/(20000\,\eps^2)$ (and $T \ge c_{\mathrm{ad}}\,kd\log(1/\delta)/\eps^2$
        with the constant $c_{\mathrm{ad}}$ of Theorem~\ref{thm:lower_adaptive}).
\end{enumerate}
\end{theorem}
In particular every fixed-design $(\eps,\delta)$-PAC algorithm with $\delta < 1/16$ has
budget at least
$\tfrac12\big( c_{\mathrm{ad}}\,kd\log(1/\delta) + kd^2/210 \big)/\eps^2$.

\begin{proof}
\leanok
\uses{lem:block_restriction, lem:lower_nonadaptive_posdef, lem:lower_nonadaptive_bayes_step, prop:width_ball, thm:lower_adaptive, lem:block_ball_basic}
(i) If $k = 0$ or $d = 0$ there is nothing to prove, so assume $k, d \ge 1$. With the notation
of Lemma~\ref{lem:block_restriction}, pick a block $i^*$ with $kT_{i^*} \le T$ and consider
the projected design $x'_t = \pi_{i^*}x_t$ on $\ball_d$, with design matrix $\Sigma'_T$, which
is $(\eps,\delta)$-PAC on $\ball_d$ with budget $T$.

\emph{Step 1: $\Sigma'_T \succ 0$.} Since $0 \in \ball_d$, $\ball_d$ is spanning and
$\delta \le 1/10 < 1/2$, Lemma~\ref{lem:lower_nonadaptive_posdef} applies to the projected
design (with any budget $T \ge 0$, including $T = 0$). In particular $T_{i^*} \ge 1$ (as
$d \ge 1$ and $\tr\Sigma'_T \le T_{i^*}$, with $\tr\Sigma'_T > 0$).

\emph{Step 2: the Bayesian step.} Lemma~\ref{lem:lower_nonadaptive_bayes_step} on $\ball_d$
gives $0.12\,\gwidth{\ball_d}{\Sigma'_T} \le \eps$.

\emph{Step 3: the width of the ball.} By Proposition~\ref{prop:width_ball},
$\gwidth{\ball_d}{\Sigma'_T} \ge \sqrt{2/\pi}\,d/\sqrt{\tr\Sigma'_T} \ge \sqrt{2/\pi}\,d/\sqrt{T_{i^*}}$.

\emph{Step 4: arithmetic.} Combining, $0.12\sqrt{2/\pi}\,d \le \eps\sqrt{T_{i^*}}$; squaring,
$0.0144\cdot(2/\pi)\,d^2 \le \eps^2T_{i^*}$, and since $\pi \le 4$, $2/\pi \ge 1/2$ and
$0.0072\,d^2 \le \eps^2 T_{i^*}$, i.e.\ $T_{i^*} \ge d^2/(139\,\eps^2) \ge d^2/(210\,\eps^2)$.
Hence $T \ge kT_{i^*} \ge kd^2/(210\eps^2)$.

(ii) $\Xs$ is a spanning action set in $\R^{kd}$ with $kd \ge 2$
(Lemma~\ref{lem:block_ball_basic}), so Theorem~\ref{thm:lower_adaptive} applies directly with
dimension $kd$.

The last claim is $T \ge \max(a,b) \ge (a+b)/2$.
\end{proof}

\section{The adaptive block algorithm}
\label{sec:separation_algorithm}

\begin{lemma}[Transporting an algorithm on $\ball_d$ to a block]\label{lem:separation_block_embedding}
\lean{Maiti2026Power.blockEmbBall, Maiti2026Power.blockProjBall_blockEmbBall, Maiti2026Power.BlockParam.F_blockEmbBall, Maiti2026Power.BlockParam.seedStep_block, Maiti2026Power.BlockParam.yω_block, Maiti2026Power.BlockParam.estN_eq, Maiti2026Power.shiftSeq, Maiti2026Power.nsMeasure_map_shiftSeq}\leanok
\uses{def:block_ball_set, def:env, def:bai_alg, def:norm_estimator, lem:seeded, def:meta_algorithm}
Let $\iota_i : \ball_d \to \Xs$ be the embedding of the unit ball on block $i$ (so
$\pi_i \circ \iota_i = \mathrm{id}$). For every $\theta \in \R^{kd}$ and $x \in \ball_d$, the
reward kernel of $\nu_\theta$ at $\iota_i x$ is that of $\nu_{\theta^{(i)}}$ at $x$:
$\ip{\iota_i x}{\theta} + \eta = \ip{x}{\theta^{(i)}} + \eta$. Consequently, the block
algorithm of Definition~\ref{def:block_algorithm} \emph{simulates} the norm estimation
meta-algorithm on each block: on the seed space, for $i < k$ and $t < N$, the step
(action, observation) of round $iN + t$ of the block algorithm against $\nu_\theta$ is
$(\iota_i x_t, y_t)$ where $(x_t, y_t)$ is the step of round $t$ of
$\mathrm{NormEst}(\eps/4, \delta/(2k))$ against $\nu_{\theta^{(i)}}$ on the seed sequence shifted
by $iN$, $(\omega_{iN + t})_t$. In particular the estimate $\hat r_i$ of block $i$ is the
estimate of the meta-algorithm against $\nu_{\theta^{(i)}}$ on the shifted seed sequence, and the
shifted seed sequence has the law $\Prob_{\mathrm{seed}}$ of the seed sequence.
\end{lemma}

\begin{proof}
\leanok
\uses{def:env, lem:block_ball_basic, lem:seeded, lem:norm_meta_plan, lem:statistic_structure}
The kernel identity is Lemma~\ref{lem:block_ball_basic} ($\ip{\iota_i x}{\theta} = \ip{x}{\pi_i\theta}$).
The simulation is a strong induction on $t < N$: the action of round $iN + t$ of the block
algorithm is $\iota_i$ applied to the next action of the meta-algorithm computed from the past
actions (projected by $\pi_i$) and observations of the rounds $iN + r$, $r \le t$, of the
block, which by the induction hypothesis are the past actions and observations of the
meta-algorithm on the shifted seed sequence; the observation follows from the kernel
identity. The shift invariance of $\Prob_{\mathrm{seed}}$ is Lemma~\ref{lem:statistic_structure}
(the coordinates $(\omega_{a + t})_t$ of an i.i.d.\ sequence form an i.i.d.\ sequence with the
same law).
\end{proof}

\textbf{Lean remark.} The block algorithm is not obtained by transporting the meta-algorithm
as an LML algorithm (which would require mapping its policy kernels by $\iota_i$ and its
histories by $\pi_i$): it is a seeded algorithm whose next-action map, in the rounds of block
$i$, applies $\iota_i$ to the next-action map of the meta-algorithm evaluated on the
$\pi_i$-projected history of the block. The simulation lemma is then a statement about two
recursively defined runs on the same seed space.

\begin{proof}
\uses{def:env, lem:block_ball_basic}
The reward kernel of $\nu_\theta$ at $\iota_i(x)$ is
$\Normal(\ip{\iota_ix}{\theta}, 1) = \Normal(\ip{x}{\theta^{(i)}}, 1)$
(Lemma~\ref{lem:block_ball_basic}), which is the reward kernel of $\nu_{\theta^{(i)}}$ at $x$.
The policy of $\iota_i\mathrm{alg}$ is the image by $\iota_i$ of the policy of $\mathrm{alg}$
applied to the projected history, and $\pi_i\circ\iota_i = \mathrm{id}$, so the projected
process satisfies the defining conditional-law conditions of an algorithm-environment
sequence for $(\mathrm{alg},\nu_{\theta^{(i)}})$ (LML \texttt{IsAlgEnvSeq}); uniqueness of the
trajectory law gives the claim about estimates.
\end{proof}

\textbf{Lean remark.} This is the analogue, for the non-invertible isometric embedding
$\iota_i$, of the transformation lemma \ref{lem:transform_preserves_pac} used for the
adaptive lower bound; it is a statement about mapping actions by a measurable map $\phi$
such that the environments correspond ($\kappa_\theta \circ \phi = \kappa_{\theta'}$), which
holds for any stationary environments and could be stated at that level of generality. The
map $\phi : \ball_d \to \Xs$, $u \mapsto \iota_i u$, is continuous hence measurable between
the subtypes, and the policy kernels of $\iota_i\mathrm{alg}$ are the images of those of
$\mathrm{alg}$ by $\phi$ (\texttt{Kernel.map}), applied to the history mapped by $\pi_i$ on
its action coordinates; the action type of $\iota_i\mathrm{alg}$ is the subtype $\Xs$, as
required to compose it with other algorithms on $\Xs$ (Lemma~\ref{lem:composition}).

\begin{definition}[Block algorithm]\label{def:block_algorithm}
\lean{Maiti2026Power.BlockParam, Maiti2026Power.BlockParam.P, Maiti2026Power.BlockParam.N, Maiti2026Power.BlockParam.n₂, Maiti2026Power.BlockParam.T₂, Maiti2026Power.BlockParam.T, Maiti2026Power.BlockParam.blockOf, Maiti2026Power.BlockParam.basisRound, Maiti2026Power.BlockParam.estN, Maiti2026Power.BlockParam.iHat, Maiti2026Power.BlockParam.thetaHat, Maiti2026Power.normalizeBall, Maiti2026Power.BlockParam.recommend, Maiti2026Power.BlockParam.nextAction, Maiti2026Power.BlockParam.alg, Maiti2026Power.BlockParam.output}\leanok
\uses{def:block_ball_set, def:meta_algorithm, lem:ball_identification, lem:seeded}
Let $\eps \in (0,1]$ and $\delta \in (0,1)$, let $N := N_{\mathrm{norm}}(\eps/4, \delta/(2k))$ be
the budget of the meta-algorithm $\mathrm{NormEst}(\eps/4, \delta/(2k))$
(Definition~\ref{def:meta_algorithm}) and $n_2 := \lceil 4(8d + 48\log(2/\delta))/\eps^2 \rceil$.
The \emph{block algorithm} on $\Xs$ is the seeded algorithm (Lemma~\ref{lem:seeded}) with
uniform sign-vector seeds and budget $N_{\mathrm{sep}}(\eps,\delta) := kN + dn_2$ defined by the
following round schedule.
\begin{enumerate}
  \item[(1)] For $i = 0, \dots, k-1$, the rounds $[iN, (i+1)N)$ run
        $\mathrm{NormEst}(\eps/4, \delta/(2k))$ on block $i$: the action of round $iN + t$ is
        $\iota_i$ applied to the action of round $t$ of the meta-algorithm computed from the
        ($\pi_i$-projected) history of the block and the seed of the round. The estimate
        $\hat r_i$ is the estimate of the meta-algorithm computed from the observations of the
        rounds of block $i$, and $\hat i := \argmax_i\hat r_i$ (smallest index in case of ties).
  \item[(2)] The rounds $kN + mn_2 + \ell$, $m < d$, $\ell < n_2$, play $\iota_{\hat i}(e_m)$
        (the basis design of Lemma~\ref{lem:ball_identification} on block $\hat i$); let
        $\hat\vartheta \in \R^d$ be the vector of the $d$ empirical means of the observations
        at each $e_m$, and recommend $\rec := \iota_{\hat i}\big(\hat\vartheta/\norm{\hat\vartheta}\big)$
        (with $\hat\vartheta/\norm{\hat\vartheta} := 0$ if $\hat\vartheta = 0$).
\end{enumerate}
\end{definition}

\textbf{Lean remark.} The parameters are gathered in a structure \texttt{BlockParam}
($k \ge 1$, $\eps$, $\delta$), with the meta-algorithm parameters
\texttt{BlockParam.P} $= (\eps/4, \delta/(2k))$. The block of a round $t < kN$ is
$\lfloor t/N \rfloor$ (\texttt{blockOf}), the basis vector of a round of phase (2) is given by
\texttt{basisRound}. The estimate of block $m$ (\texttt{estN}) is the estimate of
Definition~\ref{def:meta_algorithm} applied to the observations $y_{mN + r}$; the selected block
$\hat i$ (\texttt{iHat}, defined by \texttt{Nat.find}) and the empirical mean vector
(\texttt{thetaHat}) are measurable functions of the observation sequence, hence so is the
recommendation (\texttt{recommend}, \texttt{output}); the normalization $v \mapsto v/\norm v$
(\texttt{normalizeBall}, with $0 \mapsto 0$) is measurable but not continuous. The seeded
algorithm \texttt{alg} has next-action map \texttt{nextAction}. The recommendation $0$ when
$\hat\vartheta = 0$ replaces the paper's $e_1$; both have regret at most $2\norm{\hat\vartheta - \vartheta}$.

\textbf{Lean remark.} Phase (1) is itself a $k$-fold composition of algorithms with action
type $\Xs$ (each $\iota_i\mathrm{NormEst}$ has action type $\Xs$ by
Lemma~\ref{lem:separation_block_embedding}); the index $\hat i$ is a
measurable function of the phase-(1) history $h$ with values in $\{1,\dots,k\}$, and phase (2)
is the fixed design of Lemma~\ref{lem:ball_identification} transported by $\iota_{\hat i}$,
again with action type $\Xs$: a deterministic algorithm parametrized by $\hat i$, which is
the ``$\mathrm{alg}_2(h)$'' of Lemma~\ref{lem:composition}. The map
$h \mapsto \iota_{\hat i(h)}\mathrm{alg}$ is measurable in $h$ because $\hat i$ takes finitely
many values: on each of the measurable sets $\{\hat i = i\}$ it is the constant algorithm
$\iota_i\mathrm{alg}$, so its policy kernels are finite case distinctions of kernels
(\texttt{Kernel.piecewise}), and Lemma~\ref{lem:composition} applies.

\begin{lemma}[Selection of a good block]\label{lem:block_selection}
\lean{Maiti2026Power.BlockParam.estN_le_iHat, Maiti2026Power.norm_sub_inner_normalizeBall_le, Maiti2026Power.BlockParam.regret_le}\leanok
\uses{def:block_algorithm, lem:block_ball_basic}
Let $\theta \in \R^{kd}$ and let $y$ be an observation sequence such that
$\abs{\hat r_i - \norm{\theta^{(i)}}} \le \eps/4$ for all $i < k$, where $\hat r_i$ is the estimate
of block $i$ computed from $y$. Then, with $\hat i$ the selected block,
\[
  \max_{i<k}\norm{\theta^{(i)}} \le \norm{\theta^{(\hat i)}} + \frac{\eps}{2} ;
\]
and if moreover the empirical mean vector $\hat\vartheta$ computed from $y$ satisfies
$\norm{\hat\vartheta - \theta^{(\hat i)}} \le \eps/4$, then the recommendation
$\rec = \iota_{\hat i}(\hat\vartheta/\norm{\hat\vartheta})$ has simple regret
$r(\rec, \theta) \le \eps$ on $\Xs$.
\end{lemma}

\begin{proof}
\leanok
\uses{def:block_algorithm, lem:block_ball_basic}
With $i_* \in \argmax_i\norm{\theta^{(i)}}$ and $\hat r_{i_*} \le \hat r_{\hat i}$ (definition of
$\hat i$),
$\norm{\theta^{(\hat i)}} \ge \hat r_{\hat i} - \eps/4 \ge \hat r_{i_*} - \eps/4
\ge \norm{\theta^{(i_*)}} - \eps/2$. For the regret, by Lemma~\ref{lem:block_ball_basic} the
support function of $\Xs$ at $\theta$ is $\max_i\norm{\theta^{(i)}}$ and
$\ip{\iota_{\hat i}\rec'}{\theta} = \ip{\rec'}{\theta^{(\hat i)}}$ with
$\rec' := \hat\vartheta/\norm{\hat\vartheta}$; and for every $\vartheta, v \in \R^d$,
$\norm{\vartheta} - \ip{v/\norm v}{\vartheta} \le 2\norm{v - \vartheta}$ (if $v \ne 0$,
$\ip{v/\norm v}{\vartheta} = \norm v - \ip{v/\norm v}{v - \vartheta} \ge \norm v - \norm{v - \vartheta}
\ge \norm\vartheta - 2\norm{v - \vartheta}$; if $v = 0$ the left side is
$\norm\vartheta = \norm{v - \vartheta}$). Hence
$r(\rec,\theta) \le \big(\norm{\theta^{(\hat i)}} + \eps/2\big) - \big(\norm{\theta^{(\hat i)}} - 2\cdot\eps/4\big) = \eps$.
\end{proof}

\begin{proof}
\uses{thm:norm_estimation, lem:separation_block_embedding, lem:composition}
By Lemma~\ref{lem:separation_block_embedding} and Theorem~\ref{thm:norm_estimation}, the
run on block $i$ is a run of $\mathrm{NormEst}(\eps/4,\delta/(2k))$ against $\nu_{\theta^{(i)}}$ on
$\ball_d$, so $\Prob(|\hat r_i - \norm{\theta^{(i)}}| > \eps/4) \le \delta/(2k)$; by
Lemma~\ref{lem:composition} this bound holds conditionally on the history of the earlier
blocks, hence unconditionally. A union bound over the $k$ blocks gives the first claim with
probability at least $1 - \delta/2$. On this event, with $i_* \in \argmax_i\norm{\theta^{(i)}}$,
\[
  \norm{\theta^{(\hat i)}} \ge \hat r_{\hat i} - \frac\eps4 \ge \hat r_{i_*} - \frac\eps4
  \ge \norm{\theta^{(i_*)}} - \frac\eps2 .
\]
\end{proof}

\begin{theorem}[Polynomial separation; Theorem~7 of the paper]\label{thm:separation}
\lean{Maiti2026Power.BlockParam.seedAction_of_T₂_le, Maiti2026Power.BlockParam.thetaHat_eq, Maiti2026Power.BlockParam.recommend_yOf_seedFinHist, Maiti2026Power.BlockParam.one_sub_le_seedMeasure_real_output, Maiti2026Power.BlockParam.T_le}\leanok
\uses{def:block_algorithm, def:pac, thm:separation_lower_nonadaptive, lem:seeded}
Let $\eps \in (0,1]$ and $\delta \in (0,1)$. The block algorithm is $(\eps,\delta)$-PAC on $\Xs$
(with failure probability at most $15\delta/16$), with budget
\[
  N_{\mathrm{sep}}(\eps,\delta)
  \le \frac{1600192\,kd\log(8k/\delta) + 33\,d^2}{\eps^2}
  \le C_{\mathrm{sep}}\,\frac{kd\log(8k/\delta) + d^2}{\eps^2},
  \qquad C_{\mathrm{sep}} := 2\cdot 10^6 .
\]
If $d \le k$, this is at most $2C_{\mathrm{sep}}\,kd\big(\log(1/\delta) + \log(8k)\big)/\eps^2$,
whereas by Theorem~\ref{thm:separation_lower_nonadaptive} every non-adaptive
$(\eps,\delta)$-PAC algorithm with $\delta < 1/16$ needs at least
$\tfrac12\big(c_{\mathrm{ad}}\,kd\log(1/\delta) + kd^2/210\big)/\eps^2$ samples. Here
``non-adaptive'' means a deterministic fixed design with an arbitrary (possibly randomized)
recommendation kernel, the class of Definition~\ref{def:fixed_design} and of
Theorem~\ref{thm:lower_nonadaptive}; see Remark~\ref{rem:separation_gap}.
\end{theorem}

\begin{corollary}[Theorem~7 of the paper, existential form of the upper bound]\label{cor:separation_exists}
\lean{Maiti2026Power.exists_isPAC_blockBallSet}\leanok
\uses{thm:separation, def:pac}
Let $k \ge 1$, $\eps \in (0,1]$ and $\delta \in (0,1)$. There exist a budget
$T \le C_{\mathrm{sep}}\,(kd\log(8k/\delta) + d^2)/\eps^2$ and an identification algorithm with
budget $T$ that is $(\eps,\delta)$-PAC on the block-ball set $\Xs \subset \R^{kd}$.
\end{corollary}
\begin{proof}
\leanok
\uses{thm:separation, lem:seeded}
Take the block algorithm of Theorem~\ref{thm:separation}. (For $d = 0$ every reward vector is
$0$, every recommendation is optimal, and the budget $0$ suffices.)
\end{proof}

\begin{proof}
\leanok
\uses{lem:block_selection, lem:ball_identification, lem:separation_block_embedding, lem:seeded, lem:block_ball_basic, thm:norm_estimation, thm:separation_lower_nonadaptive, lem:statistic_structure}
\emph{Correctness.} By Lemma~\ref{lem:seeded} it suffices to prove the bound on the seed space
$(\Omega_{\mathrm{seed}}, \Prob_{\mathrm{seed}})$ of Lemma~\ref{lem:statistic_structure}, where
the recommendation of the run is the recommendation computed from the observation sequence
(it only reads the rounds $t < N_{\mathrm{sep}}$). Consider the two failure events
\[
  E_1 := \bigcup_{i<k}\Big\{ \abs{\hat r_i - \norm{\theta^{(i)}}} > \frac\eps4 \Big\},
  \qquad
  E_2 := \Big\{ \norm{\Delta} > \frac\eps4 \Big\},
  \quad \Delta := \pi'_{kN, n_2}\text{-average of the noises},
\]
where $\Delta \in \R^d$ has coordinates $\Delta_m := \frac1{n_2}\sum_{\ell<n_2}\eta_{kN + mn_2 + \ell}$.
By Lemma~\ref{lem:separation_block_embedding}, $\hat r_i$ is the estimate of
$\mathrm{NormEst}(\eps/4, \delta/(2k))$ against $\nu_{\theta^{(i)}}$ on the shifted seed sequence,
whose law is $\Prob_{\mathrm{seed}}$, so by Theorem~\ref{thm:norm_estimation} each event of the
union has probability at most $\tfrac78\cdot\delta/(2k)$ and $\Prob_{\mathrm{seed}}(E_1) \le 7\delta/16$.
In phase (2) the actions are $\iota_{\hat i}(e_m)$, so the empirical mean vector is
$\hat\vartheta = \theta^{(\hat i)} + \Delta$, where $\Delta$ does not depend on $\hat i$; the law of
the noises of the basis window is $\Normal(0,1)^{\otimes dn_2}$
(Lemma~\ref{lem:statistic_structure}), so by the chi-square tail bound of
Lemma~\ref{lem:ball_identification} with $c = \eps/4$ and the choice of $n_2$
($n_2\eps^2/16 \ge 2d + 12\log(2/\delta)$), $\Prob_{\mathrm{seed}}(E_2) \le \delta/2$. Outside
$E_1 \cup E_2$, Lemma~\ref{lem:block_selection} gives $r(\rec,\theta) \le \eps$. Hence
$\Prob_{\mathrm{seed}}(r(\rec,\theta) \le \eps) \ge 1 - 7\delta/16 - \delta/2 = 1 - 15\delta/16 \ge 1 - \delta$.
No conditioning on the first phase is needed: the second failure event only involves the
noises of the basis window, whatever the selected block.

\emph{Budget.} By Theorem~\ref{thm:norm_estimation} with $(\eps/4, \delta/(2k))$
(note $\log(4/(\delta/(2k))) = \log(8k/\delta)$), each run of
phase (1) uses at most $10^5\,d\log(8k/\delta)\cdot 16/\eps^2$ rounds, so phase (1) uses at
most $1.6\cdot10^6\,kd\log(8k/\delta)/\eps^2$ rounds. Phase (2) uses
$dn_2 \le d\big( (32d + 192\log(2/\delta))/\eps^2 + 1 \big) \le (33d^2 + 192\,d\log(2/\delta))/\eps^2$
(using $d \le d^2/\eps^2$), and $192\,d\log(2/\delta) \le 192\,kd\log(8k/\delta)$. Summing gives
the first bound, and $1600192 \le C_{\mathrm{sep}}$, $33 \le C_{\mathrm{sep}}$ give the second. If
$d \le k$ then $d^2 \le kd \le kd\log(8k/\delta)$ (as $\log(8k/\delta) \ge \log 8 > 1$), whence
the third bound with $\log(8k/\delta) = \log(1/\delta) + \log(8k)$. The lower bound is
Theorem~\ref{thm:separation_lower_nonadaptive}.
\end{proof}

\begin{remark}[The size of the gap]\label{rem:separation_gap}
For $d \le k \le d^2$ and a fixed $\delta$, the adaptive budget is
$O(kd\log k/\eps^2) = O(kd\log d/\eps^2)$ (since $\log k \le 2\log d$) while every
non-adaptive algorithm needs $\Omega(kd^2/\eps^2)$: adaptivity gains a factor
$\Omega(d/\log d)$, polynomial in the ambient dimension $kd \le d^3$. The paper phrases the
lower bound as $\Omega(kd\log(1/\delta)/\eps^2 + kd^2/\eps^2)$ using the adaptive lower bound
of Chen et al.\ \cite{chen2024assouad} on each block together with a padding argument for
expected budgets; the route above through Theorem~\ref{thm:lower_nonadaptive} gives the same
$kd^2$ term with the explicit constant $1/210$ for deterministic budgets and $\delta \le 1/10$,
without any citation.

\emph{Restriction of the comparison class.} The paper's Theorem~7 states the lower bound for
``possibly randomized non-adaptive'' algorithms. The lower bound proved here
(Theorem~\ref{thm:separation_lower_nonadaptive}(i)) covers only \emph{deterministic} fixed
designs, with an arbitrary randomized recommendation rule: this is the class of
Definition~\ref{def:fixed_design} and of Theorem~\ref{thm:lower_nonadaptive}, and it is the
class in which Theorem~\ref{thm:separation} compares the block algorithm with non-adaptive
ones. A non-adaptive algorithm whose design is drawn at random (independently of the
observations) is not covered: for such an algorithm the pigeonhole step selects a
block depending on the realized design, and the fixed-design lower bound would have to be
applied conditionally on the design and integrated, which is not done here. The adaptive
lower bound (ii) applies to every algorithm, randomized designs included, but only gives the
$kd\log(1/\delta)$ term. This restriction is recorded in the list of deviations from the
paper in the overview.

The constant $C_{\mathrm{sep}}$ inherits the (unoptimized) constant
$C_{\mathrm{norm}}$ of Theorem~\ref{thm:norm_estimation} multiplied by $16$ because the
norm-estimation accuracy is $\eps/4$.
\end{remark}


\part{Prerequisites to be added to Mathlib and LML}
\label{part:prereq}

\textbf{Overview}

The proofs of Part~\ref{part:main} rely on several results of probability theory,
information theory and bandit theory that are standard but not yet available in Mathlib or
in LML. This part develops them, at the level of generality of a library: each chapter is a
self-contained area that can be formalized (and contributed upstream) independently of the
rest.
\begin{itemize}
  \item Chapter~\ref{chap:pre_gaussian}: Gaussian vectors, linear Gaussian processes,
        suprema over compact index sets, comparison by covariance domination, expectation
        of maxima and of norms.
  \item Chapter~\ref{chap:pre_concentration}: Gaussian concentration for Lipschitz
        functions (Maurey--Pisier) and the Borell--TIS inequality for linear Gaussian
        processes.
  \item Chapter~\ref{chap:pre_subexp}: sub-exponential random variables, squares of
        Gaussian and Rademacher linear forms, Bernstein-type bounds.
  \item Chapter~\ref{chap:pre_info}: Pinsker and Bretagnolle--Huber inequalities,
        convexity of the Kullback--Leibler divergence, and the divergence decomposition for
        algorithm-environment sequences.
  \item Chapter~\ref{chap:pre_sudakov}: the Sudakov--Fernique inequality and the
        expectation of the maximum of independent Gaussians (lower bound).
  \item Chapter~\ref{chap:pre_median}: the Median Elimination algorithm for multi-armed
        bandits.
  \item Chapter~\ref{chap:pre_bayes}: likelihood ratios and Gaussian posteriors for
        adaptively collected linear-Gaussian data, and the adaptive simple-regret lower
        bound on the unit ball.
\end{itemize}

\chapter{Gaussian vectors and linear Gaussian processes}
\label{chap:pre_gaussian}

This chapter collects the facts about Gaussian vectors in $\R^n$ and about \emph{linear
Gaussian processes} $x \mapsto \ip{x}{G}$ ($G$ a centered Gaussian vector, $x$ ranging over
a compact index set $\Xs \subset \R^d$) that are used throughout Part~\ref{part:main}: closure
of Gaussian vectors under affine maps, identification of the law by mean and covariance,
rotation invariance of pairs of independent standard vectors, independence of uncorrelated
jointly Gaussian vectors, moments of Gaussian norms, tails of empirical means, measurability
and integrability of suprema over compact index sets, the expectation of a maximum of $m$
Gaussians, and comparison by covariance domination (the special case of the
Sudakov--Fernique inequality that suffices for the monotonicity of the Gaussian width,
Lemma~\ref{lem:width_mono}). Everything here is classical
\cite{vershynin2018high, boucheron2013concentration}; the point is to fix the statements at
the level of generality of a library and to name what exists.

\textbf{What Mathlib provides.} The standard Gaussian measure on a finite-dimensional inner
product space, \texttt{ProbabilityTheory.stdGaussian}, and the multivariate Gaussian
\texttt{multivariateGaussian μ S} on \texttt{EuclideanSpace ℝ ι}, defined as the image of
\texttt{stdGaussian} under $x \mapsto \mu + S^{1/2}x$ with $S^{1/2} = $ \texttt{CFC.sqrt S}
(\texttt{Mathlib/Probability/Distributions/Gaussian/Multivariate.lean}), with
\texttt{integral\_id\_multivariateGaussian}, \texttt{covarianceBilin\_multivariateGaussian},
\texttt{charFun\_multivariateGaussian}, \texttt{multivariateGaussian\_zero\_one},
\texttt{stdGaussian\_map} (invariance under linear isometry equivalences),
\texttt{map\_pi\_eq\_stdGaussian} and \texttt{stdGaussian\_eq\_map\_pi\_orthonormalBasis}
(coordinates are i.i.d.\ standard). The class \texttt{IsGaussian} of Gaussian measures
(\texttt{Gaussian/Basic.lean}: every continuous linear functional has a real Gaussian law),
closed under continuous linear maps (\texttt{isGaussian\_map}), translations, negation and
products, characterized by its characteristic function
(\texttt{IsGaussian.charFun\_eq}, \texttt{isGaussian\_iff\_charFun\_eq}); the predicate
\texttt{HasGaussianLaw X P} on random variables
(\texttt{Gaussian/HasGaussianLaw/Basic.lean}: \texttt{HasGaussianLaw.map}, \texttt{.smul},
\texttt{.neg}, \texttt{.add}, \texttt{.eval}, \texttt{.integrable}, \texttt{.memLp}), and the
independence results of \texttt{Gaussian/HasGaussianLaw/Independence.lean}
(\texttt{iIndepFun.hasGaussianLaw}, \texttt{IndepFun.hasGaussianLaw},
\texttt{HasGaussianLaw.indepFun\_of\_covariance\_eval}, \texttt{\_inner},
\texttt{\_eq\_zero}). Fernique's theorem \texttt{IsGaussian.exists\_integrable\_exp\_sq}
(\texttt{Gaussian/Fernique.lean}) and the rotation invariance
\texttt{IsGaussian.map\_rotation\_eq\_self} of $\mu \otimes \mu$ for centered $\mu$. On the real
line, \texttt{gaussianReal}, \texttt{mgf\_gaussianReal}, \texttt{integral\_id\_gaussianReal},
\texttt{variance\_id\_gaussianReal}, \texttt{gaussianReal\_map\_const\_mul},
\texttt{gaussianReal\_add\_gaussianReal\_of\_indepFun} (\texttt{Gaussian/Real.lean}).
Sub-Gaussian tails: \texttt{HasSubgaussianMGF} with \texttt{measure\_ge\_le}
(\texttt{Mathlib/Probability/Moments/SubGaussian.lean}) and \texttt{measure\_le\_le}
(LML, \texttt{ForMathlib/Probability/Moments/SubGaussian.lean}).

\textbf{What is missing.} The matrix form of the linear-image lemma (mean and covariance
matrix of $b + LV$), the uniqueness of the law given mean and covariance matrix in the
matrix language of Part~\ref{part:main}, the fourth moment of a Gaussian norm, the
lower bound $\E\norm{Y} \ge \sqrt{\tr S}/\sqrt 3$, the expectation bound for a maximum of $m$
Gaussians, the countable-dense reduction for suprema over compact index sets, and the
covariance-domination comparison lemma. All are short.

\textbf{Lean remark.} Vectors live in \texttt{EuclideanSpace ℝ (Fin d)} and matrices in
\texttt{Matrix (Fin d) (Fin d) ℝ}, acting through \texttt{Matrix.toEuclideanLin} /
\texttt{toEuclideanCLM}. The positive semidefinite square root of $S \in \PSD$ is
\texttt{CFC.sqrt S}, the one used by \texttt{multivariateGaussian}; for $S \in \PD$ it agrees
with the square root of Lemma~\ref{lem:sqrt_props}. The random vectors of this chapter are
functions \texttt{Ω → EuclideanSpace ℝ (Fin d)} with \texttt{HasLaw} a given Gaussian measure,
or with \texttt{HasGaussianLaw} together with prescribed mean and covariance.

\section{Gaussian vectors}
\label{sec:pre_gaussian_vectors}

\begin{definition}[Gaussian measures and Gaussian vectors]\label{def:pg_gaussian_vector}
Let $n \ge 1$, $\mu \in \R^n$ and $S \in \PSD$ (an $n \times n$ positive semidefinite
matrix). The \emph{Gaussian measure} $\Normal(\mu, S)$ on $\R^n$ is the image of the standard
Gaussian measure $\Normal(0, I_n)$ (Mathlib \texttt{stdGaussian}) under
$x \mapsto \mu + S^{1/2}x$, where $S^{1/2}$ is the positive semidefinite square root of $S$
(Mathlib \texttt{multivariateGaussian μ S}). A random vector $V : \Omega \to \R^n$ satisfies
$V \sim \Normal(\mu, S)$ if its law is $\Normal(\mu,S)$. More generally $V$ is a
\emph{Gaussian vector} if $\ip{t}{V}$ has a (possibly degenerate) real Gaussian law for every
$t \in \R^n$ (Mathlib \texttt{HasGaussianLaw V P}); its \emph{mean} is $\E V$ and its
\emph{covariance matrix} is $\operatorname{Cov}(V) := \E[(V - \E V)(V - \E V)^\top]$.
For $n = 1$, $\Normal(\mu, \sigma^2)$ is Mathlib's \texttt{gaussianReal μ σ²}; for
$\sigma = 0$ it is the Dirac mass at $\mu$.
\end{definition}

We recall the facts that Mathlib attaches to this definition: $\Normal(\mu,S)$ is a
probability measure with mean $\mu$ (\texttt{integral\_id\_multivariateGaussian}), covariance
bilinear form $(s,t) \mapsto s^\top S t$ (\texttt{covarianceBilin\_multivariateGaussian}, for
$S \in \PSD$), and characteristic function
$t \mapsto \exp(i\ip{t}{\mu} - \tfrac12 t^\top S t)$ (\texttt{charFun\_multivariateGaussian});
$\Normal(0, I_n) = $ \texttt{stdGaussian} (\texttt{multivariateGaussian\_zero\_one}); every
Gaussian vector is integrable with all moments (\texttt{HasGaussianLaw.memLp}).

\begin{lemma}[Standard Gaussian vectors have i.i.d.\ coordinates]\label{lem:gauss_pi}
\lean{ProbabilityTheory.hasLaw_pi_of_hasCondDistrib_const, ProbabilityTheory.iIndepFun_of_hasCondDistrib_const, ProbabilityTheory.hasLaw_of_hasCondDistrib_const}\leanok
\uses{def:pg_gaussian_vector}
Let $T \ge 1$.
\begin{enumerate}
  \item[(i)] If $\eta_0, \dots, \eta_{T-1}$ are independent real random variables, each with
        law $\Normal(0,1)$, then the random vector $\eta := (\eta_t)_{t<T} \in \R^T$
        satisfies $\eta \sim \Normal(0, I_T)$.
  \item[(ii)] Conversely, if $\eta \sim \Normal(0, I_T)$ then its coordinates
        $\eta_0,\dots,\eta_{T-1}$ are independent with law $\Normal(0,1)$.
\end{enumerate}
\end{lemma}

\begin{proof}
\uses{def:pg_gaussian_vector}
(i) The joint law of independent random variables is the product of their laws
(Mathlib \texttt{iIndepFun\_iff\_map\_fun\_eq\_pi\_map}), so the law of $\eta$, seen in
$\R^T = $ \texttt{EuclideanSpace ℝ (Fin T)}, is the image of
$\bigotimes_{t<T}\Normal(0,1)$ under the identification \texttt{WithLp.toLp 2}, which is
\texttt{stdGaussian} by \texttt{map\_pi\_eq\_stdGaussian}; conclude with
\texttt{multivariateGaussian\_zero\_one}. (ii) The same identity read backwards through the
measurable equivalence \texttt{WithLp.toLp 2}: the law of $(\eta_t)_t$ as an element of
$\mathrm{Fin}\,T \to \R$ is the product measure, which is the independence of the
coordinates with the stated marginals (\texttt{iIndepFun\_iff\_map\_fun\_eq\_pi\_map} again,
or \texttt{measurePreserving\_eval\_multivariateGaussian} for the marginals).
\end{proof}

\textbf{Lean remark.} Part (i) is Mathlib's \texttt{iIndepFun.hasLaw\_pi} together with
\texttt{map\_pi\_eq\_stdGaussian}, and (ii) is \texttt{iIndepFun\_iff\_hasLaw\_pi\_pi}. The form
in which (i) is used is sequential: if $\eta_0 \sim \nu$ and, for every $n$, the conditional law
of $\eta_{n+1}$ given $(\eta_0,\dots,\eta_n)$ (or given any variable of which this vector is a
measurable function) is the constant $\nu$, then $(\eta_0,\dots,\eta_{n-1})$ has the product law
$\nu^{\otimes n}$ and the $\eta_t$ are independent with law $\nu$
(\texttt{ProbabilityTheory.hasLaw\_pi\_of\_hasCondDistrib\_const},
\texttt{iIndepFun\_of\_hasCondDistrib\_const}; proof by induction through the measurable
equivalence \texttt{MeasurableEquiv.piFinSuccAbove}). This is what
Lemma~\ref{lem:noise_representation} provides for the noise of a linear Gaussian run.

\begin{lemma}[The law of a Gaussian vector is determined by its mean and covariance]
\label{lem:gauss_law_of_cov}
\uses{def:pg_gaussian_vector}
Let $V, V'$ be Gaussian vectors in $\R^n$ (possibly on different probability spaces).
\begin{enumerate}
  \item[(i)] If $\E V = \E V'$ and $\operatorname{Cov}(V) = \operatorname{Cov}(V')$, then $V$
        and $V'$ have the same law.
  \item[(ii)] If $\E V = \mu$ and $\operatorname{Cov}(V) = S$, then $V \sim \Normal(\mu, S)$.
  \item[(iii)] For $S \in \PSD$ and $g \sim \Normal(0, I_n)$, $S^{1/2} g \sim \Normal(0,S)$.
  \item[(iv)] If $G \sim \Normal(0, S)$ then $-G \sim \Normal(0,S)$, i.e.\ $-G \eqd G$.
\end{enumerate}
\end{lemma}

\begin{proof}
\uses{def:pg_gaussian_vector}
(i) The characteristic function of a Gaussian vector is
$t \mapsto \exp\big(i\,\E\ip{t}{V} - \tfrac12\Var\ip{t}{V}\big)$
(Mathlib \texttt{IsGaussian.charFun\_eq}, \texttt{HasGaussianLaw.charFun\_map\_eq}), and
$\E\ip{t}{V} = \ip{t}{\E V}$, $\Var\ip{t}{V} = t^\top\operatorname{Cov}(V)\,t$ (bilinearity of
the covariance, Mathlib \texttt{covarianceBilin}). Hence $V$ and $V'$ have the same
characteristic function, and a finite measure on a finite-dimensional inner product space is
determined by its characteristic function (\texttt{Measure.ext\_of\_charFun}). (ii) The
measure $\Normal(\mu,S)$ is Gaussian (\texttt{isGaussian\_multivariateGaussian}) with mean
$\mu$ and covariance $S$ (the facts recalled after Definition~\ref{def:pg_gaussian_vector});
apply (i) to $V$ and to the identity on $(\R^n, \Normal(\mu,S))$. (iii) is the definition of
$\Normal(0,S)$ (Mathlib \texttt{multivariateGaussian}, as the image of \texttt{stdGaussian}),
using that the law of $S^{1/2}g$ is the image of the law of $g$. (iv) $-G$ is a Gaussian
vector (\texttt{HasGaussianLaw.neg}) with mean $0$ and covariance
$\E[(-G)(-G)^\top] = S$; apply (ii).
\end{proof}

\begin{lemma}[Affine images of Gaussian vectors]\label{lem:gauss_linear_image}
\lean{ProbabilityTheory.multivariateGaussian_map_toEuclideanCLM, ProbabilityTheory.multivariateGaussian_map_toEuclideanLin}
\uses{def:pg_gaussian_vector}
Let $V \sim \Normal(\mu, S)$ in $\R^n$, $L \in \R^{m \times n}$ and $b \in \R^m$. Then
$b + LV \sim \Normal(b + L\mu,\ L S L^\top)$. In particular, for $g \sim \Normal(0, I_n)$,
$M \in \R^{m\times n}$ and $\mu \in \R^m$: $\mu + Mg \sim \Normal(\mu, MM^\top)$; and for
$a \in \R^n$, $\ip{a}{g} \sim \Normal(0, \norm{a}^2)$.
\end{lemma}

\begin{proof}
\uses{def:pg_gaussian_vector, lem:gauss_law_of_cov}
$b + LV$ is a Gaussian vector: $V$ is one, $L$ is a continuous linear map
(\texttt{HasGaussianLaw.map}) and translations preserve Gaussianity
(\texttt{IsGaussian} instance for \texttt{μ.map (· + c)}). Its mean is $b + L\mu$ by
linearity of the integral, and for $s, t \in \R^m$,
$\operatorname{Cov}(\ip{s}{b + LV}, \ip{t}{b+LV}) = \operatorname{Cov}(\ip{L^\top s}{V},
\ip{L^\top t}{V}) = (L^\top s)^\top S (L^\top t) = s^\top (LSL^\top) t$, i.e.\ the covariance
matrix is $LSL^\top$. Conclude with Lemma~\ref{lem:gauss_law_of_cov}(ii). The special cases
are $S = I_n$ and $L = a^\top$ (with $m = 1$, $\Normal(0, a^\top a)$ on $\R$ being
\texttt{gaussianReal 0 ‖a‖²}).
\end{proof}

\textbf{Lean remark.} With $V$ given by \texttt{HasLaw V (multivariateGaussian μ S) P}, the
conclusion is \texttt{HasLaw (fun ω ↦ b + L (V ω)) (multivariateGaussian (b + L μ) (L * S * Lᵀ)) P}.
The proof through characteristic functions is an alternative to the one through
\texttt{HasGaussianLaw}: $\E e^{i\ip{t}{b + LV}} = e^{i\ip{t}{b}}\,\E e^{i\ip{L^\top t}{V}}$
and \texttt{charFun\_multivariateGaussian} on both sides.

\begin{lemma}[Independent Gaussian vectors are jointly Gaussian]\label{lem:pg_indep_gaussian_pair}
\lean{ProbabilityTheory.multivariateGaussian_conv}
\uses{def:pg_gaussian_vector}
Let $V_1 \sim \Normal(\mu_1, S_1)$ in $\R^{n_1}$ and $V_2 \sim \Normal(\mu_2, S_2)$ in
$\R^{n_2}$ be independent, on the same probability space. Then the pair $(V_1, V_2)$, seen in
$\R^{n_1 + n_2}$, satisfies
$(V_1,V_2) \sim \Normal\big((\mu_1,\mu_2), \operatorname{diag}(S_1, S_2)\big)$.
If $n_1 = n_2$ then $V_1 + V_2 \sim \Normal(\mu_1 + \mu_2, S_1 + S_2)$.
\end{lemma}

\begin{proof}
\uses{def:pg_gaussian_vector, lem:gauss_law_of_cov, lem:gauss_linear_image}
The pair is a Gaussian vector (Mathlib \texttt{IndepFun.hasGaussianLaw}: the law of the pair
is the product measure, and a product of Gaussian measures is Gaussian). Its mean is
$(\mu_1,\mu_2)$ and its covariance matrix is block diagonal: the cross covariances
$\operatorname{Cov}(\ip{s}{V_1}, \ip{t}{V_2})$ vanish because $\ip{s}{V_1}$ and $\ip{t}{V_2}$
are independent and square-integrable (the covariance of independent variables vanishes:
Mathlib \texttt{IndepFun.covariance\_eq\_zero} in \texttt{Mathlib/Probability/Moments/Covariance.lean},
or the product formula \texttt{IndepFun.integral\_mul\_eq\_mul\_integral}). Conclude with
Lemma~\ref{lem:gauss_law_of_cov}(ii). For the sum apply Lemma~\ref{lem:gauss_linear_image}
with $L = [I_n\ I_n]$: $L\operatorname{diag}(S_1,S_2)L^\top = S_1 + S_2$.
\end{proof}

\begin{lemma}[Orthogonal invariance and rotation of pairs]\label{lem:gauss_rotation}
\uses{def:pg_gaussian_vector}
\begin{enumerate}
  \item[(i)] If $g \sim \Normal(0, I_d)$ and $U \in \R^{d\times d}$ is orthogonal
        ($U^\top U = I_d$), then $Ug \sim \Normal(0, I_d)$.
  \item[(ii)] If $g, g'$ are independent, each with law $\Normal(0, I_d)$, then for every
        $\theta \in \R$ the pair
        \[
          (g_\theta, g'_\theta) := \big(g\sin\theta + g'\cos\theta,\ g\cos\theta - g'\sin\theta\big)
        \]
        has the same law as $(g, g')$, namely $\Normal(0, I_{2d})$ on $\R^d \times \R^d$;
        in particular $g_\theta$ and $g'_\theta$ are independent $\Normal(0,I_d)$ vectors.
\end{enumerate}
\end{lemma}

\begin{proof}
\uses{def:pg_gaussian_vector, lem:pg_indep_gaussian_pair, lem:gauss_pi, lem:gauss_uncorrelated_indep}
(i) $U$ is a linear isometry equivalence of $\R^d$ and the standard Gaussian measure is
invariant under linear isometry equivalences (Mathlib \texttt{stdGaussian\_map}); the law of
$Ug$ is the image of the law of $g$. (ii) By Lemma~\ref{lem:pg_indep_gaussian_pair} the pair
$(g,g')$ has law $\Normal(0, I_{2d})$. The map
$R_\theta(u,v) := (u\sin\theta + v\cos\theta,\ u\cos\theta - v\sin\theta)$ is linear and
\[
  \norm{u\sin\theta + v\cos\theta}^2 + \norm{u\cos\theta - v\sin\theta}^2
  = \norm{u}^2(\sin^2\theta + \cos^2\theta) + \norm{v}^2(\cos^2\theta + \sin^2\theta)
    + 2\ip{u}{v}(\sin\theta\cos\theta - \cos\theta\sin\theta) = \norm{u}^2 + \norm{v}^2,
\]
so $R_\theta$ is a linear isometry of $\R^{2d}$ (bijective since it is an isometry of a
finite-dimensional space, or explicitly $R_\theta^{-1} = R_\theta^\top$). Apply (i) in
dimension $2d$. (Mathlib's \texttt{IsGaussian.map\_rotation\_eq\_self} is the same statement
for the map \texttt{ContinuousLinearMap.rotation θ} of
\texttt{Mathlib/Probability/Distributions/Fernique.lean}, namely
$(u,v)\mapsto(u\cos\theta + v\sin\theta,\ -u\sin\theta + v\cos\theta)$, applied to
$\mu\otimes\mu$ with $\mu$ a centered Gaussian measure. Since
$\mathrm{rotation}(-\theta)(u,v) = (u\cos\theta - v\sin\theta,\ u\sin\theta + v\cos\theta)$,
one has $R_\theta = \mathrm{swap}\circ\mathrm{rotation}(-\theta)$ with
$\mathrm{swap}(u,v) := (v,u)$, and the swap also preserves $\mu\otimes\mu$
(\texttt{Measure.prod\_swap}).) Independence of the components of
a $\Normal(0,I_{2d})$ vector is Lemma~\ref{lem:gauss_pi}(ii) applied blockwise, or
Lemma~\ref{lem:gauss_uncorrelated_indep}.
\end{proof}

\begin{lemma}[Joint law of two linear images]\label{lem:gauss_joint_linear}
\uses{def:pg_gaussian_vector}
Let $V \sim \Normal(\mu,S)$ in $\R^n$ and $L_1 \in \R^{m_1\times n}$, $L_2 \in \R^{m_2\times n}$.
Then $(L_1 V, L_2 V)$ is a Gaussian vector in $\R^{m_1}\times\R^{m_2}$ with
\[
  (L_1V, L_2V) \sim \Normal\Big( (L_1\mu, L_2\mu),
  \begin{pmatrix} L_1SL_1^\top & L_1SL_2^\top \\ L_2SL_1^\top & L_2SL_2^\top\end{pmatrix}\Big),
\]
and in particular the cross covariance is
$\operatorname{Cov}(L_1V, L_2V) := \E[(L_1V - L_1\mu)(L_2V - L_2\mu)^\top] = L_1 S L_2^\top$.
\end{lemma}

\begin{proof}
\uses{def:pg_gaussian_vector, lem:gauss_linear_image}
$(L_1V, L_2V) = LV$ with $L := \binom{L_1}{L_2} \in \R^{(m_1+m_2)\times n}$;
Lemma~\ref{lem:gauss_linear_image} gives $LV \sim \Normal(L\mu, LSL^\top)$ and
$LSL^\top$ is the displayed block matrix. The cross covariance is its off-diagonal block.
\end{proof}

\begin{lemma}[Uncorrelated jointly Gaussian vectors are independent]\label{lem:gauss_uncorrelated_indep}
\uses{def:pg_gaussian_vector}
Let $(V_1, V_2)$ be a Gaussian vector in $\R^{m_1}\times\R^{m_2}$ (i.e.\ the pair is jointly
Gaussian) such that $\operatorname{Cov}(V_1, V_2) = 0$, i.e.\
$\operatorname{Cov}((V_1)_i, (V_2)_j) = 0$ for all $i \le m_1$, $j \le m_2$. Then $V_1$ and
$V_2$ are independent.
\end{lemma}

\begin{proof}
\uses{def:pg_gaussian_vector}
Mathlib: \texttt{HasGaussianLaw.indepFun\_of\_covariance\_eval} (coordinates) or
\texttt{HasGaussianLaw.indepFun\_of\_covariance\_inner} (all $\operatorname{Cov}(\ip{s}{V_1},
\ip{t}{V_2}) = 0$, which follows from the coordinatewise hypothesis by bilinearity).
\end{proof}

\begin{lemma}[Expectation of an inner product with an independent centered vector]
\label{lem:indep_integral_zero}
Let $W : \Omega \to \R^d$ be integrable with $\E W = 0$, let $V : \Omega \to \mathcal V$ be a
random variable with values in a measurable space, independent of $W$, and let
$f : \mathcal V \to \R^d$ be bounded and measurable. Then $\ip{f(V)}{W}$ is integrable and
$\E\ip{f(V)}{W} = 0$.
\end{lemma}

\begin{proof}
$\ip{f(V)}{W} = \sum_{i<d} f_i(V)\,W_i$. For each $i$, $f_i(V)$ is bounded and measurable and
$W_i$ is integrable, so their product is integrable; $f_i(V)$ and $W_i$ are independent
(compositions of independent variables with measurable maps, Mathlib \texttt{IndepFun.comp}),
so $\E[f_i(V)W_i] = \E[f_i(V)]\,\E[W_i] = 0$ (product formula for independent integrable
variables, Mathlib \texttt{IndepFun.integral\_mul\_eq\_mul\_integral}). Sum over $i$.
\end{proof}

\section{Moments and tails}
\label{sec:pre_gaussian_moments}

\begin{lemma}[Moments of a centered real Gaussian]\label{lem:gauss_1d_moments}
\uses{def:pg_gaussian_vector}
Let $\sigma \ge 0$ and $Z \sim \Normal(0,\sigma^2)$. Then $Z$ has finite moments of all
orders and
\[
  \E Z = 0,\quad \E Z^2 = \sigma^2,\quad \E Z^4 = 3\sigma^4,\quad
  \E\abs{Z} = \sigma\sqrt{2/\pi},\quad \E e^{\lambda Z} = e^{\lambda^2\sigma^2/2}\ (\lambda\in\R).
\]
\end{lemma}

\begin{proof}
\uses{def:pg_gaussian_vector}
Mathlib: \texttt{memLp\_id\_gaussianReal} (all moments), \texttt{integral\_id\_gaussianReal},
\texttt{variance\_id\_gaussianReal}, \texttt{mgf\_gaussianReal} (with mean $0$ and variance
$\sigma^2$: $\E e^{\lambda Z} = \exp(0\cdot\lambda + \sigma^2\lambda^2/2)$). For the fourth
moment and the absolute moment, reduce to $\sigma = 1$ by $Z \eqd \sigma\xi$,
$\xi \sim \Normal(0,1)$ (\texttt{gaussianReal\_map\_const\_mul}; for $\sigma = 0$ both sides
vanish), and compute with the density $\varphi(x) = e^{-x^2/2}/\sqrt{2\pi}$:
$\int x^4\varphi = 3\int x^2\varphi = 3$ by integration by parts
($x^3\cdot x\varphi(x) = -x^3\varphi'(x)$, \texttt{integral\_mul\_deriv\_eq\_deriv\_mul}), and
$\int\abs{x}\varphi = 2\int_0^\infty x\varphi(x)\,dx = 2\varphi(0) = \sqrt{2/\pi}$ since
$x\varphi(x) = -\varphi'(x)$. Alternatively, $\E Z^4$ is the fourth derivative at $0$ of the
MGF $e^{\lambda^2\sigma^2/2}$ (Mathlib \texttt{iteratedDeriv\_mgf\_zero}).
\end{proof}

\begin{lemma}[Moments of the norm of a Gaussian vector]\label{lem:gauss_norm_moments}
\lean{ProbabilityTheory.integral_norm_sq_multivariateGaussian, ProbabilityTheory.integral_norm_le_sqrt_trace_multivariateGaussian, ProbabilityTheory.integral_norm_stdGaussian_le}
\uses{def:pg_gaussian_vector}
Let $S \in \PSD$ ($d \times d$) and $Y \sim \Normal(0, S)$; equivalently $Y = S^{1/2}g$ with
$g \sim \Normal(0,I_d)$. Then
\begin{enumerate}
  \item[(i)] $\E\norm{Y}^2 = \tr S$;
  \item[(ii)] $\E\norm{Y}^4 = (\tr S)^2 + 2\tr(S^2) \le 3(\tr S)^2$;
  \item[(iii)] $\E\norm{Y} \le \sqrt{\tr S}$;
  \item[(iv)] $\E\norm{Y} \ge (\tr S)^{3/2}\big/\sqrt{(\tr S)^2 + 2\tr(S^2)} \ge \sqrt{\tr S}/\sqrt3$.
\end{enumerate}
In the special case $S = I_d$: $\E\norm{g}^2 = d$, $\E\norm{g}^4 = d^2 + 2d$,
$\E\norm{g} \le \sqrt d$ and $\E\norm{g} \ge d/\sqrt{d+2} \ge \sqrt{d/3}$.
\end{lemma}

\begin{proof}
\uses{def:pg_gaussian_vector, lem:gauss_law_of_cov, lem:gauss_rotation, lem:gauss_pi, lem:gauss_1d_moments}
All quantities depend only on the law of $Y$, so by Lemma~\ref{lem:gauss_law_of_cov}(iii) we
may take $Y = S^{1/2}g$. Diagonalize $S = U\operatorname{diag}(s_1,\dots,s_d)U^\top$ with $U$
orthogonal and $s_i \ge 0$ (spectral theorem for symmetric matrices, Mathlib
\texttt{Matrix.IsHermitian.spectral\_theorem}; nonnegativity of the eigenvalues from
$S \in \PSD$). Then $\norm{Y}^2 = g^\top S g = \sum_i s_i h_i^2$ with $h := U^\top g$, and
$h \sim \Normal(0,I_d)$ by Lemma~\ref{lem:gauss_rotation}(i), so the $h_i$ are i.i.d.\
$\Normal(0,1)$ (Lemma~\ref{lem:gauss_pi}(ii)). (i) $\E\sum_i s_ih_i^2 = \sum_i s_i = \tr S$
(Lemma~\ref{lem:gauss_1d_moments}). (ii) Expanding the square and using independence,
$\E h_i^4 = 3$ and $\E h_i^2 = 1$,
\[
  \E\Big(\sum_i s_ih_i^2\Big)^2 = \sum_i s_i^2\,\E h_i^4 + \sum_{i\ne j}s_is_j\,\E h_i^2\,\E h_j^2
  = 3\sum_i s_i^2 + \Big(\big(\sum_i s_i\big)^2 - \sum_i s_i^2\Big) = (\tr S)^2 + 2\tr(S^2),
\]
using $\tr(S^2) = \sum_i s_i^2$. Since $s_i \ge 0$, $\sum_i s_i^2 \le (\sum_i s_i)^2$, whence
the bound $3(\tr S)^2$. (iii) is Jensen (or Cauchy--Schwarz): $(\E\norm{Y})^2 \le \E\norm{Y}^2$.
(iv) Hölder's inequality with exponents $3/2$ and $3$ applied to
$\norm{Y}^2 = \norm{Y}^{2/3}\cdot\norm{Y}^{4/3}$ gives
$\E\norm{Y}^2 \le (\E\norm{Y})^{2/3}(\E\norm{Y}^4)^{1/3}$, i.e.\
$(\E\norm{Y})^{2/3} \ge \tr S\,/\,((\tr S)^2 + 2\tr(S^2))^{1/3}$; raise to the power $3/2$ and
use (ii) for the last inequality. For $S = I_d$, $\tr S = d$, $\tr(S^2) = d$ and
$d^{3/2}/\sqrt{d^2+2d} = d/\sqrt{d+2}$, and $d/\sqrt{d+2} \ge \sqrt{d/3}$ since
$d \ge 1$ gives $d \ge (d+2)/3$.
\end{proof}

\begin{lemma}[Gaussian tail]\label{lem:gauss_tail}
\uses{def:pg_gaussian_vector}
Let $\mu \in \R$, $\sigma > 0$ and $Z \sim \Normal(\mu,\sigma^2)$. Then $Z - \mu$ has a
sub-Gaussian moment generating function with constant $\sigma^2$
(Mathlib \texttt{HasSubgaussianMGF (Z - μ) σ²}: $\E e^{\lambda(Z-\mu)} \le e^{\sigma^2\lambda^2/2}$
for all $\lambda$, with equality), and for every $t \ge 0$,
\[
  \Prob(Z - \mu \ge t) \le e^{-t^2/(2\sigma^2)}, \qquad
  \Prob(Z - \mu \le -t) \le e^{-t^2/(2\sigma^2)}, \qquad
  \Prob(\abs{Z - \mu} \ge t) \le 2e^{-t^2/(2\sigma^2)} .
\]
\end{lemma}

\begin{proof}
\uses{def:pg_gaussian_vector}
$Z - \mu \sim \Normal(0,\sigma^2)$ (\texttt{gaussianReal\_sub\_const}) and
\texttt{mgf\_gaussianReal} gives $\E e^{\lambda(Z-\mu)} = e^{\sigma^2\lambda^2/2}$, with
integrability from \texttt{integrable\_exp\_mul\_gaussianReal}; this is the structure
\texttt{HasSubgaussianMGF}. The tails are the Chernoff bounds
\texttt{HasSubgaussianMGF.measure\_ge\_le} (Mathlib) and
\texttt{HasSubgaussianMGF.measure\_le\_le} (LML); the two-sided bound is a union bound.
\end{proof}

\begin{lemma}[Concentration of a Gaussian empirical mean]\label{lem:gauss_mean_concentration}
\uses{def:pg_gaussian_vector, lem:gauss_tail}
Let $n \ge 1$ and $\eta_1,\dots,\eta_n$ be i.i.d.\ $\Normal(0,1)$. Then
$\bar\eta := \frac1n\sum_{i\le n}\eta_i \sim \Normal(0, 1/n)$ and for all $t \ge 0$,
$\Prob(\abs{\bar\eta} \ge t) \le 2e^{-nt^2/2}$. Consequently, for $\eps > 0$ and
$\delta \in (0,1)$, if $n \ge 8\log(2/\delta)/\eps^2$ then
$\Prob(\abs{\bar\eta} \le \eps/2) \ge 1 - \delta$. The same holds for $\bar y - \mu$ when
$y_i = \mu + \eta_i$.
\end{lemma}

\begin{proof}
\uses{lem:gauss_tail, lem:gauss_pi, lem:gauss_linear_image}
$\sum_i\eta_i \sim \Normal(0,n)$ (sums of independent real Gaussians, Mathlib
\texttt{gaussianReal\_add\_gaussianReal\_of\_indepFun} by induction; or
Lemmas~\ref{lem:gauss_pi} and~\ref{lem:gauss_linear_image} with $a = (1,\dots,1)$), and
$\bar\eta \sim \Normal(0, 1/n)$ (\texttt{gaussianReal\_map\_const\_mul}). Lemma~\ref{lem:gauss_tail}
with $\sigma^2 = 1/n$ gives $\Prob(\abs{\bar\eta}\ge t) \le 2e^{-t^2/(2/n)} = 2e^{-nt^2/2}$.
With $t = \eps/2$: $2e^{-n\eps^2/8} \le 2e^{-\log(2/\delta)} = \delta$ when
$n\eps^2/8 \ge \log(2/\delta)$.
\end{proof}

\section{Suprema of linear Gaussian processes}
\label{sec:pre_gaussian_sup}

\begin{lemma}[Suprema over a compact set reduce to a countable dense subset]
\label{lem:sup_countable_dense}
\lean{supportFn, lipschitzWith_supportFn, convexOn_supportFn, continuous_supportFn, abs_supportFn_le, exists_supportFn_eq_inner, partialSupInner, monotone_partialSupInner, abs_partialSupInner_le, exists_forall_supportFn_eq_iSup, tendsto_partialSupInner}\leanok
Let $\Xs \subset \R^d$ be nonempty and compact and $R := \max_{x\in\Xs}\norm{x}$.
\begin{enumerate}
  \item[(i)] There is a countable dense subset $D \subseteq \Xs$; fix an enumeration
        $D = \{q_n : n \in \N\}$ (repetitions allowed).
  \item[(ii)] For every $v \in \R^d$ the supremum $m(v) := \sup_{x\in\Xs}\ip{x}{v}$ is
        attained (it is a maximum), and
        $m(v) = \sup_{n\in\N}\ip{q_n}{v} = \lim_{n\to\infty}\max_{i\le n}\ip{q_i}{v}$,
        the sequence $\max_{i\le n}\ip{q_i}{v}$ being nondecreasing in $n$.
  \item[(iii)] $m : \R^d \to \R$ is convex, $R$-Lipschitz ($\abs{m(v) - m(v')} \le R\norm{v - v'}$),
        hence continuous and Borel measurable, and $\abs{m(v)} \le R\norm{v}$.
  \item[(iv)] For every random vector $G : \Omega \to \R^d$, $\sup_{x\in\Xs}\ip{x}{G} = m(G)$ is
        a real random variable with $\abs{m(G)} \le R\norm{G}$; it is integrable as soon as
        $\norm{G}$ is.
\end{enumerate}
\end{lemma}

\begin{proof}
\leanok
(i) $\Xs$ is a separable metric space (Mathlib \texttt{TopologicalSpace.exists\_countable\_dense}
for the subtype, which is second countable as a subspace of $\R^d$). (ii) For fixed $v$ the
map $x \mapsto \ip{x}{v}$ is continuous, so it attains its maximum on the compact $\Xs$
(\texttt{IsCompact.exists\_isMaxOn}). Clearly $\sup_{n}\ip{q_n}{v} \le m(v)$. Conversely, if
$x \in \Xs$ attains $m(v)$, density gives $q_{n_k} \to x$ and continuity gives
$\ip{q_{n_k}}{v} \to \ip{x}{v} = m(v)$, so $\sup_n\ip{q_n}{v} \ge m(v)$. The partial maxima are
nondecreasing and converge to the supremum of the sequence
(\texttt{tendsto\_atTop\_ciSup} for monotone bounded sequences). (iii) $m$ is a supremum of
linear functions, hence convex. For $v, v'$ and $x \in \Xs$ attaining $m(v)$,
$m(v) - m(v') \le \ip{x}{v} - \ip{x}{v'} \le \norm{x}\norm{v-v'} \le R\norm{v-v'}$; symmetrize.
$m(0) = 0$ gives $\abs{m(v)} \le R\norm{v}$. Measurability follows from continuity, or
directly from (ii): a countable supremum of measurable functions is measurable
(\texttt{Measurable.iSup}). (iv) is (ii)--(iii) composed with $G$.
\end{proof}

\textbf{Lean remark.} In Lean, $\sup_{x\in\Xs}\ip{x}{v}$ is best written as
\texttt{⨆ x : Xs, ⟪(x : E), v⟫} (a \texttt{ciSup} over the compact subtype, bounded above),
and item (ii) provides the equality with \texttt{⨆ n, ⟪q n, v⟫} for a sequence
\texttt{q : ℕ → Xs} with dense range (\texttt{TopologicalSpace.exists\_dense\_seq}). All
measure-theoretic statements about $\sup_{x\in\Xs}\ip{x}{G}$ in this blueprint are proved
through this countable form, and Chapter~\ref{chap:pre_concentration} uses the partial maxima
$\max_{i\le n}\ip{q_i}{G}$ as the finite approximations.

\begin{lemma}[Expected maximum of $m$ sub-Gaussian variables]\label{lem:gauss_max_upper}
\uses{lem:gauss_tail}
Let $m \ge 1$, $\sigma \ge 0$ and $Z_1,\dots,Z_m$ be real random variables on a common
probability space, each with a sub-Gaussian moment generating function with constant
$\sigma^2$ (Mathlib \texttt{HasSubgaussianMGF (Z i) σ²}: $e^{\beta Z_i}$ is integrable and
$\E e^{\beta Z_i}\le e^{\beta^2\sigma^2/2}$ for every $\beta\in\R$; no assumption on their
joint law). Then $\max_{i\le m}Z_i$ is integrable and
\[
  \E\max_{i\le m}Z_i \le \sigma\sqrt{2\log m}.
\]
In particular this applies to centered Gaussian variables $Z_i \sim \Normal(0,\sigma_i^2)$
with $\sigma_i^2 \le \sigma^2$ for every $i$, whatever their joint law
(Lemma~\ref{lem:gauss_tail}): $\E\max_{i\le m}Z_i \le \sigma\sqrt{2\log m}$.
\end{lemma}

\begin{proof}
\uses{lem:gauss_tail}
Integrability: $\abs{Z_i}\le e^{Z_i} + e^{-Z_i}$ is integrable ($\beta = \pm1$), and
$\abs{\max_iZ_i} \le \sum_i\abs{Z_i}$. Let $\beta > 0$. By Jensen's inequality for the convex
function $\exp$ (Mathlib \texttt{ConvexOn.map\_integral\_le} with \texttt{convexOn\_exp}; the
integrand $e^{\beta\max_iZ_i} \le \sum_ie^{\beta Z_i}$ is integrable by assumption),
\[
  \exp\Big(\beta\,\E\max_iZ_i\Big) \le \E\exp\Big(\beta\max_iZ_i\Big)
  = \E\max_i e^{\beta Z_i} \le \sum_{i\le m}\E e^{\beta Z_i}
  \le m\,e^{\beta^2\sigma^2/2},
\]
where $\max_ie^{\beta Z_i} = e^{\beta\max_iZ_i}$ because $\exp$ is increasing. Taking
logarithms, $\E\max_iZ_i \le \frac{\log m}{\beta} + \frac{\beta\sigma^2}{2}$ for every
$\beta>0$. If $m\ge2$ and $\sigma>0$, choose $\beta := \sqrt{2\log m}/\sigma$: both terms
equal $\sigma\sqrt{2\log m}/2$. If $m = 1$, let $\beta\downarrow0$: $\E Z_1\le0 = \sigma\sqrt{2\log1}$.
If $\sigma = 0$, let $\beta\to\infty$: $\E\max_iZ_i\le0$. The Gaussian case: by
Lemma~\ref{lem:gauss_tail} (with $\mu = 0$), $Z_i\sim\Normal(0,\sigma_i^2)$ with $\sigma_i>0$
has a sub-Gaussian MGF with constant $\sigma_i^2\le\sigma^2$, hence with constant $\sigma^2$
(the bound $e^{\beta^2\sigma_i^2/2}\le e^{\beta^2\sigma^2/2}$ is monotone in the constant);
if $\sigma_i = 0$ then $Z_i = 0$ a.s.\ and the constant-$\sigma^2$ bound is trivial.
\end{proof}

\begin{lemma}[Comparison by covariance domination]\label{lem:gauss_comparison}
\lean{ProbabilityTheory.gaussianWidth_multivariateGaussian_le, ProbabilityTheory.gaussianWidth_le_gaussianWidth_conv}\leanok
\uses{def:pg_gaussian_vector, lem:sup_countable_dense}
Let $\Xs \subset \R^d$ be nonempty and compact, let $S_X, S_Y \in \PSD$ with
$S_X \preceq S_Y$ (i.e.\ $S_Y - S_X \in \PSD$), and let $G_X \sim \Normal(0,S_X)$,
$G_Y \sim \Normal(0,S_Y)$ (on arbitrary probability spaces). Then
$\sup_{x\in\Xs}\ip{x}{G_X}$ and $\sup_{x\in\Xs}\ip{x}{G_Y}$ are integrable and
\[
  \E\sup_{x\in\Xs}\ip{x}{G_X} \le \E\sup_{x\in\Xs}\ip{x}{G_Y}.
\]
In particular (finite index set $\Xs = \{e_1,\dots,e_m\}$, the canonical basis of $\R^m$):
if $Z \sim \Normal(0,\Sigma_Z)$ and $W \sim \Normal(0,\Sigma_W)$ in $\R^m$ with
$\Sigma_Z \preceq \Sigma_W$, then $\E\max_{i\le m}Z_i \le \E\max_{i\le m}W_i$.
\end{lemma}

\begin{proof}
\leanok
\uses{lem:sup_countable_dense, lem:gauss_norm_moments, lem:pg_indep_gaussian_pair, lem:gauss_law_of_cov}
Write $m(v) := \sup_{x\in\Xs}\ip{x}{v}$, which is continuous with $\abs{m(v)} \le R\norm{v}$,
$R := \max_{x\in\Xs}\norm{x}$ (Lemma~\ref{lem:sup_countable_dense}); integrability of $m(G_X)$
and $m(G_Y)$ follows from $\E\norm{G} < \infty$ for Gaussian vectors
(Lemma~\ref{lem:gauss_norm_moments}(iii)). Both expectations depend only on the laws, so we may
work on the product space $(\R^d\times\R^d, \Normal(0,S_X)\otimes\Normal(0,S_Y-S_X))$ with
coordinate projections $G_X$ and $G_Z$, which are independent with
$G_Z \sim \Normal(0, S_Y - S_X)$ ($S_Y - S_X \in \PSD$ by assumption). By
Lemma~\ref{lem:pg_indep_gaussian_pair}, $G_X + G_Z \sim \Normal(0, S_X + (S_Y - S_X)) =
\Normal(0,S_Y)$, so $\E m(G_Y) = \E m(G_X + G_Z)$ (Lemma~\ref{lem:gauss_law_of_cov} is not
even needed: this is equality of laws). By Fubini on the product measure
(Mathlib \texttt{MeasureTheory.integral\_prod}; the integrand $(u,z)\mapsto m(u+z)$ is
continuous and dominated by $R(\norm{u}+\norm{z})$, which is integrable on the product),
\[
  \E m(G_X + G_Z) = \int\Big(\int m(u+z)\,\Normal(0,S_Y-S_X)(dz)\Big)\Normal(0,S_X)(du).
\]
Fix $u$. For every $x \in \Xs$ and every $z$, $\ip{x}{u} + \ip{x}{z} \le m(u+z)$; integrating in
$z$ and using $\int\ip{x}{z}\,\Normal(0,S_Y-S_X)(dz) = \ip{x}{0} = 0$ (centered), monotonicity
of the integral gives $\ip{x}{u} \le \int m(u+z)\,\Normal(0,S_Y-S_X)(dz)$. Taking the supremum
over $x\in\Xs$: $m(u) \le \int m(u+z)\,\Normal(0,S_Y-S_X)(dz)$. Integrating in $u$ yields
$\E m(G_X) \le \E m(G_X+G_Z) = \E m(G_Y)$. The finite-index statement is the case
$\Xs = \{e_1,\dots,e_m\}$, since $\ip{e_i}{G} = G_i$.
\end{proof}

\begin{remark}[Relation with Sudakov--Fernique]\label{rem:pg_comparison_vs_sf}
The Sudakov--Fernique inequality (Chapter~\ref{chap:pre_sudakov}) only requires
$\E(Z_i - Z_j)^2 \le \E(W_i - W_j)^2$ for all pairs, which is weaker than
$\Sigma_Z \preceq \Sigma_W$. The Loewner-order version above is all that the monotonicity of
the Gaussian width (Lemma~\ref{lem:width_mono}) needs, and its proof is elementary; this is
one of the deviations from the paper listed in the introduction.
\end{remark}

\chapter{Gaussian concentration and the Borell--TIS inequality}
\label{chap:pre_concentration}

This chapter proves that the supremum $M = \sup_{x\in\Xs}\ip{x}{G}$ of a linear Gaussian
process over a compact index set concentrates around its mean at the scale
$\sigma := \sup_{x\in\Xs}\sqrt{x^\top S x}$ of the largest individual standard deviation
(the Borell--TIS inequality, Theorem~\ref{thm:borell_tis}), in the form used by the
fixed-design upper bound (Corollaries~\ref{cor:sup_bound_whp}
and~\ref{cor:sup_bound_symmetric}). The route is the Maurey--Pisier argument
\cite{pisier1986probabilistic} (see also \cite[Section~5.4]{vershynin2018high} and
\cite[Chapter~5]{boucheron2013concentration}): a $C^1$ function $f$ of a standard Gaussian
vector with $\norm{\nabla f} \le L$ has a sub-Gaussian moment generating function with
constant $\cG L^2$, where
\[
  \cG := \pi^2/4 ,
\]
and the supremum of a finite family of linear forms is approximated by the smooth
``log-sum-exp'' function. The optimal constant is $1$ (Gaussian isoperimetry, or the
Gaussian log-Sobolev inequality, see \cite{adler2007random}); the factor $\cG$ is the price
of an argument that uses only Jensen's inequality, the fundamental theorem of calculus, Fubini
and the rotation invariance of Gaussian pairs, and it propagates as an explicit absolute
constant into the upper bounds of Part~\ref{part:main}.

\textbf{What Mathlib provides.} The structure \texttt{ProbabilityTheory.HasSubgaussianMGF X c μ}
(\texttt{Mathlib/Probability/Moments/SubGaussian.lean}): integrability of $e^{tX}$ for all
$t$ together with $\E e^{tX} \le e^{ct^2/2}$; its Chernoff bound
\texttt{HasSubgaussianMGF.measure\_ge\_le}, and \texttt{measure\_le\_le} in LML
(\texttt{ForMathlib/Probability/Moments/SubGaussian.lean}). Jensen's inequality
\texttt{ConvexOn.map\_integral\_le} and \texttt{ConvexOn.map\_average\_le}
(\texttt{Mathlib/Analysis/Convex/Integral.lean}) with \texttt{convexOn\_exp}. The fundamental
theorem of calculus \texttt{intervalIntegral.integral\_eq\_sub\_of\_hasDerivAt}, the chain
rule \texttt{HasFDerivAt.comp\_hasDerivAt}, the mean value inequality
\texttt{lipschitzWith\_of\_nnnorm\_fderiv\_le}, Fubini \texttt{MeasureTheory.integral\_prod}
and \texttt{integral\_integral\_swap}, dominated convergence
\texttt{tendsto\_integral\_of\_dominated\_convergence}, Fernique's theorem
\texttt{IsGaussian.exists\_integrable\_exp\_sq}, the rotation invariance of Gaussian pairs
(Lemma~\ref{lem:gauss_rotation}, Mathlib \texttt{stdGaussian\_map} and
\texttt{IsGaussian.map\_rotation\_eq\_self}), and the Gaussian MGF \texttt{mgf\_gaussianReal}.

\textbf{What is missing.} Everything specific to concentration: the Maurey--Pisier inequality
(Lemma~\ref{lem:maurey_pisier}), the log-sum-exp smoothing (Lemma~\ref{lem:logsumexp_props}),
the passage from finite to compact index sets (Lemma~\ref{lem:compact_sup_subgaussian}) and
the Borell--TIS inequality itself. Mathlib has no Gaussian concentration inequality for
Lipschitz functions at the time of writing.

\textbf{Lean remark.} Sub-Gaussianity of $M - \E M$ is stated with Mathlib's
\texttt{HasSubgaussianMGF (fun ω ↦ M ω - P[M]) (cG * σ²) P}, whose constant is an
\texttt{ℝ≥0}; all constants below are therefore nonnegative reals coerced to \texttt{ℝ≥0}.
Gradients are \texttt{gradient f x} in \texttt{EuclideanSpace ℝ (Fin d)}, with
$\ip{\nabla f(x)}{v} = f'(x)(v)$ (\texttt{HasGradientAt} / \texttt{HasFDerivAt}), and
``$\norm{\nabla f} \le L$ everywhere'' is \texttt{∀ x, ‖gradient f x‖ ≤ L} together with
differentiability; the class $C^1$ is \texttt{ContDiff ℝ 1 f}.

\section{Preliminaries}
\label{sec:pre_concentration_prelim}

\begin{definition}[The Gaussian concentration constant]\label{def:pc_gaussian_constant}
$\cG := \pi^2/4$ ($\approx 2.4674$).
\end{definition}

\begin{lemma}[MGF of a linear form of a standard Gaussian vector]\label{lem:gauss_mgf_inner}
\uses{def:pg_gaussian_vector}
Let $g' \sim \Normal(0, I_d)$ and $a \in \R^d$. Then $e^{\ip{a}{g'}}$ is integrable and
$\E e^{\ip{a}{g'}} = e^{\norm{a}^2/2}$. More generally, for $s \in \R$,
$\E e^{s\ip{a}{g'}} = e^{s^2\norm{a}^2/2}$.
\end{lemma}

\begin{proof}
\uses{lem:gauss_linear_image, lem:gauss_1d_moments}
$\ip{a}{g'} \sim \Normal(0,\norm{a}^2)$ by Lemma~\ref{lem:gauss_linear_image}, and
Lemma~\ref{lem:gauss_1d_moments} (Mathlib \texttt{mgf\_gaussianReal} at the point $s$, with
\texttt{integrable\_exp\_mul\_gaussianReal}) gives the value. Alternatively, by
Lemma~\ref{lem:gauss_pi}, $\ip{a}{g'} = \sum_ia_ig'_i$ with independent coordinates, and
$\E\prod_ie^{sa_ig'_i} = \prod_ie^{s^2a_i^2/2}$ (\texttt{iIndepFun.mgf\_sum}).
\end{proof}

\begin{lemma}[Exponential moments of Gaussian norms]\label{lem:pc_exp_norm_integrable}
\uses{def:pg_gaussian_vector}
Let $G \sim \Normal(\mu, S)$ in $\R^n$ and $c \ge 0$. Then $e^{c\norm{G}}$ is integrable.
Consequently, if $F : \R^n \to \R$ is measurable with $\abs{F(x)} \le A + B\norm{x}$ for
constants $A, B \ge 0$, then $F(G)$ is integrable and $e^{\lambda F(G)}$ is integrable for every
$\lambda \in \R$.
\end{lemma}

\begin{proof}
\uses{lem:gauss_1d_moments, lem:gauss_pi}
By Fernique's theorem (Mathlib \texttt{IsGaussian.exists\_integrable\_exp\_sq}) there is
$C > 0$ with $e^{C\norm{x}^2}$ integrable under the law of $G$; since
$c\norm{x} \le C\norm{x}^2 + c^2/(4C)$ (AM--GM), $e^{c\norm{G}} \le e^{c^2/(4C)}e^{C\norm{G}^2}$
is integrable. (Elementary alternative for $G = \mu + S^{1/2}g$: $\norm{G} \le \norm{\mu} +
\norm{S^{1/2}}_{\mathrm{op}}\sum_i\abs{g_i}$, $e^{c'\abs{g_i}} \le e^{c'g_i} + e^{-c'g_i}$, and
by independence of the coordinates (Lemma~\ref{lem:gauss_pi}) and
Lemma~\ref{lem:gauss_1d_moments} the expectation of the product is a finite product of finite
numbers.) For $F$: $\abs{F(G)} \le A + B\norm{G}$ is integrable, and
$e^{\lambda F(G)} \le e^{\abs\lambda A}e^{\abs\lambda B\norm{G}}$ is integrable by the first
part with $c = \abs{\lambda}B$.
\end{proof}

\begin{lemma}[Bounded gradient implies Lipschitz and linear growth]\label{lem:pc_lipschitz_of_grad}
Let $f : \R^d \to \R$ be differentiable with $\norm{\nabla f(x)} \le L$ for every $x$, where
$L \ge 0$. Then $\abs{f(x) - f(y)} \le L\norm{x-y}$ for all $x,y$, $\abs{f(x)} \le \abs{f(0)} +
L\norm{x}$, and $\abs{\ip{\nabla f(x)}{v}} \le L\norm{v}$ for all $x, v$. In particular, for
$g \sim \Normal(0,I_d)$, $f(g)$ is integrable and $e^{\lambda f(g)}$ is integrable for every
$\lambda\in\R$.
\end{lemma}

\begin{proof}
\uses{lem:pc_exp_norm_integrable}
The Lipschitz bound is the mean value inequality (Mathlib
\texttt{lipschitzWith\_of\_nnnorm\_fderiv\_le}, the operator norm of $f'(x)$ being
$\norm{\nabla f(x)}$ in an inner product space); $y = 0$ gives the growth bound; the last
inequality is Cauchy--Schwarz. Integrability is Lemma~\ref{lem:pc_exp_norm_integrable} with
$A = \abs{f(0)}$, $B = L$.
\end{proof}

\begin{lemma}[The log-sum-exp smoothing of a finite maximum]\label{lem:logsumexp_props}
Let $m \ge 1$, $x_1,\dots,x_m \in \R^d$, $\beta > 0$ and
\[
  f_\beta(v) := \frac1\beta\log\sum_{i\le m}e^{\beta\ip{x_i}{v}}, \qquad v \in \R^d .
\]
\begin{enumerate}
  \item[(i)] $f_\beta$ is $C^\infty$ on $\R^d$, with gradient
        $\nabla f_\beta(v) = \sum_{i\le m}p_i(v)\,x_i$ where
        $p_i(v) := e^{\beta\ip{x_i}{v}}/\sum_{j\le m}e^{\beta\ip{x_j}{v}}$ satisfies $p_i(v) \ge 0$,
        $\sum_ip_i(v) = 1$; hence $\norm{\nabla f_\beta(v)} \le \max_{i\le m}\norm{x_i}$ for all $v$.
  \item[(ii)] For every $v$,
        $\max_{i\le m}\ip{x_i}{v} \le f_\beta(v) \le \max_{i\le m}\ip{x_i}{v} + \frac{\log m}{\beta}$.
\end{enumerate}
\end{lemma}

\begin{proof}
(i) $v \mapsto \sum_ie^{\beta\ip{x_i}{v}}$ is a finite sum of compositions of $\exp$ with
linear forms, hence $C^\infty$ and strictly positive; $\log$ is $C^\infty$ on $(0,\infty)$
(Mathlib \texttt{ContDiff.log}, \texttt{contDiff\_exp}, \texttt{ContDiff.inner}). The chain
rule gives $\nabla f_\beta(v) = \frac1\beta\cdot\frac{\sum_i\beta e^{\beta\ip{x_i}{v}}x_i}
{\sum_je^{\beta\ip{x_j}{v}}} = \sum_ip_i(v)x_i$. The $p_i(v)$ are positive with sum $1$, so by
the triangle inequality $\norm{\nabla f_\beta(v)} \le \sum_ip_i(v)\norm{x_i} \le \max_i\norm{x_i}$.
(ii) Let $M(v) := \max_i\ip{x_i}{v}$. Since $\exp$ and $\log$ are increasing,
$e^{\beta M(v)} \le \sum_ie^{\beta\ip{x_i}{v}} \le m\,e^{\beta M(v)}$; take logarithms and
divide by $\beta$.
\end{proof}

\section{The Maurey--Pisier argument}
\label{sec:pre_concentration_mp}

Throughout this section $f : \R^d \to \R$ is of class $C^1$ with $\norm{\nabla f(x)} \le L$
for all $x$, and $g, g'$ are independent $\Normal(0,I_d)$ vectors on $(\Omega,\Prob)$. For
$\theta \in \R$ we write, as in Lemma~\ref{lem:gauss_rotation},
\[
  g_\theta := g\sin\theta + g'\cos\theta, \qquad g'_\theta := g\cos\theta - g'\sin\theta,
\]
so that $g_0 = g'$, $g_{\pi/2} = g$ and $\frac{d}{d\theta}g_\theta = g'_\theta$.

\begin{lemma}[Interpolation along a quarter circle]\label{lem:pc_interpolation_ftc}
Let $f : \R^d \to \R$ be of class $C^1$ and $u, v \in \R^d$. Set
$\gamma(\theta) := u\sin\theta + v\cos\theta$ and $\gamma'(\theta) := u\cos\theta - v\sin\theta$.
Then $\theta \mapsto f(\gamma(\theta))$ is differentiable on $\R$ with derivative
$\theta \mapsto \ip{\nabla f(\gamma(\theta))}{\gamma'(\theta)}$, this derivative is continuous,
and
\[
  f(u) - f(v) = \int_0^{\pi/2}\ip{\nabla f(\gamma(\theta))}{\gamma'(\theta)}\,d\theta .
\]
\end{lemma}

\begin{proof}
$\gamma$ is differentiable with $\gamma'(\theta) = u\cos\theta - v\sin\theta$
(\texttt{HasDerivAt.smul\_const} for $\sin$ and $\cos$, \texttt{HasDerivAt.add}). The chain
rule (\texttt{HasFDerivAt.comp\_hasDerivAt}) gives that $f\circ\gamma$ has derivative
$f'(\gamma(\theta))(\gamma'(\theta)) = \ip{\nabla f(\gamma(\theta))}{\gamma'(\theta)}$ at
$\theta$. This derivative is continuous in $\theta$ as a composition of continuous maps
($\nabla f$ is continuous since $f$ is $C^1$), hence interval integrable on $[0,\pi/2]$
(\texttt{ContinuousOn.intervalIntegrable}). The fundamental theorem of calculus
(\texttt{intervalIntegral.integral\_eq\_sub\_of\_hasDerivAt}) gives
$\int_0^{\pi/2}(f\circ\gamma)' = f(\gamma(\pi/2)) - f(\gamma(0)) = f(u) - f(v)$.
\end{proof}

\begin{lemma}[{Jensen's inequality for the uniform measure on $[0,\pi/2]$}]
\label{lem:pc_jensen_exp_average}
Let $h : [0,\pi/2] \to \R$ be continuous and $s \in \R$. Then
\[
  \exp\Big(s\int_0^{\pi/2}h(\theta)\,d\theta\Big)
  \le \frac2\pi\int_0^{\pi/2}\exp\Big(\frac{\pi s}{2}\,h(\theta)\Big)d\theta .
\]
\end{lemma}

\begin{proof}
Let $\nu := \frac2\pi\,\mathrm{Leb}|_{[0,\pi/2]}$, a probability measure, and
$k(\theta) := \frac{\pi s}{2}h(\theta)$, which is continuous, hence bounded on the compact
interval; so $k$ and $e^{k}$ are $\nu$-integrable. Jensen's inequality for the convex
continuous function $\exp$ on the closed convex set $\R$ (Mathlib
\texttt{ConvexOn.map\_integral\_le} with \texttt{convexOn\_exp}, \texttt{continuous\_exp},
\texttt{isClosed\_univ}) gives $\exp(\int k\,d\nu) \le \int e^{k}\,d\nu$. Finally
$\int k\,d\nu = \frac2\pi\cdot\frac{\pi s}{2}\int_0^{\pi/2}h = s\int_0^{\pi/2}h$ and
$\int e^k\,d\nu = \frac2\pi\int_0^{\pi/2}e^{k(\theta)}d\theta$.
\end{proof}

\begin{lemma}[MGF of the rotated gradient term]\label{lem:pc_rotated_mgf_bound}
\uses{def:pg_gaussian_vector, lem:pc_lipschitz_of_grad}
Let $f$ be $C^1$ with $\norm{\nabla f} \le L$ everywhere, $g, g'$ independent
$\Normal(0,I_d)$ vectors, $s \in \R$ and $\theta \in \R$. Then
$\omega \mapsto \exp(s\ip{\nabla f(g_\theta)}{g'_\theta})$ is integrable and
\[
  \E\exp\big(s\ip{\nabla f(g_\theta)}{g'_\theta}\big)
  = \E\exp\big(s\ip{\nabla f(g)}{g'}\big)
  = \E\exp\Big(\frac{s^2\norm{\nabla f(g)}^2}{2}\Big) \le \exp\Big(\frac{s^2L^2}{2}\Big).
\]
\end{lemma}

\begin{proof}
\uses{lem:gauss_rotation, lem:gauss_mgf_inner, lem:pc_lipschitz_of_grad, lem:pc_exp_norm_integrable}
Let $\Phi(u,v) := \exp(s\ip{\nabla f(u)}{v})$, a continuous function on $\R^d\times\R^d$ with
$0 < \Phi(u,v) \le \exp(\abs{s}L\norm{v})$ (Lemma~\ref{lem:pc_lipschitz_of_grad}). By
Lemma~\ref{lem:gauss_rotation}(ii), $(g_\theta, g'_\theta)$ has the same law as $(g,g')$, so
$\E\Phi(g_\theta,g'_\theta) = \E\Phi(g,g')$ (equality of laws; both sides are finite by the
domination and Lemma~\ref{lem:pc_exp_norm_integrable}). The law of $(g,g')$ is the product
$\Normal(0,I_d)\otimes\Normal(0,I_d)$ (independence), and $\Phi$ is integrable for it (dominated
by $\exp(\abs sL\norm{v})$, which is integrable in $v$ and constant in $u$), so Fubini
(\texttt{MeasureTheory.integral\_prod}) gives
\[
  \E\Phi(g,g') = \int\Big(\int\exp\big(\ip{s\nabla f(u)}{v}\big)\,\Normal(0,I_d)(dv)\Big)\Normal(0,I_d)(du)
  = \int\exp\Big(\frac{s^2\norm{\nabla f(u)}^2}{2}\Big)\Normal(0,I_d)(du),
\]
by Lemma~\ref{lem:gauss_mgf_inner} with $a = s\nabla f(u)$. The last integrand is at most
$\exp(s^2L^2/2)$ pointwise, and $\Normal(0,I_d)$ is a probability measure.
\end{proof}

\begin{lemma}[Maurey--Pisier Gaussian concentration]\label{lem:maurey_pisier}
\uses{def:pg_gaussian_vector, def:pc_gaussian_constant}
Let $L \ge 0$ and let $f : \R^d \to \R$ be of class $C^1$ with $\norm{\nabla f(x)} \le L$ for
every $x \in \R^d$. Let $g \sim \Normal(0, I_d)$. Then $f(g)$ is integrable, and for every
$\lambda \in \R$ the variable $e^{\lambda(f(g) - \E f(g))}$ is integrable with
\[
  \E\exp\big(\lambda(f(g) - \E f(g))\big) \le \exp\Big(\frac{\pi^2\lambda^2L^2}{8}\Big)
  = \exp\Big(\frac{\cG L^2\,\lambda^2}{2}\Big).
\]
In Mathlib's terms, $f(g) - \E f(g)$ satisfies \texttt{HasSubgaussianMGF} with constant
$\cG L^2$.
\end{lemma}

\begin{proof}
\uses{lem:pc_lipschitz_of_grad, lem:pc_exp_norm_integrable, lem:pc_interpolation_ftc, lem:pc_jensen_exp_average, lem:pc_rotated_mgf_bound, lem:gauss_rotation}
\emph{Step 0 (integrability).} By Lemma~\ref{lem:pc_lipschitz_of_grad}, $f$ is $L$-Lipschitz,
$\abs{f(x)} \le \abs{f(0)} + L\norm{x}$, $f(g)$ is integrable, and $e^{\lambda f(g)}$ is
integrable for every $\lambda$; hence so is $e^{\lambda(f(g)-\E f(g))}$. Fix $\lambda \in \R$
(for $\lambda = 0$ there is nothing to prove). Enlarging the probability space if necessary
(all quantities depend only on the law of $g$), let $g'$ be an independent copy of $g$; the
law of $(g,g')$ is $\gamma\otimes\gamma$ with $\gamma := \Normal(0,I_d)$.

\emph{Step 1 (symmetrization by Jensen in $g'$).} For fixed $u \in \R^d$, the variable
$\lambda(f(u) - f(g'))$ is integrable, and $e^{\lambda(f(u)-f(g'))} \le
e^{\abs\lambda(\abs{f(u)} + \abs{f(0)} + L\norm{g'})}$ is integrable
(Lemma~\ref{lem:pc_exp_norm_integrable}); Jensen's inequality for $\exp$
(\texttt{ConvexOn.map\_integral\_le}, \texttt{convexOn\_exp}) gives
\[
  \exp\big(\lambda(f(u) - \E f(g'))\big) \le \E\exp\big(\lambda(f(u) - f(g'))\big).
\]
Since $\E f(g') = \E f(g)$, evaluating at $u = g$ and integrating in $g$ (Fubini for
$\gamma\otimes\gamma$, \texttt{MeasureTheory.integral\_prod}; the function
$(u,v) \mapsto e^{\lambda(f(u)-f(v))}$ is dominated by
$e^{2\abs\lambda\abs{f(0)}}e^{\abs\lambda L\norm{u}}e^{\abs\lambda L\norm{v}}$, integrable on the
product by Lemma~\ref{lem:pc_exp_norm_integrable}),
\[
  \E\exp\big(\lambda(f(g) - \E f(g))\big) \le \E\exp\big(\lambda(f(g) - f(g'))\big).
\]

\emph{Step 2 (interpolation).} Pointwise in $\omega$, Lemma~\ref{lem:pc_interpolation_ftc}
with $u = g(\omega)$, $v = g'(\omega)$ gives
\[
  f(g) - f(g') = \int_0^{\pi/2}h_\omega(\theta)\,d\theta, \qquad
  h_\omega(\theta) := \ip{\nabla f(g_\theta)}{g'_\theta},
\]
where $\theta \mapsto h_\omega(\theta)$ is continuous.

\emph{Step 3 (Jensen for the uniform measure on $[0,\pi/2]$).} By
Lemma~\ref{lem:pc_jensen_exp_average} with $s = \lambda$, pointwise in $\omega$,
\[
  \exp\big(\lambda(f(g) - f(g'))\big)
  \le \frac2\pi\int_0^{\pi/2}\exp\Big(\frac{\pi\lambda}{2}\,h_\omega(\theta)\Big)d\theta .
\]

\emph{Step 4 (Fubini and rotation invariance).} The function
$(\omega,\theta) \mapsto \exp(\frac{\pi\lambda}{2}h_\omega(\theta))$ is measurable (a continuous
function of $(g(\omega), g'(\omega), \theta)$) and satisfies, by
Lemma~\ref{lem:pc_lipschitz_of_grad} and $\norm{g'_\theta} \le \norm{g} + \norm{g'}$,
\[
  0 < \exp\Big(\frac{\pi\lambda}{2}h_\omega(\theta)\Big)
  \le \exp\Big(\frac{\pi\abs\lambda L}{2}\big(\norm{g(\omega)} + \norm{g'(\omega)}\big)\Big) =: D(\omega),
\]
where $D$ does not depend on $\theta$ and is integrable on $\Omega$ (product of two
independent integrable variables, Lemma~\ref{lem:pc_exp_norm_integrable}); hence the function
is integrable on $\Omega\times[0,\pi/2]$ and Fubini (\texttt{integral\_integral\_swap})
applies:
\[
  \E\exp\big(\lambda(f(g) - f(g'))\big)
  \le \frac2\pi\int_0^{\pi/2}\E\exp\Big(\frac{\pi\lambda}{2}\ip{\nabla f(g_\theta)}{g'_\theta}\Big)d\theta .
\]
For each fixed $\theta$, Lemma~\ref{lem:pc_rotated_mgf_bound} with $s = \pi\lambda/2$ gives
\[
  \E\exp\Big(\frac{\pi\lambda}{2}\ip{\nabla f(g_\theta)}{g'_\theta}\Big)
  \le \exp\Big(\frac{(\pi\lambda/2)^2L^2}{2}\Big) = \exp\Big(\frac{\pi^2\lambda^2L^2}{8}\Big).
\]

\emph{Step 5 (conclusion).} Integrating the constant bound over $\theta \in [0,\pi/2]$ and
multiplying by $2/\pi$,
\[
  \E\exp\big(\lambda(f(g) - \E f(g))\big) \le \frac2\pi\cdot\frac\pi2\cdot\exp\Big(\frac{\pi^2\lambda^2L^2}{8}\Big)
  = \exp\Big(\frac{\pi^2\lambda^2L^2}{8}\Big).
\]
Finally $\frac{\pi^2L^2}{8}\lambda^2 = \frac{(\pi^2/4)L^2\,\lambda^2}{2} = \frac{\cG L^2\lambda^2}{2}$,
which together with Step~0 is the statement \texttt{HasSubgaussianMGF} with constant
$\cG L^2$.
\end{proof}

\begin{remark}[On the constant]\label{rem:pc_constant}
The factor $\pi^2/8$ in the exponent comes from Step~3: Jensen's inequality replaces
$\lambda\int_0^{\pi/2}h$ by an average of $\exp(\frac\pi2\lambda h(\theta))$, and the length
$\pi/2$ of the interval is squared in the Gaussian MGF of Step~4. With the standard
Gaussian interpolation $g_t = \sqrt t\,g + \sqrt{1-t}\,g'$ one would need the derivative
$\frac{d}{dt}g_t$, which is singular at the endpoints; the quarter-circle parametrization
avoids this at the cost of the constant. Sharper routes (Gaussian log-Sobolev + Herbst, or
the Gaussian isoperimetric inequality) give $\cG = 1$ \cite{boucheron2013concentration,
adler2007random}, and the Borell--TIS inequality with $\cG = 1$ is the form usually quoted;
none of the results of Part~\ref{part:main} is sensitive to the value of $\cG$ beyond the
absolute constants.
\end{remark}

\section{Suprema of linear Gaussian processes}
\label{sec:pre_concentration_sup}

\begin{lemma}[Concentration of a finite maximum of Gaussian linear forms]
\label{lem:finite_sup_subgaussian}
\uses{def:pg_gaussian_vector, def:pc_gaussian_constant}
Let $m \ge 1$, $x_1,\dots,x_m \in \R^d$, $S \in \PSD$ and $G \sim \Normal(0,S)$. Let
$M := \max_{i\le m}\ip{x_i}{G}$ and $\sigma^2 := \max_{i\le m}x_i^\top Sx_i$. Then $M$ is
integrable and $M - \E M$ has a sub-Gaussian MGF with constant $\cG\sigma^2$: for every
$\lambda\in\R$, $e^{\lambda(M - \E M)}$ is integrable and
$\E e^{\lambda(M-\E M)} \le \exp(\cG\sigma^2\lambda^2/2)$.
\end{lemma}

\begin{proof}
\uses{lem:gauss_law_of_cov, lem:maurey_pisier, lem:logsumexp_props, lem:pc_exp_norm_integrable}
The statement only involves the law of $M$, a measurable function of $G$, so by
Lemma~\ref{lem:gauss_law_of_cov}(iii) we may assume $G = S^{1/2}g$ with $g \sim \Normal(0,I_d)$
(Mathlib \texttt{HasSubgaussianMGF.congr\_identDistrib} / \texttt{of\_map} transport the
conclusion). Since $S^{1/2}$ is symmetric, $\ip{x_i}{S^{1/2}g} = \ip{y_i}{g}$ with
$y_i := S^{1/2}x_i$, and $\norm{y_i}^2 = x_i^\top Sx_i \le \sigma^2$. Thus $M = F(g)$ with
$F(v) := \max_{i\le m}\ip{y_i}{v}$, and $\abs{F(v)} \le \sigma\norm{v}$, so $M$ and
$e^{\lambda M}$ are integrable (Lemma~\ref{lem:pc_exp_norm_integrable}).

Fix $\lambda \in \R$ and $\beta > 0$, and let $f_\beta$ be the log-sum-exp function of
Lemma~\ref{lem:logsumexp_props} for the vectors $y_1,\dots,y_m$. It is $C^\infty$ with
$\norm{\nabla f_\beta} \le \max_i\norm{y_i} \le \sigma$, so Lemma~\ref{lem:maurey_pisier} with
$L = \sigma$ gives $\E e^{\lambda(f_\beta(g) - \E f_\beta(g))} \le \exp(\cG\sigma^2\lambda^2/2)$.
By Lemma~\ref{lem:logsumexp_props}(ii), $0 \le f_\beta(g) - M \le (\log m)/\beta$ pointwise,
hence also $0 \le \E f_\beta(g) - \E M \le (\log m)/\beta$, and therefore
\[
  \abs{(M - \E M) - (f_\beta(g) - \E f_\beta(g))} \le \frac{2\log m}{\beta}
  \quad\text{pointwise, so}\quad
  \E e^{\lambda(M - \E M)} \le e^{2\abs\lambda(\log m)/\beta}\,\E e^{\lambda(f_\beta(g) - \E f_\beta(g))}
  \le e^{2\abs\lambda(\log m)/\beta}\exp\Big(\frac{\cG\sigma^2\lambda^2}{2}\Big).
\]
The left-hand side does not depend on $\beta$; letting $\beta \to \infty$ (or taking the
infimum over $\beta > 0$: a real number bounded by $a e^{\eps}$ for every $\eps > 0$ is bounded
by $a$) gives $\E e^{\lambda(M-\E M)} \le \exp(\cG\sigma^2\lambda^2/2)$. No limit theorem for
integrals is needed here.
\end{proof}

\begin{lemma}[Concentration of the supremum over a compact index set]
\label{lem:compact_sup_subgaussian}
\uses{def:pg_gaussian_vector, def:pc_gaussian_constant, lem:sup_countable_dense}
Let $\Xs \subset \R^d$ be nonempty and compact, $S \in \PSD$ and $G \sim \Normal(0,S)$. Let
$M := \sup_{x\in\Xs}\ip{x}{G}$ (a random variable by Lemma~\ref{lem:sup_countable_dense})
and $\sigma^2 := \sup_{x\in\Xs}x^\top Sx$ (finite, attained). Then $M$ is integrable and
$M - \E M$ has a sub-Gaussian MGF with constant $\cG\sigma^2$: for every $\lambda\in\R$,
$e^{\lambda(M-\E M)}$ is integrable and $\E e^{\lambda(M - \E M)} \le \exp(\cG\sigma^2\lambda^2/2)$.
\end{lemma}

\begin{proof}
\uses{lem:sup_countable_dense, lem:finite_sup_subgaussian, lem:pc_exp_norm_integrable}
Let $R := \max_{x\in\Xs}\norm{x}$ and let $(q_n)_{n\in\N}$ enumerate a countable dense subset
of $\Xs$ (Lemma~\ref{lem:sup_countable_dense}(i)). $\sigma^2$ is the maximum of the continuous
function $x\mapsto x^\top Sx = \norm{S^{1/2}x}^2$ on the compact $\Xs$. Define the partial
maxima $M_n := \max_{i\le n}\ip{q_i}{G}$, $n \in \N$. By Lemma~\ref{lem:sup_countable_dense}(ii)
applied pointwise to $v = G(\omega)$, $M_n \uparrow M$ everywhere, and
$\abs{M_n} \le R\norm{G}$, $\abs{M} \le R\norm{G}$.

\emph{Bound for each $n$.} Lemma~\ref{lem:finite_sup_subgaussian} applied to $q_0,\dots,q_n$
gives, with $\sigma_n^2 := \max_{i\le n}q_i^\top Sq_i \le \sigma^2$, that
$\E e^{\lambda(M_n - \E M_n)} \le \exp(\cG\sigma_n^2\lambda^2/2) \le \exp(\cG\sigma^2\lambda^2/2)$
for every $\lambda$.

\emph{Passage to the limit.} Fix $\lambda\in\R$. The dominating function
$e^{\abs\lambda R\norm{G}}$ is integrable (Lemma~\ref{lem:pc_exp_norm_integrable}) and
$e^{\lambda M_n} \le e^{\abs\lambda R\norm{G}}$ for all $n$, while $e^{\lambda M_n} \to e^{\lambda M}$
pointwise (continuity of $\exp$). Dominated convergence
(\texttt{tendsto\_integral\_of\_dominated\_convergence}) gives
$\E e^{\lambda M_n} \to \E e^{\lambda M}$, and $e^{\lambda M}$ is integrable. Likewise
$\abs{M_n} \le R\norm{G}$, $M_n \to M$, and $R\norm{G}$ is integrable, so
$\E M_n \to \E M$ (dominated convergence again, or monotone convergence
\texttt{integral\_tendsto\_of\_tendsto\_of\_monotone} since $M_n$ is nondecreasing); in
particular $M$ is integrable. Therefore
\[
  \E e^{\lambda(M - \E M)} = e^{-\lambda\E M}\,\E e^{\lambda M}
  = \lim_{n\to\infty}e^{-\lambda\E M_n}\,\E e^{\lambda M_n}
  = \lim_{n\to\infty}\E e^{\lambda(M_n - \E M_n)} \le \exp\Big(\frac{\cG\sigma^2\lambda^2}{2}\Big),
\]
as a limit of numbers bounded by the right-hand side (\texttt{le\_of\_tendsto}). Integrability
of $e^{\lambda(M-\E M)} = e^{-\lambda\E M}e^{\lambda M}$ was shown above.
\end{proof}

\textbf{Lean remark.} The supremum over $\Xs$ is encoded as an \texttt{iSup} over the dense
sequence \texttt{q : ℕ → Xs} (Lemma~\ref{lem:sup_countable_dense}), so that $M_n$ is
\texttt{⨆ i : Fin (n+1), ⟪q i, G ω⟫} (or a \texttt{Finset.sup'} over \texttt{Finset.range}),
and the two limits are \texttt{Filter.Tendsto} statements along \texttt{atTop}. Lemma
\ref{lem:finite_sup_subgaussian} is stated for an indexed family \texttt{x : Fin m → E}, $m \ge 1$,
and the constant $\sigma^2$ is \texttt{⨆ i, ⟪x i, S (x i)⟫} coerced to \texttt{ℝ≥0}.

\begin{theorem}[Borell--TIS inequality for linear Gaussian processes]\label{thm:borell_tis}
\uses{def:pg_gaussian_vector, def:pc_gaussian_constant, lem:sup_countable_dense}
Let $\Xs \subset \R^d$ be nonempty and compact, $S \in \PSD$, $G \sim \Normal(0,S)$,
$M := \sup_{x\in\Xs}\ip{x}{G}$ and $\sigma^2 := \sup_{x\in\Xs}x^\top Sx$, assumed positive.
Then $M$ is integrable and for every $u \ge 0$,
\[
  \Prob\big(M - \E M \ge u\big) \le \exp\Big(-\frac{u^2}{2\cG\sigma^2}\Big), \qquad
  \Prob\big(M - \E M \le -u\big) \le \exp\Big(-\frac{u^2}{2\cG\sigma^2}\Big).
\]
(If $\sigma = 0$ then $\ip{x}{G} = 0$ a.s.\ for every $x\in\Xs$, $M = 0$ a.s., and both
probabilities vanish for $u > 0$.)
\end{theorem}

\begin{proof}
\uses{lem:compact_sup_subgaussian}
By Lemma~\ref{lem:compact_sup_subgaussian}, $M - \E M$ satisfies \texttt{HasSubgaussianMGF}
with constant $c := \cG\sigma^2$. The upper tail is the Chernoff bound
\texttt{HasSubgaussianMGF.measure\_ge\_le}: $\Prob(X \ge u) \le \exp(-u^2/(2c))$ for $u \ge 0$;
the lower tail is \texttt{HasSubgaussianMGF.measure\_le\_le} (LML; it is
\texttt{measure\_ge\_le} applied to $-X$, which is sub-Gaussian with the same constant by
\texttt{HasSubgaussianMGF.neg}). The degenerate case: $\Var\ip{x}{G} = x^\top Sx = 0$ forces
$\ip{x}{G} = 0$ a.s.\ for each $x$, hence for all $x$ in the countable dense set
simultaneously, hence $M = 0$ a.s.\ by Lemma~\ref{lem:sup_countable_dense}(ii).
\end{proof}

\begin{remark}[The classical statement]\label{rem:pc_borell_tis_classical}
The Borell--TIS inequality \cite[Theorem~2.1.1]{adler2007random} holds for every
almost surely bounded centered Gaussian process with $\cG$ replaced by $1$, and its proof
goes through the Gaussian isoperimetric inequality. Only the linear processes
$x\mapsto\ip{x}{G}$ over compact $\Xs \subset \R^d$ are needed in this blueprint, and only up
to the absolute constant $\cG$; the general statement is not pursued.
\end{remark}

\begin{corollary}[High-probability upper bound on the supremum]\label{cor:sup_bound_whp}
\uses{thm:borell_tis}
In the setting of Theorem~\ref{thm:borell_tis}, for every $\delta \in (0,1)$,
\[
  \Prob\Big(M \le \E M + \sigma\sqrt{2\cG\log(1/\delta)}\Big) \ge 1 - \delta, \qquad
  \Prob\Big(M \ge \E M - \sigma\sqrt{2\cG\log(1/\delta)}\Big) \ge 1 - \delta .
\]
\end{corollary}

\begin{proof}
\uses{thm:borell_tis}
Let $u := \sigma\sqrt{2\cG\log(1/\delta)} \ge 0$; then $u^2/(2\cG\sigma^2) = \log(1/\delta)$
and Theorem~\ref{thm:borell_tis} gives $\Prob(M - \E M \ge u) \le e^{-\log(1/\delta)} = \delta$.
The event $\{M \le \E M + u\}$ contains the complement of $\{M - \E M \ge u\}$ (in fact
equals the complement of $\{M - \E M > u\}$), so its probability is at least $1 - \delta$. The
second bound is the lower tail. If $\sigma = 0$ both events have probability one.
\end{proof}

\begin{corollary}[Symmetric index set]\label{cor:sup_bound_symmetric}
\uses{thm:borell_tis, lem:gauss_law_of_cov}
Let $\Xs \subset \R^d$ be nonempty and compact, $S \in \PSD$, $G \sim \Normal(0,S)$,
$\sigma^2 := \sup_{x\in\Xs}x^\top Sx$, and consider the difference set
$\Xs - \Xs := \{x - x' : x, x' \in \Xs\}$, which is nonempty and compact. Let
$D := \sup_{x,x'\in\Xs}\ip{x - x'}{G} = \sup_{z\in\Xs-\Xs}\ip{z}{G}$. Then
\begin{enumerate}
  \item[(i)] $\E D = 2\,\E\sup_{x\in\Xs}\ip{x}{G}$;
  \item[(ii)] $\sigma^2_{\Xs-\Xs} := \sup_{z\in\Xs-\Xs}z^\top Sz \le 4\sigma^2$;
  \item[(iii)] for every $\delta\in(0,1)$,
        $\Prob\big(D \le \E D + 2\sigma\sqrt{2\cG\log(1/\delta)}\big) \ge 1 - \delta$.
\end{enumerate}
\end{corollary}

\begin{proof}
\uses{cor:sup_bound_whp, lem:gauss_law_of_cov, lem:sup_countable_dense}
$\Xs - \Xs$ is the image of the compact $\Xs\times\Xs$ under the continuous map
$(x,x')\mapsto x - x'$, hence compact. (i) Pointwise,
$\sup_{x,x'}(\ip{x}{G} - \ip{x'}{G}) = \sup_x\ip{x}{G} + \sup_{x'}\ip{x'}{-G}$ (the two suprema
are over independent variables $x, x'$), and $-G \sim \Normal(0,S)$
(Lemma~\ref{lem:gauss_law_of_cov}(iv)), so $\E\sup_{x'}\ip{x'}{-G} = \E\sup_x\ip{x}{G}$; both
terms are integrable by Lemma~\ref{lem:sup_countable_dense}. (ii) For $z = x - x'$,
$z^\top Sz = \norm{S^{1/2}x - S^{1/2}x'}^2 \le (\norm{S^{1/2}x} + \norm{S^{1/2}x'})^2 \le (2\sigma)^2$.
(iii) Apply Corollary~\ref{cor:sup_bound_whp} to the index set $\Xs - \Xs$: with probability
at least $1 - \delta$, $D \le \E D + \sigma_{\Xs-\Xs}\sqrt{2\cG\log(1/\delta)} \le
\E D + 2\sigma\sqrt{2\cG\log(1/\delta)}$ by (ii). (If $\sigma_{\Xs - \Xs} = 0$ the bound holds
with probability one.)
\end{proof}

\textbf{Lean remark.} In the applications (Lemmas~\ref{lem:difference_process_concentration}
and~\ref{lem:region_phase1}) the index set is $\Xs - \Xs$ for the action set or for one region
of a partition, $S = \Sigma_T^{-1}$ for a fixed design, and $\sigma^2 = \max_{x}\norm{x}^2_{\Sigma_T^{-1}}$;
the compact set $\Xs - \Xs$ is \texttt{Set.image2 (· - ·) Xs Xs}, which is the image of the
product set $\Xs\times\Xs$ (\texttt{Set.prod}) under \texttt{fun p ↦ p.1 - p.2}
(\texttt{Set.image\_prod}); with \texttt{hK : IsCompact Xs} it is compact by
\texttt{(hK.prod hK).image continuous\_sub} (Mathlib \texttt{IsCompact.image}), or equivalently
\texttt{(hK.prod hK).image\_of\_continuousOn continuous\_sub.continuousOn}.

\chapter{Sub-exponential random variables}
\label{chap:pre_subexp}

This chapter develops the Bernstein-type concentration used by the $\ell_2$-norm estimation
of Chapter~\ref{chap:norm} and by the unit-ball identification algorithm
(Lemma~\ref{lem:ball_identification}): a notion of sub-exponential random variable with two
explicit parameters $(V, b)$ (Appendix~D.1 of the paper), its stability under scaling and
independent sums, the corresponding tail bound with the explicit constant $c = 1/2$, and the
sub-exponential parameters of the centered squares of Gaussian variables (the paper's
Lemma~1 in Appendix~A, with its complete-the-square proof), of chi-square variables and of
Rademacher linear forms $\ip{\eps}{\theta}$. For the last one the paper invokes the
Hanson--Wright inequality with unspecified constants; we replace it by a direct argument valid
for every sub-Gaussian variable $Z$ with $\E Z^2$ equal to its sub-Gaussian constant (a Gaussian
trick for positive $\lambda$, a fourth-moment bound for negative $\lambda$, the fourth moment
being itself controlled by the Gaussian trick), which gives explicit parameters
$(16\norm\theta^4, 4\norm\theta^2)$. All constants are explicit and tracked. References:
\cite[Section~2.7]{vershynin2018high}, \cite[Section~2.4]{boucheron2013concentration}.

\textbf{What Mathlib provides.} The sub-Gaussian structure
\texttt{ProbabilityTheory.HasSubgaussianMGF X c μ} (\texttt{Mathlib/Probability/Moments/SubGaussian.lean})
with \texttt{measure\_ge\_le}, \texttt{neg}, \texttt{add\_of\_indepFun},
\texttt{sum\_of\_iIndepFun}, \texttt{of\_map}, \texttt{congr\_identDistrib}, Hoeffding's lemma
\texttt{hasSubgaussianMGF\_of\_mem\_Icc\_of\_integral\_eq\_zero}, and LML's
\texttt{measure\_le\_le}. Moment generating functions: \texttt{ProbabilityTheory.mgf},
\texttt{iIndepFun.mgf\_sum}, \texttt{iIndepFun.integrable\_exp\_mul\_sum}, the Chernoff bound
\texttt{measure\_ge\_le\_exp\_mul\_mgf} (\texttt{Mathlib/Probability/Moments/Basic.lean}),
the derivative of the MGF \texttt{hasDerivAt\_mgf} (\texttt{MGFAnalytic.lean}), integrability
of all moments from the integrability of $e^{\lambda X}$ near $\lambda = 0$
(\texttt{integrable\_pow\_of\_mem\_interior\_integrableExpSet}, \texttt{IntegrableExpMul.lean}),
\texttt{mgf\_gaussianReal}. The Gaussian integral \texttt{integral\_gaussian}
($\int e^{-bx^2}dx = \sqrt{\pi/b}$,
\texttt{Mathlib/Analysis/SpecialFunctions/Gaussian/GaussianIntegral.lean}); the real logarithm
series bound \texttt{Real.abs\_log\_sub\_add\_sum\_range\_le}
(\texttt{Mathlib/Analysis/SpecialFunctions/Log/Deriv.lean}); \texttt{Real.add\_one\_le\_exp} and
\texttt{Real.quadratic\_le\_exp\_of\_nonneg}; independence of the coordinates of a product
measure \texttt{ProbabilityTheory.iIndepFun\_pi}; the standard Gaussian vector
\texttt{ProbabilityTheory.stdGaussian} as the image of a product of \texttt{gaussianReal 0 1}
(\texttt{map\_pi\_eq\_stdGaussian}).

\textbf{What is missing.} The sub-exponential structure itself and all its properties
(Section~\ref{sec:pre_subexp_def}), the MGF of the square of a real Gaussian
(Lemma~\ref{lem:gaussian_square_mgf}), and the sub-exponential bounds of
Sections~\ref{sec:pre_subexp_gaussian} and~\ref{sec:pre_subexp_rademacher}.

\textbf{Lean remark.} The definition below mirrors \texttt{HasSubgaussianMGF}: a structure
\texttt{HasSubexponentialMGF (X : Ω → ℝ) (V b : ℝ) (μ : Measure Ω) : Prop} with fields
\texttt{integrable\_exp\_mul : ∀ t : ℝ, b * |t| ≤ 1 → Integrable (fun ω ↦ exp (t * X ω)) μ} and
\texttt{mgf\_le : ∀ t : ℝ, b * |t| ≤ 1 → mgf X μ t ≤ exp (V * t\^{}2 / 2)}
(\texttt{COLT83/Mathlib/Probability/Moments/SubExponential.lean}). The range of $t$ is
written as $b\abs t\le1$ rather than $\abs t\le1/b$ so that \texttt{b = 0} imposes no
restriction on $t$ (the hypothesis $0\cdot\abs t\le1$ is always true):
\texttt{HasSubexponentialMGF X V 0 μ} is then equivalent to \texttt{HasSubgaussianMGF X V μ}
(same two fields, quantified over all \texttt{t}), and Mathlib's sub-Gaussian API is the
$b = 0$ case of the statements below (Lemma~\ref{lem:subexp_of_subgaussian}). Unlike
Mathlib's sub-Gaussian constant, the parameters are real numbers (not \texttt{ℝ≥0}), which
avoids coercions in the arithmetic of Chapter~\ref{chap:norm}; the lemmas that need $V > 0$
or $b > 0$ assume it explicitly. A kernel version, as for \texttt{Kernel.HasSubgaussianMGF}, is
not needed.

\section{Definition and basic properties}
\label{sec:pre_subexp_def}

\begin{definition}[Sub-exponential moment generating function]\label{def:subexp}
\lean{ProbabilityTheory.HasSubexponentialMGF}\leanok
Let $X : \Omega \to \R$ be a random variable on $(\Omega,\Prob)$ and $V, b \in \R$. We say
that $X$ \emph{has a sub-exponential MGF with parameters $(V,b)$}, written
$\mathrm{HasSubexponentialMGF}(X, V, b)$, if for every $\lambda \in \R$ with
$b\abs{\lambda} \le 1$ (i.e.\ $\abs\lambda\le1/b$ when $b>0$, and every $\lambda\in\R$ when
$b \le 0$), the variable $e^{\lambda X}$ is integrable and
\[
  \E e^{\lambda X} \le \exp\Big(\frac{\lambda^2V}{2}\Big).
\]
The case $b = 0$ is exactly Mathlib's \texttt{HasSubgaussianMGF X V}
(Lemma~\ref{lem:subexp_of_subgaussian}). No centering is assumed; it is a consequence
(Lemma~\ref{lem:ps_subexp_integral_zero}). In practice $V \ge 0$ and $b \ge 0$.
\end{definition}

\begin{lemma}[Sub-exponential variables are integrable and centered]\label{lem:ps_subexp_integral_zero}
\lean{ProbabilityTheory.HasSubexponentialMGF.zero_mem_interior_integrableExpSet, ProbabilityTheory.HasSubexponentialMGF.integrable, ProbabilityTheory.HasSubexponentialMGF.memLp, ProbabilityTheory.HasSubexponentialMGF.integrable_pow, ProbabilityTheory.HasSubexponentialMGF.integral_eq_zero}\leanok
\uses{def:subexp}
If $\mathrm{HasSubexponentialMGF}(X,V,b)$ then $X$ is integrable, all its moments are finite,
and, when $\Prob$ is a probability measure, $\E X = 0$.
\end{lemma}

\begin{proof}
\leanok
\uses{def:subexp}
The set $\{\lambda : b\abs\lambda \le 1\}$ is a neighbourhood of $0$ (for every $b$, since
$\lambda \mapsto b\abs\lambda$ is continuous and vanishes at $0$), so $0$ lies in the interior
of the set of $\lambda$ such that $e^{\lambda X}$ is integrable. Mathlib then provides the
integrability of $\abs X^p$ for every $p$
(\texttt{integrable\_pow\_of\_mem\_interior\_integrableExpSet}) and the differentiability of
$\lambda \mapsto \E e^{\lambda X}$ at $0$ with derivative $\E X$ (\texttt{hasDerivAt\_mgf}).
Centering: the function $\varphi(\lambda) := \E e^{\lambda X} - e^{\lambda^2V/2}$ is
nonpositive on a neighbourhood of $0$ and vanishes at $0$ (as $\E e^{0} = 1$ for a probability
measure), so it has a local maximum at $0$; it is differentiable there with derivative
$\E X - 0$, which must vanish (Fermat's rule, Mathlib \texttt{IsLocalMax.hasDerivAt\_eq\_zero}).
Hence $\E X = 0$.
\end{proof}

\begin{lemma}[Monotonicity in the parameters]\label{lem:ps_subexp_mono}
\lean{ProbabilityTheory.HasSubexponentialMGF.mono}\leanok
\uses{def:subexp}
If $\mathrm{HasSubexponentialMGF}(X,V,b)$, $V \le V'$ and $b \le b'$, then
$\mathrm{HasSubexponentialMGF}(X,V',b')$.
\end{lemma}

\begin{proof}
\leanok
\uses{def:subexp}
$b'\abs\lambda \le 1$ implies $b\abs\lambda \le b'\abs\lambda \le 1$, and
$e^{\lambda^2V/2} \le e^{\lambda^2V'/2}$.
\end{proof}

\begin{lemma}[Sub-Gaussian is sub-exponential with $b = 0$]\label{lem:subexp_of_subgaussian}
\lean{ProbabilityTheory.HasSubexponentialMGF.of_hasSubgaussianMGF, ProbabilityTheory.HasSubexponentialMGF.hasSubgaussianMGF}\leanok
\uses{def:subexp, lem:ps_subexp_mono}
$X$ has a sub-Gaussian MGF with constant $c \ge 0$ (Mathlib
\texttt{HasSubgaussianMGF X c}: $e^{\lambda X}$ integrable and $\E e^{\lambda X} \le e^{c\lambda^2/2}$
for all $\lambda\in\R$) if and only if $\mathrm{HasSubexponentialMGF}(X, c, 0)$. Consequently,
if $X$ has a sub-Gaussian MGF with constant $c$, then $\mathrm{HasSubexponentialMGF}(X, c, b)$
for every $b \in \R$.
\end{lemma}

\begin{proof}
\leanok
\uses{def:subexp, lem:ps_subexp_mono}
For $b = 0$ the condition $b\abs\lambda\le1$ holds for every $\lambda\in\R$, so the two
requirements of Definition~\ref{def:subexp} are literally the two fields of
\texttt{HasSubgaussianMGF}. The second claim is the same argument: the condition
$b\abs\lambda \le 1$ only restricts the set of $\lambda$ for which the two fields are required.
\end{proof}

\begin{lemma}[Scaling]\label{lem:subexp_const_mul}
\lean{ProbabilityTheory.HasSubexponentialMGF.const_mul, ProbabilityTheory.HasSubexponentialMGF.neg}\leanok
\uses{def:subexp}
If $\mathrm{HasSubexponentialMGF}(X,V,b)$ and $a \in \R$, then
$\mathrm{HasSubexponentialMGF}(aX,\ a^2V,\ \abs{a}\,b)$. In particular ($a = -1$),
$\mathrm{HasSubexponentialMGF}(-X, V, b)$.
\end{lemma}

\begin{proof}
\leanok
\uses{def:subexp}
Let $(\abs ab)\abs\lambda \le 1$, i.e.\ $b\abs{\lambda a} \le 1$; then $e^{\lambda(aX)} = e^{(\lambda a)X}$
is integrable and $\E e^{(\lambda a)X} \le e^{(\lambda a)^2V/2} = e^{\lambda^2(a^2V)/2}$
(for $a = 0$ this reads $\E e^{0} = 1 \le 1$).
\end{proof}

\begin{lemma}[Transport]\label{lem:ps_subexp_map}
\lean{ProbabilityTheory.HasSubexponentialMGF.congr, ProbabilityTheory.HasSubexponentialMGF.of_map, ProbabilityTheory.HasSubexponentialMGF.map_iff, ProbabilityTheory.HasSubexponentialMGF.id_map_iff, ProbabilityTheory.HasSubexponentialMGF.congr_identDistrib}\leanok
\uses{def:subexp}
The property $\mathrm{HasSubexponentialMGF}(X,V,b)$ depends only on the law of $X$: if
$X = Y$ almost surely, or if $X$ and $Y$ have the same law (on possibly different probability
spaces), then $X$ has parameters $(V,b)$ if and only if $Y$ does. If $Y : \Omega' \to \Omega$ is
measurable and $X : \Omega \to \R$ has parameters $(V,b)$ under the image measure
$\Prob' \circ Y^{-1}$, then $X \circ Y$ has parameters $(V,b)$ under $\Prob'$.
\end{lemma}

\begin{proof}
\leanok
\uses{def:subexp}
Both fields of Definition~\ref{def:subexp} are expressed through integrals of functions of
$X$, which are unchanged by these operations (change of variables
\texttt{integrable\_map\_measure}, \texttt{mgf\_map}). This is the same proof as for Mathlib's
\texttt{HasSubgaussianMGF.of\_map} and \texttt{congr\_identDistrib}.
\end{proof}

\begin{lemma}[Bernstein-type tail bound]\label{lem:subexp_tail}
\lean{ProbabilityTheory.HasSubexponentialMGF.measure_ge_le_exp_neg_sq, ProbabilityTheory.HasSubexponentialMGF.measure_ge_le_exp_neg_div, ProbabilityTheory.HasSubexponentialMGF.measure_ge_le, ProbabilityTheory.HasSubexponentialMGF.measure_le_le, ProbabilityTheory.HasSubexponentialMGF.measure_abs_ge_le, ProbabilityTheory.HasSubexponentialMGF.measure_abs_gt_le}\leanok
\uses{def:subexp}
Let $\mathrm{HasSubexponentialMGF}(X,V,b)$ and $t \ge 0$. Then
\[
  \Prob(X \ge t) \le
  \begin{cases}
    \exp\big(-\frac{t^2}{2V}\big) & \text{if } tb \le V \text{ and } V > 0,\\[2pt]
    \exp\big(-\frac{t}{2b}\big) & \text{if } tb \ge V \text{ and } b > 0 ,
  \end{cases}
\]
and consequently, when $V > 0$ and $b > 0$,
\[
  \Prob(X \ge t) \le \exp\Big(-\min\Big(\frac{t^2}{2V}, \frac{t}{2b}\Big)\Big), \quad
  \Prob(X \le -t) \le \exp\Big(-\min\Big(\frac{t^2}{2V}, \frac{t}{2b}\Big)\Big), \quad
  \Prob(\abs X \ge t) \le 2\exp\Big(-\min\Big(\frac{t^2}{2V}, \frac{t}{2b}\Big)\Big).
\]
If $b = 0$ (the sub-Gaussian case) and $V > 0$, the first bound applies for every $t \ge 0$
(the condition $tb \le V$ is automatic): this is Mathlib
\texttt{HasSubgaussianMGF.measure\_ge\_le} (and LML \texttt{measure\_le\_le} for the lower tail),
i.e.\ the $\min$ form with the convention $t/(2b) = +\infty$.
\end{lemma}

\begin{proof}
\leanok
\uses{def:subexp, lem:subexp_const_mul}
For $\lambda \ge 0$ with $b\lambda \le 1$, Markov's inequality applied to the nonnegative
integrable variable $e^{\lambda X}$ gives (Chernoff, Mathlib
\texttt{measure\_ge\_le\_exp\_mul\_mgf})
\[
  \Prob(X \ge t) = \Prob(e^{\lambda X} \ge e^{\lambda t}) \le e^{-\lambda t}\,\E e^{\lambda X}
  \le \exp\Big(-\lambda t + \frac{\lambda^2V}{2}\Big).
\]
\emph{Case $tb \le V$, $V > 0$:} take $\lambda := t/V \ge 0$, which satisfies $b\lambda \le 1$;
the exponent is $-t^2/V + t^2/(2V) = -t^2/(2V)$. \emph{Case $tb \ge V$, $b > 0$:} take
$\lambda := 1/b$; the exponent is $-t/b + V/(2b^2) \le -t/b + t/(2b) = -t/(2b)$ since
$V \le tb$. For the $\min$ form it is not necessary to identify which term realizes the
minimum: in each case the bound obtained is $\exp(-u)$ for one of the two terms $u$ of the
minimum, and $\exp(-u) \le \exp(-\min(\cdot,\cdot))$. The lower tail is the upper tail of
$-X$, which has the same parameters (Lemma~\ref{lem:subexp_const_mul} with $a = -1$); the
two-sided bound is a union bound over $\{X \ge t\} \cup \{X \le -t\}$.
\end{proof}

\begin{remark}[The constant $c$ of the paper]\label{rem:ps_tail_constant}
The paper's Lemma~2 (Appendix~A) states this bound with an unspecified absolute constant $c$
in the exponent; the computation above, which is the one given there, yields $c = 1/2$. This
value is used in Chapter~\ref{chap:norm} to make all sample-complexity constants explicit.
\end{remark}

\begin{lemma}[Independent sums]\label{lem:subexp_sum}
\lean{ProbabilityTheory.HasSubexponentialMGF.add_of_indepFun, ProbabilityTheory.HasSubexponentialMGF.sum_of_iIndepFun, ProbabilityTheory.HasSubexponentialMGF.fun_sum_of_iIndepFun}\leanok
\uses{def:subexp}
Let $X_1,\dots,X_K$ be independent with $\mathrm{HasSubexponentialMGF}(X_i,V_i,b)$ for every
$i \le K$ (the same $b$). Then $\mathrm{HasSubexponentialMGF}\big(\sum_{i\le K}X_i,\ \sum_{i\le K}V_i,\ b\big)$.
\end{lemma}

\begin{proof}
\leanok
\uses{def:subexp}
Let $b\abs\lambda \le 1$. Each $e^{\lambda X_i}$ is integrable and they are independent, so
$e^{\lambda\sum_iX_i} = \prod_ie^{\lambda X_i}$ is integrable (Mathlib
\texttt{iIndepFun.integrable\_exp\_mul\_sum}) and
$\E e^{\lambda\sum_iX_i} = \prod_i\E e^{\lambda X_i}$ (\texttt{iIndepFun.mgf\_sum})
$\le \prod_ie^{\lambda^2V_i/2} = e^{\lambda^2\sum_iV_i/2}$.
\end{proof}

\begin{corollary}[Tail of an average]\label{cor:subexp_average_tail}
\lean{ProbabilityTheory.HasSubexponentialMGF.average_of_iIndepFun, ProbabilityTheory.HasSubexponentialMGF.measure_abs_average_ge_le, ProbabilityTheory.HasSubexponentialMGF.measure_abs_average_ge_le_of_forall}\leanok
\uses{def:subexp}
Let $K \ge 1$ and $X_1,\dots,X_K$ be independent with $\mathrm{HasSubexponentialMGF}(X_i,V_i,b)$,
and let $\bar V := \frac1K\sum_{i\le K}V_i$. Then
$\mathrm{HasSubexponentialMGF}\big(\frac1K\sum_{i\le K}X_i,\ \bar V/K,\ b/K\big)$ and, if
$\bar V > 0$ and $b > 0$, for every $t \ge 0$,
\[
  \Prob\Big(\Big|\frac1K\sum_{i\le K}X_i\Big| \ge t\Big)
  \le 2\exp\Big(-K\min\Big(\frac{t^2}{2\bar V}, \frac{t}{2b}\Big)\Big).
\]
When all the $V_i$ are equal to $V$, $\bar V = V$.
\end{corollary}

\begin{proof}
\leanok
\uses{lem:subexp_sum, lem:subexp_const_mul, lem:subexp_tail}
Lemma~\ref{lem:subexp_sum} gives parameters $(\sum_iV_i, b) = (K\bar V, b)$ for the sum, and
Lemma~\ref{lem:subexp_const_mul} with $a = 1/K$ gives $(K\bar V/K^2, b/K) = (\bar V/K, b/K)$ for
the average. Lemma~\ref{lem:subexp_tail} (with $b/K > 0$ and $\bar V/K > 0$) then gives the
two-sided bound
$2\exp(-\min(t^2/(2\bar V/K), t/(2b/K))) = 2\exp(-K\min(t^2/(2\bar V), t/(2b)))$.
(For $b = 0$ the same computation with the $b = 0$ clause of Lemma~\ref{lem:subexp_tail} gives
$2\exp(-Kt^2/(2\bar V))$, Mathlib's \texttt{measure\_sum\_ge\_le\_of\_iIndepFun}.)
\end{proof}

\section{Squares of Gaussian variables}
\label{sec:pre_subexp_gaussian}

\begin{lemma}[MGF of the square of a real Gaussian]\label{lem:gaussian_square_mgf}
\lean{integral_exp_quadratic, integrable_exp_quadratic, ProbabilityTheory.integral_exp_mul_sq_gaussianReal, ProbabilityTheory.integrable_exp_mul_sq_gaussianReal}\leanok
\uses{def:pg_gaussian_vector}
For $b < 0$ and $c, d \in \R$,
\[
  \int_{\R} \exp(bx^2 + cx + d)\,dx = \sqrt{\frac{\pi}{-b}}\,\exp\Big(d - \frac{c^2}{4b}\Big),
\]
the integrand being integrable. Consequently, let $\mu \in \R$, $\sigma \ge 0$,
$X \sim \Normal(\mu,\sigma^2)$ and $\lambda \in \R$ with $2\sigma^2\lambda < 1$. Then
$e^{\lambda X^2}$ is integrable and
\[
  \E e^{\lambda X^2} = (1 - 2\sigma^2\lambda)^{-1/2}\exp\Big(\frac{\mu^2\lambda}{1 - 2\sigma^2\lambda}\Big).
\]
\end{lemma}

\begin{proof}
\leanok
\uses{def:pg_gaussian_vector}
Complete the square: $bx^2 + cx + d = b\,(x + \frac{c}{2b})^2 + d - \frac{c^2}{4b}$, so by the
translation invariance of Lebesgue measure (\texttt{integral\_add\_right\_eq\_self}) and
Mathlib \texttt{integral\_gaussian} ($\int e^{-\beta y^2}dy = \sqrt{\pi/\beta}$ with
$\beta = -b > 0$), the first identity follows; integrability is
\texttt{integrable\_exp\_neg\_mul\_sq} after the same translation.
If $\sigma = 0$ then $X = \mu$ a.s.\ and both sides of the second identity equal
$e^{\lambda\mu^2}$. Let $\sigma > 0$. With the density
$\frac{1}{\sqrt{2\pi\sigma^2}}e^{-(x-\mu)^2/(2\sigma^2)}$ (Mathlib \texttt{gaussianReal},
\texttt{gaussianPDFReal}, \texttt{integral\_gaussianReal\_eq\_integral\_smul}),
\[
  \E e^{\lambda X^2} = \frac{1}{\sqrt{2\pi\sigma^2}}\int_{\R}\exp\Big(\lambda x^2 - \frac{(x-\mu)^2}{2\sigma^2}\Big)dx ,
\]
and $\lambda x^2 - \frac{(x-\mu)^2}{2\sigma^2} = bx^2 + cx + d$ with
$b = \lambda - \frac{1}{2\sigma^2} = -\frac{1-2\sigma^2\lambda}{2\sigma^2} < 0$,
$c = \frac{\mu}{\sigma^2}$, $d = -\frac{\mu^2}{2\sigma^2}$. The first identity gives
$\frac{1}{\sqrt{2\pi\sigma^2}}\sqrt{\frac{2\pi\sigma^2}{1-2\sigma^2\lambda}}
= (1-2\sigma^2\lambda)^{-1/2}$ and the exponent
$d - \frac{c^2}{4b} = -\frac{\mu^2}{2\sigma^2} + \frac{\mu^2}{2\sigma^2(1-2\sigma^2\lambda)}
= \frac{\mu^2\lambda}{1-2\sigma^2\lambda}$.
\end{proof}

\begin{lemma}[Two elementary inequalities]\label{lem:log_ineq_nu}
\lean{ProbabilityTheory.neg_log_one_sub_div_two_sub_le_sq, ProbabilityTheory.one_div_one_sub_le_two}\leanok
For every $\nu \in \R$ with $\abs{\nu} \le 1/2$:
\begin{enumerate}
  \item[(i)] $-\frac12\log(1-\nu) - \frac{\nu}{2} \le \nu^2$;
  \item[(ii)] $\frac{1}{1-\nu} \le 2$.
\end{enumerate}
\end{lemma}

\begin{proof}
\leanok
(i) Mathlib's \texttt{Real.abs\_log\_sub\_add\_sum\_range\_le}, for $\abs x < 1$ and $n = 1$,
reads $\abs{x + \log(1-x)} \le \abs{x}^2/(1 - \abs x)$. With $x = \nu$ and $\abs\nu \le 1/2$,
$1 - \abs\nu \ge 1/2$, so $\abs{\nu + \log(1-\nu)} \le 2\nu^2$, hence
$-\log(1-\nu) - \nu \le 2\nu^2$ and dividing by $2$ gives (i).
(ii) $1 - \nu \ge 1 - \abs\nu \ge \frac12$.
\end{proof}

\begin{lemma}[Centered square of a Gaussian is sub-exponential]\label{lem:gaussian_square_subexp}
\lean{ProbabilityTheory.hasSubexponentialMGF_sq_sub_gaussianReal}\leanok
\uses{def:subexp, def:pg_gaussian_vector}
Let $\mu\in\R$, $\sigma \ge 0$ and $X \sim \Normal(\mu,\sigma^2)$. Then
\[
  \mathrm{HasSubexponentialMGF}\big(X^2 - (\mu^2 + \sigma^2),\ 8(\sigma^4 + \mu^2\sigma^2),\ 4\sigma^2\big),
\]
i.e.\ for all $\abs\lambda \le \frac{1}{4\sigma^2}$,
$\log\E e^{\lambda(X^2 - (\mu^2+\sigma^2))} \le 4(\sigma^4 + \mu^2\sigma^2)\lambda^2$
(and $\E X^2 = \mu^2 + \sigma^2$ by Lemma~\ref{lem:ps_subexp_integral_zero}).
\end{lemma}

\begin{proof}
\leanok
\uses{lem:gaussian_square_mgf, lem:log_ineq_nu}
Let $4\sigma^2\abs\lambda \le 1$ and $\nu := 2\sigma^2\lambda$, so that $\abs\nu \le 1/2$ and in
particular $2\sigma^2\lambda < 1$: Lemma~\ref{lem:gaussian_square_mgf} applies, $e^{\lambda X^2}$
is integrable (hence so is $e^{\lambda(X^2 - \mu^2 - \sigma^2)}$), and
\[
  \log\E e^{\lambda(X^2 - \mu^2 - \sigma^2)}
  = -\frac12\log(1-\nu) + \frac{\mu^2\lambda}{1-\nu} - \lambda(\mu^2 + \sigma^2)
  = \Big(-\frac12\log(1-\nu) - \frac\nu2\Big) + \mu^2\lambda\Big(\frac{1}{1-\nu} - 1\Big),
\]
using $\lambda\sigma^2 = \nu/2$. By Lemma~\ref{lem:log_ineq_nu}(i), the first bracket is at most
$\nu^2 = 4\sigma^4\lambda^2$. The second term equals
$\mu^2\lambda\nu/(1-\nu) = 2\mu^2\sigma^2\lambda^2\cdot\frac{1}{1-\nu} \le 4\mu^2\sigma^2\lambda^2$
by Lemma~\ref{lem:log_ineq_nu}(ii), since $\lambda\nu = 2\sigma^2\lambda^2 \ge 0$. Hence
$\log\E e^{\lambda(X^2-\mu^2-\sigma^2)} \le 4(\sigma^4 + \mu^2\sigma^2)\lambda^2 = \frac{\lambda^2}{2}\cdot8(\sigma^4 + \mu^2\sigma^2)$,
which is the claim with $V = 8(\sigma^4 + \mu^2\sigma^2)$ and $b = 4\sigma^2$.
\end{proof}

\begin{lemma}[Chi-square variables are sub-exponential]\label{lem:chi_square_subexp}
\lean{ProbabilityTheory.hasSubexponentialMGF_sum_sq_sub_one_pi_gaussianReal, ProbabilityTheory.hasSubexponentialMGF_norm_sq_sub_stdGaussian}\leanok
\uses{def:subexp, def:pg_gaussian_vector}
Let $g \sim \Normal(0, I_d)$. Then
$\mathrm{HasSubexponentialMGF}\big(\norm{g}^2 - d,\ 8d,\ 4\big)$.
\end{lemma}

\begin{proof}
\leanok
\uses{lem:gauss_pi, lem:gaussian_square_subexp, lem:subexp_sum, lem:ps_subexp_map}
$\norm{g}^2 - d = \sum_{i<d}(g_i^2 - 1)$ with $g_0,\dots,g_{d-1}$ i.i.d.\ $\Normal(0,1)$
(Lemma~\ref{lem:gauss_pi}; in Lean, $\Normal(0,I_d)$ is the image of the product measure
$\Normal(0,1)^{\otimes d}$ on $\R^d$, \texttt{map\_pi\_eq\_stdGaussian}, and the statement is
transported by Lemma~\ref{lem:ps_subexp_map}). Each $g_i^2 - 1$ has parameters $(8, 4)$ by
Lemma~\ref{lem:gaussian_square_subexp} with $\mu = 0$, $\sigma = 1$; the summands are
independent (\texttt{iIndepFun\_pi}), so Lemma~\ref{lem:subexp_sum} gives $(8d, 4)$.
\end{proof}

\begin{lemma}[Chi-square tail]\label{lem:chi_square_tail}
\lean{ProbabilityTheory.measureReal_norm_sq_sub_ge_le_stdGaussian, ProbabilityTheory.measureReal_norm_sq_ge_le_stdGaussian}\leanok
\uses{def:pg_gaussian_vector}
Let $g \sim \Normal(0,I_d)$ with $d \ge 1$. For every $t \ge 0$,
\[
  \Prob\big(\norm{g}^2 - d \ge t\big) \le \exp\Big(-\min\Big(\frac{t^2}{16d}, \frac{t}{8}\Big)\Big),
\]
and for every $\delta \in (0,1)$,
$\Prob\big(\norm{g}^2 \ge 2d + 12\log(1/\delta)\big) \le \delta$.
\end{lemma}

\begin{proof}
\leanok
\uses{lem:chi_square_subexp, lem:subexp_tail}
The first bound is Lemma~\ref{lem:subexp_tail} with $(V,b) = (8d,4)$ from
Lemma~\ref{lem:chi_square_subexp}: $t^2/(2V) = t^2/(16d)$ and $t/(2b) = t/8$. For the second,
let $L := \log(1/\delta) > 0$ and $t := \max(4\sqrt{dL}, 8L)$. Then $t^2/(16d) \ge 16dL/(16d) = L$
and $t/8 \ge L$, so $\Prob(\norm{g}^2 - d \ge t) \le e^{-L} = \delta$. Moreover, by AM--GM,
$4\sqrt{dL} = 2\sqrt{d\cdot4L} \le d + 4L$, so $t \le \max(d + 4L, 8L) \le d + 8L \le d + 12L$.
Hence $\{\norm{g}^2 \ge 2d + 12L\} = \{\norm{g}^2 - d \ge d + 12L\} \subseteq \{\norm{g}^2 - d \ge t\}$
has probability at most $\delta$.
\end{proof}

\begin{lemma}[Gaussian quadratic forms are sub-exponential]\label{lem:gauss_quadratic_form_subexp}
\uses{def:subexp, def:pg_gaussian_vector}
Let $S \in \PSD$ ($d\times d$, $S \ne 0$) and $g \sim \Normal(0,I_d)$. Then
$\E[g^\top Sg] = \tr S$ and
\[
  \mathrm{HasSubexponentialMGF}\big(g^\top Sg - \tr S,\ 8\tr(S^2),\ 4\norm{S}_{\mathrm{op}}\big),
\]
where $\norm{S}_{\mathrm{op}}$ is the largest eigenvalue of $S$.
\end{lemma}

\begin{proof}
\uses{lem:gauss_rotation, lem:gauss_pi, lem:gaussian_square_subexp, lem:subexp_const_mul, lem:ps_subexp_mono, lem:subexp_sum, lem:gauss_norm_moments}
Diagonalize $S = U\operatorname{diag}(s_1,\dots,s_d)U^\top$ with $U$ orthogonal and
$s_i \ge 0$ (Mathlib \texttt{Matrix.IsHermitian.spectral\_theorem}), so that
$\tr S = \sum_is_i$, $\tr(S^2) = \sum_is_i^2$ and $\norm{S}_{\mathrm{op}} = s_{\max} := \max_is_i > 0$.
Let $h := U^\top g$; then $h \sim \Normal(0,I_d)$ (Lemma~\ref{lem:gauss_rotation}(i)), its
coordinates are i.i.d.\ $\Normal(0,1)$ (Lemma~\ref{lem:gauss_pi}), and
$g^\top Sg - \tr S = \sum_is_i(h_i^2 - 1)$. The expectation is
Lemma~\ref{lem:gauss_norm_moments}(i). For each $i$ with $s_i > 0$,
Lemma~\ref{lem:gaussian_square_subexp} gives $(8,4)$ for $h_i^2 - 1$ and
Lemma~\ref{lem:subexp_const_mul} gives $(8s_i^2, 4s_i)$ for $s_i(h_i^2-1)$, hence
$(8s_i^2, 4s_{\max})$ by Lemma~\ref{lem:ps_subexp_mono}; for $s_i = 0$ the summand is $0$,
which has parameters $(0, 4s_{\max})$. The summands are independent, so
Lemma~\ref{lem:subexp_sum} gives $(8\sum_is_i^2, 4s_{\max}) = (8\tr(S^2), 4\norm{S}_{\mathrm{op}})$.
\end{proof}

\begin{remark}[Not used]\label{rem:ps_quadratic_form_unused}
Lemma~\ref{lem:gauss_quadratic_form_subexp} is recorded for completeness (it is the general
form of Lemma~\ref{lem:chi_square_subexp}, which is the case $S = I_d$) but none of the main
results uses it: the large-norm regime of the norm estimation uses a deterministic basis
design, for which only the chi-square case is needed. It is not formalized.
\end{remark}

\section{Squares of sub-Gaussian variables and of Rademacher linear forms}
\label{sec:pre_subexp_rademacher}

\begin{lemma}[MGF of the square of a sub-Gaussian variable, positive parameter]
\label{lem:subgaussian_square_mgf_pos}
\lean{ProbabilityTheory.HasSubgaussianMGF.lintegral_exp_mul_sq_le, ProbabilityTheory.HasSubgaussianMGF.integrable_exp_mul_sq, ProbabilityTheory.HasSubgaussianMGF.integral_exp_mul_sq_le}\leanok
\uses{def:pg_gaussian_vector}
Let $Z$ be a real random variable with a sub-Gaussian MGF with constant $c \ge 0$
(Mathlib \texttt{HasSubgaussianMGF Z c}), and let $0 \le \lambda$ with $2c\lambda < 1$
(any $\lambda \ge 0$ if $c = 0$). Then $e^{\lambda Z^2}$ is integrable and
\[
  \E e^{\lambda Z^2} \le (1 - 2c\lambda)^{-1/2}.
\]
\end{lemma}

\begin{proof}
\leanok
\uses{lem:gaussian_square_mgf, lem:gauss_1d_moments}
\emph{Gaussian trick:} for every real $z$, by the MGF of $g' \sim \Normal(0,1)$ at
$s = \sqrt{2\lambda}\,z$ (Mathlib \texttt{mgf\_gaussianReal}),
\[
  e^{\lambda z^2} = e^{(\sqrt{2\lambda}\,z)^2/2} = \E_{g'}\exp\big(\sqrt{2\lambda}\,z\,g'\big).
\]
Hence, by Tonelli's theorem for the nonnegative measurable function
$(\omega, v) \mapsto \exp(\sqrt{2\lambda}\,Z(\omega)\,v)$ on the product of $(\Omega,\Prob)$
with $(\R, \Normal(0,1))$ (\texttt{MeasureTheory.lintegral\_lintegral\_swap}),
\[
  \E e^{\lambda Z^2} = \int\Big(\int\exp\big(\sqrt{2\lambda}\,Z(\omega)\,v\big)\,\Normal(0,1)(dv)\Big)\Prob(d\omega)
  = \int\Big(\int\exp\big((\sqrt{2\lambda}\,v)\,Z(\omega)\big)\Prob(d\omega)\Big)\Normal(0,1)(dv).
\]
For fixed $v$, the inner integral is the MGF of $Z$ at $s := \sqrt{2\lambda}\,v$, which is at
most $e^{cs^2/2} = e^{c\lambda v^2}$ by sub-Gaussianity. Therefore
\[
  \E e^{\lambda Z^2} \le \int e^{c\lambda v^2}\,\Normal(0,1)(dv) = (1 - 2c\lambda)^{-1/2}
\]
by Lemma~\ref{lem:gaussian_square_mgf} with $\mu = 0$, $\sigma = 1$ and the parameter
$c\lambda < 1/2$. Since the right-hand side is finite, $e^{\lambda Z^2}$ is integrable: all
integrals above are Lebesgue integrals in $[0,\infty]$ of nonnegative functions, and
\texttt{Integrable} is the conjunction of \texttt{AEStronglyMeasurable} (clear) and
\texttt{HasFiniteIntegral}, which for a nonnegative function $f$ is
$\int^-\mathrm{ENNReal.ofReal}(f)\,d\Prob < \infty$ (Mathlib \texttt{hasFiniteIntegral\_iff\_ofReal}).
\end{proof}

\begin{lemma}[Fourth moment of a sub-Gaussian variable]\label{lem:subgaussian_fourth_moment}
\lean{ProbabilityTheory.HasSubgaussianMGF.integral_pow_four_le}\leanok
Let $Z$ have a sub-Gaussian MGF with constant $c \ge 0$ under a probability measure. Then
$\E Z^4 \le 14c^2$.
\end{lemma}

\begin{proof}
\leanok
\uses{lem:subgaussian_square_mgf_pos}
If $c = 0$ then $Z = 0$ a.s.\ (Mathlib \texttt{ae\_eq\_zero\_of\_hasSubgaussianMGF\_zero}). Let
$c > 0$ and $\lambda := 1/(4c)$, so that $2c\lambda = 1/2$ and
Lemma~\ref{lem:subgaussian_square_mgf_pos} gives $\E e^{\lambda Z^2} \le \sqrt2$. For $x \ge 0$,
$1 + x + x^2/2 \le e^x$ (Mathlib \texttt{Real.quadratic\_le\_exp\_of\_nonneg}), so with
$x = \lambda Z^2 \ge 0$, pointwise $\frac{\lambda^2}{2}Z^4 \le e^{\lambda Z^2} - 1$. Taking
expectations, $\frac{\lambda^2}{2}\E Z^4 \le \sqrt2 - 1$, i.e.\
$\E Z^4 \le 32c^2(\sqrt2 - 1) \le 14c^2$ since $\sqrt2 \le 23/16$.
\end{proof}

\begin{lemma}[A quadratic upper bound for the exponential on the negative half-line]
\label{lem:ps_exp_le_quadratic}
\lean{Real.exp_le_one_add_add_sq_div_two_of_nonpos}\leanok
For every $x \le 0$, $e^x \le 1 + x + \frac{x^2}{2}$.
\end{lemma}

\begin{proof}
\leanok
Let $y := -x \ge 0$. By Mathlib \texttt{Real.quadratic\_le\_exp\_of\_nonneg},
$e^{y} \ge 1 + y + y^2/2 > 0$, so $e^x = 1/e^y \le 1/(1 + y + y^2/2)$. Finally
$1/(1 + y + y^2/2) \le 1 - y + y^2/2$, because
$(1 + y + y^2/2)(1 - y + y^2/2) = (1 + y^2/2)^2 - y^2 = 1 + y^4/4 \ge 1$ and
$1 - y + y^2/2 = ((1-y)^2 + 1)/2 > 0$. Since $1 - y + y^2/2 = 1 + x + x^2/2$, this is the claim.
\end{proof}

\begin{lemma}[Centered square of a sub-Gaussian variable is sub-exponential]
\label{lem:subgaussian_square_subexp}
\lean{ProbabilityTheory.HasSubgaussianMGF.hasSubexponentialMGF_sq_sub}\leanok
\uses{def:subexp}
Let $Z$ have a sub-Gaussian MGF with constant $c \ge 0$ under a probability measure, and
assume $\E Z^2 = c$. Then
\[
  \mathrm{HasSubexponentialMGF}\big(Z^2 - c,\ 16c^2,\ 4c\big).
\]
\end{lemma}

\begin{proof}
\leanok
\uses{lem:subgaussian_square_mgf_pos, lem:subgaussian_fourth_moment, lem:log_ineq_nu, lem:ps_exp_le_quadratic}
$Z$ has all moments (Mathlib \texttt{integrable\_pow\_of\_mem\_interior\_integrableExpSet}).
Let $\abs\lambda \le 1/(4c)$; we show that $e^{\lambda(Z^2 - c)}$ is integrable and
$\E e^{\lambda(Z^2 - c)} \le e^{8c^2\lambda^2} = e^{\lambda^2\cdot16c^2/2}$.

\emph{Case $0 \le \lambda \le 1/(4c)$.} Let $\nu := 2c\lambda \in [0, 1/2]$, so that
$2c\lambda < 1$ and Lemma~\ref{lem:subgaussian_square_mgf_pos} gives the integrability of
$e^{\lambda Z^2}$ (hence of $e^{\lambda(Z^2-c)} = e^{-\lambda c}e^{\lambda Z^2}$) and
$\E e^{\lambda Z^2} \le (1-\nu)^{-1/2}$. Hence, using $\lambda c = \nu/2$ and
Lemma~\ref{lem:log_ineq_nu}(i),
\[
  \log\E e^{\lambda(Z^2 - c)} \le -\frac12\log(1-\nu) - \frac{\nu}{2} \le \nu^2 = 4c^2\lambda^2 \le 8c^2\lambda^2 .
\]

\emph{Case $-1/(4c) \le \lambda < 0$.} Pointwise $\lambda Z^2 \le 0$, so
$e^{\lambda(Z^2-c)} \le e^{-\lambda c}$ is bounded, hence integrable, and
Lemma~\ref{lem:ps_exp_le_quadratic} gives $e^{\lambda Z^2} \le 1 + \lambda Z^2 + \frac{\lambda^2Z^4}{2}$.
Taking expectations and using $\E Z^2 = c$ and Lemma~\ref{lem:subgaussian_fourth_moment}
($\E Z^4 \le 14c^2$),
\[
  \E e^{\lambda Z^2} \le 1 + \lambda c + 7c^2\lambda^2
  \le \exp\big(\lambda c + 7c^2\lambda^2\big),
\]
by $1 + y \le e^y$. Multiplying by $e^{-\lambda c}$:
$\E e^{\lambda(Z^2 - c)} \le e^{7c^2\lambda^2} \le e^{8c^2\lambda^2}$.

In both cases $\E e^{\lambda(Z^2-c)} \le \exp(\lambda^2\cdot16c^2/2)$, which is the claim with
$V = 16c^2$ and $b = 4c$.
\end{proof}

\begin{definition}[Rademacher vector]\label{def:ps_rademacher}
\lean{ProbabilityTheory.signMeasure, ProbabilityTheory.rademacherMeasure, ProbabilityTheory.integral_signMeasure, ProbabilityTheory.iIndepFun_eval_rademacherMeasure, ProbabilityTheory.map_eval_rademacherMeasure, ProbabilityTheory.ae_rademacherMeasure_mem_Icc, ProbabilityTheory.integral_eval_rademacherMeasure, ProbabilityTheory.integral_eval_sq_rademacherMeasure, ProbabilityTheory.inner_toLp_eq_sum}\leanok
A \emph{Rademacher vector} in $\R^d$ is a random vector $\eps = (\eps_1,\dots,\eps_d)$ whose
coordinates are independent with $\Prob(\eps_i = 1) = \Prob(\eps_i = -1) = 1/2$; its law is the
product $\big(\frac12\delta_{1} + \frac12\delta_{-1}\big)^{\otimes d}$
(Lean: \texttt{signMeasure := 2⁻¹ • (dirac 1 + dirac (-1))} and
\texttt{rademacherMeasure ι := Measure.pi (fun \_ ↦ signMeasure)} on \texttt{ι → ℝ}; the
coordinates are independent by \texttt{iIndepFun\_pi}, each with law \texttt{signMeasure},
$\E\eps_i = 0$, $\E\eps_i^2 = 1$ and $\abs{\eps_i} \le 1$ a.s.).
For $\theta \in \R^d$ we consider the linear form $Z := \ip{\eps}{\theta} = \sum_i\eps_i\theta_i$,
which satisfies $\abs Z \le \sum_i\abs{\theta_i}$.
\end{definition}

\begin{lemma}[Rademacher linear forms are sub-Gaussian]\label{lem:rademacher_linear_subgaussian}
\lean{ProbabilityTheory.hasSubgaussianMGF_sum_mul_rademacherMeasure, ProbabilityTheory.hasSubgaussianMGF_sum_mul_rademacherMeasure'}\leanok
\uses{def:ps_rademacher}
Let $\eps$ be a Rademacher vector in $\R^d$ and $\theta \in \R^d$. Then $\ip{\eps}{\theta}$ has a
sub-Gaussian MGF with constant $\norm{\theta}^2$ (Mathlib
\texttt{HasSubgaussianMGF ⟪ε, θ⟫ ‖θ‖²}): $\E e^{\lambda\ip{\eps}{\theta}} \le e^{\norm\theta^2\lambda^2/2}$
for all $\lambda\in\R$.
\end{lemma}

\begin{proof}
\leanok
\uses{def:ps_rademacher}
Each $\theta_i\eps_i$ takes values in $[-\abs{\theta_i}, \abs{\theta_i}]$ and has mean zero, so
by Hoeffding's lemma (Mathlib \texttt{hasSubgaussianMGF\_of\_mem\_Icc\_of\_integral\_eq\_zero})
it has a sub-Gaussian MGF with constant $((2\abs{\theta_i})/2)^2 = \theta_i^2$. The
$\theta_i\eps_i$ are independent, so \texttt{HasSubgaussianMGF.sum\_of\_iIndepFun} gives the
constant $\sum_i\theta_i^2 = \norm\theta^2$ for the sum.
\end{proof}

\begin{lemma}[First and second moments of a Rademacher linear form]\label{lem:rademacher_fourth_moment}
\lean{ProbabilityTheory.integral_eval_mul_eval_rademacherMeasure, ProbabilityTheory.integral_sum_mul_rademacherMeasure, ProbabilityTheory.integral_sq_sum_mul_rademacherMeasure}\leanok
\uses{def:ps_rademacher}
Let $\eps$ be a Rademacher vector in $\R^d$, $\theta\in\R^d$ and $Z := \ip{\eps}{\theta}$. Then
$\E[\eps_i\eps_j] = \indic\{i = j\}$, $\E Z = 0$ and $\E Z^2 = \norm{\theta}^2$.
\end{lemma}

\begin{proof}
\leanok
\uses{def:ps_rademacher}
All the variables involved are bounded, hence integrable. $\E[\eps_i^2] = 1$ and, for
$i \ne j$, $\E[\eps_i\eps_j] = \E\eps_i\,\E\eps_j = 0$ by independence. Then
$\E Z = \sum_i\theta_i\E\eps_i = 0$ and, expanding the square,
$\E Z^2 = \sum_{i,j}\theta_i\theta_j\E[\eps_i\eps_j] = \sum_i\theta_i^2$.
\end{proof}

\begin{lemma}[Centered square of a Rademacher linear form is sub-exponential]
\label{lem:rademacher_square_subexp}
\lean{ProbabilityTheory.hasSubexponentialMGF_sq_sum_mul_sub_rademacherMeasure}\leanok
\uses{def:subexp, def:ps_rademacher}
Let $\eps$ be a Rademacher vector in $\R^d$, $\theta \in \R^d$ with $r := \norm{\theta}$, and
$Z := \ip{\eps}{\theta}$. Then
\[
  \mathrm{HasSubexponentialMGF}\big(Z^2 - r^2,\ 16r^4,\ 4r^2\big).
\]
\end{lemma}

\begin{proof}
\leanok
\uses{lem:rademacher_linear_subgaussian, lem:rademacher_fourth_moment, lem:subgaussian_square_subexp}
By Lemma~\ref{lem:rademacher_linear_subgaussian}, $Z$ is sub-Gaussian with constant
$c = r^2$, and $\E Z^2 = r^2 = c$ by Lemma~\ref{lem:rademacher_fourth_moment}. This is exactly
the hypothesis of Lemma~\ref{lem:subgaussian_square_subexp}, which gives the parameters
$(16c^2, 4c) = (16r^4, 4r^2)$.
\end{proof}

\begin{remark}[Constants, and the case $r = 0$]\label{rem:ps_rademacher_constants}
The proof of Lemma~\ref{lem:subgaussian_square_subexp} shows the stronger parameters
$(8c^2, 4c)$ (the bound $e^{7c^2\lambda^2}$ for negative $\lambda$ could be improved by using
the exact fourth moment $\E Z^4 = 3r^4 - 2\sum_i\theta_i^4$ of a Rademacher form instead of
the general bound $14c^2$, but this is not needed); the value $V = 16r^4$ is kept because it is
the one used by Chapter~\ref{chap:norm} (Lemma~\ref{lem:term1_subexp}). The paper obtains
parameters $(Cr^4, Cr^2)$ with an unspecified $C$ from the Hanson--Wright inequality (its
Lemma \emph{lem:hanson-wright}, Appendix~A) and a tail-to-MGF conversion; the argument above
is self-contained and explicit. If $r = 0$ then $Z = 0$ a.s.\ and $Z^2 - r^2 = 0$, which has
parameters $(0, 0)$, consistent with the statement; the norm-estimation chapter treats this
case separately anyway.
\end{remark}

\textbf{Lean remark.} The Gaussian trick of Lemma~\ref{lem:subgaussian_square_mgf_pos} needs an
auxiliary standard Gaussian independent of $Z$; it is encoded by Tonelli's theorem on the
product measure \texttt{μ.prod (gaussianReal 0 1)}, without introducing a new random variable.
The Rademacher results are stated for the product measure \texttt{rademacherMeasure ι} on
\texttt{ι → ℝ} and the coordinate sum \texttt{∑ i, ε i * θ i}; the identification with the
inner product on \texttt{EuclideanSpace ℝ ι} is \texttt{inner\_toLp\_eq\_sum}.

\chapter{Information-theoretic tools and the divergence decomposition}
\label{chap:pre_info}

This chapter collects the information-theoretic facts used by the lower bounds of
Chapters~\ref{chap:lower_adaptive} and~\ref{chap:log_gains}: the Kullback--Leibler divergence
of two Gaussians, the Pinsker and Bretagnolle--Huber inequalities, convexity of the divergence
in the pair of measures (hence in its second argument), and, most importantly, the \emph{divergence decomposition} for
algorithm-environment sequences (the ``chain rule'' of the paper, Appendix~A): the divergence
between the laws of the histories produced by one algorithm against two stationary
environments is the expected sum of the one-step divergences along the trajectory. The
decomposition is stated for the LML framework, for arbitrary reward kernels, and the linear
Gaussian case of the paper is a corollary.

\textbf{What Mathlib/LML provides.} Mathlib defines the Kullback--Leibler divergence
\texttt{InformationTheory.klDiv μ ν} as an \texttt{ℝ≥0∞}-valued \texttt{irreducible\_def}
(\texttt{Mathlib/InformationTheory/KullbackLeibler/Basic.lean}): it equals
$\int \mathrm{llr}(\mu,\nu)\,d\mu + \nu(\Omega) - \mu(\Omega)$ when $\mu \ll \nu$ and the
log-likelihood ratio $\mathrm{llr}(\mu,\nu) = \log \frac{d\mu}{d\nu}$ is $\mu$-integrable, and
$\infty$ otherwise (\texttt{klDiv\_of\_ac\_of\_integrable}, \texttt{klDiv\_of\_not\_ac},
\texttt{klDiv\_of\_not\_integrable}). It is also the $f$-divergence
$\int \mathrm{klFun}(\frac{d\mu}{d\nu})\,d\nu$ with $\mathrm{klFun}(x) = x\log x + 1 - x$
(\texttt{klDiv\_eq\_lintegral\_klFun\_of\_ac}). The chain rule is available in the form
\texttt{klDiv\_compProd\_eq\_add}:
$\KL(\mu \otimes \kappa \,\|\, \nu \otimes \eta) = \KL(\mu\|\nu) + \KL(\mu\otimes\kappa\,\|\,\mu\otimes\eta)$,
together with \texttt{klDiv\_compProd\_left} ($\KL(\mu\otimes\kappa\,\|\,\nu\otimes\kappa) = \KL(\mu\|\nu)$),
and the data-processing inequalities \texttt{klDiv\_map\_le}, \texttt{klDiv\_trim\_le},
\texttt{klDiv\_comp\_right\_le}. The lower bound \texttt{mul\_log\_le\_klDiv} and the converse
Gibbs inequality \texttt{klDiv\_eq\_zero\_iff} are also there. Radon--Nikodym derivatives of
composition-products are in \texttt{Mathlib/Probability/Kernel/Composition/RadonNikodym.lean}
(\texttt{rnDeriv\_measure\_compProd}) and kernel derivatives \texttt{Kernel.rnDeriv} are
measurable (\texttt{Mathlib/Probability/Kernel/RadonNikodym.lean}, under the hypothesis
\texttt{MeasurableSpace.CountableOrCountablyGenerated}, which is implied by a countably
generated target space); the same file proves that $\{x : \kappa_x \ll \eta_x\}$ is
measurable (\texttt{Kernel.measurableSet\_absolutelyContinuous}, via
\texttt{Kernel.singularPart\_eq\_zero\_iff\_absolutelyContinuous}). Real Gaussians are
\texttt{ProbabilityTheory.gaussianReal μ v} with density \texttt{gaussianPDFReal}
(\texttt{Mathlib/Probability/Distributions/Gaussian/Real.lean}), with
\texttt{gaussianReal\_absolutelyContinuous}, \texttt{rnDeriv\_gaussianReal},
\texttt{integral\_id\_gaussianReal}. LML provides the trajectory measure
\texttt{Learning.trajMeasure alg env}, the step decomposition of histories
(\texttt{history\_succ}, \texttt{hasCondDistrib\_step}, \texttt{hasCondDistrib\_action}),
stationary environments \texttt{stationaryEnv} (feedback kernel
\texttt{feedback\_stationaryEnv}: $(h,a) \mapsto \nu(a)$), and uniqueness of the law of the
history (\texttt{IsAlgEnvSeqUntil.map\_history}).

\textbf{What is missing.} Mathlib has neither the divergence of two Gaussians, nor Pinsker's
inequality, nor the Bretagnolle--Huber inequality, nor convexity of $\KL$ in the pair of
measures, nor the ``integrated'' form of the chain rule
$\KL(\mu\otimes\kappa\,\|\,\mu\otimes\eta) = \int \KL(\kappa_x\|\eta_x)\,\mu(dx)$ (the Mathlib file
explicitly avoids it because $x \mapsto \KL(\kappa_x\|\eta_x)$ need not be measurable in
general; it is measurable when the target space is countably generated, which is all we need).
The integrated chain rule and the invariance of $\KL$ under measurable embeddings now live in LML
(\texttt{LeanMachineLearning/ForMathlib/InformationTheory/KullbackLeibler/ChainRule.lean}); the
policy/reward specialization of Lemma~\ref{lem:pi_kl_step} and the divergence decomposition for
algorithm-environment sequences are new to LML. Mathlib also lacks the behaviour of $\KL$ under
restrictions (Lemma~\ref{lem:kl_restrict}), the invariance under sufficient statistics
(Lemma~\ref{lem:kl_sufficient_statistic}) and the continuity of $\KL$ along a generating
sequence of $\sigma$-algebras (Lemma~\ref{lem:kl_iSup_map}), which the stopping-time and
trajectory versions of the decomposition (Section~\ref{sec:pre_info_stopping}) rely on.

\section{Basic inequalities}
\label{sec:pre_info_basic}

\begin{lemma}[Invariance under measurable embeddings]\label{lem:pi_kl_measurable_equiv}
\lean{InformationTheory.klDiv_map_measurableEmbedding, InformationTheory.klDiv_map_measurableEquiv}\leanok
Let $\mu,\nu$ be finite measures on $\Omega$ and $f : \Omega \to \Omega'$ a measurable embedding
(for example a measurable equivalence). Then $\KL(f_*\mu \,\|\, f_*\nu) = \KL(\mu\|\nu)$.
\end{lemma}

\begin{proof}
\leanok
Two applications of the data-processing inequality (Mathlib \texttt{klDiv\_map\_le}): with
$f$ we get $\le$, and with a measurable left inverse $g$ of $f$ applied to $f_*\mu, f_*\nu$ we get
$\ge$, since $g_*f_*\mu = (g \circ f)_*\mu = \mu$. If $\Omega$ is empty then $\mu = \nu = 0$ and
both sides vanish; otherwise such a $g$ is given by \texttt{MeasurableEmbedding.invFun}.
\end{proof}

\begin{lemma}[Divergence of two Gaussians with the same variance]\label{lem:kl_gaussian}
\lean{ProbabilityTheory.gaussianReal_absolutelyContinuous_gaussianReal, ProbabilityTheory.llr_gaussianReal, ProbabilityTheory.integrable_llr_gaussianReal, ProbabilityTheory.klDiv_gaussianReal, ProbabilityTheory.klDiv_gaussianReal_one}\leanok
\uses{lem:gauss_1d_moments}
Let $m_1, m_2 \in \R$ and $v > 0$. Then $\Normal(m_1, v) \ll \Normal(m_2,v)$, the log-likelihood
ratio is $\mathrm{llr}(y) = \frac{(m_1 - m_2)}{v}\big(y - \frac{m_1+m_2}{2}\big)$ for Lebesgue-a.e.\
$y$, it is integrable, and
\[
  \KL\big(\Normal(m_1,v)\,\big\|\,\Normal(m_2,v)\big) = \frac{(m_1-m_2)^2}{2v} < \infty .
\]
In particular $\KL(\Normal(m_1,1)\|\Normal(m_2,1)) = (m_1-m_2)^2/2$.
\end{lemma}

\begin{proof}
\leanok
\uses{lem:gauss_1d_moments}
Both measures are $\mathrm{Leb}$ with the densities
$p_i(y) = (2\pi v)^{-1/2}\exp(-(y-m_i)^2/(2v))$ (Mathlib \texttt{gaussianReal\_of\_var\_ne\_zero}),
which are positive everywhere, so each is absolutely continuous with respect to the other and
$\frac{d\Normal(m_1,v)}{d\Normal(m_2,v)} = p_1/p_2$ a.e.\ (Mathlib
\texttt{Measure.rnDeriv\_withDensity} and the chain rule \texttt{Measure.rnDeriv\_mul\_rnDeriv}
for $\mathrm{Leb} \to \Normal(m_2,v) \to \Normal(m_1,v)$, or directly the uniqueness of the
Radon--Nikodym derivative, \texttt{Measure.eq\_rnDeriv}). Taking logarithms,
\[
  \log\frac{p_1(y)}{p_2(y)} = \frac{(y-m_2)^2 - (y-m_1)^2}{2v}
  = \frac{(m_1-m_2)\,(2y - m_1 - m_2)}{2v} ,
\]
an affine function of $y$, hence integrable under $\Normal(m_1,v)$ (Gaussians have a first
moment, Lemma~\ref{lem:gauss_1d_moments}; Mathlib \texttt{integral\_id\_gaussianReal}). Its
integral under $\Normal(m_1,v)$ is obtained by replacing $y$ by its mean $m_1$:
$\frac{(m_1-m_2)(2m_1 - m_1 - m_2)}{2v} = \frac{(m_1-m_2)^2}{2v}$. Conclude with Mathlib
\texttt{klDiv\_of\_ac\_of\_integrable} (both measures are probability measures so the correction
term $\nu(\Omega)-\mu(\Omega)$ vanishes).
\end{proof}

\begin{lemma}[Divergence of two Bernoulli measures]\label{lem:pi_kl_bernoulli}
\lean{InformationTheory.klBin, InformationTheory.klDiv_bernoulliMeasure, InformationTheory.ofReal_le_klDiv_bernoulliMeasure, ProbabilityTheory.lintegral_bernoulliMeasure, ProbabilityTheory.bernoulliMeasure_eq_withDensity, ProbabilityTheory.rnDeriv_bernoulliMeasure}\leanok
For $p \in [0,1]$ write $\mathrm{B}(p) := p\,\delta_1 + (1-p)\,\delta_0$ for the Bernoulli
measure on $\{0,1\}$. Let $p \in [0,1]$ and $q \in (0,1)$. Then
\[
  \KL(\mathrm{B}(p)\,\|\,\mathrm{B}(q)) = p\log\frac{p}{q} + (1-p)\log\frac{1-p}{1-q}
  =: \mathrm{kl}(p,q),
\]
with the convention $0 \log 0 = 0$. If $q \in \{0,1\}$ and $p \ne q$ then
$\KL(\mathrm{B}(p)\|\mathrm{B}(q)) = \infty$.
\end{lemma}

\begin{proof}
\leanok
For $q \in (0,1)$, $\mathrm{B}(q)$ charges both points so $\mathrm{B}(p) \ll \mathrm{B}(q)$, with
Radon--Nikodym derivative $p/q$ at $1$ and $(1-p)/(1-q)$ at $0$; every function on a finite set is
integrable. By Mathlib \texttt{klDiv\_eq\_lintegral\_klFun\_of\_ac}, the divergence is
$q\,\mathrm{klFun}(p/q) + (1-q)\,\mathrm{klFun}((1-p)/(1-q))$ with
$\mathrm{klFun}(x) = x\log x + 1 - x$, which expands to
$p\log(p/q) + (1-p)\log((1-p)/(1-q)) + (q - p) + ((1-q) - (1-p)) = \mathrm{kl}(p,q)$; the
convention $0\log 0 = 0$ is exactly $\mathrm{klFun}(0) = 1$ (Mathlib \texttt{klFun\_zero}). If
$q = 0$ and $p > 0$ then $\mathrm{B}(p)(\{1\}) > 0 = \mathrm{B}(0)(\{1\})$ so $\mathrm{B}(p)
\not\ll \mathrm{B}(q)$ and the divergence is $\infty$ (\texttt{klDiv\_of\_not\_ac}); the case
$q = 1$ is symmetric.
\end{proof}

\begin{lemma}[Pinsker's inequality for Bernoulli measures]\label{lem:kl_bernoulli_pinsker}
\lean{InformationTheory.klBin_pinsker}\leanok
\uses{lem:pi_kl_bernoulli}
For $p \in [0,1]$ and $q \in (0,1)$, $\mathrm{kl}(p,q) \ge 2(p-q)^2$.
\end{lemma}

\begin{proof}
\leanok
\uses{lem:pi_kl_bernoulli}
Fix $p \in (0,1)$ first and let $g(q) := \mathrm{kl}(p,q) - 2(p-q)^2$ on $(0,1)$. Then
$g(p) = 0$ and
\[
  g'(q) = -\frac{p}{q} + \frac{1-p}{1-q} + 4(p-q) = \frac{q-p}{q(1-q)} - 4(q-p)
  = (q-p)\Big(\frac{1}{q(1-q)} - 4\Big).
\]
Since $q(1-q) \le 1/4$, the last factor is nonnegative, so $g' \le 0$ on $(0,p]$ and $g' \ge 0$
on $[p,1)$: $g$ is nonincreasing then nondecreasing (Mathlib
\texttt{antitoneOn\_of\_deriv\_nonpos}, \texttt{monotoneOn\_of\_deriv\_nonneg}) and its minimum
is $g(p) = 0$. For $p = 1$ the claim is $-\log q \ge 2(1-q)^2$: the function
$h(q) := -\log q - 2(1-q)^2$ satisfies $h(1) = 0$ and $h'(q) = -1/q + 4(1-q) = -(2q-1)^2/q \le 0$,
so $h \ge 0$ on $(0,1]$. The case $p = 0$ is the case $p = 1$ with $q$ replaced by $1-q$.
\end{proof}

\begin{lemma}[Pinsker's inequality]\label{lem:pinsker}
\lean{InformationTheory.pinsker_measureReal, InformationTheory.abs_sub_le_sqrt_klDiv}\leanok
\uses{lem:pi_kl_bernoulli, lem:kl_bernoulli_pinsker}
Let $\mu, \nu$ be probability measures on a measurable space and $A$ a measurable set. Then
\[
  2\,\big(\mu(A) - \nu(A)\big)^2 \le \KL(\mu\|\nu)
\]
(as elements of $[0,\infty]$). In particular, if $\KL(\mu\|\nu) < \infty$ then
$\abs{\mu(A) - \nu(A)} \le \sqrt{\KL(\mu\|\nu)/2}$, and if $\KL(\mu\|\nu) = \infty$ the
statement is empty.
\end{lemma}

\begin{proof}
\leanok
\uses{lem:pi_kl_bernoulli, lem:kl_bernoulli_pinsker}
Let $g := \indic_A : \Omega \to \{0,1\}$, measurable. Then $g_*\mu = \mathrm{B}(p)$ and
$g_*\nu = \mathrm{B}(q)$ with $p := \mu(A)$, $q := \nu(A)$, and the data-processing inequality
(Mathlib \texttt{klDiv\_map\_le}) gives $\KL(\mathrm{B}(p)\|\mathrm{B}(q)) \le \KL(\mu\|\nu)$.
If $p = q$ there is nothing to prove. If $q \in \{0,1\}$ and $p \ne q$, then
$\KL(\mathrm{B}(p)\|\mathrm{B}(q)) = \infty$ by Lemma~\ref{lem:pi_kl_bernoulli}, so
$\KL(\mu\|\nu) = \infty$ and the inequality holds trivially. Otherwise $q \in (0,1)$ and
Lemmas~\ref{lem:pi_kl_bernoulli} and~\ref{lem:kl_bernoulli_pinsker} give
$2(p-q)^2 \le \mathrm{kl}(p,q) = \KL(\mathrm{B}(p)\|\mathrm{B}(q)) \le \KL(\mu\|\nu)$.
\end{proof}

\begin{lemma}[Bretagnolle--Huber inequality]\label{lem:bretagnolle_huber}
\lean{InformationTheory.bretagnolle_huber}\leanok
Let $\mu, \nu$ be probability measures on a measurable space and $A$ a measurable set. Then
\[
  \mu(A) + \nu(A^c) \ge \tfrac12 \exp\big(-\KL(\mu\|\nu)\big),
\]
where the right-hand side is $0$ if $\KL(\mu\|\nu) = \infty$.
\end{lemma}

\begin{proof}
\leanok
If $\KL(\mu\|\nu) = \infty$ there is nothing to prove; assume $\mu \ll \nu$ and that
$\mathrm{llr}(\mu,\nu)$ is $\mu$-integrable, so that $\KL(\mu\|\nu) = \int \mathrm{llr}\,d\mu$
(Mathlib \texttt{klDiv\_of\_ac\_of\_integrable}). Let $\lambda := \mu + \nu$ and
$p := d\mu/d\lambda$, $q := d\nu/d\lambda$, which exist (Mathlib
\texttt{Measure.rnDeriv}, \texttt{withDensity\_rnDeriv\_eq} since $\mu,\nu \ll \lambda$) and
satisfy $p + q = 1$ $\lambda$-a.e.\ (\texttt{Measure.rnDeriv\_add}, \texttt{rnDeriv\_self}), in
particular $p, q \le 1$. The proof has three steps.
\begin{enumerate}
  \item[(i)] $\mu(A) + \nu(A^c) = \int_A p\,d\lambda + \int_{A^c} q\,d\lambda \ge \int \min(p,q)\,d\lambda$.
  \item[(ii)] By the Cauchy--Schwarz inequality in $L^2(\lambda)$ applied to
        $\sqrt{\min(p,q)}$ and $\sqrt{\max(p,q)}$, whose product is $\sqrt{pq}$,
        \[
          \Big(\int\sqrt{pq}\,d\lambda\Big)^2 \le \int\min(p,q)\,d\lambda\cdot\int\max(p,q)\,d\lambda
          \le 2\int \min(p,q)\,d\lambda ,
        \]
        since $\max(p,q) \le p + q$ and $\int (p+q)\,d\lambda = \mu(\Omega) + \nu(\Omega) = 2$.
  \item[(iii)] Since $\mu \ll \nu$, we have $p = 0$ $\lambda$-a.e.\ on $\{q = 0\}$, and
        $\frac{d\mu}{d\nu} = p/q$ $\mu$-a.e.\ (Mathlib \texttt{Measure.rnDeriv\_mul\_rnDeriv}:
        $\frac{d\mu}{d\nu}\cdot\frac{d\nu}{d\lambda} = \frac{d\mu}{d\lambda}$ $\lambda$-a.e.). Hence
        \[
          \int \sqrt{pq}\,d\lambda = \int_{\{p>0\}} \sqrt{q/p}\;p\,d\lambda
          = \int \sqrt{q/p}\,d\mu = \int \exp\big(-\tfrac12\,\mathrm{llr}(\mu,\nu)\big)\,d\mu
          \ge \exp\Big(-\tfrac12\int \mathrm{llr}(\mu,\nu)\,d\mu\Big)
          = e^{-\KL(\mu\|\nu)/2},
        \]
        by Jensen's inequality for the convex function $\exp$ and the probability measure
        $\mu$ (Mathlib \texttt{ConvexOn.map\_integral\_le}, with $\mathrm{llr}$ integrable).
\end{enumerate}
Combining, $\mu(A) + \nu(A^c) \ge \int\min(p,q)\,d\lambda \ge \tfrac12 (\int\sqrt{pq}\,d\lambda)^2
\ge \tfrac12 e^{-\KL(\mu\|\nu)}$. This is Theorem~14.2 of~\cite{lattimore2020bandit}; see also
\cite{tsybakov2009introduction}.
\end{proof}

\begin{lemma}[Convexity of the divergence in the pair of measures]\label{lem:kl_joint_convex}
\lean{InformationTheory.klDiv_finsetSum_smul_le, InformationTheory.klDiv_sum_smul_le, InformationTheory.klDiv_smul_add_smul_le}\leanok
Let $\mu_1,\dots,\mu_J$ and $\nu_1,\dots,\nu_J$ be finite measures on the same space and let
$c_1,\dots,c_J \ge 0$. Then
\[
  \KL\Big(\sum_{j=1}^J c_j\mu_j \,\Big\|\, \sum_{j=1}^J c_j\nu_j\Big)
  \le \sum_{j=1}^J c_j\,\KL(\mu_j\|\nu_j) \quad\text{in } [0,\infty].
\]
With $\sum_j c_j = 1$ this is the convexity of $\KL$ in the pair of measures; that hypothesis is
not needed, since $\KL$ is positively homogeneous.
\end{lemma}

\begin{proof}
\leanok
\uses{lem:kl_compProd_kernel}
Let $\beta := \sum_j c_j\delta_j$ be the measure with weights $c$ on the index set
$\{1,\dots,J\}$, equipped with the discrete $\sigma$-algebra, and let $\kappa,\eta$ be the kernels
from the index set with $\kappa_j = \mu_j$ and $\eta_j = \nu_j$ (any function on a countable space
with measurable singletons is a kernel, Mathlib \texttt{Kernel.ofFunOfCountable}; they are finite
kernels because $\kappa_j(\Omega) \le \sum_i \mu_i(\Omega) < \infty$). Then
$\kappa\circ\beta = \sum_j c_j\mu_j$ and $\eta\circ\beta = \sum_j c_j\nu_j$, and $\kappa\circ\beta$
is the image of $\beta\otimes\kappa$ under the second projection (Mathlib
\texttt{Measure.snd\_compProd}). The data-processing inequality (Mathlib \texttt{klDiv\_map\_le})
therefore gives
\[
  \KL(\kappa\circ\beta \,\|\, \eta\circ\beta) \le \KL(\beta\otimes\kappa \,\|\, \beta\otimes\eta) ,
\]
and, the index set being countable, Lemma~\ref{lem:kl_compProd_kernel} (stated in the next
section) identifies the right-hand side with
\[
  \int^{-} \KL(\kappa_j\|\eta_j)\,\beta(dj) = \sum_{j=1}^J c_j\,\KL(\mu_j\|\nu_j) .
\]
\end{proof}

\begin{lemma}[Convexity of the divergence in its second argument]\label{lem:kl_mixture_convex}
\lean{MeasureTheory.isProbabilityMeasure_sum_smul, InformationTheory.klDiv_finsetSum_smul_right_le, InformationTheory.klDiv_sum_smul_right_le}\leanok
Let $\mu$ be a finite measure, $\nu_1,\dots,\nu_J$ finite measures on the same space,
$c_1,\dots,c_J \ge 0$ with $\sum_j c_j = 1$, and $\nu := \sum_j c_j\nu_j$. Then
\[
  \KL(\mu\|\nu) \le \sum_{j=1}^J c_j\,\KL(\mu\|\nu_j) \quad\text{in } [0,\infty].
\]
If moreover the $\nu_j$ are probability measures, then so is $\nu$.
\end{lemma}

\begin{proof}
\leanok
\uses{lem:kl_joint_convex}
Apply Lemma~\ref{lem:kl_joint_convex} with $\mu_j := \mu$ for every $j$: the mixture on the left
is $\sum_j c_j\mu = \big(\sum_j c_j\big)\mu = \mu$. The mass of $\nu$ is $\sum_j c_j = 1$.
\end{proof}

\textbf{Lean remark.} In Mathlib, $\KL$ takes values in \texttt{ℝ≥0∞}; the statements above
are therefore stated in $[0,\infty]$ with \texttt{ENNReal.ofReal} on the real side (e.g.\
\texttt{ENNReal.ofReal (2 * (μ.real A - ν.real A)\^{}2) ≤ klDiv μ ν}), and their real versions are
obtained through \texttt{klDiv\_ne\_top\_iff} and \texttt{ENNReal.toReal}. The mixture in
Lemma~\ref{lem:kl_mixture_convex} is the measure \texttt{∑ j, (c j : ℝ≥0) • ν j}.

\section{The integrated chain rule}
\label{sec:pre_info_chain}

\begin{lemma}[Chain rule with an integrated conditional divergence]\label{lem:kl_compProd_kernel}
\lean{InformationTheory.measurable_klDiv_kernel, InformationTheory.klDiv_compProd_right_eq_lintegral, InformationTheory.klDiv_compProd_eq_add_lintegral}\leanok
Let $\mu$ be a finite measure on $\mathcal{H}$ and $\kappa,\eta$ be Markov kernels from
$\mathcal{H}$ to a countably generated measurable space $\mathcal{Y}$. Then
$x \mapsto \KL(\kappa_x\|\eta_x)$ is measurable and
\[
  \KL(\mu\otimes\kappa\,\|\,\mu\otimes\eta) = \int^{-} \KL(\kappa_x\|\eta_x)\,\mu(dx) ,
\]
where $\int^-$ denotes the Lebesgue integral of a $[0,\infty]$-valued function. Consequently,
for any finite measure $\nu$ on $\mathcal{H}$,
$\KL(\mu\otimes\kappa\,\|\,\nu\otimes\eta) = \KL(\mu\|\nu) + \int^-\KL(\kappa_x\|\eta_x)\,\mu(dx)$.
\end{lemma}

\begin{proof}
\leanok
Mathlib \texttt{Measure.absolutelyContinuous\_compProd\_right\_iff} gives
$\mu\otimes\kappa \ll \mu\otimes\eta$ iff $\kappa_x \ll \eta_x$ for $\mu$-a.e.\ $x$.

\emph{Measurability.} Mathlib's kernel Radon--Nikodym theory
(\texttt{Mathlib/Probability/Kernel/RadonNikodym.lean}) is developed under the hypothesis
\texttt{MeasurableSpace.CountableOrCountablyGenerated 𝓗 𝓨}, which is implied by
\texttt{CountablyGenerated 𝓨} (an instance). Under it, the kernel Radon--Nikodym
derivative $\frac{d\kappa}{d\eta} : \mathcal{H}\times\mathcal{Y}\to[0,\infty]$ is jointly
measurable (\texttt{Kernel.rnDeriv}, \texttt{Kernel.measurable\_rnDeriv}) and
$\frac{d\kappa}{d\eta}(x,\cdot) = \frac{d\kappa_x}{d\eta_x}$ $\eta_x$-a.e.\ for every $x$
(\texttt{Kernel.rnDeriv\_eq\_rnDeriv\_measure}). The set $\{x : \kappa_x \ll \eta_x\}$ is
measurable: this is exactly Mathlib \texttt{Kernel.measurableSet\_absolutelyContinuous} (for
finite kernels), whose proof rewrites the condition as
$\kappa.\mathrm{singularPart}(\eta)(x) = 0$ by
\texttt{Kernel.singularPart\_eq\_zero\_iff\_absolutelyContinuous} and uses the measurability
of $x \mapsto \kappa.\mathrm{singularPart}(\eta)(x)(\mathcal{Y})$. On this set,
$\KL(\kappa_x\|\eta_x) = \int^-\mathrm{klFun}\big(\tfrac{d\kappa}{d\eta}(x,y)\big)\,\eta_x(dy)$
(\texttt{klDiv\_eq\_lintegral\_klFun\_of\_ac}), which is measurable in $x$ by
\texttt{Measurable.lintegral\_kernel\_prod\_right}; off this set it is the constant $\infty$.

\emph{Equality.} If $\mu\otimes\kappa \not\ll \mu\otimes\eta$, then $\kappa_x\not\ll\eta_x$ on a
set of positive $\mu$-measure, where $\KL(\kappa_x\|\eta_x) = \infty$: both sides are $\infty$.
Otherwise, by \texttt{klDiv\_eq\_lintegral\_klFun\_of\_ac} and the chain rule for derivatives
of composition-products (\texttt{rnDeriv\_measure\_compProd}, again under
\texttt{CountableOrCountablyGenerated}:
$\frac{d(\mu\otimes\kappa)}{d(\mu\otimes\eta)}(x,y) = \frac{d\kappa}{d\eta}(x,y)$ for
$(\mu\otimes\eta)$-a.e.\ $(x,y)$),
\[
  \KL(\mu\otimes\kappa\,\|\,\mu\otimes\eta)
  = \int^- \mathrm{klFun}\Big(\frac{d\kappa}{d\eta}(x,y)\Big)\,(\mu\otimes\eta)(d(x,y))
  = \int^-\!\!\int^- \mathrm{klFun}\Big(\frac{d\kappa}{d\eta}(x,y)\Big)\,\eta_x(dy)\,\mu(dx)
\]
by Tonelli for composition-products (\texttt{Measure.lintegral\_compProd}), and the inner
integral is $\KL(\kappa_x\|\eta_x)$ for $\mu$-a.e.\ $x$ (those with $\kappa_x\ll\eta_x$). The
last claim is Mathlib \texttt{klDiv\_compProd\_eq\_add} followed by the equality just proved.
\end{proof}

\begin{lemma}[Sufficient statistics and the divergence]\label{lem:kl_sufficient_statistic}
\lean{InformationTheory.klDiv_map_of_leftInverse, InformationTheory.klDiv_map_of_eq_withDensity_comp, InformationTheory.klDiv_compProd_comap, MeasureTheory.Measure.map_compProd_comap, MeasureTheory.Measure.map_compProd_of_forall_map_eq, ProbabilityTheory.HasCondDistrib.map_of_forall_map_eq, ProbabilityTheory.HasCondDistrib.prodMk_const_left, ProbabilityTheory.HasCondDistrib.restrict_preimage, MeasureTheory.Measure.restrict_compProd_prod_univ}\leanok
\uses{lem:pi_kl_measurable_equiv}
Let $\mu,\nu$ be finite measures on $\Omega$.
\begin{enumerate}
\item[(i)] If $g : \Omega \to \Omega'$ is measurable with a measurable left inverse, then
  $\KL(g_*\mu\|g_*\nu) = \KL(\mu\|\nu)$.
\item[(ii)] If $\mu = (f\circ g)\cdot\nu$ for a measurable $g : \Omega\to\Omega'$ and a measurable
  $f : \Omega' \to [0,\infty]$ (the density factors through $g$, i.e.\ $g$ is a sufficient
  statistic), then $\KL(g_*\mu\|g_*\nu) = \KL(\mu\|\nu)$.
\item[(iii)] Let $\kappa,\eta$ be finite kernels from $\mathcal{B}$ to $\mathcal{C}$ and
  $f : \Omega\to\mathcal{B}$ measurable. Then
  $\KL(\mu\otimes(\kappa\circ f)\,\|\,\mu\otimes(\eta\circ f)) = \KL(f_*\mu\otimes\kappa\,\|\,f_*\mu\otimes\eta)$:
  the conditional divergence of two kernels which depend on the conditioning variable only through
  $f$ is the conditional divergence given $f$.
\end{enumerate}
No countable generation hypothesis is needed.
\end{lemma}

\begin{proof}
\leanok
\uses{lem:pi_kl_measurable_equiv}
(i) is two applications of the data-processing inequality (as for measurable embeddings).
(ii): $g_*\mu = f\cdot g_*\nu$ (change of variables), so both pairs are absolutely continuous
with densities $f\circ g$ and $f$, and the $\mathrm{klFun}$ integrals agree by the change of
variables $g$. (iii): if $f_*\mu\otimes\kappa\not\ll f_*\mu\otimes\eta$ then, transporting along
$(\omega,c)\mapsto(f(\omega),c)$, neither is the left-hand pair, and both divergences are
infinite. Otherwise let $D$ be the density of $f_*\mu\otimes\kappa$ with respect to
$f_*\mu\otimes\eta$; the sections of $D$ integrate to $\kappa(b,t)$ for $f_*\mu$-almost every $b$,
for each fixed measurable $t$, so $\mu\otimes(\kappa\circ f)$ and $(D\circ(f\times\mathrm{id}))\cdot(\mu\otimes(\eta\circ f))$
agree on rectangles, hence everywhere, and (ii) applies with $g = f\times\mathrm{id}$.
\end{proof}

\begin{lemma}[One step of an algorithm-environment pair]\label{lem:pi_kl_step}
\lean{ProbabilityTheory.Kernel.compProd_prodMkLeft_apply, InformationTheory.klDiv_compProd_compProd_prodMkLeft_eq_klDiv_comp_compProd, InformationTheory.klDiv_compProd_compProd_prodMkLeft}\leanok
\uses{lem:kl_compProd_kernel, lem:kl_sufficient_statistic}
Let $P$ be a probability measure on a measurable space $\mathcal{H}$ (``histories''), $\pi$ a
Markov kernel from $\mathcal{H}$ to a countably generated space $\mathcal{A}$ (``policy''),
and $\kappa,\kappa'$ Markov kernels from $\mathcal{A}$ to a countably generated space
$\mathcal{Y}$ (``reward kernels'').
Write $\tilde\kappa$ for the kernel $(h,a)\mapsto\kappa(a)$ from $\mathcal{H}\times\mathcal{A}$
to $\mathcal{Y}$ (LML/Mathlib \texttt{Kernel.prodMkLeft}), and similarly $\tilde\kappa'$. Then
\[
  \KL\big(P\otimes(\pi\otimes\tilde\kappa)\,\big\|\,P\otimes(\pi\otimes\tilde\kappa')\big)
  = \int^- \KL(\kappa(a)\|\kappa'(a))\,(\pi\circ P)(da)
  = \int^-\!\!\int^- \KL(\kappa(a)\|\kappa'(a))\,\pi(h,da)\,P(dh),
\]
where $\pi\circ P := \int \pi(h,\cdot)\,P(dh)$ is the law of the next action. In
composition-product form, without any countable generation hypothesis,
$\KL\big(P\otimes(\pi\otimes\tilde\kappa)\,\big\|\,P\otimes(\pi\otimes\tilde\kappa')\big)
  = \KL\big((\pi\circ P)\otimes\kappa\,\big\|\,(\pi\circ P)\otimes\kappa'\big)$.
\end{lemma}

\begin{proof}
\leanok
\uses{lem:kl_compProd_kernel}
Apply Lemma~\ref{lem:kl_compProd_kernel} to $P$ and the kernels $\pi\otimes\tilde\kappa$,
$\pi\otimes\tilde\kappa'$, Markov kernels from $\mathcal{H}$ to the countably generated space
$\mathcal{A}\times\mathcal{Y}$ (Mathlib \texttt{MeasurableSpace.CountablyGenerated} is stable
under products). For each $h$, the measure
$(\pi\otimes\tilde\kappa)(h)$ is the composition-product $\pi(h)\otimes\kappa$ of the measure
$\pi(h)$ on $\mathcal{A}$ with the kernel $\kappa$ (Mathlib
\texttt{Kernel.compProd\_apply} and \texttt{Kernel.prodMkLeft\_apply}). By
\texttt{klDiv\_compProd\_eq\_add}, \texttt{klDiv\_self} and Lemma~\ref{lem:kl_compProd_kernel}
again,
\[
  \KL\big(\pi(h)\otimes\kappa\,\big\|\,\pi(h)\otimes\kappa'\big)
  = \KL(\pi(h)\|\pi(h)) + \int^-\KL(\kappa(a)\|\kappa'(a))\,\pi(h,da)
  = \int^-\KL(\kappa(a)\|\kappa'(a))\,\pi(h,da).
\]
Integrating in $h$ and using $\int^-\!\int^- f(a)\,\pi(h,da)\,P(dh) = \int^- f\,d(\pi\circ P)$
(Mathlib \texttt{Measure.lintegral\_bind} / \texttt{Kernel.lintegral\_comp}) gives the claim.
\end{proof}

\section{The divergence decomposition}
\label{sec:pre_info_decomposition}

\begin{definition}[A pair of stationary environments and the laws of histories]\label{def:env_pair}
Let $\mathcal{A}$ (actions) and $\mathcal{Y}$ (observations) be countably generated measurable
spaces (for instance Borel subsets of $\R^d$, $\R$, or finite types; in particular
$\mathcal{A}$ may be any action set $\Xs$ in the sense of Definition~\ref{def:arm_set}, i.e.\
a nonempty compact subset of $\R^d$, spanning or not). Let $\kappa, \kappa'$ be Markov kernels from $\mathcal{A}$ to
$\mathcal{Y}$ and let $\nu := \mathrm{stationaryEnv}(\kappa)$, $\nu' := \mathrm{stationaryEnv}(\kappa')$
be the corresponding stationary environments (LML \texttt{Learning.stationaryEnv}: the
observation at every round has law $\kappa(x_t)$ given the past and the action $x_t$). Let
$\mathrm{alg}$ be an LML algorithm with actions in $\mathcal{A}$ and observations in $\mathcal{Y}$,
with policy kernels $\pi_n$ (from histories $(x_t,y_t)_{t< n}$ of the first $n$ rounds to
$\mathcal{A}$; the initial law is $p_0 = \pi_0(\emptyset)$). For $n \ge 0$ let
\[
  \Prob^{(n)} := \mathrm{law}\big((x_t,y_t)_{t < n}\big)
  = (\mathrm{trajMeasure}(\mathrm{alg},\nu))_*\mathrm{hist}_n,
  \qquad
  \Prob'^{(n)} := (\mathrm{trajMeasure}(\mathrm{alg},\nu'))_*\mathrm{hist}_n ,
\]
the laws of the history of the first $n$ rounds (a history of \emph{length} $n$) under the two
environments. We write $\E^{(n)}$ for the expectation under $\Prob^{(n)}$ and also
$\Prob^T := \Prob^{(T)}$ when we want to emphasize the length. By LML
\texttt{IsAlgEnvSeqUntil.map\_history}, $\Prob^{(n)}$ is the law of
$(x_t,y_t)_{t< n}$ for \emph{any} algorithm-environment sequence of $(\mathrm{alg},\nu)$ on any
probability space.
\end{definition}

\begin{theorem}[Divergence decomposition]\label{thm:divergence_decomposition}
\lean{Learning.IsAlgEnvSeq.map_history_succ, Learning.IsAlgEnvSeq.klDiv_map_history, Learning.IsAlgEnvSeq.klDiv_map_history_compProd, Learning.IsAlgEnvSeq.klDiv_map_history_stepKernel}\leanok
\uses{def:env_pair, lem:stopped_history_chain_rule}
In the setting of Definition~\ref{def:env_pair}, for every $T \ge 0$,
\[
  \KL\big(\Prob^{(T)}\,\big\|\,\Prob'^{(T)}\big)
  = \sum_{t<T} \int^- \KL\big(\kappa(x_t)\,\big\|\,\kappa'(x_t)\big)\,d\Prob^{(T)}
  \qquad\text{in } [0,\infty],
\]
where $x_t$ denotes the $t$-th action of the history, that is
$\KL(\Prob^T\|\Prob'^T) = \sum_{t<T}\E^{(T)}[\KL(\kappa(x_t)\|\kappa'(x_t))]$.
If moreover $\KL(\kappa(a)\|\kappa'(a)) \le C < \infty$ for all $a \in \mathcal{A}$, then
$\KL(\Prob^T\|\Prob'^T) \le TC < \infty$ and the identity holds in $\R$ with ordinary
expectations.
\end{theorem}

\begin{proof}
\leanok
\uses{lem:pi_kl_measurable_equiv, lem:kl_compProd_kernel, lem:pi_kl_step, lem:stopped_history_chain_rule}
In the formalization the theorem is the special case $S = \emptyset$ (no stopping) of the chain
rule for stopped histories, Lemma~\ref{lem:stopped_history_chain_rule}: the stopping time is
$+\infty$, the history stopped at $\min(\tau, T)$ is the history of the first $T$ rounds, and the
one-step terms are those of Lemma~\ref{lem:pi_kl_step}. We nevertheless record the direct
induction, which is the argument of the paper.

Write $\mathcal{H}_n := (\{0,\dots,n-1\} \to \mathcal{A}\times\mathcal{Y})$ for the space of
histories of the first $n$ rounds. By definition of the stationary environment, the feedback
kernel at every round is $\tilde\kappa : (h,a)\mapsto\kappa(a)$ (LML
\texttt{feedback\_stationaryEnv}). We prove the identity by induction on $T$.

\emph{Base case $T = 0$.} $\mathcal{H}_0$ is a one-point space, so $\Prob^{(0)} = \Prob'^{(0)}$
is the Dirac mass at the empty history (LML \texttt{hasLaw\_history\_zero}) and both sides
vanish (\texttt{klDiv\_self} and the empty sum).

\emph{Induction step.} Assume the identity for $T$. By LML \texttt{history\_succ} and
\texttt{hasCondDistrib\_step}, $\mathrm{hist}_{T+1} = e^{-1}(\mathrm{hist}_T, (x_T,y_T))$
for the measurable equivalence $e : \mathcal{H}_{T+1}\simeq\mathcal{H}_T\times(\mathcal{A}\times\mathcal{Y})$
(\texttt{MeasurableEquiv.finSuccProd}), and the law of $(\mathrm{hist}_T,(x_T,y_T))$ is
the composition-product $\Prob^{(T)}\otimes(\pi_T\otimes\tilde\kappa)$ (the step kernel
\texttt{stepKernel alg env n = alg.policy n ⊗ₖ env.feedback n}); similarly with primes, with
the \emph{same} policy kernel $\pi_T$. Hence, by Lemma~\ref{lem:pi_kl_measurable_equiv},
\texttt{klDiv\_compProd\_eq\_add} and Lemma~\ref{lem:pi_kl_step},
\begin{align*}
  \KL(\Prob^{(T+1)}\|\Prob'^{(T+1)})
  &= \KL\big(\Prob^{(T)}\otimes(\pi_T\otimes\tilde\kappa)\,\big\|\,\Prob'^{(T)}\otimes(\pi_T\otimes\tilde\kappa')\big)\\
  &= \KL(\Prob^{(T)}\|\Prob'^{(T)})
     + \KL\big(\Prob^{(T)}\otimes(\pi_T\otimes\tilde\kappa)\,\big\|\,\Prob^{(T)}\otimes(\pi_T\otimes\tilde\kappa')\big)\\
  &= \sum_{t<T}\int^-\KL(\kappa(x_t)\|\kappa'(x_t))\,d\Prob^{(T)}
     + \int^-\KL(\kappa(a)\|\kappa'(a))\,(\pi_T\circ\Prob^{(T)})(da).
\end{align*}
Finally $\pi_T\circ\Prob^{(T)}$ is the law of $x_T$ under $\Prob^{(T+1)}$ (LML
\texttt{hasCondDistrib\_action}: the conditional law of $x_T$ given $\mathrm{hist}_T$ is
$\pi_T(\mathrm{hist}_T)$), and $\Prob^{(T)}$ is the image of $\Prob^{(T+1)}$ under the restriction
$\mathcal{H}_{T+1}\to\mathcal{H}_T$ (consistency of the laws of histories, LML
\texttt{history\_eq\_comp\_history}), so every term is an integral under $\Prob^{(T+1)}$ of a
function of the history, and the sum runs over $t < T+1$.

\emph{Finite case.} If the one-step divergences are bounded by $C$, the right-hand side is at
most $(n+1)C$, so all quantities are finite and \texttt{ENNReal.toReal} turns the identity into
a real identity, the integrals becoming Bochner integrals of bounded measurable functions
(\texttt{integral\_eq\_lintegral\_of\_nonneg\_ae}).
\end{proof}

\textbf{Lean remark.} The formalization (\texttt{Learning.IsAlgEnvSeq.klDiv\_map\_history})
states the theorem for two algorithm-environment sequences \texttt{IsAlgEnvSeq X Y alg
(stationaryEnv κ) P} and \texttt{IsAlgEnvSeq X' Y' alg (stationaryEnv κ') P'} on arbitrary
probability spaces, as \texttt{klDiv (P.map (history X Y T)) (P'.map (history X' Y' T)) =
∑ t ∈ range T, ∫⁻ ω, klDiv (κ (X t ω)) (κ' (X t ω)) ∂P}; this covers the trajectory
measures (which are such sequences) and avoids the uniqueness-of-law step. The finite form is
only recorded for the Gaussian case (Corollary~\ref{cor:divergence_decomposition_gaussian}).
Every conditional divergence is recorded in two forms: the \emph{integral form} above, which
needs the hypothesis \texttt{MeasurableSpace.CountablyGenerated 𝓨} for the measurability of the
integrands (Lemma~\ref{lem:kl_compProd_kernel}; it holds for $\mathcal{Y} = \R$), and the
\emph{composition-product form} \texttt{∑ t ∈ range T, klDiv (P.map (X t) ⊗ₘ κ) (P.map
(X t) ⊗ₘ κ')} (\texttt{klDiv\_map\_history\_compProd}), valid on arbitrary measurable spaces.
The chain rule for two arbitrary environments, with the conditional divergences of the step
kernels, is \texttt{klDiv\_map\_history\_stepKernel}. Only the environment changes: the policy
kernels are shared, which is what makes \texttt{klDiv\_compProd\_left} (inside
Lemma~\ref{lem:pi_kl_step}) applicable. This is the environment analogue of LML's
\texttt{Algorithm.density}, and of Lemma~\ref{lem:env_likelihood_ratio} below, which give the
Radon--Nikodym derivative rather than its expected logarithm.

\begin{lemma}[Appending the recommendation does not change the divergence]
\label{lem:kl_data_processing_recommendation}
\lean{Learning.IdentAlg.IsFixedBudget.stoppingTime_eq, Learning.IdentAlg.IsFixedBudget.stoppedHist_eq, Learning.IdentAlg.IsFixedBudget.stoppingTime_ne_top, Learning.IdentAlg.IsRun.klDiv_map_stoppedHist_out_of_isFixedBudget, Learning.IdentAlg.IsRun.klDiv_map_out_le, Learning.LinearBandit.IsRun.klDiv_map_out_le, Learning.LinearBandit.IsRun.klDiv_map_out_le_of_sq_le, Learning.LinearBandit.IsRun.exp_neg_le_measureReal_add_of_sq_le}\leanok
\uses{def:env_pair, def:bai_alg, lem:pinsker, lem:bretagnolle_huber}
In the setting of Definition~\ref{def:env_pair}, let $T \ge 1$ and let $\rho$ be a Markov
kernel from histories of length $T$ to a measurable space $\mathcal{X}$ (a recommendation rule).
Then
\[
  \KL\big(\Prob^T\otimes\rho\,\big\|\,\Prob'^T\otimes\rho\big) = \KL(\Prob^T\|\Prob'^T),
\]
and for every measurable set $E$ of pairs (history, recommendation),
\[
  (\Prob^T\otimes\rho)(E) - (\Prob'^T\otimes\rho)(E) \le \sqrt{\KL(\Prob^T\|\Prob'^T)/2},
  \qquad
  (\Prob^T\otimes\rho)(E) + (\Prob'^T\otimes\rho)(E^c) \ge \tfrac12\exp\big(-\KL(\Prob^T\|\Prob'^T)\big).
\]
When the environments are those of Definition~\ref{def:env} on an arbitrary action set
$\Xs$ (nonempty compact, not assumed spanning) and $(\mathrm{alg},\rho)$ is an
identification algorithm (Definition~\ref{def:bai_alg}), $\Prob^T\otimes\rho$ is the law
$\Prob_\theta$ of $(\hist_T,\rec)$.
\end{lemma}

\begin{proof}
\leanok
\uses{lem:pinsker, lem:bretagnolle_huber}
The equality is Mathlib \texttt{klDiv\_compProd\_left} (same kernel on both sides). The two
inequalities are Lemma~\ref{lem:pinsker} and Lemma~\ref{lem:bretagnolle_huber} applied to the
probability measures $\Prob^T\otimes\rho$, $\Prob'^T\otimes\rho$ (trivial if the divergence is
infinite). The last sentence is the definition of $\Prob_\theta$ in
Definition~\ref{def:bai_alg} together with the uniqueness of the law of the history.
\end{proof}

\textbf{Lean remark.} In the formalization the recommendation is the output of a run
(Definition~\ref{def:bai_alg}, \texttt{IsRun}): for a fixed-budget algorithm the stopping time is
$T$ and the history at the stopping time is the history of the $T$ rounds
(\texttt{IsFixedBudget.stoppingTime\_eq}, \texttt{IsFixedBudget.stoppedHist\_eq}); the law of
(history, output) is the composition-product with the output kernel, which is a probability kernel
on the stopping set, so \texttt{klDiv\_compProd\_left} applies in the form
\texttt{klDiv\_compProd\_left\_of\_ae}. Only the inequality
$\KL(\mathrm{law}(\rec)\|\mathrm{law}'(\rec)) \le \KL(\Prob^T\|\Prob'^T)$ is recorded
(\texttt{IsRun.klDiv\_map\_out\_le}, by data processing), which is what the lower bounds use, and
the Bretagnolle--Huber consequence for linear Gaussian environments,
$\Prob_\theta(\rec\in B) + \Prob_{\theta'}(\rec\notin B) \ge \tfrac12\exp(-TR^2\norm{\theta-\theta'}^2/2)$,
is \texttt{LinearBandit.IsRun.exp\_neg\_le\_measureReal\_add}.

\subsection{Stopped histories and infinite trajectories}
\label{sec:pre_info_stopping}

The paper only needs the divergence decomposition for a fixed number of rounds. For the library
we also prove it for histories stopped at a stopping time (the change-of-measure lemma of
fixed-confidence lower bounds, Kaufmann, Capp\'e and Garivier 2016) and for whole trajectories.
The results below are formalized in \texttt{Maiti2026Power/LeanMachineLearning/\{HistoryLaw,
StoppedHistory, DivergenceDecomposition, RunDivergence\}.lean} and
\texttt{Maiti2026Power/Mathlib/InformationTheory/\{KLCompProd, KLRestrict, KLFiltration\}.lean}.

\begin{definition}[Histories of the first $n$ rounds, step kernels]\label{def:fin_history}
\lean{Learning.history, Learning.stepKernel, Learning.IsAlgEnvSeq.hasCondDistrib_action, Learning.IsAlgEnvSeq.hasCondDistrib_feedback, Learning.IsAlgEnvSeq.hasCondDistrib_step, Learning.IsAlgEnvSeq.map_history_succ, Learning.map_history_succ_of_hasCondDistrib, Learning.feedback_stationaryEnv, Learning.stepKernel_stationaryEnv}\leanok
\uses{def:env_pair}
For action and feedback processes $X, Y$ on $\Omega$, the \emph{history of the first $n$ rounds}
is $\mathrm{hist}_n := (X_t, Y_t)_{t < n} : \Omega \to (\{0,\dots,n-1\}\to\mathcal{A}\times\mathcal{Y})$
(LML \texttt{Learning.history}; $n = 0$ is allowed and gives the one-point space), and
$\mathrm{hist}_{n+1} = (\mathrm{hist}_n, (X_n, Y_n))$ (\texttt{Fin.snoc}). For an algorithm
$\mathrm{alg}$ with policy kernels $\pi_n$ and an environment $\nu$ with feedback kernels
$\nu_n$, the \emph{step kernel at round $n$} is $\sigma_n := \pi_n\otimes\nu_n$, a Markov kernel
from histories of $n$ rounds to $\mathcal{A}\times\mathcal{Y}$ (LML \texttt{Learning.stepKernel}).
For an algorithm-environment sequence, $X_n$, $Y_n$ and $(X_n, Y_n)$ have conditional laws
$\pi_n$, $\nu_n$ and $\sigma_n$ given $\mathrm{hist}_n$ (and $X_n$), and the law of
$\mathrm{hist}_{n+1}$ is $\mathrm{law}(\mathrm{hist}_n)\otimes\sigma_n$. For a stationary
environment with reward kernel $\kappa$, $\nu_n = \tilde\kappa : (h, a)\mapsto\kappa(a)$ at
every round, so $\sigma_n = \pi_n\otimes\tilde\kappa$.
\end{definition}

\begin{definition}[Stopping rules, stopping times, stopped histories]\label{def:stopping_rule}
\lean{Learning.stoppingTime, Learning.stoppedHist, Learning.stoppingTime_le_iff, Learning.lt_stoppingTime_iff, Learning.stoppingTime_eq_coe_iff, Learning.stoppingTime_eq_top_iff, Learning.stoppingTime_empty, Learning.notMem_of_lt_stoppingTime, Learning.stoppedHist_mem_of_ne_top, Learning.stoppedHist_min_of_le, Learning.stoppedHist_min_of_lt, Learning.measurable_stoppingTime, Learning.measurable_stoppedHist, Learning.IsAlgEnvSeq.isStoppingTime_stoppingTime, Learning.exists_measurableSet_preimage_lt_stoppingTime, Learning.truncHist, Learning.truncHist_stoppedHist, Learning.map_stoppedHist_min_eq_map_truncHist, Learning.restrict_map_truncHist, Learning.IdentAlg.stoppingTime, Learning.IdentAlg.stoppedHist}\leanok
\uses{def:fin_history, def:bai_alg}
A \emph{stopping rule} is a measurable set $S$ of histories of variable length,
$S \subseteq \bigsqcup_n (\{0,\dots,n-1\}\to\mathcal{A}\times\mathcal{Y})$: the interaction stops after
$n$ rounds if the history of these $n$ rounds belongs to $S$. Its \emph{stopping time}
$\tau_S := \inf\{n : \mathrm{hist}_n \in S\} \in \N\cup\{\infty\}$ (Mathlib \texttt{hittingAfter})
is the number of rounds played; it is a stopping time of the history filtration of any
algorithm-environment sequence, $\{n < \tau_S\}$ is determined by $\mathrm{hist}_n$, and
$\tau_S = \infty$ for $S = \emptyset$. For a random time $\tau$, the \emph{stopped history}
$\mathrm{hist}_\tau := \langle\tau, \mathrm{hist}_\tau\rangle$ is the history of the first $\tau$
rounds as a history of variable length (of length $0$ if $\tau = \infty$); it belongs to $S$ when
$\tau = \tau_S < \infty$. The \emph{truncation} $\mathrm{trunc}_M$ keeps the first $M$ rounds of a
history of variable length: $\mathrm{trunc}_M(\mathrm{hist}_\tau) = \mathrm{hist}_{\min(\tau,M)}$
when $\tau < \infty$, so that the law of $\mathrm{hist}_{\min(\tau,M)}$ is the image of the law of
$\mathrm{hist}_\tau$ by $\mathrm{trunc}_M$ when $\tau$ is almost surely finite, and both laws
agree on histories of length $< M$. The stopping time and stopped history of an identification
algorithm (Definition~\ref{def:bai_alg}) are those of its stopping rule.
\end{definition}

\begin{lemma}[Restrictions and the divergence]\label{lem:kl_restrict}
\lean{MeasureTheory.Measure.rnDeriv_restrict_restrict, InformationTheory.klDiv_restrict_of_ac, InformationTheory.klDiv_restrict_add_restrict_compl, InformationTheory.klDiv_restrict_le, InformationTheory.klDiv_add_add_of_measure_eq_zero, InformationTheory.klDiv_eq_iSup_restrict}\leanok
Let $\mu,\nu$ be finite measures on $\Omega$ and $s$ a measurable set. Then
$\KL(\mu\|\nu) = \KL(\mu|_s\,\|\,\nu|_s) + \KL(\mu|_{s^c}\,\|\,\nu|_{s^c})$; in particular
$\KL(\mu|_s\|\nu|_s) \le \KL(\mu\|\nu)$, and the divergence of two sums of measures carried by
$s$ and $s^c$ respectively is the sum of the divergences of the two parts. If $s_n \uparrow \Omega$
is an increasing sequence of measurable sets, $\KL(\mu\|\nu) = \sup_n \KL(\mu|_{s_n}\|\nu|_{s_n})$.
\end{lemma}

\begin{proof}
\leanok
The Radon--Nikodym derivative of $\mu|_s$ with respect to $\nu|_s$ is that of $\mu$ with respect
to $\nu$ (both are $\nu|_s$-densities of $\mu|_s = (\nu|_s)\cdot\frac{d\mu}{d\nu}$), so when
$\mu\ll\nu$ the identities are those of the integral $\int\mathrm{klFun}(d\mu/d\nu)\,d\nu$ over
$s$, $s^c$ and $s_n$ (monotone convergence for the last one, $\mathrm{klFun}\ge 0$). When
$\mu\not\ll\nu$, a set $A$ with $\nu(A) = 0 < \mu(A)$ meets $s$ or $s^c$, and some $s_n$, with
positive $\mu$-measure, so the corresponding restricted divergences are infinite too.
\end{proof}

\begin{lemma}[Chain rule for stopped histories]\label{lem:stopped_history_chain_rule}
\lean{Learning.map_stoppedHist_min_eq_add, Learning.map_stoppedHist_min_succ_eq_add, Learning.map_stoppedHist_min_zero, Learning.IsAlgEnvSeq.map_history_succ_restrict_lt_stoppingTime, Learning.IsAlgEnvSeq.map_action_restrict_lt_stoppingTime, Learning.IsAlgEnvSeq.klDiv_map_stoppedHist_min_stepKernel, Learning.IsAlgEnvSeq.klDiv_map_stoppedHist_min_compProd, Learning.IsAlgEnvSeq.klDiv_map_stoppedHist_min, Learning.sum_ite_lt_eq_sum_range_toNat_min}\leanok
\uses{def:env_pair, def:fin_history, def:stopping_rule, lem:kl_restrict, lem:pi_kl_measurable_equiv, lem:pi_kl_step, lem:kl_sufficient_statistic}
Let $\mathrm{alg}$ be an algorithm, $\nu,\nu'$ two environments, $(X,Y)$ and $(X',Y')$
algorithm-environment sequences of $(\mathrm{alg},\nu)$ and $(\mathrm{alg},\nu')$ on probability
spaces $(\Omega,\Prob)$ and $(\Omega',\Prob')$, $S$ a stopping rule with stopping times $\tau$,
$\tau'$ on the two spaces, and $M \ge 0$. Then
\[
  \KL\big(\mathrm{law}(\mathrm{hist}_{\min(\tau,M)})\,\big\|\,\mathrm{law}(\mathrm{hist}'_{\min(\tau',M)})\big)
  = \sum_{t < M} \KL\big(\Prob|_{\{t<\tau\}}\circ\mathrm{hist}_t^{-1}\otimes\sigma_t\,\big\|\,
      \Prob|_{\{t<\tau\}}\circ\mathrm{hist}_t^{-1}\otimes\sigma'_t\big),
\]
the sum over the rounds of the conditional divergences of the step kernels $\sigma_t,\sigma'_t$
of the two environments given the first $t$ rounds, on the event $\{t < \tau\}$ (the round is
played). For two stationary environments with reward kernels $\kappa,\kappa'$ the $t$-th term is
$\KL(\Prob|_{\{t<\tau\}}\circ X_t^{-1}\otimes\kappa\,\|\,\Prob|_{\{t<\tau\}}\circ X_t^{-1}\otimes\kappa')$
(composition-product form) and, when $\mathcal{Y}$ is countably generated, the right-hand side is
$\E\big[\sum_{t<\min(\tau,M)}\KL(\kappa(X_t)\|\kappa'(X_t))\big]$ (integral form).
\end{lemma}

\begin{proof}
\leanok
\uses{lem:kl_restrict, lem:pi_kl_measurable_equiv, lem:pi_kl_step, lem:kl_sufficient_statistic, lem:kl_compProd_kernel}
Induction on $M$. For $M = 0$ both laws are the Dirac mass at the empty history. For the step
$M\to M+1$, write $Z_M := \mathrm{hist}_{\min(\tau,M)}$. The law of $Z_{M+1}$ splits according
to whether $\tau \le M$: on $\{\tau\le M\}$, $Z_{M+1} = Z_M$, and on $\{M < \tau\}$, $Z_{M+1}$ is the
history of the first $M+1$ rounds; the two parts are carried by the disjoint sets of histories of
length $\le M$ and of length $M+1$. Likewise the law of $Z_M$ splits into its restriction to
$\{\tau\le M\}$ (carried by $S$, as $Z_M = \mathrm{hist}_\tau \in S$ there) and the law of
$\mathrm{hist}_M$ on $\{M<\tau\}$ (carried by $S^c$). By the additivity of
Lemma~\ref{lem:kl_restrict} over disjoint supports, the invariance under the embeddings
$h\mapsto\langle n,h\rangle$ (Lemma~\ref{lem:pi_kl_measurable_equiv}) and the chain rule
\texttt{klDiv\_compProd\_eq\_add}, since on $\{M<\tau\}$ (an event determined by
$\mathrm{hist}_M$) the law of $\mathrm{hist}_{M+1}$ is $\mathrm{law}(\mathrm{hist}_M)\otimes\sigma_M$,
\[
  \KL(Z_{M+1}) = \KL(Z_M|_{\tau\le M}) + \KL(\rho_M) + \KL(\rho_M\otimes\sigma_M\,\|\,\rho_M\otimes\sigma'_M)
  = \KL(Z_M) + \KL(\rho_M\otimes\sigma_M\,\|\,\rho_M\otimes\sigma'_M),
\]
where $\rho_M$ is the law of $\mathrm{hist}_M$ on $\{M<\tau\}$ (primes omitted on the second
arguments). The stationary case is Lemma~\ref{lem:pi_kl_step} in composition-product form, the
law of $X_t$ on $\{t<\tau\}$ being the policy applied to the law of $\mathrm{hist}_t$ on that
event; the integral form is Lemma~\ref{lem:kl_compProd_kernel} and the exchange of the finite sum
with the integral.
\end{proof}

\begin{theorem}[Divergence decomposition at a stopping time]\label{thm:divergence_decomposition_stopping}
\lean{Learning.IsAlgEnvSeq.klDiv_map_stoppedHist_stepKernel, Learning.IsAlgEnvSeq.klDiv_map_stoppedHist_compProd, Learning.IsAlgEnvSeq.klDiv_map_stoppedHist}\leanok
\uses{def:env_pair, def:stopping_rule, lem:stopped_history_chain_rule, lem:kl_restrict}
In the setting of Lemma~\ref{lem:stopped_history_chain_rule}, assume that the stopping time is
almost surely finite under both $\Prob$ and $\Prob'$. Then
\[
  \KL\big(\mathrm{law}(\mathrm{hist}_\tau)\,\big\|\,\mathrm{law}(\mathrm{hist}'_{\tau'})\big)
  = \sum_{t=0}^{\infty} \KL\big(\Prob|_{\{t<\tau\}}\circ\mathrm{hist}_t^{-1}\otimes\sigma_t\,\big\|\,
      \Prob|_{\{t<\tau\}}\circ\mathrm{hist}_t^{-1}\otimes\sigma'_t\big),
\]
and for two stationary environments with reward kernels $\kappa,\kappa'$,
\[
  \KL\big(\mathrm{law}(\mathrm{hist}_\tau)\,\big\|\,\mathrm{law}(\mathrm{hist}'_{\tau'})\big)
  = \sum_{t=0}^{\infty}\KL\big(\Prob|_{\{t<\tau\}}\circ X_t^{-1}\otimes\kappa\,\big\|\,\Prob|_{\{t<\tau\}}\circ X_t^{-1}\otimes\kappa'\big)
  = \E\Big[\sum_{t<\tau}\KL(\kappa(X_t)\|\kappa'(X_t))\Big],
\]
the last form when $\mathcal{Y}$ is countably generated.
\end{theorem}

\begin{proof}
\leanok
\uses{lem:stopped_history_chain_rule, lem:kl_restrict}
The series is the supremum over $M$ of the partial sums, which are the divergences of the laws
of $\mathrm{hist}_{\min(\tau,M)}$ by Lemma~\ref{lem:stopped_history_chain_rule}. Since the
stopping times are almost surely finite, these laws are the images of the laws
$\mu,\mu'$ of the stopped histories by the truncation $\mathrm{trunc}_M$, so their divergences
are at most $\KL(\mu\|\mu')$ (data processing). Conversely, the images agree with $\mu,\mu'$ on
histories of length $< M$, and these sets increase to the whole space, so by
Lemma~\ref{lem:kl_restrict} (monotone convergence, then monotonicity under restriction)
$\KL(\mu\|\mu') = \sup_M\KL(\mu|_{\mathrm{length}<M}\|\mu'|_{\mathrm{length}<M}) \le \sup_M\KL(\mathrm{trunc}_M{}_*\mu\|\mathrm{trunc}_M{}_*\mu')$.
The integral form follows from the composition-product form as in
Lemma~\ref{lem:stopped_history_chain_rule}, the series and the integral being exchanged by
monotone convergence, and $\sum_t\mathbf 1_{t<\tau}(\cdot) = \sum_{t<\tau}(\cdot)$ on $\{\tau<\infty\}$.
\end{proof}

\begin{lemma}[Change of measure at a stopping time for runs]\label{lem:run_change_of_measure}
\lean{Learning.IdentAlg.IsRun.ae_stoppedHist_mem_stopSet, Learning.IdentAlg.IsRun.klDiv_map_stoppedHist_out, Learning.IdentAlg.IsRun.klDiv_map_out_le_stoppedHist, Learning.IdentAlg.IsRun.klDiv_map_out_le_tsum_compProd, Learning.IdentAlg.IsRun.klDiv_map_out_le_lintegral}\leanok
\uses{def:bai_alg, def:stopping_rule, thm:divergence_decomposition_stopping, lem:kl_data_processing_recommendation}
Let $A$ be an identification algorithm and consider two runs of $A$, in two environments, on
probability spaces $(\Omega,\Prob)$ and $(\Omega',\Prob')$, whose stopping times are almost surely
finite. Then the divergence between the laws of (stopped history, output) is the divergence
between the laws of the stopped histories, and the divergence between the laws of the outputs
is at most the latter. For two stationary environments with reward kernels $\kappa,\kappa'$,
\[
  \KL\big(\mathrm{law}(\rec)\,\big\|\,\mathrm{law}'(\rec')\big)
  \le \E\Big[\sum_{t<\tau}\KL(\kappa(X_t)\|\kappa'(X_t))\Big]
\]
(when $\mathcal{Y}$ is countably generated; the composition-product form holds in general).
\end{lemma}

\begin{proof}
\leanok
\uses{thm:divergence_decomposition_stopping, lem:kl_data_processing_recommendation}
As in Lemma~\ref{lem:kl_data_processing_recommendation}: the stopped history belongs to the
stopping set almost surely, on which the output rule is a probability kernel, so
\texttt{klDiv\_compProd\_left\_of\_ae} applies; then data processing along the projection on the
output, and Theorem~\ref{thm:divergence_decomposition_stopping}.
\end{proof}

\begin{lemma}[The divergence along a generating sequence of maps]\label{lem:kl_iSup_map}
\lean{InformationTheory.klDiv_map_eq_lintegral_klFun_condExp, InformationTheory.klDiv_eq_iSup_map, MeasurableSpace.iSup_comap_restrictFin}\leanok
\uses{lem:kl_restrict}
Let $\mu,\nu$ be finite measures on $\Omega$ and $g_n : \Omega\to\Omega_n$ be measurable maps whose
$\sigma$-algebras $\sigma(g_n)$ increase to the $\sigma$-algebra of $\Omega$. Then
$\KL(\mu\|\nu) = \sup_n \KL({g_n}_*\mu\,\|\,{g_n}_*\nu)$. The projections of a sequence space
$\N\to E$ on the coordinates $\le n$ satisfy the hypothesis.
\end{lemma}

\begin{proof}
\leanok
\uses{lem:kl_restrict}
The inequality $\ge$ is the data-processing inequality. For $\le$, if $\mu\ll\nu$ with density
$f$, then $\KL({g_n}_*\mu\|{g_n}_*\nu) = \int\mathrm{klFun}(\E_\nu[f\mid\sigma(g_n)])\,d\nu$ (Mathlib
\texttt{toReal\_rnDeriv\_map}), the conditional expectations converge $\nu$-almost everywhere to
$f$ by L\'evy's upward theorem (Mathlib \texttt{Integrable.tendsto\_ae\_condExp}), and Fatou's
lemma gives $\int\mathrm{klFun}(f)\,d\nu \le \liminf_n\int\mathrm{klFun}(\E_\nu[f\mid\sigma(g_n)])\,d\nu$.
If $\mu\not\ll\nu$, pick $A$ with $\nu(A) = 0 < \mu(A)$; by L\'evy's theorem under $\mu+\nu$
applied to $\mathbf 1_A$, the $\sigma(g_n)$-measurable sets $B_n := \{\E_{\mu+\nu}[\mathbf 1_A\mid\sigma(g_n)] > 1/2\}$
satisfy $\mu(B_n)\to\mu(A)$ and $\nu(B_n)\to 0$, and the lower bound
$\mu(B)\log(\mu(B)/\nu(B)) + \nu(B) - \mu(B) \le \KL(\mu|_B\|\nu|_B) \le \KL(\mu\|\nu)$ (Mathlib
\texttt{mul\_log\_le\_klDiv} and Lemma~\ref{lem:kl_restrict}), applied to the images and the
sets $B_n = g_n^{-1}(C_n)$, shows that the divergences of the images tend to infinity.
\end{proof}

\begin{theorem}[Divergence decomposition for trajectories]\label{thm:divergence_decomposition_trajectory}
\lean{Learning.klDiv_map_trajectory_eq_iSup, Learning.IsAlgEnvSeq.klDiv_map_trajectory_stepKernel, Learning.IsAlgEnvSeq.klDiv_map_trajectory_compProd, Learning.IsAlgEnvSeq.klDiv_map_trajectory}\leanok
\uses{def:env_pair, thm:divergence_decomposition, lem:kl_iSup_map}
In the setting of Definition~\ref{def:env_pair}, for two algorithm-environment sequences on
$(\Omega,\Prob)$, $(\Omega',\Prob')$, the divergence between the laws of the whole trajectories
$(X_t,Y_t)_{t\in\N}$ is
\[
  \sum_{t=0}^{\infty}\KL\big(\Prob\circ X_t^{-1}\otimes\kappa\,\big\|\,\Prob\circ X_t^{-1}\otimes\kappa'\big)
  = \sum_{t=0}^{\infty}\E\big[\KL(\kappa(X_t)\|\kappa'(X_t))\big]
\]
(the last form when $\mathcal{Y}$ is countably generated), and for two arbitrary environments it
is the series of the conditional divergences of the step kernels given the past.
\end{theorem}

\begin{proof}
\leanok
\uses{thm:divergence_decomposition, lem:kl_iSup_map}
By Lemma~\ref{lem:kl_iSup_map} applied to the projections on the first $n+1$ coordinates, the
divergence between the laws of the trajectories is the supremum over $n$ of the divergences
between the laws of the histories up to time $n$, which are the partial sums of the series by
Theorem~\ref{thm:divergence_decomposition}.
\end{proof}

\begin{corollary}[Divergence decomposition for linear Gaussian environments]
\label{cor:divergence_decomposition_gaussian}
\lean{Learning.LinearBandit.klDiv_linearGaussianKernel, Learning.LinearBandit.klDiv_map_history, Learning.LinearBandit.klDiv_map_history_le}\leanok
\uses{def:env, def:env_pair, thm:divergence_decomposition, lem:kl_gaussian}
Let $\Xs$ be an action set (Definition~\ref{def:arm_set}: nonempty compact, \emph{not}
assumed spanning; Part~\ref{part:main} applies this corollary to non-spanning sets) and
$\theta,\theta'\in\R^d$. For every
algorithm with actions in $\Xs$ and every $T \ge 1$, writing $\Prob_\theta^T$ for the law of the
history of length $T$ under $\nu_\theta$ (Definition~\ref{def:env}; this is $\Prob^T$ of
Definition~\ref{def:env_pair} with $\mathcal{A} = \Xs$, $\mathcal{Y} = \R$ and
$\kappa(x) = \Normal(\ip{x}{\theta},1)$),
\[
  \KL\big(\Prob_\theta^T\,\big\|\,\Prob_{\theta'}^T\big)
  = \frac12\sum_{t<T}\E_\theta\big[\ip{x_t}{\theta-\theta'}^2\big]
  \le \frac{T}{2}\Big(\sup_{x\in\Xs}\norm{x}\Big)^2\norm{\theta-\theta'}^2 < \infty .
\]
\end{corollary}

\begin{proof}
\leanok
\uses{thm:divergence_decomposition, lem:kl_gaussian}
The subtype $\Xs$ of $\R^d$ is countably generated, so Definition~\ref{def:env_pair}
applies with $\mathcal{A} = \Xs$ (compactness and spanning play no role there). The reward
kernels are $\kappa(x) = \Normal(\ip{x}{\theta},1)$ and
$\kappa'(x) = \Normal(\ip{x}{\theta'},1)$, so by Lemma~\ref{lem:kl_gaussian}
$\KL(\kappa(x)\|\kappa'(x)) = \tfrac12\ip{x}{\theta-\theta'}^2 \le \tfrac12 R^2\norm{\theta-\theta'}^2$
with $R := \sup_{x\in\Xs}\norm{x} < \infty$ (compactness). Apply
Theorem~\ref{thm:divergence_decomposition} in its finite form.
\end{proof}

\textbf{Lean remark.} The action type is the subtype $\Xs$ of \texttt{EuclideanSpace ℝ (Fin d)}
and the kernel is \texttt{fun x ↦ gaussianReal ⟪x, θ⟫ 1}, measurable by
\texttt{measurable\_gaussianReal} composed with the continuous inner product. The expectation
$\E_\theta[\ip{x_t}{\theta-\theta'}^2]$ is an integral of a bounded measurable function of the
history, so no integrability side condition arises.

\chapter{The Sudakov--Fernique inequality}
\label{chap:pre_sudakov}

This chapter proves the Sudakov--Fernique comparison inequality for finite families of
centered Gaussian random variables, by the Gaussian interpolation method
(\cite[Theorem~7.2.11]{vershynin2018high}), and a lower bound on the expected maximum of
independent standard Gaussians. Both are only used for the secondary lower bound
$\gw(\Xs) \ge c\sqrt{d\log d}$ (Proposition~\ref{prop:width_lower_sqrt_dlogd}); the
covariance-domination special case needed everywhere else has a direct proof
(Lemma~\ref{lem:gauss_comparison}). The chapter is written so that each step is a small
declaration: Gaussian integration by parts (Stein's identity), the derivative of an
interpolated expectation, the Hessian of the soft-max, and elementary tail estimates.

\textbf{What Mathlib/LML provides.} Gaussian vectors: \texttt{ProbabilityTheory.stdGaussian},
\texttt{multivariateGaussian}, \texttt{IsGaussian}, and their linear images (see
Chapter~\ref{chap:pre_gaussian}). Differentiation under the integral sign:
\texttt{hasDerivAt\_integral\_of\_dominated\_loc\_of\_deriv\_le}
(\texttt{Mathlib/Analysis/Calculus/ParametricIntegral.lean}). Integration by parts on $\R$ with
limits at $\pm\infty$: \texttt{MeasureTheory.integral\_mul\_deriv\_eq\_deriv\_mul} and
\texttt{integral\_deriv\_mul\_eq\_sub} in \texttt{MeasureTheory/Integral/IntegralEqImproper.lean}.
Monotonicity from the sign of the derivative: \texttt{antitoneOn\_of\_deriv\_nonpos}. The
Gaussian density \texttt{gaussianPDFReal} and \texttt{gaussianReal\_apply\_eq\_integral}.
Numerical bounds on $\pi$, $e$ and $\log 2$: \texttt{Real.pi\_lt\_d2}, \texttt{Real.pi\_gt\_d2},
\texttt{Real.exp\_one\_lt\_d9}, \texttt{Real.log\_two\_lt\_d9}, \texttt{Real.log\_two\_gt\_d9}.

\textbf{What is missing.} Stein's identity for Gaussian vectors, the interpolation formula, the
Sudakov--Fernique inequality itself, the lower Mills-ratio bound and the lower bound on the
expected maximum are all absent.

\section{Gaussian integration by parts}
\label{sec:pre_sudakov_stein}

Throughout, ``$F : \R^m\to\R$ is $C^1$ with $\norm{\nabla F}\le L$'' means that $F$ is
continuously differentiable and $\norm{\nabla F(x)}\le L$ for every $x$; such an $F$ is
$L$-Lipschitz with $\abs{F(x)}\le\abs{F(0)}+L\norm x$ (Lemma~\ref{lem:pc_lipschitz_of_grad}),
so that for a Gaussian vector $X$ the variables $F(X)$, $X_iF(X)$ and $\partial_jF(X)$ are
integrable (Lemma~\ref{lem:pc_exp_norm_integrable}, or Lemma~\ref{lem:gauss_norm_moments}
for the second moments needed for $X_iF(X)$). Similarly, ``$F$ is $C^2$ with
$\norm{\mathrm{Hess}\,F}_{\mathrm{op}}\le L$'' means that the Hessian matrix of $F$ has
operator norm at most $L$ at every point; for a $C^2$ function this is equivalent to $\nabla F$
being $L$-Lipschitz (mean value inequality, Mathlib
\texttt{lipschitzWith\_of\_nnnorm\_fderiv\_le}).

\begin{lemma}[One-dimensional Stein identity]\label{lem:psf_ibp_1d}
\uses{def:pg_gaussian_vector}
Let $\xi\sim\Normal(0,1)$, $L\ge0$, and $f : \R\to\R$ be $C^1$ with $\abs{f'(x)}\le L$ for
every $x$. Then $\xi f(\xi)$ and $f'(\xi)$ are integrable and $\E[\xi f(\xi)] = \E[f'(\xi)]$.
\end{lemma}

\begin{proof}
\uses{lem:gauss_1d_moments, lem:pc_lipschitz_of_grad}
Let $\varphi(x) := (2\pi)^{-1/2}e^{-x^2/2}$, so that $\varphi'(x) = -x\varphi(x)$ and
$\E[\xi f(\xi)] = \int x f(x)\varphi(x)\,dx = -\int f(x)\varphi'(x)\,dx$. By
Lemma~\ref{lem:pc_lipschitz_of_grad} (in dimension one), $\abs{f(x)}\le\abs{f(0)}+L\abs x$.
Integration by parts on $\R$ (Mathlib \texttt{integral\_mul\_deriv\_eq\_deriv\_mul}) requires:
$f\varphi'$ and $f'\varphi$ integrable, which holds since
$\abs{f(x)\varphi'(x)}\le(\abs{f(0)}+L\abs x)\abs x\varphi(x)$ and $\abs{f'\varphi}\le L\varphi$,
and $\xi$ has a second moment (Lemma~\ref{lem:gauss_1d_moments}); and
$f(x)\varphi(x)\to 0$ as $x\to\pm\infty$, which holds since
$\abs{f(x)\varphi(x)}\le(\abs{f(0)}+L\abs x)\varphi(x)\to0$. It gives
$-\int f\varphi' = \int f'\varphi = \E[f'(\xi)]$.
\end{proof}

\begin{lemma}[Multivariate Stein identity]\label{lem:gauss_ibp}
\uses{def:pg_gaussian_vector}
Let $S \in \PSD$ (of size $m$), $X \sim \Normal(0,S)$ on $\R^m$, $L\ge0$, and
$F : \R^m\to\R$ be $C^1$ with $\norm{\nabla F(x)}\le L$ for every $x$. Then $X_iF(X)$ and
$\partial_jF(X)$ are integrable and, for every $i \le m$,
\[
  \E\big[X_i F(X)\big] = \sum_{j=1}^m S_{ij}\,\E\big[\partial_j F(X)\big].
\]
\end{lemma}

\begin{proof}
\uses{lem:gauss_pi, lem:gauss_linear_image, lem:gauss_law_of_cov, lem:psf_ibp_1d, lem:pc_lipschitz_of_grad, lem:pc_exp_norm_integrable, lem:gauss_norm_moments}
\emph{Integrability.} By Lemma~\ref{lem:pc_lipschitz_of_grad}, $\abs{F(x)}\le\abs{F(0)}+L\norm x$,
so $\abs{X_iF(X)}\le\norm X(\abs{F(0)}+L\norm X)$ is integrable by
Lemma~\ref{lem:gauss_norm_moments} (first and second moments; or
Lemma~\ref{lem:pc_exp_norm_integrable}), and $\abs{\partial_jF(X)}\le L$ is bounded.

\emph{Standard case $S = I_m$.} By Lemma~\ref{lem:gauss_pi} the coordinates of $X = g$ are
i.i.d.\ standard Gaussians, so the law of $g$ is the product measure. Fix $i$ and write
$g = (g_i, g_{-i})$. By Fubini (justified by the integrability just shown),
$\E[g_iF(g)] = \E_{g_{-i}}\big[\E_{g_i}[g_i F(g_i,g_{-i})]\big]$, and for fixed $g_{-i}$ the
function $x\mapsto F(x,g_{-i})$ is $C^1$ with derivative $\partial_iF(x,g_{-i})$ bounded by
$L$ in absolute value (as $\abs{\partial_iF}\le\norm{\nabla F}\le L$). Lemma~\ref{lem:psf_ibp_1d}
gives $\E_{g_i}[g_iF(g_i,g_{-i})] = \E_{g_i}[\partial_iF(g_i,g_{-i})]$; integrate in $g_{-i}$.

\emph{General case.} Let $M := S^{1/2}$, so that $MM^\top = S$ and $Mg\sim\Normal(0,S)$ has the
law of $X$ (Lemmas~\ref{lem:gauss_linear_image} and~\ref{lem:gauss_law_of_cov}); expectations
of measurable functions only depend on the law, so we may take $X = Mg$. Then
$X_i = \sum_k M_{ik}g_k$ and $G := F\circ M$ is $C^1$ with
$\nabla G(g) = M^\top\nabla F(Mg)$, i.e.\ $\partial_kG(g) = \sum_j M_{jk}(\partial_jF)(Mg)$
(chain rule), and $\norm{\nabla G}\le\norm{M}_{\mathrm{op}}L$: $G$ has a bounded gradient. By
the standard case,
\[
  \E[X_iF(X)] = \sum_k M_{ik}\E[g_kG(g)] = \sum_k M_{ik}\E[\partial_kG(g)]
  = \sum_k\sum_j M_{ik}M_{jk}\E[\partial_jF(Mg)] = \sum_j S_{ij}\E[\partial_jF(X)].
\]
\end{proof}

\section{Gaussian interpolation}
\label{sec:pre_sudakov_interp}

\begin{lemma}[Derivative of an interpolated expectation]\label{lem:interpolation_derivative}
\uses{def:pg_gaussian_vector}
Let $S_X, S_Y \in \PSD$ (size $m$), $X\sim\Normal(0,S_X)$ and $Y\sim\Normal(0,S_Y)$ be
independent, $L\ge0$, and $F : \R^m\to\R$ be $C^2$ with
$\norm{\mathrm{Hess}\,F(z)}_{\mathrm{op}}\le L$ for every $z$ (so that $\nabla F$ is
$L$-Lipschitz). For $t\in[0,1]$ set
$Z_t := \sqrt t\,X + \sqrt{1-t}\,Y$ and $\psi(t) := \E[F(Z_t)]$. Then $\psi$ is continuous on
$[0,1]$, differentiable on $(0,1)$, and for $t\in(0,1)$,
\[
  \psi'(t) = \frac12\sum_{i,j=1}^m (S_X - S_Y)_{ij}\,\E\big[\partial_{ij}F(Z_t)\big].
\]
\end{lemma}

\begin{proof}
\uses{lem:gauss_pi, lem:gauss_linear_image, lem:gauss_law_of_cov, lem:gauss_ibp, lem:gauss_norm_moments}
Since $\norm{\mathrm{Hess}\,F}_{\mathrm{op}}\le L$ everywhere, $\nabla F$ is $L$-Lipschitz
(mean value inequality applied to the $C^1$ map $\nabla F$, whose derivative is the Hessian;
Mathlib \texttt{lipschitzWith\_of\_nnnorm\_fderiv\_le}), so
$\norm{\nabla F(z)} \le \norm{\nabla F(0)} + L\norm{z}$ and $\abs{F(z)} \le \abs{F(0)} + \norm{\nabla F(0)}\norm{z} + L\norm{z}^2/2$.
Hence $F(Z_t)$, $\partial_iF(Z_t)$ and $X_i\partial_iF(Z_t)$ are integrable for every $t$
(Lemma~\ref{lem:gauss_norm_moments}: $X$ and $Y$ have second moments), and each
$\partial_iF$ is $C^1$ with $\norm{\nabla\partial_iF(z)} = \norm{\mathrm{Hess}\,F(z)\,e_i}\le L$.

\emph{Continuity.} $\abs{F(Z_t) - F(Z_s)} \le (\norm{\nabla F(0)} + L(\norm{X}+\norm{Y}))\,(\abs{\sqrt t-\sqrt s}\norm{X} + \abs{\sqrt{1-t}-\sqrt{1-s}}\norm{Y})$,
whose expectation tends to $0$ as $s\to t$ (second moments, Lemma~\ref{lem:gauss_norm_moments}).

\emph{Differentiability.} For $t \in (0,1)$, $\frac{d}{dt}F(Z_t) = \big\langle\nabla F(Z_t),\ \tfrac{X}{2\sqrt t} - \tfrac{Y}{2\sqrt{1-t}}\big\rangle$,
which on any interval $[a,b]\subset(0,1)$ is dominated by
$(\norm{\nabla F(0)} + L(\norm{X}+\norm{Y}))(\norm{X}+\norm{Y})/(2\min(\sqrt a,\sqrt{1-b}))$, an
integrable function. Differentiation under the integral sign (Mathlib
\texttt{hasDerivAt\_integral\_of\_dominated\_loc\_of\_deriv\_le}) gives
\[
  \psi'(t) = \frac{1}{2\sqrt t}\sum_i\E\big[X_i\,\partial_iF(Z_t)\big]
  - \frac{1}{2\sqrt{1-t}}\sum_i\E\big[Y_i\,\partial_iF(Z_t)\big].
\]

\emph{Stein's identity.} Write $X = S_X^{1/2}g$, $Y = S_Y^{1/2}g'$ with $(g,g')$ a standard
Gaussian vector of $\R^{2m}$ (Lemmas~\ref{lem:gauss_pi}, \ref{lem:gauss_law_of_cov}; independence
of $X$ and $Y$ is the product structure of the law of $(g,g')$). Then $W := (X,Y)$ is the linear
image of $(g,g')$ by the block-diagonal matrix $\mathrm{diag}(S_X^{1/2},S_Y^{1/2})$, hence
$W\sim\Normal(0,\mathrm{diag}(S_X,S_Y))$ (Lemma~\ref{lem:gauss_linear_image}). Fix $i$ and apply
Lemma~\ref{lem:gauss_ibp} to $W$ and the $C^1$ function
$G(x,y) := \partial_iF(\sqrt t\,x + \sqrt{1-t}\,y)$, whose gradient is
$\nabla G(x,y) = (\sqrt t\,\nabla\partial_iF(Z), \sqrt{1-t}\,\nabla\partial_iF(Z))$ with
$Z = \sqrt t x+\sqrt{1-t}y$, so that
$\norm{\nabla G(x,y)}^2 = (t + (1-t))\norm{\mathrm{Hess}\,F(Z)\,e_i}^2\le L^2$: $G$ has a
gradient bounded by $L$, as Lemma~\ref{lem:gauss_ibp} requires.
The covariance of $W$ is block-diagonal, so
\[
  \E[X_iG(W)] = \sum_j (S_X)_{ij}\E[\partial_{x_j}G(W)] = \sqrt t\sum_j(S_X)_{ij}\E[\partial_{ij}F(Z_t)],
  \qquad
  \E[Y_iG(W)] = \sqrt{1-t}\sum_j(S_Y)_{ij}\E[\partial_{ij}F(Z_t)].
\]
Substituting in the formula for $\psi'(t)$, the factors $\sqrt t$ and $\sqrt{1-t}$ cancel and
we obtain $\psi'(t) = \frac12\sum_{ij}(S_X)_{ij}\E[\partial_{ij}F(Z_t)] - \frac12\sum_{ij}(S_Y)_{ij}\E[\partial_{ij}F(Z_t)]$.
\end{proof}

\begin{lemma}[The soft-max and its Hessian]\label{lem:logsumexp_hessian}
For $\beta > 0$ and $z\in\R^m$ let $F_\beta(z) := \beta^{-1}\log\sum_{i=1}^m e^{\beta z_i}$ and
$p_i(z) := e^{\beta z_i}/\sum_k e^{\beta z_k}$, so that $p(z)$ is a probability vector. Then:
\begin{enumerate}
  \item[(i)] $\max_i z_i \le F_\beta(z) \le \max_i z_i + (\log m)/\beta$;
  \item[(ii)] $F_\beta$ is $C^\infty$, $\partial_iF_\beta = p_i$, hence $\norm{\nabla F_\beta(z)}\le 1$
        and $\abs{F_\beta(z) - F_\beta(z')} \le \norm{z-z'}$;
  \item[(iii)] $\partial_{ij}F_\beta = \beta(\delta_{ij}p_i - p_ip_j)$, i.e.\
        $\mathrm{Hess}\,F_\beta(z) = \beta(\operatorname{diag}(p) - pp^\top)$, a positive
        semidefinite matrix with $\norm{\mathrm{Hess}\,F_\beta(z)}_{\mathrm{op}}\le\beta$ (and
        all entries bounded by $\beta$ in absolute value);
  \item[(iv)] for every symmetric matrix $\Gamma\in\R^{m\times m}$ and every $z$,
        \[
          \sum_{i,j}\Gamma_{ij}\,\partial_{ij}F_\beta(z)
          = \frac{\beta}{2}\sum_{i,j}p_i(z)p_j(z)\,\big(\Gamma_{ii} + \Gamma_{jj} - 2\Gamma_{ij}\big).
        \]
\end{enumerate}
\end{lemma}

\begin{proof}
(i): $e^{\beta\max_i z_i} \le \sum_i e^{\beta z_i} \le m\,e^{\beta\max_i z_i}$; take logarithms
and divide by $\beta$. (ii): $\partial_i\log\sum_ke^{\beta z_k} = \beta e^{\beta z_i}/\sum_ke^{\beta z_k}$;
$\norm{p}_2 \le \norm{p}_1 = 1$; the Lipschitz bound is the mean value inequality
(Mathlib \texttt{Convex.norm\_image\_sub\_le\_of\_norm\_fderiv\_le}). (iii): differentiate
$p_i = e^{\beta z_i}/\sum_k e^{\beta z_k}$: $\partial_jp_i = \beta\delta_{ij}p_i - \beta p_ip_j$,
and $\abs{\delta_{ij}p_i - p_ip_j}\le 1$. For the operator norm, for every $v\in\R^m$, with
$\bar v := \sum_ip_iv_i$,
$v^\top(\operatorname{diag}(p) - pp^\top)v = \sum_ip_iv_i^2 - \bar v^2 = \sum_ip_i(v_i-\bar v)^2$,
the variance of $v$ under the probability vector $p$, which lies in
$[0, \sum_ip_iv_i^2]\subseteq[0,\norm v^2]$; so the symmetric matrix
$\mathrm{Hess}\,F_\beta(z)$ has its eigenvalues in $[0,\beta]$. (iv): by (iii),
$\sum_{ij}\Gamma_{ij}\partial_{ij}F_\beta = \beta\big(\sum_i\Gamma_{ii}p_i - \sum_{ij}\Gamma_{ij}p_ip_j\big)$.
On the other hand, using $\sum_jp_j = 1$,
\[
  \sum_{i,j}p_ip_j(\Gamma_{ii}+\Gamma_{jj}-2\Gamma_{ij})
  = \sum_ip_i\Gamma_{ii} + \sum_jp_j\Gamma_{jj} - 2\sum_{ij}p_ip_j\Gamma_{ij}
  = 2\Big(\sum_i\Gamma_{ii}p_i - \sum_{ij}\Gamma_{ij}p_ip_j\Big),
\]
and multiplying by $\beta/2$ gives the claim.
\end{proof}

\begin{theorem}[Sudakov--Fernique inequality]\label{thm:sudakov_fernique}
\uses{def:pg_gaussian_vector}
Let $X$ and $Y$ be centered Gaussian vectors in $\R^m$, with laws $\Normal(0,S_X)$ and
$\Normal(0,S_Y)$, such that
\[
  \E\big[(X_i - X_j)^2\big] \le \E\big[(Y_i - Y_j)^2\big] \qquad\text{for all } i,j\le m .
\]
Then $\E\big[\max_{i\le m}X_i\big] \le \E\big[\max_{i\le m}Y_i\big]$.
\end{theorem}

\begin{proof}
\uses{lem:gauss_law_of_cov, lem:interpolation_derivative, lem:logsumexp_hessian}
Both sides only depend on the laws of $X$ and $Y$, so we may assume that $X$ and $Y$ are
independent (realize them on a product space; Lemma~\ref{lem:gauss_law_of_cov}). Fix
$\beta > 0$, let $F_\beta$ be the soft-max of Lemma~\ref{lem:logsumexp_hessian}, and
$Z_t := \sqrt t\,X + \sqrt{1-t}\,Y$, $\psi(t) := \E[F_\beta(Z_t)]$, so that $Z_0 = Y$ and
$Z_1 = X$. By Lemma~\ref{lem:logsumexp_hessian}(iii) $F_\beta$ is $C^2$ with
$\norm{\mathrm{Hess}\,F_\beta}_{\mathrm{op}}\le\beta$, so Lemma~\ref{lem:interpolation_derivative}
applies with $L = \beta$: $\psi$ is continuous on
$[0,1]$, differentiable on $(0,1)$, and with $\Gamma := S_X - S_Y$ (symmetric) and
Lemma~\ref{lem:logsumexp_hessian}(iv),
\[
  \psi'(t) = \frac12\,\E\Big[\sum_{i,j}\Gamma_{ij}\partial_{ij}F_\beta(Z_t)\Big]
  = \frac{\beta}{4}\,\E\Big[\sum_{i,j}p_i(Z_t)p_j(Z_t)\big(\Gamma_{ii}+\Gamma_{jj}-2\Gamma_{ij}\big)\Big].
\]
Now $\Gamma_{ii}+\Gamma_{jj}-2\Gamma_{ij} = \big((S_X)_{ii}+(S_X)_{jj}-2(S_X)_{ij}\big) - \big((S_Y)_{ii}+(S_Y)_{jj}-2(S_Y)_{ij}\big)
= \E[(X_i-X_j)^2] - \E[(Y_i-Y_j)^2] \le 0$ by hypothesis, and $p_ip_j \ge 0$, so $\psi'(t)\le 0$
on $(0,1)$. Hence $\psi$ is nonincreasing on $[0,1]$ (Mathlib
\texttt{antitoneOn\_of\_deriv\_nonpos} with continuity on the closed interval), and
$\E[F_\beta(X)] = \psi(1) \le \psi(0) = \E[F_\beta(Y)]$. Finally, by
Lemma~\ref{lem:logsumexp_hessian}(i),
\[
  \E\big[\max_iX_i\big] \le \E[F_\beta(X)] \le \E[F_\beta(Y)] \le \E\big[\max_iY_i\big] + \frac{\log m}{\beta},
\]
and letting $\beta\to\infty$ concludes (all expectations are finite since $\abs{\max_iX_i}\le\norm{X}$).
\end{proof}

\begin{remark}[On the sign of the interpolation derivative]\label{rem:psf_sign}
The outline of this chapter hesitated about the sign of $\psi'$. The computation above is the
correct one: with $Z_t = \sqrt t X + \sqrt{1-t}Y$ the derivative is
$\frac12\sum_{ij}(S_X-S_Y)_{ij}\E[\partial_{ij}F_\beta(Z_t)]$ (no minus sign), and the Hessian
identity of Lemma~\ref{lem:logsumexp_hessian}(iv) carries the factor $+\beta/2$ (the sign error in
the outline came from writing $-\beta/2$ there). Smaller increments for $X$ therefore make
$\psi$ \emph{nonincreasing} from $\psi(0) = \E F_\beta(Y)$ to $\psi(1) = \E F_\beta(X)$, which is
exactly $\E\max X\le\E\max Y$, the standard statement (\cite[Theorem~7.2.11]{vershynin2018high}).
The compact-index-set version (suprema over $\Xs$) follows by density,
Lemma~\ref{lem:sup_countable_dense}, but is not needed.
\end{remark}

\section{Expected maximum of independent Gaussians}
\label{sec:pre_sudakov_max}

\begin{lemma}[Lower Mills-ratio bound]\label{lem:gauss_tail_lower}
Let $\xi\sim\Normal(0,1)$ and $\varphi(t) := (2\pi)^{-1/2}e^{-t^2/2}$. For every $t \ge 0$,
\[
  \Prob(\xi\ge t) \ge \frac{t}{1+t^2}\,\varphi(t) = \frac{t}{1+t^2}\cdot\frac{e^{-t^2/2}}{\sqrt{2\pi}} .
\]
\end{lemma}

\begin{proof}
Let $h(t) := \Prob(\xi\ge t) - \frac{t}{1+t^2}\varphi(t)$ for $t\in\R$. The tail
$t\mapsto\Prob(\xi\ge t) = \int_t^\infty\varphi$ is differentiable with derivative $-\varphi(t)$
(fundamental theorem of calculus, Mathlib \texttt{intervalIntegral.integral\_hasDerivAt\_right}
after writing the tail as $1 - \int_{-\infty}^t\varphi$; the tail is
\texttt{gaussianReal\_apply\_eq\_integral}). Using $\varphi'(t) = -t\varphi(t)$,
\[
  \frac{d}{dt}\Big[\frac{t}{1+t^2}\varphi(t)\Big]
  = \varphi(t)\Big[\frac{1-t^2}{(1+t^2)^2} - \frac{t^2}{1+t^2}\Big],
  \qquad
  h'(t) = -\varphi(t)\Big[1 + \frac{1-t^2}{(1+t^2)^2} - \frac{t^2}{1+t^2}\Big]
  = -\frac{2\,\varphi(t)}{(1+t^2)^2} < 0,
\]
since $1 - \frac{t^2}{1+t^2} = \frac{1}{1+t^2}$ and $\frac{1+t^2}{(1+t^2)^2} + \frac{1-t^2}{(1+t^2)^2} = \frac{2}{(1+t^2)^2}$.
Thus $h$ is decreasing (Mathlib \texttt{antitone\_of\_deriv\_nonpos}) and
$h(t)\to 0$ as $t\to\infty$ (both terms tend to $0$: the tail by dominated convergence, the
second term since $\varphi(t)\to 0$). Hence $h(t) \ge \lim_{s\to\infty}h(s) = 0$ for every $t$.
\end{proof}

\begin{lemma}[Expected maximum of two standard Gaussians]\label{lem:psf_max_two}
\uses{lem:gauss_pi, lem:gauss_linear_image, lem:gauss_1d_moments}
If $\xi_1,\xi_2$ are i.i.d.\ $\Normal(0,1)$ then $\E[\max(\xi_1,\xi_2)] = 1/\sqrt\pi \ge 0.56$.
\end{lemma}

\begin{proof}
\uses{lem:gauss_pi, lem:gauss_linear_image, lem:gauss_1d_moments}
$\max(a,b) = \frac{a+b}{2} + \frac{\abs{a-b}}{2}$. The first term has expectation $0$. By
Lemmas~\ref{lem:gauss_pi} and~\ref{lem:gauss_linear_image}, $\xi_1-\xi_2\sim\Normal(0,2)$, so
$\E\abs{\xi_1-\xi_2} = \sqrt2\cdot\sqrt{2/\pi} = 2/\sqrt\pi$ (Lemma~\ref{lem:gauss_1d_moments}),
and $\E\max(\xi_1,\xi_2) = 1/\sqrt\pi$; finally $\pi < 3.15$ gives $1/\sqrt\pi > 0.563$.
\end{proof}

\begin{lemma}[Monotonicity in the number of variables]\label{lem:psf_max_mono}
Let $(\xi_i)_{i\ge1}$ be real random variables with $\E\abs{\xi_i}<\infty$. For $1\le m\le m'$,
$\E[\max_{i\le m}\xi_i] \le \E[\max_{i\le m'}\xi_i]$.
\end{lemma}

\begin{proof}
Pointwise, $\max_{i\le m}\xi_i\le\max_{i\le m'}\xi_i$; both are integrable since
$\abs{\max_{i\le m}\xi_i}\le\sum_{i\le m}\abs{\xi_i}$.
\end{proof}

\begin{lemma}[Lower bound on the expected maximum, large $m$]\label{lem:psf_max_lower_large}
\uses{lem:gauss_tail_lower, lem:gauss_1d_moments}
Let $\xi_1,\dots,\xi_m$ be i.i.d.\ $\Normal(0,1)$ with $m \ge 128$. Then
$\E[\max_{i\le m}\xi_i] \ge \frac14\sqrt{\log m}$.
\end{lemma}

\begin{proof}
\uses{lem:gauss_tail_lower, lem:gauss_1d_moments}
Let $t := \sqrt{\log m}$ and $p_t := \Prob(\xi_1\ge t)$. We proceed in steps.
\begin{enumerate}
  \item[(i)] Since $m\ge128 = 2^7$, $t^2 = \log m \ge 7\log 2 > 4.85$ (Mathlib
        \texttt{Real.log\_two\_gt\_d9}), so $t > 2.2 > 1$.
  \item[(ii)] By Lemma~\ref{lem:gauss_tail_lower} and $e^{-t^2/2} = m^{-1/2}$,
        $p_t \ge \frac{t}{1+t^2}\cdot\frac{m^{-1/2}}{\sqrt{2\pi}} \ge \frac{m^{-1/2}}{2t\sqrt{2\pi}}$,
        using $\frac{t}{1+t^2}\ge\frac{1}{2t}$ for $t\ge1$ (equivalent to $t^2\ge1$).
  \item[(iii)] Hence $mp_t \ge \frac{\sqrt m}{2\sqrt{2\pi}\sqrt{\log m}} \ge 1$: the inequality
        $\sqrt m \ge 2\sqrt{2\pi}\sqrt{\log m}$ is equivalent to $m \ge 8\pi\log m$; at
        $m = 128$ it reads $128\ge 8\pi\cdot7\log2$, true since $8\pi\cdot7\log 2 < 8\cdot3.15\cdot7\cdot0.6932 < 122.3$;
        and $m\mapsto m/\log m$ is nondecreasing on $[e,\infty)$ (its derivative is
        $(\log m - 1)/(\log m)^2\ge0$), so $m\ge 8\pi\log m$ for all $m\ge128$.
  \item[(iv)] By independence, $\Prob(\max_i\xi_i < t) = (1-p_t)^m \le e^{-mp_t} \le e^{-1}$,
        so $\Prob(\max_i\xi_i\ge t) \ge 1 - e^{-1} \ge \frac12$ (Mathlib
        \texttt{Real.add\_one\_le\_exp}, \texttt{Real.exp\_one\_gt\_d9}).
  \item[(v)] Write $M := \max_i\xi_i = M^+ - M^-$. Since $t > 0$, $M^+ \ge t\,\indic\{M\ge t\}$, so
        $\E[M^+] \ge t\,\Prob(M\ge t) \ge t/2$. Since $M\ge\xi_1$, $M^- \le \xi_1^-$ and
        $\E[M^-]\le\E[\xi_1^-] = \frac12\E\abs{\xi_1} = \frac{1}{\sqrt{2\pi}} < 0.4$
        (Lemma~\ref{lem:gauss_1d_moments} and symmetry of $\xi_1$).
  \item[(vi)] Therefore $\E[M] \ge \frac t2 - 0.4 \ge \frac t4$, because $\frac t4 \ge 0.4$ is
        implied by $t > 2.2$. That is, $\E[\max_i\xi_i]\ge\frac14\sqrt{\log m}$.
\end{enumerate}
\end{proof}

\begin{lemma}[Lower bound on the expected maximum of i.i.d.\ Gaussians]\label{lem:gauss_max_lower}
\uses{lem:psf_max_two, lem:psf_max_mono, lem:psf_max_lower_large}
Let $\xi_1,\dots,\xi_m$ be i.i.d.\ $\Normal(0,1)$ with $m\ge2$. Then
\[
  \E\Big[\max_{i\le m}\xi_i\Big] \ge c_{\max}\sqrt{\log m},\qquad c_{\max} := \frac14 .
\]
\end{lemma}

\begin{proof}
\uses{lem:psf_max_two, lem:psf_max_mono, lem:psf_max_lower_large}
If $m\ge128$ this is Lemma~\ref{lem:psf_max_lower_large}. If $2\le m<128$, then by
Lemmas~\ref{lem:psf_max_mono} and~\ref{lem:psf_max_two},
$\E[\max_{i\le m}\xi_i]\ge\E[\max(\xi_1,\xi_2)] = 1/\sqrt\pi$, and
$\frac14\sqrt{\log m} \le \frac14\sqrt{\log 128} = \frac14\sqrt{7\log2}$. It remains to check
$1/\sqrt\pi\ge\frac14\sqrt{7\log2}$, i.e.\ $16/\pi\ge7\log2$: indeed $16/\pi > 16/3.15 > 5.07$
and $7\log 2 < 7\cdot0.6932 < 4.86$ (Mathlib \texttt{Real.pi\_lt\_d2},
\texttt{Real.log\_two\_lt\_d9}).
\end{proof}

\begin{remark}[Constants]\label{rem:psf_constants}
The constant $c_{\max} = 1/4$ holds for every $m\ge2$; the argument is not sharp
($\E\max_i\xi_i\sim\sqrt{2\log m}$), but only an absolute constant is needed for
Proposition~\ref{prop:width_lower_sqrt_dlogd}. The splitting at $m = 128$ is chosen so that
both branches only require the crude numerical bounds $\pi<3.15$, $0.693<\log2<0.6932$ and
$e^{-1}<1/2$, which are available in Mathlib. The outline suggested the threshold
$t = \sqrt{\log m}$ and a Paley--Zygmund argument on the number of exceedances; the direct
bound $\Prob(\max\ge t)\ge1-e^{-mp_t}$ used above is simpler and gives the same order.
\end{remark}

\textbf{Lean remark.} The i.i.d.\ family is \texttt{Fin m → ℝ} with law
\texttt{Measure.pi (fun \_ ↦ gaussianReal 0 1)}; the maximum is \texttt{Finset.sup'} or
\texttt{⨆ i, ξ i}. In Theorem~\ref{thm:sudakov_fernique} the vectors are
\texttt{EuclideanSpace ℝ (Fin m)}-valued with laws \texttt{multivariateGaussian 0 S}, and the
soft-max $F_\beta$ is \texttt{fun z ↦ β⁻¹ * Real.log (∑ i, Real.exp (β * z i))}, whose
derivatives are best handled through \texttt{ContDiff} and \texttt{fderiv} of the composition
$\log\circ\sum\circ\exp$.

\chapter{Median Elimination}
\label{chap:pre_median}

This chapter formalizes the Median Elimination algorithm of Even-Dar, Mannor and
Mansour~\cite{even2002pac,evendar2006action} for $\eps$-best arm identification in a
$K$-armed bandit with sub-Gaussian rewards, with a fixed (deterministic) budget of
$O(K\log(1/\delta)/\eps^2)$ samples. It is the second phase of the algorithm of
Chapter~\ref{chap:log_gains} (paper, Appendix~C), where the arms are $d$ candidate vectors of
$\Xs$ selected from a first, non-adaptive phase. The algorithm is written as an LML
\texttt{Algorithm} with a deterministic policy, and its guarantee is proved through the
sequential composition of Lemma~\ref{lem:composition}, one elimination round at a time.

\textbf{What Mathlib/LML provides.} Sub-Gaussian random variables
\texttt{ProbabilityTheory.HasSubgaussianMGF X c μ} with the tail bounds \texttt{measure\_ge\_le}
(Mathlib) and \texttt{measure\_le\_le} (LML \texttt{ForMathlib/Probability/Moments/SubGaussian.lean}),
closure under sums of independent variables \texttt{HasSubgaussianMGF.sum\_of\_iIndepFun} and scaling
\texttt{HasSubgaussianMGF.const\_mul}, and Hoeffding's inequality
\texttt{measure\_sum\_ge\_le\_of\_iIndepFun}. LML provides stationary environments
\texttt{stationaryEnv ν} for a kernel \texttt{ν : Kernel (Fin K) ℝ}, deterministic algorithms
\texttt{detAlgorithm nextA} and the class \texttt{IsDeterministicAlg}, the round-robin
algorithm \texttt{roundRobinAlgorithm} (\texttt{SequentialLearning/Algorithms/RoundRobin.lean}),
pull counts \texttt{pullCount}, cumulative sums and empirical means \texttt{sumRewards},
\texttt{empMean} (\texttt{SequentialLearning/SumRewards.lean}), the conditional law of the
feedback given the history and the current action (the structure field
\texttt{IsAlgEnvSeq.hasCondDistrib\_feedback n : HasCondDistrib (Y (n+1)) (history A Y n, A (n+1)) (env.feedback n) P},
and \texttt{IsObliviousEnv.hasCondDistrib\_feedback\_history\_action} in
\texttt{SequentialLearning/StationaryEnv.lean}; the lemma
\texttt{hasCondDistrib\_feedback\_stationaryEnv} conditions on the current action only), and
the sequential composition of Lemma~\ref{lem:composition}.

\textbf{What is missing.} The Median Elimination algorithm, its per-round analysis and its
sample complexity are not in LML. The empirical means used by the algorithm are per-round
averages (not the cumulative \texttt{empMean}), so a small variant of \texttt{empMean}
restricted to a time window is needed.

\section{Sub-Gaussian bandits and fixed designs}
\label{sec:pre_median_setting}

\begin{definition}[Sub-Gaussian $K$-armed bandit]\label{def:pm_subgaussian_bandit}
Let $K\ge1$. A \emph{$1$-sub-Gaussian $K$-armed bandit} is a Markov kernel
$\nu : \mathrm{Fin}\,K\to\R$ (LML \texttt{Kernel (Fin K) ℝ}) such that, for every arm $a$, the
mean $\mu_a := \int y\,\nu(a,dy)$ exists and the centered reward $y - \mu_a$ under $\nu(a)$ has
sub-Gaussian moment generating function with constant $1$: $\int e^{\lambda(y-\mu_a)}\nu(a,dy)\le e^{\lambda^2/2}$
for all $\lambda\in\R$ (Mathlib \texttt{HasSubgaussianMGF (fun y ↦ y - μ a) 1 (ν a)}). The
associated environment is $\mathrm{stationaryEnv}(\nu)$. The Gaussian bandit
$\nu(a) = \Normal(\mu_a,1)$ is $1$-sub-Gaussian (Mathlib \texttt{mgf\_gaussianReal}).
\end{definition}

\begin{lemma}[Rewards of a history-independent design are independent]\label{lem:pm_fixed_design_rewards}
\uses{def:pm_subgaussian_bandit}
Let $\nu$ be a Markov kernel from $\mathcal{A}$ to $\mathcal{Y}$, $a_0,a_1,\dots\in\mathcal{A}$ a
deterministic sequence, and $\mathrm{alg}$ the deterministic algorithm playing $a_t$ at round $t$
regardless of the history (LML \texttt{detAlgorithm} with a history-independent next action).
Under any algorithm-environment sequence for $(\mathrm{alg},\mathrm{stationaryEnv}(\nu))$, the
observations $(y_t)_{t\ge0}$ are independent and $y_t\sim\nu(a_t)$ for every $t$.
\end{lemma}

\begin{proof}
By induction on $t$. The structure field \texttt{IsAlgEnvSeq.hasCondDistrib\_feedback n} of
LML states that the conditional law of $y_{n+1}$ given the \emph{pair} (history up to time
$n$, $x_{n+1}$) is the feedback kernel \texttt{env.feedback n} evaluated at that pair; for a
stationary environment this kernel is $(h,a)\mapsto\nu(a)$ (\texttt{feedback\_stationaryEnv};
equivalently \texttt{IsObliviousEnv.hasCondDistrib\_feedback\_history\_action} in
\texttt{SequentialLearning/StationaryEnv.lean}, whose kernel is
\texttt{(feedbackCondAction env (n+1)).prodMkLeft \_}). Since the algorithm is deterministic
and history-independent, $x_{n+1} = a_{n+1}$ almost surely (\texttt{action\_ae\_eq}), so the
conditional law of $y_{n+1}$ given $(\mathrm{hist}_n, x_{n+1})$ is the \emph{constant} kernel
$\nu(a_{n+1})$: it depends neither on the past observations nor on the action. A conditional
distribution given $(\mathrm{hist}_n,x_{n+1})$ that is a constant kernel means that $y_{n+1}$
is independent of $(\mathrm{hist}_n,x_{n+1})$, hence of $(y_s)_{s\le n}$, with law
$\nu(a_{n+1})$ (Mathlib \texttt{HasCondDistrib} with a constant kernel gives \texttt{IndepFun}
and the marginal law). Note that \texttt{hasCondDistrib\_feedback\_stationaryEnv} conditions
on the current action only and would \emph{not} give independence from the past. The case
$t = 0$ is \texttt{hasCondDistrib\_feedback\_zero} (law $\nu(x_0) = \nu(a_0)$). Independence of
$y_{n+1}$ from $(y_s)_{s\le n}$ for every $n$ gives joint independence of the whole sequence
(\texttt{iIndepFun}, by induction on finite subfamilies). This is the argument of
Lemma~\ref{lem:noise_representation} for a general kernel.
\end{proof}

\begin{lemma}[Tail of an empirical mean]\label{lem:pm_empmean_tail}
\uses{def:pm_subgaussian_bandit}
Let $Y_1,\dots,Y_n$ be independent real random variables such that each $Y_i - \mu$ has
sub-Gaussian moment generating function with constant $1$, and let $\hat\mu := \frac1n\sum_iY_i$.
Then $\hat\mu - \mu$ has sub-Gaussian moment generating function with constant $1/n$ and, for
every $u\ge0$,
\[
  \Prob(\hat\mu - \mu\ge u)\le e^{-nu^2/2},\qquad \Prob(\hat\mu-\mu\le -u)\le e^{-nu^2/2}.
\]
In particular, if $n\ge 8\log(3/\delta)/\eps^2$ then each of $\Prob(\hat\mu-\mu\ge\eps/2)$ and
$\Prob(\hat\mu-\mu\le-\eps/2)$ is at most $\delta/3$.
\end{lemma}

\begin{proof}
Mathlib \texttt{HasSubgaussianMGF.sum\_of\_iIndepFun} (independent summands, constants add up to $n$) and
\texttt{HasSubgaussianMGF.const\_mul} with the factor $1/n$ (constant $n\cdot(1/n)^2 = 1/n$),
then \texttt{measure\_ge\_le} / \texttt{measure\_le\_le}: $\Prob(X\ge u)\le\exp(-u^2/(2c))$ with
$c = 1/n$. With $u = \eps/2$ the bound is $\exp(-n\eps^2/8)\le\exp(-\log(3/\delta)) = \delta/3$.
\end{proof}

\section{The algorithm}
\label{sec:pre_median_alg}

\begin{definition}[Schedule of Median Elimination]\label{def:me_schedule}
For $\eps>0$, $\delta\in(0,1)$ and $\ell\in\N$ (rounds are $0$-indexed) let
\[
  \eps_\ell := \frac{\eps}{4}\Big(\frac34\Big)^\ell,\qquad
  \delta_\ell := \frac{\delta}{2^{\ell+1}},\qquad
  n_\ell := \Big\lceil\frac{8\log(3/\delta_\ell)}{\eps_\ell^2}\Big\rceil ,
\]
so that $\eps_0 = \eps/4$, $\delta_0 = \delta/2$, $\eps_{\ell+1} = 3\eps_\ell/4$ and
$\delta_{\ell+1} = \delta_\ell/2$.
\end{definition}

\begin{lemma}[Sums of the schedule]\label{lem:me_schedule_sums}
\uses{def:me_schedule}
$\sum_{\ell\ge0}\eps_\ell = \eps$ and $\sum_{\ell\ge0}\delta_\ell = \delta$; in particular every
finite partial sum is at most $\eps$, resp.\ $\delta$. Moreover $n_\ell\ge 8\log(3/\delta_\ell)/\eps_\ell^2$
and $\log(3/\delta_\ell) = \log(3/\delta) + (\ell+1)\log 2$.
\end{lemma}

\begin{proof}
\uses{def:me_schedule}
Geometric series: $\frac\eps4\sum_\ell(3/4)^\ell = \frac\eps4\cdot4$ and
$\frac\delta2\sum_\ell 2^{-\ell} = \delta$ (Mathlib \texttt{tsum\_geometric\_of\_lt\_one}); the
rest is the definition and $\log(2^{\ell+1}) = (\ell+1)\log2$.
\end{proof}

\begin{definition}[Halving sizes and number of rounds]\label{def:pm_halving}
For $K\ge1$ and $\ell\in\N$ let $s_\ell := \lceil K/2^\ell\rceil$ and
$L := \min\{\ell : s_\ell = 1\} = \lceil\log_2K\rceil$. Then $s_0 = K$,
$s_{\ell+1} = \lceil s_\ell/2\rceil$, $s_\ell\ge2$ for $\ell<L$, and $s_\ell\le 2K/2^\ell$ for
$\ell<L$ (since $\lceil x\rceil\le x+1\le 2x$ when $x\ge1$, and $K/2^\ell>1$ when $s_\ell\ge2$).
\end{definition}

\begin{definition}[Top-half selection]\label{def:pm_top}
Let $S\subseteq\mathrm{Fin}\,K$ be finite with $\abs S = s\ge2$ and $v : S\to\R$. Order $S$ by the
strict total order $a\succ b$ iff $v_a>v_b$, or $v_a = v_b$ and $a<b$ (ties broken by index).
$\mathrm{Top}(S,v)$ is the set of the $\lceil s/2\rceil$ largest elements of $S$ for $\succ$.
It has the property that if $a\in S\setminus\mathrm{Top}(S,v)$ then $v_b\ge v_a$ for every
$b\in\mathrm{Top}(S,v)$.
\end{definition}

\begin{definition}[Median Elimination]\label{def:median_elimination}
\uses{def:pm_subgaussian_bandit, def:me_schedule, def:pm_halving, def:pm_top}
Let $K\ge1$, $\eps>0$, $\delta\in(0,1)$. Median Elimination $\mathrm{ME}(K,\eps,\delta)$ is the
deterministic algorithm with actions in $\mathrm{Fin}\,K$ and observations in $\R$ defined round
by round. Set $S_0 := \mathrm{Fin}\,K$ and $\tau_0 := 0$. For $\ell<L$, given the surviving set
$S_\ell$ (with $\abs{S_\ell} = s_\ell$) and the starting time $\tau_\ell$:
\begin{enumerate}
  \item[(i)] \emph{Round $\ell$} occupies the times $\tau_\ell\le t<\tau_{\ell+1} := \tau_\ell + s_\ell n_\ell$;
        the action at time $\tau_\ell + j$ is the $(j\bmod s_\ell)$-th element of $S_\ell$ in
        increasing order (round-robin over $S_\ell$), so each arm of $S_\ell$ is pulled exactly
        $n_\ell$ times during the round, at times that do not depend on the observations.
  \item[(ii)] For $a\in S_\ell$, $\hat\mu^{(\ell)}_a$ is the average of the $n_\ell$ observations
        received when pulling $a$ during round $\ell$ (a measurable function of the history of
        length $\tau_{\ell+1}$).
  \item[(iii)] $S_{\ell+1} := \mathrm{Top}(S_\ell,\hat\mu^{(\ell)})$, of size $s_{\ell+1}$.
\end{enumerate}
The \emph{budget} of the algorithm is the deterministic integer
$N_{\mathrm{ME}}(K,\eps,\delta) := \tau_L = \sum_{\ell<L}s_\ell n_\ell$, and its
\emph{recommendation} is the unique element $\hat a$ of $S_L$, a measurable function of the
history of length $N_{\mathrm{ME}}$. After time $\tau_L$ the algorithm plays an arbitrary fixed
arm (its actions are irrelevant). Formally, $\mathrm{ME}(K,\eps,\delta)$ \emph{is} the LML
deterministic algorithm \texttt{detAlgorithm nextA} whose next-action function
$\mathrm{nextA}(n,h)$ recomputes $S_0,\dots,S_\ell$ from the history $h$ of length $n$ by
(ii)--(iii), where $\ell$ is the round with $\tau_\ell\le n<\tau_{\ell+1}$, and returns the
action prescribed by (i) (see the Lean remark below). Its relation to the sequential
composition of Lemma~\ref{lem:composition} is the content of
Lemma~\ref{lem:pm_policy_composition}; it is not part of the definition.
\end{definition}

\textbf{Lean remark.} The algorithm is \texttt{detAlgorithm nextA} where
\texttt{nextA n h} decodes from $n$ the round $\ell$ with $\tau_\ell\le n<\tau_{\ell+1}$ and the
offset $j = n-\tau_\ell$ (all deterministic functions of $(K,\eps,\delta,n)$), recomputes
$S_0,\dots,S_\ell$ from the history \texttt{h : Iic (n-1) → Fin K × ℝ} by the recursion
(ii)--(iii), and returns the $(j\bmod s_\ell)$-th element of $S_\ell$ (\texttt{Finset.sort} or
\texttt{Finset.orderEmbOfFin}). Measurability of \texttt{nextA n} holds because it is a function
of finitely many real coordinates of \texttt{h} through comparisons: the map
$h\mapsto S_\ell(h)$ takes finitely many values and each fiber is a finite Boolean combination
of measurable sets $\{\hat\mu^{(k)}_a\ge\hat\mu^{(k)}_b\}$. The per-round empirical mean
$\hat\mu^{(\ell)}_a$ is \texttt{(∑ t ∈ Finset.Ico (τ ℓ) (τ (ℓ+1)) with (h t).1 = a, (h t).2) / n ℓ}.
The definition deliberately does not go through the composition operation of
Lemma~\ref{lem:composition} (an existence lemma cannot be part of a definition); the
identification with the composition, round by round, is Lemma~\ref{lem:pm_policy_composition}.

\begin{lemma}[Median Elimination as a sequential composition, round by round]
\label{lem:pm_policy_composition}
\uses{def:median_elimination, lem:composition}
Let $\ell<L$ and $T_1 := \tau_\ell$. Let $\mathrm{alg}_1 := \mathrm{ME}(K,\eps,\delta)$ and, for a
history $h$ of length $T_1$, let $\mathrm{alg}_2(h)$ be the history-independent deterministic
algorithm playing the round-robin design of Definition~\ref{def:median_elimination}(i) over
$S_\ell(h)$ (the set computed from $h$ by (ii)--(iii)) for $s_\ell n_\ell$ rounds, then an
arbitrary fixed arm. Then $h\mapsto\mathrm{alg}_2(h)$ is measurable in the sense of
Lemma~\ref{lem:composition}, and the policy of $\mathrm{ME}(K,\eps,\delta)$ at every time
$t<\tau_{\ell+1}$ coincides with the policy of the sequential composition of $\mathrm{alg}_1$ and
$\mathrm{alg}_2$ (Lemma~\ref{lem:composition}) with $T_1 = \tau_\ell$. Consequently, against
every environment $\nu$, the law of the history of length $\tau_{\ell+1}$ under
$\mathrm{ME}(K,\eps,\delta)$ is its law under the composition; in particular, conditionally
on the history $h$ of length $\tau_\ell$, the actions and observations of round $\ell$ form an
algorithm-environment sequence for $(\mathrm{alg}_2(h),\nu)$, and the probability bound of
Lemma~\ref{lem:composition} applies to events of the history of length $\tau_{\ell+1}$.
\end{lemma}

\begin{proof}
\uses{def:median_elimination, lem:composition}
For $t<\tau_\ell$ both algorithms follow the policy of $\mathrm{alg}_1 = \mathrm{ME}$. For
$\tau_\ell\le t<\tau_{\ell+1}$ and a history $h'$ of length $t$ with prefix $h$ of length
$\tau_\ell$, $\mathrm{nextA}(t,h')$ decodes the round $\ell$ and the offset $j = t-\tau_\ell$,
recomputes $S_0,\dots,S_\ell$ from $h'$, and plays the $(j\bmod s_\ell)$-th element of
$S_\ell$. The sets $S_k$, $k\le\ell$, only depend on the observations received before time
$\tau_\ell$ (Definition~\ref{def:median_elimination}(ii)--(iii)), i.e.\ on $h$, so this action is
the action of $\mathrm{alg}_2(h)$ at its time $j$: the two policy kernels at time $t$ are the
same deterministic kernel of the same measurable function of $h'$. Measurability of
$h\mapsto\mathrm{alg}_2(h)$: $\mathrm{alg}_2(h)$ is the fixed round-robin algorithm
reparametrized by $S_\ell(h)$, a function with finitely many values whose fibers are finite
Boolean combinations of the measurable sets $\{\hat\mu^{(k)}_a\ge\hat\mu^{(k)}_b\}$, $k<\ell$,
so the policy of $\mathrm{alg}_2(h)$ at time $j$, as a function of $(h,\cdot)$, is a
deterministic kernel of a measurable function. Since the trajectory measure is determined by
the initial law and the policy kernels (\texttt{Learning.trajMeasure}) and the two algorithms
agree at all times $t<\tau_{\ell+1}$, the laws of the histories of length $\tau_{\ell+1}$
coincide (the policies at later times do not affect this law, by the consistency of
\texttt{trajMeasure} under \texttt{frestrictLe}). The conditional statement and the
probability bound are then those of Lemma~\ref{lem:composition}.
\end{proof}

\section{Analysis}
\label{sec:pre_median_analysis}

\begin{lemma}[The median counting argument]\label{lem:pm_selection}
\uses{def:pm_top}
Let $S$ be finite with $\abs S = s\ge2$, $\mu,v : S\to\R$, $\eps>0$, and $a_*\in S$ with
$\mu_{a_*} = \max_{a\in S}\mu_a$. Define the set of \emph{bad} arms
\[
  B := \{a\in S : \mu_a < \mu_{a_*}-\eps \text{ and } v_a\ge v_{a_*}\}.
\]
If $\abs B < s/2$ then $\max_{a\in\mathrm{Top}(S,v)}\mu_a\ge\mu_{a_*}-\eps$.
\end{lemma}

\begin{proof}
\uses{def:pm_top}
Let $S' := \mathrm{Top}(S,v)$, of size $\lceil s/2\rceil\ge s/2>\abs B$. If $a_*\in S'$ the
conclusion holds since $\mu_{a_*}\ge\mu_{a_*}-\eps$. Otherwise $a_*\in S\setminus S'$, so by the
property of Definition~\ref{def:pm_top} every $b\in S'$ satisfies $v_b\ge v_{a_*}$. Since
$\abs{S'}>\abs B$, there is $b\in S'\setminus B$; as $v_b\ge v_{a_*}$ and $b\notin B$, the first
condition defining $B$ fails: $\mu_b\ge\mu_{a_*}-\eps$. Hence $\max_{a\in S'}\mu_a\ge\mu_b\ge\mu_{a_*}-\eps$.
\end{proof}

\begin{lemma}[One round of Median Elimination]\label{lem:me_round}
\uses{def:pm_top, lem:pm_selection, lem:pm_empmean_tail, lem:pm_fixed_design_rewards, def:median_elimination, def:me_schedule}
Let $S$ be finite with $\abs S = s\ge2$, $\eps>0$, $\delta\in(0,1)$, and let $(\hat\mu_a)_{a\in S}$
be real random variables (not necessarily independent) and $(\mu_a)_{a\in S}$ real numbers such
that for every $a\in S$,
\[
  \Prob(\hat\mu_a-\mu_a\ge\eps/2)\le\delta/3\qquad\text{and}\qquad\Prob(\hat\mu_a-\mu_a\le-\eps/2)\le\delta/3 .
\]
Then
\[
  \Prob\Big(\max_{a\in\mathrm{Top}(S,\hat\mu)}\mu_a\ge\max_{a\in S}\mu_a-\eps\Big)\ge1-\delta .
\]
In particular this applies, with $\delta = \delta_\ell$ and $\eps = \eps_\ell$, to the empirical
means $\hat\mu^{(\ell)}$ of round $\ell$ of Median Elimination conditionally on the history
before the round, by Lemmas~\ref{lem:pm_fixed_design_rewards} and~\ref{lem:pm_empmean_tail}.
\end{lemma}

\begin{proof}
\uses{lem:pm_selection, lem:pm_empmean_tail, lem:pm_fixed_design_rewards, def:median_elimination, lem:me_schedule_sums, lem:pm_policy_composition}
Fix $a_*\in S$ with $\mu_{a_*} = \max_S\mu$ and let $B$ be the (random) set of bad arms of
Lemma~\ref{lem:pm_selection} with $v = \hat\mu$. Consider the events
\[
  E_1 := \{\hat\mu_{a_*}<\mu_{a_*}-\eps/2\},\qquad
  N := \sum_{a\in S}\indic\{\hat\mu_a-\mu_a\ge\eps/2\},\qquad E_2 := \{N\ge s/2\}.
\]
By hypothesis $\Prob(E_1)\le\delta/3$ and $\E[N]\le s\delta/3$, so by Markov's inequality
$\Prob(E_2)\le\frac{s\delta/3}{s/2} = 2\delta/3$; hence $\Prob(E_1\cup E_2)\le\delta$. On
$E_1^c$, every bad arm $a\in B$ satisfies $\hat\mu_a\ge\hat\mu_{a_*}\ge\mu_{a_*}-\eps/2>\mu_a+\eps/2$,
so $\indic_B(a)\le\indic\{\hat\mu_a-\mu_a\ge\eps/2\}$ and $\abs B\le N$. Therefore on
$E_1^c\cap E_2^c$ we have $\abs B\le N<s/2$, and Lemma~\ref{lem:pm_selection} gives
$\max_{\mathrm{Top}(S,\hat\mu)}\mu\ge\mu_{a_*}-\eps$. The last sentence: conditionally on the
history before round $\ell$, the set $S_\ell$ is fixed and the actions of the round form a
history-independent design (Definition~\ref{def:median_elimination}(i)), i.e.\ round $\ell$ is
an algorithm-environment sequence for the fixed design $\mathrm{alg}_2(h)$ of
Lemma~\ref{lem:pm_policy_composition}; so by
Lemma~\ref{lem:pm_fixed_design_rewards} the $n_\ell$ observations of each arm $a\in S_\ell$ are
independent with law $\nu(a)$, and Lemma~\ref{lem:pm_empmean_tail} with
$n_\ell\ge8\log(3/\delta_\ell)/\eps_\ell^2$ (Lemma~\ref{lem:me_schedule_sums}) gives the two tail
bounds with $\delta_\ell/3$.
\end{proof}

\begin{theorem}[Median Elimination is $(\eps,\delta)$-PAC]\label{thm:median_elimination}
\uses{def:pm_subgaussian_bandit, def:median_elimination, lem:me_round, lem:me_schedule_sums, lem:composition, lem:pm_policy_composition}
Let $\nu$ be a $1$-sub-Gaussian $K$-armed bandit with means $(\mu_a)_a$, $\eps>0$ and
$\delta\in(0,1)$. Run against $\mathrm{stationaryEnv}(\nu)$, Median Elimination
$\mathrm{ME}(K,\eps,\delta)$ uses the deterministic budget $N_{\mathrm{ME}}(K,\eps,\delta)$ and
its recommendation $\hat a$ satisfies
\[
  \Prob\Big(\mu_{\hat a}\ge\max_{a}\mu_a-\eps\Big)\ge1-\delta .
\]
\end{theorem}

\begin{proof}
\uses{lem:me_round, lem:me_schedule_sums, lem:composition, lem:pm_fixed_design_rewards, lem:pm_policy_composition}
If $K = 1$ then $L = 0$, $\hat a$ is the unique arm and there is nothing to prove. Let $K\ge2$.
For $\ell<L$ let $G_\ell := \{\max_{a\in S_\ell}\mu_a\le\max_{a\in S_{\ell+1}}\mu_a+\eps_\ell\}$,
an event of the history of length $\tau_{\ell+1}$. By Lemma~\ref{lem:pm_policy_composition},
the history of length $\tau_{\ell+1}$ under $\mathrm{ME}$ has the law of the composition
(Lemma~\ref{lem:composition}) of $\mathrm{ME}$ up to time $\tau_\ell$ with the fixed design
of round $\ell$ parametrized by $S_\ell$, so by Lemma~\ref{lem:me_round} the conditional
probability of $G_\ell$ given the history of length $\tau_\ell$ is at least $1-\delta_\ell$
for every value of that history.
Iterating the last clause of Lemma~\ref{lem:composition} over $\ell = 0,\dots,L-1$ (with
$E_h := G_\ell$ and $\delta' := \sum_{k<\ell}\delta_k$ at step $\ell$),
$\Prob(\bigcap_{\ell<L}G_\ell)\ge1-\sum_{\ell<L}\delta_\ell\ge1-\delta$ by
Lemma~\ref{lem:me_schedule_sums}. On $\bigcap_{\ell<L}G_\ell$, telescoping,
\[
  \max_{a}\mu_a = \max_{a\in S_0}\mu_a\le\max_{a\in S_L}\mu_a+\sum_{\ell<L}\eps_\ell
  = \mu_{\hat a}+\sum_{\ell<L}\eps_\ell\le\mu_{\hat a}+\eps ,
\]
again by Lemma~\ref{lem:me_schedule_sums}, since $S_L = \{\hat a\}$.
\end{proof}

\begin{lemma}[Budget of Median Elimination]\label{lem:me_budget}
\uses{def:median_elimination, def:me_schedule, def:pm_halving, lem:me_schedule_sums}
For all $K\ge1$, $\eps>0$, $\delta\in(0,1)$,
\[
  N_{\mathrm{ME}}(K,\eps,\delta)\le\frac{256K}{\eps^2}\Big(9\log\frac3\delta+81\log2\Big)+4K .
\]
If moreover $\eps\le1$, then $N_{\mathrm{ME}}(K,\eps,\delta)\le C_{\mathrm{ME}}\,K\log(3/\delta)/\eps^2$
with $C_{\mathrm{ME}} := 16000$.
\end{lemma}

\begin{proof}
\uses{lem:me_schedule_sums}
By Definition~\ref{def:pm_halving}, $s_\ell\le2K/2^\ell$ for $\ell<L$, and
$n_\ell\le8\log(3/\delta_\ell)/\eps_\ell^2+1$ with $1/\eps_\ell^2 = 16(16/9)^\ell/\eps^2$ and
$\log(3/\delta_\ell) = \log(3/\delta)+(\ell+1)\log2$ (Lemma~\ref{lem:me_schedule_sums}). Hence
\[
  s_\ell n_\ell\le\frac{2K}{2^\ell}\cdot\frac{128\,(16/9)^\ell}{\eps^2}\Big(\log\frac3\delta+(\ell+1)\log2\Big)+\frac{2K}{2^\ell}
  = \frac{256K}{\eps^2}\Big(\frac89\Big)^\ell\Big(\log\frac3\delta+(\ell+1)\log2\Big)+\frac{2K}{2^\ell}.
\]
Summing over $\ell<L$ and extending the sums to all $\ell\ge0$ (nonnegative terms), with
$\sum_{\ell\ge0}(8/9)^\ell = 9$, $\sum_{\ell\ge0}(\ell+1)(8/9)^\ell = 1/(1-8/9)^2 = 81$ (Mathlib
\texttt{tsum\_geometric\_of\_lt\_one}, \texttt{tsum\_coe\_mul\_geometric\_of\_norm\_lt\_one}) and
$\sum_{\ell\ge0}2^{-\ell} = 2$,
\[
  N_{\mathrm{ME}} = \sum_{\ell<L}s_\ell n_\ell\le\frac{256K}{\eps^2}\Big(9\log\frac3\delta+81\log2\Big)+4K .
\]
For the second form, use $\log(3/\delta)\ge\log3>1.09$ (as $\delta<1$) and $\eps\le1$:
$81\log2 = \frac{81\log2}{\log3}\log3\le 51.2\log(3/\delta)$ (since $81\cdot0.6932/1.0986<51.2$),
so $9\log(3/\delta)+81\log2\le 60.2\log(3/\delta)$ and $256\cdot60.2<15412$; and
$4K\le 4K\log(3/\delta)/(\eps^2\log3)\le 3.7K\log(3/\delta)/\eps^2$. Altogether
$N_{\mathrm{ME}}\le(15412+3.7)K\log(3/\delta)/\eps^2\le16000\,K\log(3/\delta)/\eps^2$.
\end{proof}

\begin{remark}[Constants of the budget]\label{rem:pm_budget_constants}
The outline announced the bound $(128K/\eps^2)(9\log(3/\delta)+81\log2)+2K$; the factor $128$
there implicitly uses $s_\ell\le K/2^\ell$, which fails because of the ceilings
($s_\ell = \lceil K/2^\ell\rceil$). The correct bookkeeping uses $s_\ell\le2K/2^\ell$ (valid while
$s_\ell\ge2$), which doubles the constants: $256$ and $4K$. The resulting
$C_{\mathrm{ME}} = 16000$ is far from optimal (the schedule of Definition~\ref{def:me_schedule}
is the one of \cite{even2002pac} and was not tuned), but only its existence matters for
Theorem~\ref{thm:log_gains}.
\end{remark}

\begin{corollary}[Median Elimination on candidate arms of a linear bandit]\label{cor:median_elimination_linear}
\uses{def:env, def:median_elimination, thm:median_elimination, lem:me_budget, lem:composition, lem:gauss_tail}
Let $\Xs$ be an action set, $\theta\in\R^d$, $K\ge1$, and $x^{(1)},\dots,x^{(K)}\in\Xs$. Consider
the algorithm with actions in $\Xs$ that runs $\mathrm{ME}(K,\eps,\delta)$ on the arms
$\mathrm{Fin}\,K$ and plays $x^{(a)}$ whenever $\mathrm{ME}$ pulls $a$ (its next action is
$x^{(\mathrm{nextA}(n,\cdot))}$ applied to the sequence of past observations). Against
$\nu_\theta$ (Definition~\ref{def:env}) it uses the budget $N_{\mathrm{ME}}(K,\eps,\delta)\le C_{\mathrm{ME}}K\log(3/\delta)/\eps^2$
(for $\eps\le1$) and its recommendation $\hat a$ satisfies
\[
  \Prob_\theta\Big(\ip{x^{(\hat a)}}{\theta}\ge\max_{a\le K}\ip{x^{(a)}}{\theta}-\eps\Big)\ge1-\delta .
\]
The same holds when the candidates $x^{(a)} = x^{(a)}(h)$ are measurable functions of a history $h$
of a previous phase, conditionally on $h$, so that by Lemma~\ref{lem:composition} the composed
algorithm satisfies the conclusion with $\delta$ replaced by $\delta+\delta'$ if the previous
phase fails with probability at most $\delta'$.
\end{corollary}

\begin{proof}
\uses{thm:median_elimination, lem:me_budget, lem:composition, lem:gauss_tail}
Let $\nu'(a) := \Normal(\ip{x^{(a)}}{\theta},1)$, a $1$-sub-Gaussian $K$-armed bandit with means
$\mu_a = \ip{x^{(a)}}{\theta}$ (Definition~\ref{def:pm_subgaussian_bandit}; the Gaussian moment
generating function is \texttt{mgf\_gaussianReal}, see also Lemma~\ref{lem:gauss_tail}). The
sequence of arms pulled and observations received by the $\Xs$-valued algorithm against
$\nu_\theta$ is an algorithm-environment sequence for $(\mathrm{ME},\mathrm{stationaryEnv}(\nu'))$:
the next arm is the same deterministic function of past observations, and the conditional law
of the observation given the past and the arm $a$ is $\Normal(\ip{x^{(a)}}{\theta},1) = \nu'(a)$
(the structure field \texttt{IsAlgEnvSeq.hasCondDistrib\_feedback} for $\nu_\theta$, whose
kernel is $(h,x)\mapsto\Normal(\ip{x}{\theta},1)$, composed with the map $a\mapsto x^{(a)}$). Hence Theorem~\ref{thm:median_elimination} applies and gives the
probability bound, and Lemma~\ref{lem:me_budget} the budget. The conditional version is
Lemma~\ref{lem:composition} with $\mathrm{alg}_2(h)$ the algorithm just described for the
candidates $x^{(a)}(h)$; measurability in $h$ holds because $h\mapsto(x^{(a)}(h))_a$ is
measurable and the policy of $\mathrm{alg}_2(h)$ is the deterministic kernel of a measurable
function of $(h,\text{observations})$.
\end{proof}

\textbf{Lean remark.} The $\Xs$-valued algorithm is \texttt{detAlgorithm} with next action
\texttt{fun n h ↦ x (nextA n (fun t ↦ (h t).2))}, where \texttt{nextA} only reads observations;
this avoids having to recover the arm index from the played vector $x^{(a)}$ (the map
$a\mapsto x^{(a)}$ need not be injective). The identification of the two
algorithm-environment sequences is an instance of the general fact that mapping the actions of
a deterministic algorithm through a measurable map $f$ turns an algorithm-environment sequence
for the environment $\nu\circ f$ into one for $\nu$, to be contributed to LML.

\chapter{Gaussian posteriors for adaptively collected data and the unit-ball lower bound}
\label{chap:pre_bayes}

This chapter replaces the citation of \cite{chen2024assouad} in Appendix~H of the paper by a
self-contained argument: for every (adaptive, possibly randomized) identification algorithm
with budget $T$ on the unit ball $\ball_d$, there is a reward vector $\theta$ of norm
$O(d/\sqrt T)$ under which the expected simple regret is $\Omega(d/\sqrt T)$; consequently an
$(\eps,\delta)$-PAC algorithm on $\ball_d$ needs $T = \Omega(d^2/\eps^2)$ samples
(Corollary~\ref{cor:ball_pac_lower}, used in Chapter~\ref{chap:log_gains}). The proof is
Bayesian: under a Gaussian prior $\theta\sim\Normal(0,\sigma^2I_d)$, the posterior mean of
$\theta$ given the history is $m_T = (\sigma^{-2}I+S_T)^{-1}b_T$, and the recommendation cannot
correlate with $\theta$ better than with $m_T$. We never use abstract conditional expectations:
the key identity $\E\ip{\varphi}{\theta-m_T} = 0$ is proved by Fubini, the likelihood ratio of
the history with respect to the noise-only environment, and an explicit completion of the
square in $\R^d$. The likelihood ratio is first established for general stationary environments
whose kernels have densities, as the environment analogue of LML's \texttt{Algorithm.density}.

\textbf{What Mathlib/LML provides.} Measures with densities \texttt{Measure.withDensity},
\texttt{withDensity\_mul}, \texttt{Measure.compProd\_withDensity}
(\texttt{Mathlib/Probability/Kernel/CompProdEqIff.lean}: $\mu\otimes(\kappa.\mathrm{withDensity}\,f) = (\mu\otimes\kappa).\mathrm{withDensity}(f)$),
and in LML \texttt{ForMathlib/Probability/WithDensity.lean}: \texttt{compProd\_withDensity\_left},
\texttt{compProd\_withDensity\_withDensity}, \texttt{map\_equiv\_withDensity},
\texttt{Kernel.compProd\_withDensity\_left} (a density on the \emph{left} factor of a kernel
composition-product). The algorithm density
\texttt{Algorithm.density} and \texttt{IsAlgEnvSeq.hasLaw\_history\_withDensity}
(\texttt{SequentialLearning/AlgorithmDensity.lean}) are the model for
Lemma~\ref{lem:env_likelihood_ratio}. Bayesian sequences: \texttt{IsBayesAlgEnvSeq},
\texttt{bayesStationaryEnv}, \texttt{IT.bayesTrajMeasure}, \texttt{IsBayesAlgEnvSeq.ae\_IsAlgEnvSeq}
(\texttt{SequentialLearning/BayesStationaryEnv.lean}) give the joint law of a parameter drawn
from a prior and the trajectory. Gaussians: \texttt{gaussianReal\_of\_var\_ne\_zero},
\texttt{gaussianPDFReal}, \texttt{integral\_gaussian}, \texttt{multivariateGaussian},
\texttt{stdGaussian}, the linear change of variables for the Lebesgue measure
\texttt{MeasureTheory.Measure.map\_linearMap\_addHaar\_eq\_smul\_addHaar} and
\texttt{MeasureTheory.integral\_map\_equiv}.

\textbf{Formalization note.} The chapter is formalized along the plan above, with two
deviations. The posterior identities of Section~\ref{sec:pre_bayes_posterior} are obtained
from Stein's identity (Gaussian integration by parts, Lemma~\ref{lem:gauss_ibp}, extended to
functions of exponential growth in Lemma~\ref{lem:gauss_ibp_exp}) applied to the tilt
$\theta \mapsto \exp(\ip{b}{\theta} - \tfrac12\theta^\top S\theta)$, instead of the Lebesgue
density of the prior and a change of variables; no density with respect to the Lebesgue
measure on $\R^d$ is ever used. And the lower bound $\E\norm\theta \ge \sigma d/\sqrt{d+2}$ of
the outline is replaced by the sharper $\E\norm\theta \ge \sqrt{2/\pi}\,\sigma\sqrt d$ of
Lemma~\ref{lem:gauss_norm_lower}, which improves the constants (we keep the outline's
headline constants).

\section{Likelihood ratio of the history with respect to a reference environment}
\label{sec:pre_bayes_lr}

\begin{definition}[Environment density along a history]\label{def:pb_env_density}
\lean{Learning.envDensity, Learning.measurable_envDensity, Learning.envDensity_zero, Learning.envDensity_snoc}\leanok
Let $\mathcal{A},\mathcal{Y}$ be measurable spaces and $\rho : \mathcal{A}\times\mathcal{Y}\to[0,\infty]$
be measurable. For $n\ge0$ and a history $h = (x_t,y_t)_{t\le n}$ define
$D_n^\rho(h) := \prod_{t=0}^n\rho(x_t,y_t)$, a measurable function of $h$. For a history of
length $T\ge1$ we write $D_T^\rho := D_{T-1}^\rho$.
\end{definition}

\begin{lemma}[Composition-product of a kernel with a kernel having a density]
\label{lem:pb_kernel_compProd_withDensity_right}
\lean{ProbabilityTheory.Kernel.compProd_withDensity_right, ProbabilityTheory.Kernel.prodMkLeft_withDensity}\leanok
Let $\pi$ be an s-finite kernel from $\mathcal{H}$ to $\mathcal{A}$, $\eta$ an s-finite kernel
from $\mathcal{H}\times\mathcal{A}$ to $\mathcal{Y}$, and
$f : (\mathcal{H}\times\mathcal{A})\times\mathcal{Y}\to[0,\infty]$ measurable such that the
kernel $\eta.\mathrm{withDensity}(f)$ is s-finite (in all applications below it is a Markov
kernel). Then, as kernels from $\mathcal{H}$ to $\mathcal{A}\times\mathcal{Y}$,
\[
  \pi\otimes\big(\eta.\mathrm{withDensity}(f)\big)
  = (\pi\otimes\eta).\mathrm{withDensity}\big((h,(a,y))\mapsto f((h,a),y)\big).
\]
In particular, if $f((h,a),y) = \rho(a,y)$ does not depend on $h$, the density of the
composition-product is $(h,(a,y))\mapsto\rho(a,y)$, i.e.\ $\rho\circ\mathrm{snd}$.
\end{lemma}

\begin{proof}
\leanok
(The Lean proof is not fiberwise: both sides are evaluated on a measurable set through
\texttt{Kernel.compProd\_apply}, \texttt{Kernel.withDensity\_apply'} and
\texttt{Kernel.lintegral\_compProd}, exactly as Mathlib's measure-level
\texttt{Measure.compProd\_withDensity}.) Fiberwise. Fix $h\in\mathcal{H}$. For every kernel $\eta'$ from $\mathcal{H}\times\mathcal{A}$
to $\mathcal{Y}$, $(\pi\otimes\eta')(h) = \pi(h)\otimes\eta'_h$, where
$\eta'_h := \eta'((h,\cdot))$ is the section kernel from $\mathcal{A}$ to $\mathcal{Y}$
(Mathlib \texttt{Kernel.compProd\_apply\_eq\_compProd\_sectR}). The section of a kernel with
a density is the section with the sectioned density:
$(\eta.\mathrm{withDensity}(f))_h = \eta_h.\mathrm{withDensity}\big((a,y)\mapsto f((h,a),y)\big)$,
both sides being the measure $\eta((h,a)).\mathrm{withDensity}(f((h,a),\cdot))$ at each $a$
(\texttt{Kernel.withDensity\_apply}). The measure-level identity
\texttt{Measure.compProd\_withDensity} (\texttt{Mathlib/Probability/Kernel/CompProdEqIff.lean})
then gives
\[
  \pi(h)\otimes\big(\eta_h.\mathrm{withDensity}(f((h,\cdot),\cdot))\big)
  = \big(\pi(h)\otimes\eta_h\big).\mathrm{withDensity}\big((a,y)\mapsto f((h,a),y)\big),
\]
and the right-hand side is $\big((\pi\otimes\eta).\mathrm{withDensity}(\tilde f)\big)(h)$
with $\tilde f(h,(a,y)) := f((h,a),y)$, by \texttt{Kernel.withDensity\_apply} again
($\tilde f$ is measurable, being $f$ composed with a measurable rearrangement of the
variables) and \texttt{Kernel.compProd\_apply\_eq\_compProd\_sectR}. Two kernels that agree
at every point are equal (\texttt{Kernel.ext}).
\end{proof}

\begin{lemma}[Environment likelihood ratio]\label{lem:env_likelihood_ratio}
\lean{Learning.IsAlgEnvSeq.map_history_eq_withDensity_env, Learning.IsAlgEnvSeq.map_finHistory_eq_withDensity_env, Learning.toFinHistory_IicSuccProd_symm}\leanok
\uses{def:pb_env_density, lem:pb_kernel_compProd_withDensity_right}
Let $\kappa_0,\kappa$ be Markov kernels from $\mathcal{A}$ to $\mathcal{Y}$ and
$\rho : \mathcal{A}\times\mathcal{Y}\to[0,\infty]$ measurable with
$\kappa(a) = \kappa_0(a).\mathrm{withDensity}(\rho(a,\cdot))$ for every $a\in\mathcal{A}$. Let
$\mathrm{alg}$ be any algorithm and, for $n\ge0$, let $\Prob^{(n)}$ and $\Prob_0^{(n)}$ be the
laws of the history $(x_t,y_t)_{t\le n}$ under $\mathrm{stationaryEnv}(\kappa)$, resp.\
$\mathrm{stationaryEnv}(\kappa_0)$ (Definition~\ref{def:env_pair}). Then
\[
  \Prob^{(n)} = \Prob_0^{(n)}.\mathrm{withDensity}\big(D_n^\rho\big),
\]
so $\Prob^{(n)}\ll\Prob_0^{(n)}$ and, for every measurable $f\ge0$ (or $f$ integrable under
$\Prob^{(n)}$), $\int f\,d\Prob^{(n)} = \int f\,D_n^\rho\,d\Prob_0^{(n)}$.
\end{lemma}

\begin{proof}
\leanok
\uses{def:pb_env_density, lem:pb_kernel_compProd_withDensity_right}
Induction on $n$, exactly as LML's \texttt{hasLaw\_history\_withDensity} for the algorithm
density. Write $\tilde\kappa$, $\tilde\kappa_0$ for the kernels $(h,a)\mapsto\kappa(a)$,
$(h,a)\mapsto\kappa_0(a)$ (the feedback kernels of the stationary environments, LML
\texttt{feedback\_stationaryEnv}), so that
$\tilde\kappa = \tilde\kappa_0.\mathrm{withDensity}((h,a),y)\mapsto\rho(a,y))$.

\emph{Base case.} $(x_0,y_0)$ has law $p_0\otimes\kappa = p_0\otimes(\kappa_0.\mathrm{withDensity}\,\rho)
= (p_0\otimes\kappa_0).\mathrm{withDensity}((a,y)\mapsto\rho(a,y))$ (Mathlib
\texttt{Measure.compProd\_withDensity}), and $\mathrm{hist}_0$ is its image by the measurable
equivalence $\mathcal{A}\times\mathcal{Y}\simeq\mathcal{H}_0$, which commutes with
\texttt{withDensity} (LML \texttt{map\_equiv\_withDensity}); $D_0^\rho(h) = \rho(x_0,y_0)$.

\emph{Induction step.} By LML \texttt{history\_succ} and \texttt{hasCondDistrib\_step},
$\Prob^{(n+1)}$ is the image under $e^{-1} : \mathcal{H}_n\times(\mathcal{A}\times\mathcal{Y})\simeq\mathcal{H}_{n+1}$
of $\Prob^{(n)}\otimes(\pi_n\otimes\tilde\kappa)$, where $\pi_n$ is the policy kernel. Now
$\pi_n\otimes\tilde\kappa = (\pi_n\otimes\tilde\kappa_0).\mathrm{withDensity}((h,(a,y))\mapsto\rho(a,y))$
(Lemma~\ref{lem:pb_kernel_compProd_withDensity_right} with $f((h,a),y) := \rho(a,y)$; here
$\tilde\kappa = \tilde\kappa_0.\mathrm{withDensity}(f)$ is a Markov kernel), and by the
induction hypothesis
$\Prob^{(n)} = \Prob_0^{(n)}.\mathrm{withDensity}(D_n^\rho)$. A composition-product of a measure
with a density and a kernel with a density is the composition-product with the product
density (LML \texttt{compProd\_withDensity\_withDensity}):
\[
  \Prob^{(n)}\otimes(\pi_n\otimes\tilde\kappa)
  = \big(\Prob_0^{(n)}\otimes(\pi_n\otimes\tilde\kappa_0)\big).\mathrm{withDensity}\big((h,(a,y))\mapsto D_n^\rho(h)\rho(a,y)\big).
\]
Pushing forward by $e^{-1}$ (\texttt{map\_equiv\_withDensity}) gives
$\Prob^{(n+1)} = \Prob_0^{(n+1)}.\mathrm{withDensity}(D_{n+1}^\rho)$ since
$D_{n+1}^\rho(e^{-1}(h,(a,y))) = D_n^\rho(h)\rho(a,y)$. The integral formulas are
\texttt{lintegral\_withDensity\_eq\_lintegral\_mul} and
\texttt{integral\_withDensity\_eq\_integral\_smul}.
\end{proof}

\textbf{Lean remark.} \texttt{Learning.envDensity ρ n : (Fin n → 𝓐 × 𝓨) → ℝ≥0∞} is the
product $\prod_{t<n}\rho(x_t,y_t)$ over \texttt{Fin n} (so that it applies to the histories
\texttt{finHistory} used throughout), with the recursion \texttt{envDensity\_snoc}; the law
identity is proved first for LML's \texttt{history X Y n : Iic n → 𝓐 × 𝓨} by the induction of
\texttt{hasLaw\_history\_withDensity} (\texttt{map\_history\_eq\_withDensity\_env}) and
transported to \texttt{finHistory} through \texttt{toFinHistory}
(\texttt{map\_finHistory\_eq\_withDensity\_env}). The kernel-level identity
$\pi\otimes(\tilde\kappa_0.\mathrm{withDensity}\,\rho) = (\pi\otimes\tilde\kappa_0).\mathrm{withDensity}(\rho\circ\mathrm{snd})$
is Lemma~\ref{lem:pb_kernel_compProd_withDensity_right}, the right-factor analogue of LML's
\texttt{Kernel.compProd\_withDensity\_left}; it is not in Mathlib or LML and is proved
fiberwise from \texttt{Measure.compProd\_withDensity}.

\begin{lemma}[Gaussian density ratio]\label{lem:pb_gaussian_density_ratio}
\lean{ProbabilityTheory.gaussianReal_eq_withDensity_exp, ProbabilityTheory.gaussianPDFReal_eq_mul_exp, ProbabilityTheory.measurable_gaussianDensityRatio}\leanok
For $m\in\R$, $\Normal(m,1) = \Normal(0,1).\mathrm{withDensity}\big(y\mapsto e^{ym-m^2/2}\big)$.
\end{lemma}

\begin{proof}
\leanok
Both are Lebesgue measure with densities (Mathlib \texttt{gaussianReal\_of\_var\_ne\_zero}):
$\Normal(m,1) = \mathrm{Leb}.\mathrm{withDensity}(\varphi(\cdot-m))$ with
$\varphi(y) = (2\pi)^{-1/2}e^{-y^2/2}$, and
$\varphi(y-m) = \varphi(y)\,e^{ym-m^2/2}$ since $-(y-m)^2/2 = -y^2/2+ym-m^2/2$. Conclude with
\texttt{withDensity\_mul}: $(\mathrm{Leb}.\mathrm{withDensity}\,\varphi).\mathrm{withDensity}(e^{ym-m^2/2}) = \mathrm{Leb}.\mathrm{withDensity}(\varphi\cdot e^{ym-m^2/2})$.
\end{proof}

\begin{corollary}[Likelihood ratio of a linear Gaussian history]\label{cor:pb_gaussian_likelihood}
\lean{Learning.LinearBandit.stepLogLR, Learning.LinearBandit.stepLR, Learning.LinearBandit.measurable_uncurry_stepLR, Learning.LinearBandit.linearGaussianKernel_eq_withDensity, Learning.LinearBandit.likelihood, Learning.LinearBandit.likelihood_pos, Learning.LinearBandit.envDensity_stepLR, Learning.LinearBandit.measurable_likelihood, Learning.LinearBandit.measurable_likelihood_prod, Learning.IsAlgEnvSeq.map_finHistory_eq_withDensity_likelihood, Maiti2026Power.histB, Maiti2026Power.histS, Maiti2026Power.likelihood_eq_gaussTilt}\leanok
\uses{def:env, lem:env_likelihood_ratio, lem:pb_gaussian_density_ratio}
Let $\Xs$ be an action set, $\theta\in\R^d$, $T\ge1$, and for a history $h = (x_t,y_t)_{t<T}$
define
\[
  S_T(h) := \sum_{t<T}x_tx_t^\top\in\PSD,\qquad b_T(h) := \sum_{t<T}y_tx_t\in\R^d,\qquad
  \ell_\theta(h) := \exp\Big(\ip{b_T(h)}{\theta}-\tfrac12\theta^\top S_T(h)\theta\Big) .
\]
For every algorithm with actions in $\Xs$, the law $\Prob_\theta^T$ of the history of length
$T$ under $\nu_\theta$ (Definition~\ref{def:env}) satisfies
$\Prob_\theta^T = \Prob_0^T.\mathrm{withDensity}(\ell_\theta)$, where $\Prob_0^T$ is the law
under $\nu_0$ (pure noise: $y_t\sim\Normal(0,1)$). Moreover $(\theta,h)\mapsto\ell_\theta(h)$ is
jointly measurable and positive. (In Lean, \texttt{withDensity} takes an
\texttt{ℝ≥0∞}-valued density, so the density is \texttt{fun h ↦ ENNReal.ofReal (ℓ θ h)}, and
likewise $\mathrm{ENNReal.ofReal}\circ\rho_\theta$ in the proof; joint measurability of
$(\theta,h)\mapsto\mathrm{ENNReal.ofReal}(\ell_\theta(h))$ follows from that of $\ell$ and
\texttt{ENNReal.measurable\_ofReal}.) In Lean, \texttt{likelihood θ h} is the sum form
$\exp(\sum_{t<T}(y_t\ip{x_t}{\theta} - \tfrac12\ip{x_t}{\theta}^2))$, valid on any action
set; the identification with $\ip{b_T}{\theta} - \tfrac12\theta^\top S_T\theta$ is
\texttt{Maiti2026Power.likelihood\_eq\_gaussTilt} (stated on the unit ball, with
\texttt{histB}, \texttt{histS} for $b_T$, $S_T$).
\end{corollary}

\begin{proof}
\leanok
\uses{lem:env_likelihood_ratio, lem:pb_gaussian_density_ratio}
Lemma~\ref{lem:pb_gaussian_density_ratio} with $m = \ip{x}{\theta}$ gives
$\Normal(\ip{x}{\theta},1) = \Normal(0,1).\mathrm{withDensity}(\rho_\theta(x,\cdot))$ with
$\rho_\theta(x,y) := \exp(y\ip{x}{\theta}-\tfrac12\ip{x}{\theta}^2)$, jointly continuous in
$(\theta,x,y)$. Lemma~\ref{lem:env_likelihood_ratio} then gives the density
$D_T^{\rho_\theta}(h) = \prod_{t<T}\exp(y_t\ip{x_t}{\theta}-\tfrac12\ip{x_t}{\theta}^2) = \ell_\theta(h)$,
since $\sum_ty_t\ip{x_t}{\theta} = \ip{b_T}{\theta}$ and
$\sum_t\ip{x_t}{\theta}^2 = \theta^\top S_T\theta$. Since $\ell_\theta > 0$ is real-valued, the
\texttt{ℝ≥0∞} density $\mathrm{ENNReal.ofReal}(\ell_\theta)$ is finite everywhere and
$\mathrm{ENNReal.ofReal}(\prod_t\rho_\theta(x_t,y_t)) = \prod_t\mathrm{ENNReal.ofReal}(\rho_\theta(x_t,y_t))$
(\texttt{ENNReal.ofReal\_prod\_of\_nonneg}).
\end{proof}

\begin{lemma}[The law of the run as a Markov kernel in $\theta$]\label{lem:pb_kernel_of_likelihood}
\lean{Learning.LinearBandit.noiseHistLaw, Learning.LinearBandit.histKernel, Learning.LinearBandit.histKernel_apply, Learning.IsAlgEnvSeq.map_finHistory_eq_histKernel, Learning.IsAlgEnvSeq.hasLaw_finHistory_histKernel, Learning.LinearBandit.pairKernel, Learning.LinearBandit.pairKernel_apply, Learning.IdentAlg.IsRun.hasLaw_finHistory_out_pairKernel, Learning.LinearBandit.IsPAC.le_measureReal_pairKernel, Learning.LinearBandit.IsPAC.integral_simpleRegret_pairKernel_le, Maiti2026Power.expRegret, Maiti2026Power.stronglyMeasurable_expRegret}\leanok
\uses{def:bai_alg, cor:pb_gaussian_likelihood}
Let $\Xs$ be an action set and $\mathsf{Alg} = (\mathrm{alg},\rho)$ an identification algorithm
with budget $T\ge1$ on $\Xs$ (Definition~\ref{def:bai_alg}). Define, for $\theta\in\R^d$,
\[
  K(\theta) := \Prob_0^T.\mathrm{withDensity}\big(\mathrm{ENNReal.ofReal}\circ\ell_\theta\big),
  \qquad \kappa(\theta) := K(\theta)\otimes\rho .
\]
Then $K$ is a Markov kernel from $\R^d$ to histories of length $T$ with $K(\theta) = \Prob_\theta^T$
for every $\theta$, and $\kappa$ is a Markov kernel from $\R^d$ to pairs (history,
recommendation) with $\kappa(\theta) = \Prob_\theta^T\otimes\rho = \Prob_\theta$, the law of
$(\hist_T,\rec)$ of Definition~\ref{def:bai_alg}, for \emph{every} $\theta$. In particular
$\kappa$ satisfies the hypothesis of Lemma~\ref{lem:pac_expected_regret} for every prior
$\pi$, and for every bounded measurable $g$ of $(\theta,\hist_T,\rec)$ the map
$\theta\mapsto\E_\theta[g(\theta,\hist_T,\rec)] = \int g(\theta,\cdot)\,d\kappa(\theta)$ is
measurable.
\end{lemma}

\begin{proof}
\leanok
\uses{cor:pb_gaussian_likelihood}
$K$ is \texttt{Kernel.withDensity} of the constant kernel $\theta\mapsto\Prob_0^T$ with the
density $(\theta,h)\mapsto\mathrm{ENNReal.ofReal}(\ell_\theta(h))$, which is jointly measurable
by Corollary~\ref{cor:pb_gaussian_likelihood}; this joint measurability is exactly what
\texttt{Kernel.withDensity} requires to produce a kernel, and $K(\theta) = \Prob_\theta^T$ by
the same corollary (\texttt{Kernel.withDensity\_apply}). Each $K(\theta)$ is a probability
measure, so $K$ is a Markov kernel. Then $\kappa := K\otimes\rho'$ with
$\rho' := \rho.\mathrm{prodMkLeft}$ (the kernel $(\theta,h)\mapsto\rho(h)$) is a Markov kernel
(\texttt{Kernel.compProd} of Markov kernels) with
$\kappa(\theta) = K(\theta)\otimes\rho = \Prob_\theta^T\otimes\rho$
(\texttt{Kernel.compProd\_apply\_eq\_compProd\_sectR}, \texttt{Kernel.prodMkLeft\_apply}),
which is $\Prob_\theta$ by Definition~\ref{def:bai_alg} and the uniqueness of the law of the
history. The measurability of $\theta\mapsto\int g\,d\kappa(\theta)$ is the measurability of
integrals against a kernel (\texttt{Kernel.measurable\_lintegral},
\texttt{Measurable.lintegral\_kernel\_prod\_right}, and
\texttt{StronglyMeasurable.integral\_kernel\_prod\_right} for real-valued bounded $g$).
\end{proof}

\textbf{Lean remark.} This is the construction announced in the Lean remark of
Lemma~\ref{lem:pac_expected_regret}: the measurability of
$\theta\mapsto\mathrm{trajMeasure}(\mathrm{alg},\nu_\theta)$ is never needed, because the kernel
is built from the single reference law $\Prob_0^T$ and the explicit density $\ell_\theta$.

\section{The joint law of the parameter, the history and the recommendation}
\label{sec:pre_bayes_joint}

\begin{definition}[Bayesian joint law and posterior quantities]\label{def:bayes_joint}
\lean{Maiti2026Power.noisePairLaw, Maiti2026Power.pairKernel_eq_withDensity, Maiti2026Power.bayesJointLaw, Maiti2026Power.bayesJointLaw_eq_withDensity, ProbabilityTheory.gaussPrior, ProbabilityTheory.postPrec, ProbabilityTheory.postCov, ProbabilityTheory.postMean, ProbabilityTheory.posDef_postPrec, ProbabilityTheory.posDef_postCov, ProbabilityTheory.toEuclideanCLM_postPrec_postMean, ProbabilityTheory.sub_toEuclideanCLM_postMean}\leanok
\uses{def:bai_alg, cor:pb_gaussian_likelihood, lem:pb_kernel_of_likelihood, lem:loewner_inv, lem:eigen_quadratic_bounds}
Let $d\ge1$, $\Xs$ an action set, $\mathsf{Alg} = (\mathrm{alg},\rho)$ an identification algorithm
with budget $T\ge1$ on $\Xs$ (Definition~\ref{def:bai_alg}) and $\sigma>0$. Let
$\pi := \Normal(0,\sigma^2I_d)$ (the prior). The \emph{joint law} of $(\theta,\hist_T,\rec)$ is
the probability measure on $\R^d\times\mathcal{H}_T\times\Xs$
\[
  Q := \pi\otimes\kappa,\qquad \kappa(\theta) = \Prob_\theta^T\otimes\rho = \Prob_\theta ,
\]
where $\kappa$ is the Markov kernel of Lemma~\ref{lem:pb_kernel_of_likelihood} (built from
$\Prob_\theta^T = \Prob_0^T.\mathrm{withDensity}(\ell_\theta)$,
Corollary~\ref{cor:pb_gaussian_likelihood}). Equivalently, with $Q_0 := \Prob_0^T\otimes\rho$
the law of (history, recommendation) under pure noise,
$\kappa(\theta) = Q_0.\mathrm{withDensity}((h,x)\mapsto\ell_\theta(h))$ and
\[
  Q = (\pi\times Q_0).\mathrm{withDensity}\big((\theta,h,x)\mapsto\ell_\theta(h)\big) .
\]
We write $\E$ for the expectation under $Q$ and, for
a history $h$ of length $T$,
\[
  V_T(h) := \big(\sigma^{-2}I_d+S_T(h)\big)^{-1}\in\PD,\qquad m_T(h) := V_T(h)\,b_T(h)\in\R^d ,
\]
so that $(\sigma^{-2}I + S_T)m_T = b_T$ and $b_T - S_Tm_T = \sigma^{-2}m_T$.
\end{definition}

\textbf{Lean remark.} The pure-noise laws are \texttt{noiseHistLaw} and \texttt{noisePairLaw},
the kernel is \texttt{pairKernel}, the joint law is \texttt{bayesJointLaw A T σ} and its
product-with-density form is \texttt{bayesJointLaw\_eq\_withDensity}. The posterior quantities
are defined for a general tilt $(b, S)$: \texttt{postPrec σ S} $= \sigma^{-2}I + S$,
\texttt{postCov σ S} its inverse and \texttt{postMean σ b S} $= V b$; they are measurable
functions of the history through $b_T$, $S_T$, but this measurability is never needed (see the
Lean remark after Lemma~\ref{lem:posterior_second_moment}).

\begin{lemma}[Fubini for the joint law]\label{lem:pb_joint_fubini}
\lean{Maiti2026Power.integral_bayesJointLaw, Maiti2026Power.integrable_bayesJointLaw_iff, Maiti2026Power.integral_likelihood_noisePairLaw, Maiti2026Power.integrable_likelihood_noisePairLaw, Maiti2026Power.integrable_likelihood_mul_prod, Maiti2026Power.integral_likelihood_mul_prod, Maiti2026Power.integral_indicator_mul_simpleRegret_bayesJointLaw}\leanok
\uses{def:bayes_joint, cor:pb_gaussian_likelihood}
For every measurable $f : \R^d\times\mathcal{H}_T\times\Xs\to[0,\infty]$,
\[
  \E[f(\theta,\hist_T,\rec)]
  = \int\!\!\int\!\!\int f(\theta,h,x)\,\rho(h,dx)\,\Prob_\theta^T(dh)\,\pi(d\theta)
  = \int\!\!\int\ell_\theta(h)\,f(\theta,h,x)\,Q_0(dh,dx)\,\pi(d\theta)
  = \int\!\!\int\ell_\theta(h)\,f(\theta,h,x)\,\pi(d\theta)\,Q_0(dh,dx),
\]
and the same holds for real-valued $f$ that are $Q$-integrable, i.e.\ such that
$\ell\cdot f$ is $\pi\times Q_0$-integrable. In particular $\int\ell_\theta\,dQ_0 = 1$ for every
$\theta$, the marginal law of $\theta$ is $\pi$, the conditional law
of $(\hist_T,\rec)$ given $\theta$ is $\Prob_\theta$, and for a function $g$ of $\theta$ alone,
$\int\ell_\theta(h)g(\theta)\,d(\pi\times Q_0) = \int g\,d\pi$.
\end{lemma}

\begin{proof}
\leanok
\uses{cor:pb_gaussian_likelihood}
The first equality is Tonelli/Fubini for composition-products (Mathlib
\texttt{Measure.lintegral\_compProd}, \texttt{Measure.integral\_compProd}), the second is the
density of Corollary~\ref{cor:pb_gaussian_likelihood} and the product form of $Q$
(\texttt{integral\_withDensity\_eq\_integral\_toReal\_smul}), the third is Fubini for the product
$\pi\times Q_0$ (\texttt{integral\_prod}, \texttt{integral\_prod\_symm}). Since
$\kappa(\theta)$ is a probability measure, $\int\ell_\theta\,dQ_0 = \kappa(\theta)(\text{all}) = 1$.
\end{proof}

\section{Posterior mean and posterior second moment}
\label{sec:pre_bayes_posterior}

\begin{lemma}[Stein's identity for functions of exponential growth]\label{lem:gauss_ibp_exp}
\lean{ProbabilityTheory.integral_sub_mul_gaussianReal_of_integrable, ProbabilityTheory.integral_mul_gaussianReal_zero_one_of_integrable, ProbabilityTheory.integrable_exp_mul_abs_gaussianReal, ProbabilityTheory.IsGaussian.integrable_of_abs_le_mul_exp, ProbabilityTheory.IsGaussian.integrable_norm_mul_of_abs_le_mul_exp, ProbabilityTheory.IsGaussian.integrable_inner_mul_of_abs_le_mul_exp, ProbabilityTheory.IsGaussian.integrable_fderiv_apply_of_norm_fderiv_le_mul_exp, ProbabilityTheory.integral_apply_mul_stdGaussian_of_le_exp, ProbabilityTheory.integral_inner_mul_stdGaussian_euclidean_of_le_exp, ProbabilityTheory.integral_inner_mul_stdGaussian_of_le_exp, ProbabilityTheory.integral_inner_mul_multivariateGaussian_of_le_exp, ProbabilityTheory.norm_toLp_insertNth_le, ProbabilityTheory.abs_apply_toLp_insertNth_le}\leanok
\uses{lem:gauss_ibp, lem:psf_ibp_1d, def:pg_gaussian_vector}
Let $F : \R^d\to\R$ be $C^1$ with $\abs{F(x)} \le Ae^{B\norm x}$ and
$\norm{DF(x)} \le Ae^{B\norm x}$ for all $x$ (some constants $A, B$). Then for
$g\sim\Normal(0,I_d)$ and every $a\in\R^d$, $\E[\ip{a}{g}F(g)] = \E[DF(g)\,a]$, and for
$X\sim\Normal(0,S)$ with $S\in\PSD$, $\E[\ip{a}{X}F(X)] = \E[DF(X)\,(Sa)]$.
\end{lemma}

\begin{proof}
\leanok
\uses{lem:gauss_ibp, lem:psf_ibp_1d}
The one-dimensional identity $\E[\xi f(\xi)] = \E[f'(\xi)]$ (Lemma~\ref{lem:psf_ibp_1d}) only
needs the integrability of $f$, $\xi f(\xi)$ and $f'$, which holds for functions of exponential
growth since $e^{B\abs t}$ has all Gaussian moments. The multivariate identity follows by the
coordinate-wise Fubini argument of Lemma~\ref{lem:gauss_ibp}: the section
$t\mapsto F(x + te_i)$ has growth $Ae^{B\norm{x}}e^{B\abs t}$, and the integrability of
$\ip{a}{g}F(g)$ and $DF(g)a$ under a Gaussian measure is Fernique's theorem.
\end{proof}

\begin{lemma}[Gaussian tilt: derivative and growth]\label{lem:pb_complete_square_identity}
\lean{ProbabilityTheory.gaussTilt, ProbabilityTheory.gaussTilt_pos, ProbabilityTheory.gaussTilt_le, ProbabilityTheory.hasFDerivAt_gaussTilt, ProbabilityTheory.contDiff_gaussTilt, ProbabilityTheory.hasFDerivAt_affine_mul_gaussTilt, ProbabilityTheory.abs_affine_mul_gaussTilt_le, ProbabilityTheory.norm_fderiv_affine_mul_gaussTilt_le, ProbabilityTheory.inner_toEuclideanCLM_symm, ProbabilityTheory.inner_toEuclideanCLM_nonneg}\leanok
Let $S\in\PSD$, $b\in\R^d$ and $\ell(\theta) := \exp(\ip{b}{\theta}-\tfrac12\theta^\top S\theta)$.
Then $\ell > 0$, $\ell(\theta) \le e^{\norm b\norm\theta}$, $\ell$ is $C^1$ with
$D\ell(\theta) = \ell(\theta)\,\ip{b - S\theta}{\cdot}$, and for $c\in\R^d$, $r\in\R$ the
function $\theta\mapsto(\ip{c}{\theta}+r)\ell(\theta)$ and its derivative have exponential
growth: they are bounded by $A e^{(\norm b+2)\norm\theta}$ with
$A = (\norm c+\abs r)(\norm b+\norm S+1)+\norm c$.
\end{lemma}

\begin{proof}
\leanok
$\theta^\top S\theta\ge0$ gives the bound; the derivative is the chain rule, and the growth
bounds use $\norm\theta \le e^{\norm\theta}$.
\end{proof}

\begin{lemma}[Stein's identity for the affine tilt]\label{lem:pb_gaussian_integrals}
\lean{ProbabilityTheory.integral_inner_mul_affine_mul_gaussTilt, ProbabilityTheory.integrable_affine_mul_affine_mul_gaussTilt, ProbabilityTheory.integrable_affine_mul_gaussTilt, ProbabilityTheory.integrable_gaussTilt, ProbabilityTheory.integrable_gaussTilt_smul, ProbabilityTheory.integral_inner_mul_gaussTilt, ProbabilityTheory.integral_gaussTilt_pos}\leanok
\uses{lem:gauss_ibp_exp, lem:pb_complete_square_identity}
Let $\sigma>0$, $\pi = \Normal(0,\sigma^2I_d)$, $S\in\PSD$, $b\in\R^d$ and $\ell$ as above.
For all $a, c\in\R^d$ and $r\in\R$,
\[
  \int\ip{a}{\theta}\,(\ip{c}{\theta}+r)\,\ell(\theta)\,\pi(d\theta)
  = \sigma^2\int\big(\ip{c}{a} + (\ip{c}{\theta}+r)\ip{b-S\theta}{a}\big)\ell(\theta)\,\pi(d\theta).
\]
All the functions $\theta\mapsto(\ip{a}{\theta}+r)(\ip{c}{\theta}+s)\ell(\theta)$ are
$\pi$-integrable, $Z := \int\ell\,d\pi\in(0,\infty)$, and
$\int\ip{a}{\theta}\ell(\theta)\,\pi(d\theta) = \ip{a}{\int\ell(\theta)\theta\,\pi(d\theta)}$.
\end{lemma}

\begin{proof}
\leanok
\uses{lem:gauss_ibp_exp, lem:pb_complete_square_identity}
Lemma~\ref{lem:gauss_ibp_exp} for $\Normal(0,\sigma^2I)$ applied to
$F = (\ip{c}{\cdot}+r)\ell$, whose derivative is computed in
Lemma~\ref{lem:pb_complete_square_identity}; $Sa = \sigma^2a$ here.
\end{proof}

\begin{lemma}[Gaussian posterior identities]\label{lem:gaussian_complete_square}
\lean{ProbabilityTheory.integral_gaussTilt_smul_eq, ProbabilityTheory.integral_inner_sub_postMean_mul_gaussTilt, ProbabilityTheory.integral_inner_sub_mul_inner_sub_mul_gaussTilt_add, ProbabilityTheory.integral_inner_single_sub_mul_inner_sub_mul_gaussTilt, ProbabilityTheory.integral_norm_sq_sub_postMean_mul_gaussTilt, ProbabilityTheory.integral_norm_sq_mul_gaussTilt, ProbabilityTheory.integrable_inner_sub_mul_inner_sub_mul_gaussTilt, ProbabilityTheory.integrable_inner_sub_mul_gaussTilt}\leanok
\uses{lem:pb_gaussian_integrals, def:bayes_joint}
Let $\sigma>0$, $S\in\PSD$, $b\in\R^d$, $V := (\sigma^{-2}I_d+S)^{-1}$, $m := Vb$,
$\pi := \Normal(0,\sigma^2I_d)$, $\ell(\theta) := \exp(\ip{b}{\theta}-\tfrac12\theta^\top S\theta)$
and $Z := \int\ell\,d\pi > 0$. Then
\[
  \int\ell(\theta)\,\theta\,\pi(d\theta) = Z\,m,\qquad
  \int\ip{c}{\theta-m}\,\ell(\theta)\,\pi(d\theta) = 0\ \ (c\in\R^d),
\]
\[
  \int\norm{\theta-m}^2\ell(\theta)\,\pi(d\theta) = \tr(V)\,Z,\qquad
  \int\norm{\theta}^2\ell(\theta)\,\pi(d\theta) = (\tr V + \norm m^2)\,Z .
\]
(Equivalently, $\ell\,d\pi/Z$ is the Gaussian $\Normal(m,V)$; we only need these moments.)
\end{lemma}

\begin{proof}
\leanok
\uses{lem:pb_gaussian_integrals, def:bayes_joint}
\emph{Mean.} Lemma~\ref{lem:pb_gaussian_integrals} with $c = 0$, $r = 1$: writing
$u := \int\ell(\theta)\theta\,d\pi$, for every $a$,
$\ip{a}{u} = \sigma^2(\ip{b}{a}Z - \ip{a}{Su})$ (using $\ip{S\theta}{a} = \ip{\theta}{Sa}$).
Hence $u + \sigma^2Su = \sigma^2Zb$, i.e.\ $(\sigma^{-2}I+S)u = Zb$ and $u = Zm$. The second
identity is $\ip{c}{u} - \ip{c}{m}Z = 0$.

\emph{Second moment.} Let $\Phi(a,c) := \int\ip{a}{\theta-m}\ip{c}{\theta-m}\ell\,d\pi$.
Lemma~\ref{lem:pb_gaussian_integrals} with $r = -\ip{c}{m}$ (so that
$\ip{c}{\theta}+r = \ip{c}{\theta-m}$), together with
$b - S\theta = \sigma^{-2}m - S(\theta-m)$ (Definition~\ref{def:bayes_joint}) and the mean
identity, gives $\Phi(a,c) + \sigma^2\Phi(Sa,c) = \sigma^2\ip{c}{a}Z$ for all $a, c$. With
$a := \sigma^{-2}Ve_i$ one has $a + \sigma^2Sa = \sigma^2(\sigma^{-2}I+S)a = e_i$, so
$\Phi(e_i,c) = \ip{c}{Ve_i}Z$; taking $c = e_i$ and summing over $i$ gives
$\int\norm{\theta-m}^2\ell\,d\pi = \tr(V)Z$. Finally
$\norm\theta^2 = \norm{\theta-m}^2 + 2\ip{m}{\theta-m} + \norm m^2$ and the cross term
integrates to $0$.
\end{proof}

\textbf{Lean remark.} \texttt{PosteriorGaussian.lean} contains this section for the general tilt
$(b, S)$ (\texttt{gaussTilt b S}, \texttt{postCov σ S}, \texttt{postMean σ b S}); it is applied
to $b = b_T(h)$, $S = S_T(h)$ for each fixed history $h$ through
\texttt{Maiti2026Power.likelihood\_eq\_gaussTilt}.

\begin{lemma}[Posterior mean identity]\label{lem:posterior_mean}
\lean{Maiti2026Power.postU, Maiti2026Power.stronglyMeasurable_postU, Maiti2026Power.integrable_likelihood_mul_inner_prod, Maiti2026Power.integrable_norm_postU, Maiti2026Power.integral_inner_bayesJointLaw_le}\leanok
\uses{def:bayes_joint, lem:pb_joint_fubini, lem:gaussian_complete_square}
Let $\Xs \subseteq \ball_d$ and define, for a history $h$,
$u_T(h) := \int\ell_\theta(h)\,\theta\,\pi(d\theta) = Z_T(h)\,m_T(h)$ with
$Z_T(h) := \int\ell_\theta(h)\,\pi(d\theta)$; $u_T$ is a measurable function of $h$ with
$\int\norm{u_T}\,dQ_0 < \infty$. Then
\[
  \E\ip{\rec}{\theta} = \int\ip{x}{u_T(h)}\,Q_0(dh,dx) \le \int\norm{u_T(h)}\,Q_0(dh,dx) .
\]
\end{lemma}

\begin{proof}
\leanok
\uses{lem:pb_joint_fubini, lem:gaussian_complete_square}
By Lemma~\ref{lem:pb_joint_fubini} with the $\theta$-integral innermost (the integrand
$\ell_\theta(h)\ip{x}{\theta}$ is $\pi\times Q_0$-integrable, being dominated by
$\ell_\theta(h)\norm\theta$ whose integral is $\E_\pi\norm\theta$),
$\E\ip{\rec}{\theta} = \int\int\ell_\theta(h)\ip{x}{\theta}\,\pi(d\theta)\,Q_0(dh,dx)
= \int\ip{x}{u_T(h)}\,Q_0(dh,dx)$ (Lemma~\ref{lem:pb_gaussian_integrals}), and
$\ip{x}{u_T(h)} \le \norm{u_T(h)}$ since $\norm x\le1$. Measurability of $u_T$ is that of a
parametric integral (Mathlib \texttt{StronglyMeasurable.integral\_prod\_right'}).
\end{proof}

\begin{lemma}[Posterior second moment]\label{lem:posterior_second_moment}
\lean{Maiti2026Power.postZ, Maiti2026Power.postN, Maiti2026Power.stronglyMeasurable_postZ, Maiti2026Power.stronglyMeasurable_postN, Maiti2026Power.postZ_pos, Maiti2026Power.norm_postU_sq_le, Maiti2026Power.norm_postU_le, Maiti2026Power.integrable_postZ, Maiti2026Power.integral_postZ, Maiti2026Power.integrable_postN, Maiti2026Power.integral_postN, Maiti2026Power.integral_norm_postU_le}\leanok
\uses{def:bayes_joint, lem:pb_joint_fubini, lem:gaussian_complete_square, lem:trace_inv_lower}
Let $\Xs \subseteq \ball_d$ and $N_T(h) := \int\norm\theta^2\ell_\theta(h)\,\pi(d\theta)$. For
every history $h$, with $c := d^2/(d\sigma^{-2}+T)$,
\[
  \norm{u_T(h)}^2 = Z_T(h)^2\norm{m_T(h)}^2 \le Z_T(h)\big(N_T(h) - c\,Z_T(h)\big),
\]
and $\int Z_T\,dQ_0 = 1$, $\int N_T\,dQ_0 = \E_\pi\norm\theta^2 = \sigma^2d$. Consequently
\[
  \int\norm{u_T(h)}\,Q_0(dh,dx)\ \le\ \sqrt{\sigma^2d - \frac{d^2}{d\sigma^{-2}+T}} .
\]
\end{lemma}

\begin{proof}
\leanok
\uses{lem:pb_joint_fubini, lem:gaussian_complete_square, lem:trace_inv_lower}
By Lemma~\ref{lem:gaussian_complete_square} for the fixed history $h$,
$u_T = Z_Tm_T$ and $N_T = (\tr V_T + \norm{m_T}^2)Z_T$, so
$Z_T\norm{m_T}^2 = N_T - \tr(V_T)Z_T \le N_T - cZ_T$ by Lemma~\ref{lem:trace_inv_lower}; multiply
by $Z_T$. The integrals of $Z_T$ and $N_T$ are Lemma~\ref{lem:pb_joint_fubini} (Fubini, the
$\theta$-integral outermost). Finally, by the Cauchy--Schwarz inequality in $L^2(Q_0)$ applied
to $\sqrt{Z_T}$ and $\sqrt{N_T - cZ_T}$,
$\int\norm{u_T}\,dQ_0 \le \int\sqrt{Z_T}\sqrt{N_T-cZ_T}\,dQ_0
\le \big(\int Z_T\,dQ_0\big)^{1/2}\big(\int(N_T - cZ_T)\,dQ_0\big)^{1/2} = \sqrt{\sigma^2d - c}$.
\end{proof}

\textbf{Lean remark.} This route never integrates $m_T$ or $\tr V_T$ over histories, only
$Z_T$, $N_T$ and $u_T$, which are parametric integrals in $h$ (measurable by
\texttt{StronglyMeasurable.integral\_prod\_right'}); the outline's integrability lemma
$\E\norm{b_T}^2 < \infty$ and the measurability of $h\mapsto V_T(h)$ (matrix inversion) are
not needed.

\section{The unit-ball lower bound}
\label{sec:pre_bayes_ball}

\begin{lemma}[Lower bound on the posterior trace]\label{lem:trace_inv_lower}
\lean{Maiti2026Power.le_trace_postCov_histS, Maiti2026Power.trace_histS_le, Maiti2026Power.posSemidef_histS}\leanok
\uses{def:bayes_joint, lem:trace_inv_amhm}
If $\Xs\subseteq\ball_d$, then for every history $h$, $\tr V_T(h)\ge\frac{d^2}{d\sigma^{-2}+T}$.
\end{lemma}

\begin{proof}
\leanok
\uses{lem:trace_inv_amhm}
$\tr S_T(h) = \sum_{t<T}\norm{x_t}^2\le T$, so $\tr(\sigma^{-2}I+S_T(h))\le d\sigma^{-2}+T$, and
Lemma~\ref{lem:trace_inv_amhm} applied to the positive definite matrix $\sigma^{-2}I+S_T(h)$
gives $\tr V_T(h)\ge d^2/\tr(\sigma^{-2}I+S_T(h))\ge d^2/(d\sigma^{-2}+T)$.
\end{proof}

\begin{lemma}[Expected norm of a Gaussian vector, lower bound]\label{lem:pb_norm_lower}
\lean{Maiti2026Power.le_integral_norm_gaussPrior, Maiti2026Power.integral_norm_sq_gaussPrior, Maiti2026Power.integrable_norm_gaussPrior, Maiti2026Power.integrable_norm_sq_gaussPrior}\leanok
\uses{lem:gauss_norm_lower, lem:gauss_norm_moments}
If $\theta\sim\Normal(0,\sigma^2I_d)$ then $\E\norm\theta^2 = \sigma^2d$ and
$\E\norm\theta\ge\sqrt{2/\pi}\,\sigma\sqrt d$.
\end{lemma}

\begin{proof}
\leanok
\uses{lem:gauss_norm_lower, lem:gauss_norm_moments}
Lemma~\ref{lem:gauss_norm_moments}(i) and Lemma~\ref{lem:gauss_norm_lower} with the weights
$c_i = 1/\sqrt d$ and $S = \sigma^2I_d$.
\end{proof}

\begin{lemma}[Bayesian simple regret on the ball]\label{lem:bayes_regret_ball}
\lean{Maiti2026Power.integral_simpleRegret_bayesJointLaw_ge}\leanok
\uses{def:bayes_joint, def:regret, lem:posterior_mean, lem:posterior_second_moment, lem:pb_norm_lower}
Let $\Xs = \ball_d$. Under the joint law $Q$ of Definition~\ref{def:bayes_joint},
\[
  \E[r(\rec,\theta)] = \E\norm\theta-\E\ip{\rec}{\theta}
  \ \ge\ \sqrt{2/\pi}\,\sigma\sqrt d-\sqrt{\sigma^2d-\frac{d^2}{d\sigma^{-2}+T}} .
\]
\end{lemma}

\begin{proof}
\leanok
\uses{lem:posterior_mean, lem:posterior_second_moment, lem:pb_norm_lower}
On the unit ball, $\max_{x\in\ball_d}\ip{x}{\theta} = \norm\theta$, so
$r(\rec,\theta) = \norm\theta-\ip{\rec}{\theta}$ (Lemma~\ref{lem:block_ball_basic}). The marginal
of $\theta$ is $\pi$, so $\E\norm\theta \ge \sqrt{2/\pi}\sigma\sqrt d$
(Lemma~\ref{lem:pb_norm_lower}), and
$\E\ip{\rec}{\theta} \le \int\norm{u_T}\,dQ_0 \le \sqrt{\sigma^2d - d^2/(d\sigma^{-2}+T)}$
by Lemmas~\ref{lem:posterior_mean} and~\ref{lem:posterior_second_moment}.
\end{proof}

\begin{theorem}[Simple-regret lower bound on the unit ball]\label{thm:ball_simple_regret_lower}
\lean{Maiti2026Power.simpleRegret_unitBall_nonneg, Maiti2026Power.simpleRegret_unitBall_le, Maiti2026Power.expRegret_nonneg, Maiti2026Power.expRegret_le, Maiti2026Power.integrable_expRegret, Maiti2026Power.integrable_simpleRegret_bayesJointLaw, Maiti2026Power.integral_norm_sq_bayesJointLaw, Maiti2026Power.integral_indicator_mul_simpleRegret_bayesJointLaw_ge, Maiti2026Power.le_measureReal_norm_le_gaussPrior, Maiti2026Power.exists_norm_le_and_le_expRegret, Maiti2026Power.exists_norm_le_and_le_expRegret_unitBall}\leanok
\uses{def:bayes_joint, def:regret, lem:bayes_regret_ball, lem:pb_joint_fubini, lem:gauss_norm_moments}
Let $d\ge1$, $T\ge1$ and $\mathsf{Alg}$ be any identification algorithm with budget $T$ on
$\Xs = \ball_d$. With the prior $\pi = \Normal(0,\sigma^2I_d)$, $\sigma^2 := d/(4T)$:
\begin{enumerate}
  \item[(i)] $\E[r(\rec,\theta)]\ge\dfrac{d}{8\sqrt T}$ under the joint law $Q$;
  \item[(ii)] there exists $\theta\in\R^d$ with $\norm\theta\le\dfrac{10\,d}{\sqrt T}$ and
        $\E_\theta[r(\rec,\theta)]\ge\dfrac{3d}{40\sqrt T}$.
\end{enumerate}
\end{theorem}

\begin{proof}
\leanok
\uses{lem:bayes_regret_ball, lem:pb_joint_fubini, lem:gauss_norm_moments, lem:pb_kernel_of_likelihood}
(i) With $\sigma^2 = d/(4T)$ we have $d\sigma^{-2} = 4T$, hence
$\frac{d^2}{d\sigma^{-2}+T} = \frac{d^2}{5T}$ and $\sigma^2d-\frac{d^2}{5T} = \frac{d^2}{4T}-\frac{d^2}{5T} = \frac{d^2}{20T}$,
while $\sqrt{2/\pi}\,\sigma\sqrt d = \sqrt{2/\pi}\,\frac{d}{2\sqrt T} \ge 0.79\cdot\frac{d}{2\sqrt T}$
(as $\pi < 3.15$). Lemma~\ref{lem:bayes_regret_ball} gives
\[
  \E[r(\rec,\theta)]\ge\frac{d}{\sqrt T}\Big(\frac{0.79}{2}-\frac{1}{4.47}\Big)
  \ge\frac{d}{\sqrt T}\,(0.395-0.2238) \ge \frac{d}{8\sqrt T} ,
\]
using $4.47^2 \le 20$.

(ii) Let $R := \frac{10d}{\sqrt T} = 20\sigma\sqrt d$.
Since $0 \le r(\rec,\theta)\le2\norm\theta$ on the ball and
$\norm\theta\indic\{\norm\theta>R\}\le\norm\theta^2/R$,
\[
  \E\big[r(\rec,\theta)\indic\{\norm\theta>R\}\big]\le\frac{2\,\E\norm\theta^2}{R}
  = \frac{2\sigma^2d}{R} = \frac{d}{20\sqrt T} ,
\]
so $\E[r\indic\{\norm\theta\le R\}] \ge \frac{d}{8\sqrt T}-\frac{d}{20\sqrt T} = \frac{3d}{40\sqrt T}$.
By Lemma~\ref{lem:pb_joint_fubini}, $\E[r\indic\{\norm\theta\le R\}] = \int_{\norm\theta\le R}\E_\theta[r(\rec,\theta)]\,\pi(d\theta)$,
where $\theta\mapsto\E_\theta[r(\rec,\theta)]$ is measurable
(Lemma~\ref{lem:pb_kernel_of_likelihood}, $r$ being bounded by $2\norm\theta$).
If every $\theta$ with $\norm\theta\le R$ had $\E_\theta[r]<\frac{3d}{40\sqrt T}$, the
nonnegative function $g := \indic\{\norm\theta\le R\}\,(\frac{3d}{40\sqrt T} - \E_\theta[r])$
would be positive on $\{\norm\theta\le R\}$, a set of positive $\pi$-measure (by Markov's
inequality $\pi(\norm\theta>R)\le\sigma^2d/R^2 = 1/400$), so $\int g\,d\pi > 0$; but
$\int g\,d\pi = \frac{3d}{40\sqrt T}\pi(\norm\theta\le R) - \E[r\indic\{\norm\theta\le R\}] \le 0$,
a contradiction. Hence some $\theta$ with $\norm\theta\le R$ satisfies
$\E_\theta[r(\rec,\theta)]\ge\frac{3d}{40\sqrt T}$.
\end{proof}

\begin{remark}[Constants]\label{rem:pb_constants}
The outline claimed the constant $0.13$ in (i), from
$\frac{1}{2\sqrt2}-\frac{1}{\sqrt{20}} = 0.12995\ldots$, which is slightly \emph{below} $0.13$; we
use $1/8$ instead (the sharper Lemma~\ref{lem:pb_norm_lower} would even give $0.17$). The
truncation radius was changed from $16\sigma\sqrt d$ to $20\sigma\sqrt d$
so that Corollary~\ref{cor:ball_pac_lower} keeps the outline's headline constants
($\delta\le1/500$, $T\ge d^2/(1000\eps^2)$): with $R = k\sigma\sqrt d$ the final coefficient is
$\frac18-\frac1k-k\delta$, which for $\delta = 1/500$ is maximized near $k = 22$; $k = 20$ gives
$7/200$. The choice $\sigma^2 = d/(4T)$ is not optimal either but keeps the arithmetic simple.
\end{remark}

\begin{corollary}[PAC lower bound on the unit ball]\label{cor:ball_pac_lower}
\lean{Maiti2026Power.unitBall_le_budget_of_isPAC, Maiti2026Power.pairKernel_zero_eq, Maiti2026Power.not_isPAC_unitBall_of_isFixedBudget_zero}\leanok
\uses{def:pac, lem:pac_expected_regret, lem:pb_kernel_of_likelihood, thm:ball_simple_regret_lower}
Let $d\ge2$, $\eps>0$ and $\delta\le1/500$. Every $(\eps,\delta)$-PAC identification algorithm
on $\ball_d$ with budget $T$ satisfies
\[
  T\ge\frac{49\,d^2}{40000\,\eps^2}\ge\frac{d^2}{1000\,\eps^2} .
\]
\end{corollary}

\begin{proof}
\leanok
\uses{lem:pac_expected_regret, lem:pb_kernel_of_likelihood, thm:ball_simple_regret_lower}
If $T = 0$ the law of the recommendation does not depend on $\theta$, and the two instances
$\pm2\eps e_1$ have regrets summing to $4\eps$ at every arm, so the algorithm cannot be
$(\eps,\delta)$-PAC for $\delta<1/2$ (as in Lemma~\ref{lem:singular_design_fails}).
Let $T\ge1$ and let $\theta$ be given by Theorem~\ref{thm:ball_simple_regret_lower}(ii). On
the ball the regret lies in $[0, 2\norm\theta]$, so Lemma~\ref{lem:pac_expected_regret}
applied to the law $\Prob_\theta = \kappa(\theta)$ of Lemma~\ref{lem:pb_kernel_of_likelihood}
gives $\E_\theta[r(\rec,\theta)]\le\eps+2\delta\norm\theta\le\eps+\frac{20\delta d}{\sqrt T}$
(if $2\norm\theta < \eps$ this is trivial). Combining
with $\E_\theta[r]\ge\frac{3d}{40\sqrt T}$,
\[
  \Big(\frac{3}{40}-20\delta\Big)\frac{d}{\sqrt T}\le\eps,\qquad
  \frac{3}{40}-20\delta\ge\frac{3}{40}-\frac{1}{25} = \frac{7}{200},
\]
hence $\sqrt T\ge\frac{7d}{200\,\eps}$ and $T\ge\frac{49d^2}{40000\eps^2}$; finally
$40000/49<1000$.
\end{proof}

\textbf{Lean remark.} The unit ball is \texttt{unitBall ι = Metric.closedBall 0 1} in
\texttt{EuclideanSpace ℝ ι}; it is compact and spans $\R^d$, hence an action set. The
expectation $\E_\theta[r(\rec,\theta)]$ is \texttt{expRegret T A θ}, the integral of a bounded
function under $\Prob_\theta^T\otimes\rho = \kappa(\theta)$ (\texttt{pairKernel A T θ}), and its
measurability in $\theta$ is Mathlib's \texttt{StronglyMeasurable.integral\_kernel\_prod\_right'}.

